\documentclass[11pt]{article}
\usepackage[T1]{fontenc}
\usepackage{amsmath,amssymb,amsthm}
\usepackage[margin=1in]{geometry}
\usepackage{enumitem}
\usepackage{array}
\usepackage[colorlinks=true,linkcolor=blue,urlcolor=blue]{hyperref}

\newtheorem{theorem}{Theorem}[section]
\newtheorem{proposition}[theorem]{Proposition}
\newtheorem{lemma}[theorem]{Lemma}
\newtheorem{corollary}[theorem]{Corollary}
\newtheorem{definition}[theorem]{Definition}
\newtheorem{warning}[theorem]{Warning}
\newtheorem{firewall}[theorem]{Claim firewall}
\newtheorem{decision}[theorem]{Next decision}
\newtheorem{assessment}[theorem]{Assessment}
\newtheorem{ledger}[theorem]{Acquisition ledger}
\newtheorem{record}[theorem]{Record}
\newtheorem{audit}[theorem]{Audit}
\newtheorem{criterion}[theorem]{Criterion}
\theoremstyle{remark}
\newtheorem{remark}[theorem]{Remark}

\title{Renewal Standard Model--Einstein Continuum Research Notes\\
Twentieth Cycle, V1}
\author{}
\date{12 September 2026}

\begin{document}
\maketitle

\begin{abstract}
This compact restart continues the Renewal Standard Model--Einstein
continuum programme after the nineteenth cycle reached V100 and was
frozen.  We first give context and a theorem-level ledger of the frozen
cycle without reproducing its proofs.  We then open the non-abelian
sector in the only setting where the programme has exact tools: the
plane-symmetric \(SU(2)\) Einstein--Yang--Mills system with one
transverse colour component in each of two colours, whose diagonal
metric ansatz is proved consistent.  Its abelian limit is two Maxwell
fields; its first non-abelian effect, computed exactly at linear order
in the coupling squared, is an effective mass \(\tfrac9{16}e^2\) that
each colour wave acquires from the other, a resonant third harmonic,
and a colour-magnetic flux whose stress the emergent geometry reads as
a transverse relief and a longitudinal contraction.  No full Standard
Model--Einstein continuum is claimed.
\end{abstract}

\tableofcontents

\section{Context, lineage, and target}
\label{sec:v20v1}

The intended endpoint is a single cutoff removal for gauge, Higgs,
chiral fermion, ghost, and gravitational variables, with a positive
Euclidean physical state, the required gauge and BV identities, a
conserved limiting stress tensor, Osterwalder--Schrader reconstruction,
the correct spectral sectors, and a well-posed Einstein evolution
sourced by the limiting state.  Each item is a separate mathematical
gate.

Nineteen cycles are frozen as
\texttt{Renewal\_SM\_Einstein\_Continuum\_Research\_Notes\_V100\_f}\(k\)\texttt{.tex},
\(k=1,\dots,19\).  The nineteenth, frozen with SHA--256
\texttt{DA1C78C077C7712D1FB871A85CA09C40DD2C601DF73576C5E0CB9D6E0CFCAF11}
(854326 bytes, 20645 lines), is the immediate predecessor of this
file.  Its discipline is kept: one appended section per version, an
exact-rational or symbolic verifier per numeric claim with both
SHA--256 values quoted, companion audits every fifth version, sweeps
every tenth, and a firewall, decision, assessment, and record in each
section.  Verifiers of this cycle are named
\texttt{scripts/verify\_sm\_v20\_vNN\_*.py} and emit
\texttt{artifacts/sm\_v20\_vNN\_*.json}.

\section{Major results proved before this restart}
\label{sec:v20v1-summary}

\begin{ledger}[Cycles one to eighteen]
\label{led:v20v1-legacy}
The frozen record separates finite regulator identities from continuum
claims: exact Schur and determinant-line reductions, local anomaly
cancellations with global-anomaly firewalls, finite Korn and
coarse-floor estimates, exact Cartan and Einstein-divergence defect
identities, the Euclidean Einstein--Hilbert conformal instability,
conditional covariance decay, subsequential tightness, reflection
positivity closed under limits, controlled shell cumulant expansions,
the Higgs conformal response on trees, and an exact kinematic
projective link law shown at the close of the eighteenth cycle to miss
the smooth small-curvature chart.
\end{ledger}

\begin{ledger}[Nineteenth cycle, V1--V31: transfer and constraint gates]
\label{led:v20v1-early}
Exact two-dimensional heat-kernel refinement and a four-dimensional
shared-link intertwiner; positive Haar projectors and a finite
reflection-positive transfer kernel; the Fock-positive Higgs hinge;
reflected Grassmann Gram matrices, determinant-line descent, spectral-flow
parity, and multiplet cancellation of mod-two determinant holonomy; a
finite hybrid gauge--Higgs--fermion transfer square; quadratic-curvature
stabilization and its two-time reflection counterexample; a sharp
reflection-positive conformal stiffness threshold; the constraint
projection that removes a negative transfer mode; the nonzero-momentum
ADM reduction with positive TT transfer and the infrared prescriptions
for the homogeneous mode; quadratic ADM compatibility and the three
translation charges on the flat torus; and exact nonlinear and smooth
small-amplitude CMC branches through flat data.
\end{ledger}

\begin{ledger}[Nineteenth cycle, V32--V88: positivity transport and the slice constraint]
\label{led:v20v1-middle}
Reflection positivity is transported through the nonlinear constraint
map, with odd parity of \(K\) and a sharp-time no-go; link seeds on a
Euclidean time lattice; the Euclidean-signature slice constraint is
multivalued; stationary chains violate the Einstein identities;
toroidal rigidity forces \(\Lambda<0\) on time-symmetric data, with the
analytic branch \(\Lambda_2=-\langle|\nabla h|^2\rangle/8\),
\(w_2=\tfrac7{128}\cos2z\), \(\Lambda_4=-145/2048\); the corrected rate
\(\dot\tau=|K|^2-\Lambda\); the Ornstein--Uhlenbeck seed as an exact
oscillator; the operator-level averaged constraint
\(\tfrac23\hat\tau^2-\widehat E-2\Lambda\) with its induced two-branch
physical space, its structural singular locus, and its branch symmetry;
the fourth-order quartic forms of the plane class, resonant mixing,
self-energy growth, and the graviton mass counterterm; the plane-class
gauge completeness; the three-dimensional exact curvature engine, the
orthogonal pair coefficient \(-\tfrac5{64}(l^2+1)\), the York momentum
sector with the erratum that superseded the plane-class momentum
completion; and the conclusion that no sector of the averaged slice
constraint cancels the graviton mass term.
\end{ledger}

\begin{ledger}[Nineteenth cycle, V89--V100: the emergent cosmology]
\label{led:v20v1-late}
The exact second- and fourth-order plane-symmetric dynamics of the
balanced standing wave, with \(\Lambda_4\) recovered from the dynamics
and both constraints preserved; the Isaacson identification
\(8\pi\rho=-\Lambda\); the homogeneous two-scale reduction confirmed
transversally and refuted longitudinally by the inhomogeneous secular
modulation of a standing wave; the wave's growth \(\tfrac5{64}\) and
frequency shift \(\tfrac3{64}\); the focusing theorem: any
time-symmetric datum on a compact slice with \(\Lambda<0\) focuses
within \(\tfrac\pi2\sqrt{3/|\Lambda|}\); for finite-mode data,
additive null dust per mode, the Friedmann limit
\(\ddot a/a=\tfrac23\Lambda\), the focusing time
\(\pi\sqrt6/\sqrt{\langle|\nabla h|^2\rangle}\), the certified quartic
form of \(\Lambda_4\) with the universal pair value
\(-\tfrac1{16}(|k_a|^2+|k_b|^2)\) and the Sobolev bound
\(|\Lambda_4|\leq C\|h\|_{H^2}^4\); the \(\Lambda=0\) sectors as
radiation universes with \(H^2=e/6\) and the redshift law; the massless
scalar as an abelian graviton polarization; the massive scalar as null
dust plus dust in the ratio \(k^2:m^2\); and the Maxwell field with its
induced polarization \(\tfrac14\sin^2z\sin^2t\) and no second-order
gravitational radiation.
\end{ledger}

The five gaps recorded at the close of the nineteenth cycle are: no
gauge-invariant three-dimensional constrained dynamics beyond plane
symmetry; no number-changing terms beyond one block; no renormalized
mode-number limit beyond normal ordering; no derivation of the
classical dynamics from the constraint operator; and no non-abelian
or fermionic sector.  This file begins on the last.

\section{The plane-symmetric \texorpdfstring{$SU(2)$}{SU(2)} Einstein--Yang--Mills system}
\label{sec:v20v1-system}

\begin{proposition}[Consistent colour ansatz and stress]
\label{prop:v20v1-ansatz}
Let \(ds^2=-dt^2+e^{2a}(e^{2b}dx^2+e^{-2b}dy^2)+e^{2c}dz^2\) with
\(a,b,c\) functions of \((z,t)\), and let the \(SU(2)\) potential be
\(A^1=f(z,t)\,dx\), \(A^2=g(z,t)\,dy\), \(A^3=0\), with coupling \(e\)
and \(8\pi\) absorbed.  Then \(F^3_{xy}=efg\) is the only non-abelian
component, the stress tensor is traceless with
\(T_{xy}=T_{tx}=T_{ty}=T_{zx}=T_{zy}=0\), so the diagonal ansatz is
consistent, and with \(W_\pm=e^{-2a\pm2b}\),
\(\mu=\tfrac12e^2f^2g^2e^{-4a}\),

\[
\begin{aligned}
 T_{tt}&=\rho=\tfrac12W_-(f_t^2+e^{-2c}f_z^2)+\tfrac12W_+(g_t^2+e^{-2c}g_z^2)+\mu,\qquad
 e^{-2c}T_{zz}=\rho-2\mu,\\
 W_-T_{xx}&=\mu-\pi,\qquad W_+T_{yy}=\mu+\pi,\qquad
 \pi=\tfrac12W_-(f_t^2-e^{-2c}f_z^2)-\tfrac12W_+(g_t^2-e^{-2c}g_z^2),\\
 T_{tz}&=W_-f_tf_z+W_+g_tg_z.
\end{aligned}                                                          \tag{3.1}
\]

The colour-magnetic energy \(\mu\) has positive transverse and negative
longitudinal pressure.
\end{proposition}

\begin{proof}
Symbolic computation of \(F^a=dA^a+e\,\epsilon^{abc}A^b\wedge A^c\) and
of \(T=\sum_aF^a\cdot F^a-\tfrac14g|F|^2\); the vanishing of the five
off-diagonal components, the tracelessness, and the five displayed
forms are certified.
\end{proof}

\begin{theorem}[Field equations]
\label{thm:v20v1-equations}
The Einstein equations \(\operatorname{Ric}=\Lambda g+T\) are the
constraints and metric equations of the Maxwell case (V99 of the
nineteenth cycle) with \(\rho\), \(\pi\), \(T_{tz}\) of (3.1), namely

\[
\begin{aligned}
 &a_t^2+2a_tc_t-b_t^2-e^{-2c}\bigl(3a_z^2-2a_zc_z+2a_{zz}+b_z^2\bigr)-\rho=\Lambda,\\
 &a_{tz}+a_ta_z-a_zc_t+b_tb_z+\tfrac12T_{tz}=0,\\
 &b_{tt}+b_t(2a_t+c_t)-e^{-2c}\bigl(b_{zz}+b_z(2a_z-c_z)\bigr)+\pi=0,\\
 &a_{tt}+2a_t^2+a_tc_t-e^{-2c}\bigl(a_{zz}+2a_z^2-a_zc_z\bigr)-\mu=\Lambda,\\
 &c_{tt}+c_t^2+2a_tc_t-e^{-2c}\bigl(2a_z^2-2a_zc_z+2a_{zz}+2b_z^2\bigr)-\rho+2\mu=\Lambda,
\end{aligned}                                                          \tag{3.2}
\]

and the Yang--Mills equations \(D_\mu F^{a\mu\nu}=0\) reduce to the two
coupled wave equations

\[
\begin{aligned}
 &f_{tt}+f_t(c_t-2b_t)-e^{-2c}\bigl(f_{zz}-f_z(2b_z+c_z)\bigr)+e^2g^2f\,W_+=0,\\
 &g_{tt}+g_t(c_t+2b_t)-e^{-2c}\bigl(g_{zz}+g_z(2b_z-c_z)\bigr)+e^2f^2g\,W_-=0,
\end{aligned}                                                          \tag{3.3}
\]

all other colour and spacetime components vanishing identically.  At
\(e=0\) the system is two decoupled Maxwell fields, \(f\) in the \(x\)
and \(g\) in the \(y\) polarization.
\end{theorem}

\begin{proof}
Certified symbolically: the six forms against the components of
\(\operatorname{Ric}-\Lambda g-T\) and against
\(\nabla_\mu F^{a\mu\nu}+e\epsilon^{abc}A^b_\mu F^{c\mu\nu}\); the
remaining ten Yang--Mills components are certified to vanish.  The
non-abelian terms are the potential \(e^2g^2f\) of colour one in the
field of colour two and vice versa, together with the magnetic energy
\(\mu\) in the metric equations.
\end{proof}

\section{The first non-abelian effect}
\label{sec:v20v1-effect}

Let \(f=\varepsilon f_1+\varepsilon^3f_3+\dots\), \(g\) likewise, with
\(f_1=g_1=\cos z\cos t\), and let the metric be expanded as in the
nineteenth cycle.  The coupling enters the equations only through
\(e^2\), and at orders \(\varepsilon^3\) and \(\varepsilon^4\) it enters
linearly: the second-order metric is \(e\)-independent, the third-order
colour waves carry \(e^2\) only through \(-e^2g_1^2f_1\), and the
fourth-order metric carries \(e^2\) only through \(\mu\) and through the
abelian stress evaluated on \((f_1,f_3)\).  The \(e^2\)-linear parts
are therefore exactly computable from these terms alone.

\begin{theorem}[Effective mass and third harmonic]
\label{thm:v20v1-mass}
The \(e^2\)-part of the third-order colour wave is

\[
\begin{aligned}
 f_3^{\rm NA}&=-\frac3{128}\cos z\cos t-\frac9{32}\,t\cos z\sin t+\frac3{128}\cos z\cos3t\\
 &\quad-\frac3{128}\cos3z\cos t+\frac3{128}\cos3z\cos3t-\frac1{96}\,t\cos3z\sin3t,
\end{aligned}                                                          \tag{4.1}
\]

and the same for \(g_3^{\rm NA}\).  Its resonant part is the frequency
shift \(\delta=\tfrac9{32}e^2\) of the colour wave, that is, an
effective mass \(m_{\rm eff}^2=\tfrac9{16}e^2\) acquired from the other
colour's field, larger than the naive average \(e^2\langle g_1^2\rangle=\tfrac14e^2\)
by the self-phase factor \(\tfrac94\), with no amplitude growth; and a
resonant third harmonic \(\cos3z\sin3t\) growing linearly in time.
\end{theorem}

\begin{proof}
\(f_3^{\rm NA}\) solves \(f_{tt}-f_{zz}=-g_1^2f_1\) with zero data; the
source \(\cos^3z\cos^3t=\tfrac1{16}(3\cos z+\cos3z)(3\cos t+\cos3t)\)
is resonant in its \((1,1)\) and \((3,3)\) components.  The
coefficient of \(t\cos z\sin t\) is \(-\delta\); certified.
\end{proof}

\begin{theorem}[The geometry of the colour-magnetic flux]
\label{thm:v20v1-geometry}
With the time-symmetric conformal datum and zero fourth-order data for
the colour waves, the \(e^2\)-parts of the fourth-order quantities are

\[
 \Lambda_4^{(e^2)}=-\frac3{16},\qquad
 \langle a_4\rangle^{(e^2)}_{\rm sec}=\frac{51}{2048}-\frac{15}{256}t^2,\qquad
 \langle c_4\rangle^{(e^2)}_{\rm sec}=-\frac{51}{1024}-\frac9{128}t^2,\qquad
 b_4^{(e^2)}=0,                                                         \tag{4.2}
\]

with both constraints and the rate law
\(\dot\tau_4=-\Lambda_4+\rho_4\) holding identically in their
\(e^2\)-parts.  The secular accelerations decompose as

\[
 2\langle\ddot a_4\rangle^{(e^2)}=\Lambda_4^{(e^2)}+\langle\mu_4\rangle=-\frac3{16}+\frac9{128},\qquad
 2\langle\ddot c_4\rangle^{(e^2)}=\Lambda_4^{(e^2)}-\langle\mu_4\rangle+\langle\rho^{\rm ab}_4\rangle
 =-\frac3{16}-\frac9{128}+\frac{15}{128},                               \tag{4.3}
\]

where \(\langle\mu_4\rangle=\tfrac9{128}\) is the mean colour-magnetic
energy and \(\langle\rho^{\rm ab}_4\rangle=\tfrac{15}{128}\) the mean
abelian stress induced by the effective mass: the colour-magnetic flux
relieves the transverse contraction and deepens the longitudinal one,
and the net non-abelian effect contracts the longitudinal direction
more than the transverse, in the ratio \(6:5\).
\end{theorem}

\begin{proof}
The \(e^2\)-parts of the fourth-order sources are \(\Lambda_4+\mu_4\)
for \(a_4\), \(\Lambda_4+2a_{4,zz}+\rho_4-2\mu_4\) for \(c_4\), and
\(-\pi_4=0\) for \(b_4\) (the two colours contribute opposite
anisotropic stresses of equal size), with
\(\rho_4=2(f_{1,t}f_{3,t}+f_{1,z}f_{3,z})+\mu_4\) and
\(\mu_4=\tfrac12f_1^2g_1^2\); \(\Lambda_4^{(e^2)}=-\langle\rho_4\rangle(0)=-\tfrac12\langle\cos^4z\rangle\)
from the constraint at \(t=0\).  All certified, including the
identities (4.3) between the certified secular coefficients and the
certified means.
\end{proof}

\begin{proposition}[What the geometry reads off the gauge coupling]
\label{prop:v20v1-reading}
At the first non-abelian order the emergent geometry reads the
coupling \(e^2\) in three ways: through the effective mass of the
colour waves, which by the massive-scalar law of the nineteenth cycle
converts part of the null dust into dust; through the colour-magnetic
energy \(\mu\), which is the only source of the transverse metric
equation and the only negative longitudinal pressure in the system;
and through the resonant third harmonic, a genuine non-abelian
mode-generation absent for Maxwell and scalar fields.  The abelian
polarization symmetry of V97 of the nineteenth cycle is broken by
the coupling, since \(\pi\) and the Yang--Mills potentials
distinguish the colours' spatial directions.
\end{proposition}

\begin{proof}
Theorems \ref{thm:v20v1-mass} and \ref{thm:v20v1-geometry}; the
statement about \(\mu\) is read from (3.2), whose fourth equation has
no other source.
\end{proof}

\subsection{Executable certificate}

The verifier

\[
 \texttt{scripts/verify\_sm\_v20\_v1\_yangmills\_sector.py}           \tag{4.4}
\]

derives (3.1)--(3.3) symbolically, checks the consistency of the
ansatz and the abelian limit, integrates the \(e^2\)-linear parts of
orders three and four with the exact engine of the nineteenth cycle,
and checks (4.1)--(4.3), both constraints, and the rate law.  It emits

\[
 \texttt{artifacts/sm\_v20\_v1\_yangmills\_sector.json}.              \tag{4.5}
\]

The hashes are

\[
\begin{array}{c|l}
\hbox{object}&\hbox{SHA--256}\\ \hline
\hbox{verifier}&
\texttt{6907B7A9B5DF504B4909B8CB6853E9FD9EAC73CCF34D977EF22011E1E4377B0A}\\
\hbox{canonical payload}&
\texttt{B4414897D6014EF627045EBE983E74E10E81D9E23AAD720616C6FCCC0023259D}.
\end{array}                                                     \tag{4.6}
\]

\begin{firewall}[One colour ansatz, equal amplitudes, first order in the coupling]
\label{fire:v20v1}
The ansatz has one transverse component per colour and no colour
three; equal amplitudes and equal phases are chosen, which makes
\(b_4^{(e^2)}=0\); only the \(e^2\)-linear parts are computed, the
\(e^0\) fourth-order dynamics of two Maxwell fields being left aside;
nothing is said about the quantum Yang--Mills field, about ghosts, or
about the Standard Model's gauge group.  The reading of the
effective mass through the massive-scalar law is by analogy of
averaged stresses, not by a computation of the full fourth-order
metric.
\end{firewall}

\begin{decision}[V2 acquisition]
\label{dec:v20v1-next}
V2 will break the equal-amplitude degeneracy, \(f_1=A\cos z\cos t\),
\(g_1=B\cos z\cos t\) with \(A\neq B\), to obtain the non-abelian
polarization source \(\pi_4^{(e^2)}\) and the induced gravitational
polarization, and will add the colour-three component to test whether
the resonant third harmonic feeds back into the metric at the next
order.  V5 is the first companion audit of the cycle.
\end{decision}

\begin{assessment}[V1 acquisition]
\label{ass:v20v1}
The non-abelian sector is opened with an exact and consistent
plane-symmetric \(SU(2)\) Einstein--Yang--Mills system and its first
effect computed exactly: an effective mass \(\tfrac9{16}e^2\) of each
colour wave from the other, a resonant third harmonic, and the
colour-magnetic flux that shifts \(\Lambda_4\) by \(-\tfrac3{16}e^2\)
and contracts the longitudinal direction more than the transverse.
The emergent geometry reads the gauge coupling at the first order
where it appears.
\end{assessment}

\begin{record}[V1 record]
\label{rec:v20v1}
This file opens the twentieth cycle after the nineteenth was frozen
as \texttt{Renewal\_SM\_Einstein\_Continuum\_Research\_Notes\_V100\_f19.tex}
with SHA--256
\texttt{DA1C78C077C7712D1FB871A85CA09C40DD2C601DF73576C5E0CB9D6E0CFCAF11}.
It is the only active numeric SM note; frozen archives and unrelated
notes are unchanged.
\end{record}


\section{V2: the relative phase of the colour waves, and non-abelian generation of a gravitational wave}
\label{sec:v20v2}

V1 computed the first non-abelian effect for two colour waves in
phase.  This note varies the only remaining datum of the ansatz at
this order, the relative phase: \(f_1=\cos z\cos t\) and
\(g_1=\cos z\cos(t-\varphi)\) with \(\varphi=0\) (in phase) and
\(\varphi=\tfrac\pi2\) (quadrature, \(g_1=\cos z\sin t\)).  The
\(e^2\)-linear analysis of V1 applies to any first-order data, since
the coupling enters the third-order colour waves only through
\(-e^2g_1^2f_1\) and the fourth-order metric only through the
colour-magnetic energy and the abelian stress on \((f_1,f_3)\).  The
effective mass, the colour-magnetic energy, the fourth-order
cosmological constant, and the secular accelerations all depend on
the phase, and in quadrature the anisotropic non-abelian stress
acquires a resonant component: the colour pair generates a
gravitational wave at twice its wavenumber and frequency, growing
linearly in time, together with a growing homogeneous anisotropy.
Colour waves in phase generate no gravitational wave at this order;
abelian fields never do.

\subsection{Phase dependence of the scalar quantities}

\begin{theorem}[Effective mass and energies at the two phases]
\label{thm:v20v2-phase}
For the two phases the \(e^2\)-linear quantities of V1 are

\[
\begin{array}{l|c|c|c|c|c|c}
\varphi&\delta&m_{\rm eff}^2&\langle\mu_4\rangle&\langle\rho^{\rm ab}_4\rangle&\Lambda_4^{(e^2)}&(2\langle\ddot a_4\rangle,\,2\langle\ddot c_4\rangle)^{(e^2)}\\ \hline
0&9/32&9/16&9/128&15/128&-3/16&(-15/128,\,-9/64)\\
\pi/2&3/32&3/16&3/128&-3/128&0&(3/128,\,-3/64)
\end{array}                                                            \tag{5.1}
\]

with \(\delta=\delta_f=\delta_g\) in both cases and no amplitude growth
of the colour waves.  In quadrature the colour-magnetic flux is one
third of its in-phase value, the effective mass one third, the
fourth-order cosmological constant vanishes because the datum carries
no colour-magnetic field at \(t=0\), and the non-abelian effect
expands the transverse directions while contracting the longitudinal
one.
\end{theorem}

\begin{proof}
Same computation as Theorems \ref{thm:v20v1-mass} and
\ref{thm:v20v1-geometry} with \(g_1=\cos z\sin t\); the resonant part
of \(-g_1^2f_1=-\cos^3z\sin^2t\cos t\) is \(-\tfrac34\cdot\tfrac14\cos z\cos t\)
against \(-\tfrac34\cdot\tfrac34\) in phase, which is the ratio \(1:3\)
of the effective masses; the frequency shift of \(g\) is read from the
coefficient of \(t\cos z\cos t\) in \(g_3\).  The constraints and the
rate law hold identically in their \(e^2\)-parts at both phases; all
values are certified.
\end{proof}

\subsection{The induced gravitational wave}

\begin{theorem}[Non-abelian generation of a gravitational wave]
\label{thm:v20v2-gw}
In quadrature the anisotropic stress has the \(e^2\)-part

\[
 \pi_4^{(e^2)}=\frac3{32}\cos2z-\frac1{12}\cos2z\cos2t-\frac1{96}\cos2z\cos4t
 +\frac3{32}\,t\sin2t+\frac1{32}\,t\cos4z\sin2t,                        \tag{5.2}
\]

whose \(\cos2z\cos2t\) component is resonant for the polarization
equation, and the induced gravitational polarization is

\[
\begin{aligned}
 b_4^{(e^2)}&=\frac1{48}\,t\cos2z\sin2t+\frac3{128}\,t\sin2t-\frac1{384}\,t\cos4z\sin2t\\
 &\quad-\frac3{128}+\frac3{128}\cos2t-\frac3{128}\cos2z+\frac7{288}\cos2z\cos2t
 -\frac1{1152}\cos2z\cos4t+\frac1{1152}\cos4z\cos2t-\frac1{1152}\cos4z\cos4t,
\end{aligned}                                                          \tag{5.3}
\]

with \(b_4^{(e^2)}(0)=\dot b_4^{(e^2)}(0)=0\): a gravitational wave of
wavenumber and frequency two, growing at the rate \(\tfrac1{48}e^2\)
per unit time, a homogeneous anisotropy growing at \(\tfrac3{128}e^2\),
and a driven non-resonant \(\cos4z\) term with the same linear growth.
In phase \(\pi_4^{(e^2)}=0\) and \(b_4^{(e^2)}=0\); for the abelian
fields of the nineteenth cycle the induced polarization is bounded.
\end{theorem}

\begin{proof}
\(\pi_4^{(e^2)}=(f_{1,t}f_{3,t}-f_{1,z}f_{3,z})-(g_{1,t}g_{3,t}-g_{1,z}g_{3,z})\)
with \(f_3,g_3\) the \(e^2\)-parts of the third-order colour waves,
which have different time structure at the two phases; in phase the
two brackets are equal.  The polarization equation
\(b_{4,tt}-b_{4,zz}=-\pi_4\) with zero data is solved mode by mode: the
\((2,2)\) component of \(-\pi_4\), \(\tfrac1{12}\cos2z\cos2t\), gives
\(\tfrac1{48}t\cos2z\sin2t\), and the \(t\sin2t\) components of
\(\pi_4\), themselves inherited from the secular terms of \(f_3,g_3\),
give the remaining linear terms.  Certified coefficient by
coefficient, including the absence of any growth faster than linear.
\end{proof}

\begin{proposition}[Reading of the phase]
\label{prop:v20v2-reading}
The emergent geometry reads the relative phase of the colour waves at
the first non-abelian order through the sign of the transverse secular
acceleration, through the presence or absence of a fourth-order
cosmological constant, and through the generation of gravitational
radiation, which requires a nonzero colour-electric field
\(F^3_{xy}\) that oscillates out of phase with the colour waves'
own stresses.  The three signatures are exact at order \(e^2\) and
distinguish the quadrature configuration from both the in-phase one
and from any abelian configuration of equal energy.
\end{proposition}

\begin{proof}
Theorems \ref{thm:v20v2-phase} and \ref{thm:v20v2-gw}, and
Theorem 4.3 of the nineteenth cycle's V99 for the abelian bound.
\end{proof}

\subsection{Executable certificate}

The verifier

\[
 \texttt{scripts/verify\_sm\_v20\_v2\_colour\_phase.py}               \tag{5.4}
\]

runs the \(e^2\)-linear computation at both phases, checks the V1
values in phase, the quadrature values of (5.1), the resonant
coefficients of (5.3), the linear bound on the growth, and both
constraints and the rate law.  It emits

\[
 \texttt{artifacts/sm\_v20\_v2\_colour\_phase.json}.                  \tag{5.5}
\]

The hashes are

\[
\begin{array}{c|l}
\hbox{object}&\hbox{SHA--256}\\ \hline
\hbox{verifier}&
\texttt{07B7F4648251C615023F7B616069B8B3B38C2A0D80203FFBA870D3F599FF4F9C}\\
\hbox{canonical payload}&
\texttt{0365746C0665A5E50A7E9080F4B02873E051C96ADB9B3D2E7AEAF1A334ADBA6E}.
\end{array}                                                     \tag{5.6}
\]

\begin{firewall}[Two phases, equal amplitudes, first order in the coupling]
\label{fire:v20v2}
Only the phases \(0\) and \(\tfrac\pi2\) are computed, the general phase
being outside exact rational arithmetic; the amplitudes are equal;
the \(e^0\) fourth-order dynamics is not computed; the linear growth of
the induced wave is a secular term of the perturbative solution and
says nothing about late times.  The reading of Proposition
\ref{prop:v20v2-reading} refers to these two configurations only.
\end{firewall}

\begin{decision}[V3 acquisition]
\label{dec:v20v2-next}
V3 will compute the \(e^0\) fourth-order dynamics of the two Maxwell
fields underlying the ansatz, completing the fourth order of the
Einstein--Yang--Mills solution at both phases, so that the full
fourth-order secular structure, and not only its \(e^2\)-part, is
known; this is needed before the effective-mass reading of V1 can be
compared with the massive-scalar law as an identity rather than an
analogy.
\end{decision}

\begin{assessment}[V2 acquisition]
\label{ass:v20v2}
The relative phase of the colour waves is read by the geometry at the
first non-abelian order: the effective mass and the colour-magnetic
flux are three times smaller in quadrature, the fourth-order
cosmological constant vanishes there, and the anisotropic non-abelian
stress becomes resonant, generating a gravitational wave at twice the
wavenumber and frequency with growth rate \(\tfrac1{48}e^2\), a
non-abelian effect absent for abelian fields and for colour waves in
phase.
\end{assessment}

\begin{record}[V2 append record]
\label{rec:v20v2}
V2 is appended to the validated V1 whose installed SHA--256 was
\texttt{8F89CD95CDABEA544AF5720C78328077F0F7BD7FA1CC232FE163E846DA2D20EE}.
After installation only V2 is the active numeric SM note; frozen
archives and unrelated notes are unchanged.
\end{record}


\section{V3: the abelian fourth order, and the exactness of the Bianchi--I reduction for electromagnetic radiation}
\label{sec:v20v3}

V1 and V2 computed the \(e^2\)-linear parts of the fourth-order
Einstein--Yang--Mills solution.  This note computes the remaining
\(e^0\) part: the fourth-order plane-symmetric Einstein--Maxwell
dynamics of the two fields of the ansatz, \(f\) in the \(x\) and \(g\)
in the \(y\) polarization, at both relative phases.  The perturbative
sources are no longer expanded by hand: they are generated
symbolically from the exact system by substituting the
\(\varepsilon\)-expansions of all fields and extracting coefficients,
and each monomial is then evaluated with the exact \((z,t)\) engine.
Both constraints and the rate law are verified identically at orders
two, three, and four.  Two results stand out.  For electromagnetic
radiation the homogeneous Bianchi--I reduction of the nineteenth cycle
is exact through fourth order in both the transverse and the
longitudinal sector, at both phases, in contrast with the graviton,
whose longitudinal sector it fails.  And the gravitational
self-renormalization of the light waves' frequency vanishes in phase
and equals \(\tfrac18\) in quadrature, where the homogeneous breathing
anisotropy of the metric parametrically shifts both waves.

\subsection{Symbolic generation of the perturbative sources}

\begin{proposition}[Order-by-order system from the exact equations]
\label{prop:v20v3-generation}
Substituting \(a=\varepsilon^2a_2+\varepsilon^4a_4\),
\(b=\varepsilon^2b_2+\varepsilon^4b_4\), \(c=\varepsilon^2c_2+\varepsilon^4c_4\),
\(f=\varepsilon f_1+\varepsilon^3f_3\), \(g=\varepsilon g_1+\varepsilon^3g_3\),
\(\Lambda=\varepsilon^2\Lambda_2+\varepsilon^4\Lambda_4\) into the
\(e=0\) system (3.2)--(3.3) and expanding to order \(\varepsilon^4\)
gives, at each order, a linear equation for the new fields with a
polynomial source in the lower ones.  The equations are solved in the
order \(a_n,b_n,c_n\) (the \(c\)-equation contains \(a_{n,zz}\)),
with the Hamiltonian constraint at \(t=0\) fixing \(\Lambda_n\) and the
conformal datum \(a_n(0)=c_n(0)\), \(b_n(0)=0\), and zero data for the
third-order fields.
\end{proposition}

\begin{proof}
This is the construction of the verifier: the coefficients are
computed by symbolic series expansion, and their monomials in the
functions and their derivatives are evaluated by the engine.  Its
correctness is certified by the identical vanishing of the
Hamiltonian and momentum constraints and of the rate law
\(\dot\tau=2a_t^2+2b_t^2+c_t^2-\Lambda+\rho\) at every order on the
computed solution, which the construction does not impose.
\end{proof}

\subsection{The two phases}

\begin{theorem}[Abelian fourth order of the two-field datum]
\label{thm:v20v3-fourth}
With \(f_1=\cos z\cos t\), \(g_1=\cos z\cos(t-\varphi)\), and the
conformally flat time-symmetric metric datum:

\[
\begin{array}{l|c|c}
&\varphi=0&\varphi=\pi/2\\ \hline
\Lambda_2&-1/2&-1/2\\
w_2&-\tfrac1{32}\cos2z&0\\
a_2&-\tfrac{t^2}4-\tfrac1{16}\cos2z\cos2t&-\tfrac{t^2}4\\
b_2&0&\tfrac18-\tfrac18\cos2t\\
c_2&-\tfrac1{16}\cos2z&0\\
\Lambda_4&3/128&0\\
w_4&\tfrac1{128}\cos2z-\tfrac3{8192}\cos4z&0\\
\langle a_4\rangle_{\rm sec}&\tfrac5{2048}+\tfrac3{256}t^2-\tfrac{t^4}{24}&-\tfrac{t^4}{24}\\
\langle c_4\rangle_{\rm sec}&-\tfrac1{1024}+\tfrac{t^4}{48}&\tfrac1{512}-\tfrac{t^2}{64}+\tfrac{t^4}{48}\\
\langle b_4\rangle&0&-\tfrac1{16}+\tfrac3{32}t\sin2t-\tfrac1{16}t^2\cos2t+\dots\\
\delta_f=\delta_g&0&1/8
\end{array}                                                            \tag{6.1}
\]

where \(\delta\) is the frequency shift of the light waves and the dots
in \(\langle b_4\rangle\) are bounded oscillatory terms; no field has
amplitude growth at third order.  Both constraints and the rate law
vanish identically at orders two, three, and four.
\end{theorem}

\begin{proof}
Certified by the generated system.  In phase the \(x\) and \(y\)
stresses cancel in \(\pi\), so \(b\) is never sourced; in quadrature
the energy density at \(t=0\) is homogeneous, so the conformal factor
and \(\Lambda_4\) vanish, while the anisotropic stress
\(\pi_2=\tfrac14\cos2t\) drives the homogeneous breathing \(b_2\), which
in turn shifts both light waves through \(\mp2b_tf_t\) at the same
rate \(\tfrac18\) and drives the growing homogeneous anisotropy of
\(b_4\).
\end{proof}

\begin{theorem}[Exactness of the Bianchi--I reduction for electromagnetic radiation]
\label{thm:v20v3-bianchi}
The secular \(t^4\) coefficients of (6.1), \(-\tfrac1{24}\)
transverse and \(\tfrac1{48}\) longitudinal at both phases, are
exactly the \(T^4\) coefficients of the homogeneous Bianchi--I
reduction of the nineteenth cycle (equation (94.3) of its V91) with
\(\Lambda_2=-\tfrac12\), obtained from its \(\Lambda_2=-\tfrac18\)
series by \(T\to2T\): \(-\tfrac{16}{384}=-\tfrac1{24}\) and
\(\tfrac{16}{768}=\tfrac1{48}\).  For electromagnetic radiation the
reduction therefore holds through fourth order in both sectors,
whereas for the graviton it fails in the longitudinal sector.
\end{theorem}

\begin{proof}
The reduced system is homogeneous of degree two under
\((\Lambda,\rho,T^{-2})\to\lambda(\Lambda,\rho,T^{-2})\), which gives the
rescaling; the coefficients are read from (6.1) and from the
nineteenth cycle's (94.3).  The mechanism is that of its V99: the
Maxwell longitudinal Ricci source \(e^{2c}\rho\) has an unmodulated
time average, so \(c_2\) carries no secular \(\cos2z\) term at either
phase, and the inhomogeneous feedback that spoiled the graviton's
longitudinal sector is absent.
\end{proof}

\begin{corollary}[Complete fourth order of the Yang--Mills solution]
\label{cor:v20v3-complete}
Through order \(\varepsilon^4\), the plane-symmetric \(SU(2)\)
solution of V1--V2 is the sum of (6.1) and \(e^2\) times the
\(e^2\)-linear parts of (5.1)--(5.3): in particular the secular
accelerations are

\[
\begin{aligned}
 \varphi=0:&\quad 2\langle\ddot a_4\rangle=\frac3{64}-\frac{15}{128}e^2,\qquad
           2\langle\ddot c_4\rangle=-\frac9{64}e^2,\\
 \varphi=\tfrac\pi2:&\quad 2\langle\ddot a_4\rangle=\frac3{128}e^2,\qquad
           2\langle\ddot c_4\rangle=-\frac1{16}-\frac3{64}e^2,
\end{aligned}                                                          \tag{6.2}
\]

for the \(t^2\) parts, and \(\Lambda_4=\tfrac3{128}-\tfrac3{16}e^2\) in
phase, \(\Lambda_4=0\) in quadrature.
\end{corollary}

\begin{proof}
Linearity of the fourth-order equations in their sources and the
exact linearity in \(e^2\) of the sources at this order (V1).
\end{proof}

\subsection{Executable certificate}

The verifier

\[
 \texttt{scripts/verify\_sm\_v20\_v3\_abelian\_fourth\_order.py}       \tag{6.3}
\]

generates the order-by-order system symbolically, solves it at both
phases with the engine, and certifies the constraints, the rate law,
and (6.1).  It emits

\[
 \texttt{artifacts/sm\_v20\_v3\_abelian\_fourth\_order.json}.          \tag{6.4}
\]

The hashes are

\[
\begin{array}{c|l}
\hbox{object}&\hbox{SHA--256}\\ \hline
\hbox{verifier}&
\texttt{C71F65816A5AB15FF1BEEB75C44C5F2A4AF56C48F9C2A41B02AA7A926EC5FD7D}\\
\hbox{canonical payload}&
\texttt{E0112677A80CE52CC0E6C2031DEC204FFEE9BB1A9B67EDD0F2DFE0C58AD10367}.
\end{array}                                                     \tag{6.5}
\]

\begin{firewall}[Two phases, equal amplitudes, conformally flat datum]
\label{fire:v20v3}
The exactness statement of Theorem \ref{thm:v20v3-bianchi} is for the
secular \(t^4\) coefficients of the two-field datum at two phases and
is not a theorem about general electromagnetic data; the fourth-order
\(e^2\) and \(e^0\) parts are combined by linearity but the
\(O(e^4)\) terms of the same order in the amplitude are not computed;
the quantum field is absent.
\end{firewall}

\begin{decision}[V4 acquisition]
\label{dec:v20v3-next}
V4 will treat the general relative phase by rational points on the
circle, \((\cos\varphi,\sin\varphi)=(\tfrac35,\tfrac45),(\tfrac45,\tfrac35),(\tfrac5{13},\tfrac{12}{13})\),
to determine the exact trigonometric polynomials in \(\varphi\) of the
\(e^2\)-linear quantities, and in particular the non-abelian
exchange of energy between the colour waves, which vanishes at both
phases computed so far and is expected to be maximal at
\(\varphi=\tfrac\pi4\).  V5 is the first companion audit of the cycle.
\end{decision}

\begin{assessment}[V3 acquisition]
\label{ass:v20v3}
The abelian fourth order of the two-field datum is computed at both
phases with symbolically generated sources and fully verified
constraints, completing the fourth order of the Yang--Mills solution.
The Bianchi--I reduction is exact through fourth order for
electromagnetic radiation, and the light waves' gravitational
frequency shift is \(0\) in phase and \(\tfrac18\) in quadrature.
\end{assessment}

\begin{record}[V3 append record]
\label{rec:v20v3}
V3 is appended to the validated V2 whose installed SHA--256 was
\texttt{68F8E2BC4DBBC293465CB999F4B238B00C2A8089DA62E1059EE9C88743B76A8F}.
After installation only V3 is the active numeric SM note; frozen
archives and unrelated notes are unchanged.
\end{record}


\section{V4: the general relative phase: colour energy exchange and the exact phase laws}
\label{sec:v20v4}

V2 computed the first non-abelian effects at the phases \(0\) and
\(\tfrac\pi2\).  This note obtains them for every relative phase
\(\varphi\) of the colour waves, \(g_1=\cos z\cos(t-\varphi)\), by
evaluating the \(e^2\)-linear solution at rational points of the
circle and fitting each quantity exactly as a trigonometric polynomial
of degree two in \(\varphi\), the degree forced by the quartic
dependence on the fields; the fits are then verified at a sixth
rational point that was not used to construct them.  The result is a
set of exact phase laws.  The two colours exchange energy at the rate
\(\tfrac3{32}e^2\sin2\varphi\), antisymmetrically; their effective
masses are \(\tfrac3{16}+\tfrac3{32}\cos2\varphi\); the fourth-order
cosmological constant is \(-\tfrac3{16}e^2\cos^2\varphi\); and the
generated gravitational wave grows at the rate \(\tfrac1{48}e^2|\sin\varphi|\),
maximal in quadrature and absent in phase.

\subsection{The phase laws}

\begin{theorem}[Exact phase dependence of the \(e^2\)-linear quantities]
\label{thm:v20v4-laws}
For \(f_1=\cos z\cos t\), \(g_1=\cos z\cos(t-\varphi)\), the
\(e^2\)-linear parts of the third- and fourth-order solution satisfy,
with \(\alpha\) the coefficient of \(t\cos z\cos(t-\hbox{phase})\)
(amplitude growth) and \(\delta\) the frequency shift of each colour
wave,

\[
\begin{aligned}
 \alpha_f&=-\alpha_g=\frac3{32}\sin2\varphi,\qquad
 \delta_f=\delta_g=\frac3{16}+\frac3{32}\cos2\varphi,\qquad
 \Lambda_4^{(e^2)}=-\frac3{16}\cos^2\varphi,\\
 \langle\mu_4\rangle&=\frac3{128}(2+\cos2\varphi),\qquad
 \langle\rho^{\rm ab}_4\rangle=\frac3{64}+\frac9{128}\cos2\varphi,\\
 2\langle\ddot a_4\rangle^{(e^2)}&=-\frac3{64}-\frac9{128}\cos2\varphi,\qquad
 2\langle\ddot c_4\rangle^{(e^2)}=-\frac3{32}-\frac3{64}\cos2\varphi,
\end{aligned}                                                          \tag{7.1}
\]

and the linearly growing parts of the induced polarization are

\[
 b_4^{(e^2)}\ni\frac{t\cos2z}{96}\bigl[(1-\cos2\varphi)\sin2t+\sin2\varphi\cos2t\bigr]
 +\frac{3t}{256}\bigl[(1-\cos2\varphi)\sin2t+\sin2\varphi\cos2t\bigr],  \tag{7.2}
\]

so that the generated gravitational wave of wavenumber and frequency
two has growth rate \(\tfrac1{48}|\sin\varphi|\) and the homogeneous
anisotropy growth rate \(\tfrac3{128}|\sin\varphi|\), both with the
time phase \(2t+\chi\), \(\tan\chi=\sin2\varphi/(1-\cos2\varphi)=\cot\varphi\).
Both constraints and the rate law hold identically in their
\(e^2\)-parts at every phase computed.
\end{theorem}

\begin{proof}
Each quantity is a polynomial of degree four in the first-order
fields, hence a trigonometric polynomial of degree two in \(\varphi\),
\(c_0+c_1\cos2\varphi+s_1\sin2\varphi+c_2\cos4\varphi+s_2\sin4\varphi\);
its five coefficients are determined exactly from the five phases
\((\cos\varphi,\sin\varphi)=(1,0),(0,1),(\tfrac35,\tfrac45),(\tfrac45,\tfrac35),(\tfrac5{13},\tfrac{12}{13})\)
and the resulting form is certified to reproduce the independently
computed values at \((\tfrac{12}{13},\tfrac5{13})\).  The degree-four
coefficients \(c_2,s_2\) vanish for every quantity.  The exchange
antisymmetry \(\alpha_f+\alpha_g=0\) holds at all six points.  The
frequency shift and growth of \(g\) are read from the coefficient of
\(t\cos z\,e^{i(t-\varphi)}\) in \(g_3\) by the rotation by \(\varphi\).
\end{proof}

\begin{corollary}[Colour energy exchange]
\label{cor:v20v4-exchange}
At the first non-abelian order the colour waves exchange energy
without loss, colour one growing and colour two decaying when
\(\sin2\varphi>0\), at the rate \(\tfrac3{32}e^2\sin2\varphi\) per unit
time, maximal at \(\varphi=\tfrac\pi4\); the exchange vanishes for
colours in phase or in quadrature, where the coupling acts only as a
mass.  The mean colour-magnetic energy and the frequency shifts are
even in \(\varphi\) and take their extremal values at those two phases.
\end{corollary}

\begin{proof}
Read from (7.1): \(\alpha_f=-\alpha_g\) is the statement of exchange;
the resonant source \(-g_1^2f_1\) contains
\(\tfrac3{16}\sin2\varphi\cos z\sin t\), in quadrature with \(f_1\),
which drives the amplitude, and \(\tfrac3{16}(2+\cos2\varphi)\cos z\cos t\),
in phase, which drives the frequency.
\end{proof}

\begin{proposition}[Reading of the coupling at general phase]
\label{prop:v20v4-reading}
The emergent geometry reads \(e^2\) and \(\varphi\) at fourth order
through three independent even functions,
\(\Lambda_4^{(e^2)}\propto\cos^2\varphi\),
\(\langle\mu_4\rangle\propto2+\cos2\varphi\), and the secular
accelerations, and through the odd function \(\sin\varphi\) in the
growth rate of the generated gravitational wave; the energy exchange
itself, being a redistribution between colours of equal stress, is
invisible to the metric at this order except through the abelian
cross term \(\langle\rho^{\rm ab}_4\rangle\).
\end{proposition}

\begin{proof}
Theorem \ref{thm:v20v4-laws}.
\end{proof}

\subsection{Executable certificate}

The verifier

\[
 \texttt{scripts/verify\_sm\_v20\_v4\_colour\_exchange.py}            \tag{7.3}
\]

evaluates the \(e^2\)-linear solution at the six rational phases,
checks constraints and rate law at each, fits and verifies the phase
laws, and checks the V1--V2 values and the exchange antisymmetry.  It
emits

\[
 \texttt{artifacts/sm\_v20\_v4\_colour\_exchange.json}.               \tag{7.4}
\]

The hashes are

\[
\begin{array}{c|l}
\hbox{object}&\hbox{SHA--256}\\ \hline
\hbox{verifier}&
\texttt{BCF899CE71690E071BEE11DC42FC8DF5C1A362DC19D200EA2DAC93559B3F7E18}\\
\hbox{canonical payload}&
\texttt{2E758D05AC7229D8516A93B6B30A719518EB6604A90F090F10A8763C8F051A5F}.
\end{array}                                                     \tag{7.5}
\]

\begin{firewall}[Equal amplitudes, first order in the coupling, secular terms only]
\label{fire:v20v4}
The laws (7.1)--(7.2) are exact for the \(e^2\)-linear parts of the
perturbative solution with equal colour amplitudes; the exchange rate
and the growth rates are secular coefficients, valid for
\(e^2t\ll1\), and the late-time exchange dynamics, which classically
is the chaotic Yang--Mills mechanics, is not addressed.  The degree
bound of the fit is a theorem, but the six-point verification is the
certificate only for the listed quantities.
\end{firewall}

\begin{decision}[V5 acquisition]
\label{dec:v20v4-next}
V5 is the first compulsory companion audit of the cycle.  V6 will
lift the equal-amplitude restriction, \(f_1=A\cos z\cos t\),
\(g_1=B\cos z\cos(t-\varphi)\), whose \(e^2\)-quantities are
polynomials in \(A^2,B^2\) of total degree two, to obtain the complete
first-order non-abelian phase and amplitude laws, and will then ask
which combination is the adiabatic invariant of the colour exchange.
\end{decision}

\begin{assessment}[V4 acquisition]
\label{ass:v20v4}
The first non-abelian order of the plane-symmetric Einstein--Yang--Mills
solution is now known for every relative phase of the colour waves
as exact trigonometric polynomials, certified by an out-of-sample
rational point: energy exchange at \(\tfrac3{32}e^2\sin2\varphi\),
effective masses \(\tfrac3{16}+\tfrac3{32}\cos2\varphi\), the
cosmological shift \(-\tfrac3{16}e^2\cos^2\varphi\), and gravitational
radiation generated at \(\tfrac1{48}e^2|\sin\varphi|\).
\end{assessment}

\begin{record}[V4 append record]
\label{rec:v20v4}
V4 is appended to the validated V3 whose installed SHA--256 was
\texttt{95264CD1FD77DE114A18EA01622A1FCBAC90A5F5C5756197685A819A08EB66A8}.
After installation only V4 is the active numeric SM note; frozen
archives and unrelated notes are unchanged.
\end{record}


\section{V5: first compulsory companion audit of the twentieth cycle}
\label{sec:v20v5}

This is the required audit at the fifth version.  Every active
companion note present in the shared folder was inspected read-only
before the amplitude and averaged-dynamics work announced in V4
begins.

\begin{record}[V5 companion snapshot]
\label{rec:v20v5-snapshot}
The files present were

\[
\begin{array}{c|r|r}
\hbox{file and SHA--256}&\hbox{bytes}&\hbox{lines}\\ \hline
\texttt{Renewal\_NS\_Stability\_Research\_Notes\_V26.tex}
&278283&5319\\
\multicolumn{3}{l}{
\texttt{82C887F8C2C69E4F28FAD7C3DC8DE5B9099CB5A4BF431EBA1FFF3E91935978AD}}
\\
\texttt{Renewal\_GRH\_Cofinal\_Visibility\_Attack\_V61.tex}
&591674&17536\\
\multicolumn{3}{l}{
\texttt{F8A1EDDCD98C5184FB88D501D0B17ABE6A923AFC03A87797BC6D9C818E527D41}}
\\
\texttt{Renewal\_Yang\_Mills\_Mass\_Gap\_Research\_Notes\_V35.tex}
&520002&12860\\
\multicolumn{3}{l}{
\texttt{6C014195DC1E1A52F7943D3840A669EB813C202A77B90B80F31E251AB1C1BB9D}}
\\
\texttt{Renewal\_YM\_Parallel\_Research\_Notes\_V5.tex}
&54174&1351\\
\multicolumn{3}{l}{
\texttt{306F1BB84CFA8A8D47EA569EC17E9545F9867C8D32B548EC9D9C444B4F3C0512}}.
\end{array}                                                     \tag{8.1}
\]

Since the close of the nineteenth cycle the GRH note advanced from
V60 to V61 (the three-pair determinant for at least seven real
padding zeros, up to six residual inequalities); the NS, Yang--Mills
mass-gap, and YM parallel notes are unchanged.
\end{record}

\begin{decision}[V5 audit decision]
\label{dec:v20v5-audit}
No companion theorem, numerical constant, or executable artifact is
imported.  The GRH advance is a determinant inequality on a disk and
has no object in common with the plane-symmetric Einstein--Yang--Mills
system of this cycle.  The V4 direction stands.
\end{decision}

\begin{proof}
The only changed file is the GRH note; its V61 section concerns
sum-of-squares and determinant certificates for a polynomial
condition, none of whose hypotheses or conclusions involve a
constrained Lorentzian datum, a gauge field, or a mode sum.
\end{proof}

\begin{assessment}[V5 acquisition]
\label{ass:v20v5}
The audit is current.  In four versions the cycle has established
the consistent plane-symmetric \(SU(2)\) ansatz, its exact system,
the first non-abelian effects at every relative phase (effective
masses, colour energy exchange, cosmological shift, gravitational
radiation), the complete abelian fourth order with symbolically
generated sources, and the exactness of the Bianchi--I reduction for
electromagnetic radiation.  No Standard Model--Einstein continuum is
claimed.
\end{assessment}

\begin{record}[V5 append record]
\label{rec:v20v5}
V5 is appended to the validated V4 whose installed SHA--256 was
\texttt{B71FDABBEFA33D2272114D1D549E8F01DAEBE163AA507D51A9AAC33D69913CBF}.
After installation only V5 is the active numeric SM note; the
companions, frozen archives, and all unrelated notes are unchanged.
The next compulsory companion audit is V10, which is also the first
ten-version sweep of the cycle.
\end{record}


\section{V6: the averaged colour dynamics is Hamiltonian and integrable}
\label{sec:v20v6}

V4 obtained the first non-abelian effects as secular coefficients of
the perturbative solution.  This note replaces the secular expansion
by the averaged equations that generate it: the resonant projection of
the cubic Yang--Mills terms onto the two colour modes gives a flow on
the complex amplitudes \((u,v)\) in the slow time \(e^2\varepsilon^2t\),
which is Hamiltonian with the quartic Hamiltonian
\(H=\tfrac3{32}(2|u|^2|v|^2+\operatorname{Re}u^2\bar v^2)\), conserves
the total colour energy \(N=|u|^2+|v|^2\), and reduces to a
one-dimensional motion of the energy difference between two turning
points.  The colour exchange of V4 is thus periodic and lossless at
this order, and the whole set of phase and amplitude laws of V1--V4 is
the linearization of this flow at the initial datum, certified by
exact agreement with the engine at six rational points of amplitude
and phase.

\subsection{The averaged system}

\begin{theorem}[Averaged colour equations]
\label{thm:v20v6-averaged}
Write the colour waves as \(f=\varepsilon\operatorname{Re}(u\,e^{it})\cos z+\dots\),
\(g=\varepsilon\operatorname{Re}(v\,e^{it})\cos z+\dots\).  The resonant
projection of the non-abelian terms \(-e^2g^2f\), \(-e^2f^2g\) of (3.3)
onto the mode \(\cos z\,e^{it}\) gives, in the slow time
\(\tau=e^2\varepsilon^2t\),

\[
 \dot u=i\,\frac3{32}\bigl(2|v|^2u+v^2\bar u\bigr),\qquad
 \dot v=i\,\frac3{32}\bigl(2|u|^2v+u^2\bar v\bigr),                     \tag{9.1}
\]

that is, \(\dot u=i\,\partial H/\partial\bar u\), \(\dot v=i\,\partial H/\partial\bar v\)
with

\[
 H=\frac3{32}\Bigl(2|u|^2|v|^2+\operatorname{Re}\,u^2\bar v^2\Bigr),      \tag{9.2}
\]

and \(N=|u|^2+|v|^2\) and \(H\) are conserved.  The certified secular
coefficients of V1--V4 are the values of (9.1) at the datum:
\(\alpha_f+i\delta_f=\dot u\) at \(u=A\), \(v=Be^{-i\varphi}\), and
likewise for \(g\).
\end{theorem}

\begin{proof}
The coefficient of \(e^{it}\) in \(g_1^2f_1\) is
\(\tfrac18(2|v|^2u+v^2\bar u)\) and the \(\cos z\) content of
\(\cos^3z\) is \(\tfrac34\); the resonant response of \(f_{tt}-f_{zz}\)
to a source \(\operatorname{Re}(\sigma e^{it})\cos z\) is
\(\operatorname{Re}(-i\sigma t\,e^{it})\cos z\), which identifies
\(\dot u\).  The Hamiltonian form and the two invariants are
certified symbolically in real coordinates; the agreement with the
engine's third-order resonant coefficients is certified at
\((A,B,\cos\varphi,\sin\varphi)=(1,1,1,0)\), \((1,1,0,1)\),
\((2,1,\tfrac35,\tfrac45)\), \((1,3,\tfrac5{13},\tfrac{12}{13})\),
\((\tfrac32,\tfrac12,\tfrac45,-\tfrac35)\), \((1,2,-\tfrac5{13},\tfrac{12}{13})\).
\end{proof}

\subsection{Integrability and the exchange cycle}

\begin{theorem}[One-dimensional reduction]
\label{thm:v20v6-reduction}
Let \(P=|u|^2-|v|^2\) and \(W=\bar u^2v^2\).  Along (9.1),
\(\dot P=-\tfrac38\operatorname{Im}W\), \(|W|=|u|^2|v|^2=\tfrac14(N^2-P^2)\),
and \(\operatorname{Re}W=\tfrac{32}3H-\tfrac12(N^2-P^2)\), so that

\[
 \dot P^2=\frac9{64}\Bigl[\Bigl(\frac{N^2-P^2}4\Bigr)^2-\Bigl(\frac{32}3H-\frac{N^2-P^2}2\Bigr)^2\Bigr],
                                                                    \tag{9.3}
\]

a quartic potential in \(P\).  Hence the energy difference oscillates
periodically between the two roots of the bracket that enclose its
initial value, and the colour exchange is lossless and periodic.  For
the equal-amplitude datum at \(\varphi=\tfrac\pi4\), \(N=2\),
\(H=\tfrac3{16}\), the bracket is \(\tfrac1{16}(4-3P^2)(4+P^2)\), so the
exchange carries \(|u|^2\) between \(1\mp\tfrac1{\sqrt3}\), with
\(\dot P(0)=\tfrac3{16}\sin2\varphi\cdot2=\tfrac38\) in agreement with
the exchange rate of V4.
\end{theorem}

\begin{proof}
The three identities are certified symbolically; (9.3) follows by
eliminating \(\operatorname{Im}W\).  The bracket is nonnegative on
the interval between the turning points and the motion is that of a
particle in the potential \(-\tfrac9{64}[\dots]\) at zero energy, hence
periodic with period given by an elliptic integral.  The example is
arithmetic.
\end{proof}

\begin{corollary}[What the secular laws are]
\label{cor:v20v6-linearization}
The phase laws (7.1) are the first-order Taylor coefficients of the
flow (9.1): the exchange rate \(\tfrac3{32}\sin2\varphi\) is
\(-\tfrac3{16}\operatorname{Im}W(0)/(2A)\), the frequency shifts are
the real parts of \(\dot u/u\) and \(\dot v/v\), and the vanishing of
the exchange at \(\varphi=0,\tfrac\pi2\) is the vanishing of
\(\operatorname{Im}W\) there, which are the two families of fixed
points of the reduced motion, stable and unstable respectively.
\end{corollary}

\begin{proof}
Evaluate (9.1) at the datum; the fixed points of (9.3) are the roots
of \(\operatorname{Im}W=0\) with \(P\) at an extremum of the bracket.
\end{proof}

\subsection{Executable certificate}

The verifier

\[
 \texttt{scripts/verify\_sm\_v20\_v6\_averaged\_colour\_dynamics.py}  \tag{9.4}
\]

derives (9.1) symbolically by resonant projection, certifies the
Hamiltonian form, the invariants, and the identities behind (9.3),
and checks the flow against the engine at the six rational points.  It
emits

\[
 \texttt{artifacts/sm\_v20\_v6\_averaged\_colour\_dynamics.json}.     \tag{9.5}
\]

The hashes are

\[
\begin{array}{c|l}
\hbox{object}&\hbox{SHA--256}\\ \hline
\hbox{verifier}&
\texttt{D099F8DD29E9432FCFE60AD256A44C96127B0955F1D4BEEC0539B792DE9378F9}\\
\hbox{canonical payload}&
\texttt{C8919B09B8C8C3AB3F76BC04820C879225846FE6DF6B54B968F9D5B6AD99DF20}.
\end{array}                                                     \tag{9.6}
\]

\begin{firewall}[Averaging, two modes, no metric feedback]
\label{fire:v20v6}
The averaged system is the leading term of a two-scale expansion
whose validity for times of order \(1/(e^2\varepsilon^2)\) is the
standard averaging theorem and is not proved here; the third
harmonic and the non-resonant terms are dropped; the metric responds
to the colour waves but does not feed back into (9.1) at this order;
the classical Yang--Mills mechanics of more modes is not integrable.
\end{firewall}

\begin{decision}[V7 acquisition]
\label{dec:v20v6-next}
V7 will quantize the averaged colour Hamiltonian (9.2) in the number
basis of the two colour oscillators of the seed, where the
pair-exchange term \(u^2\bar v^2\) couples \(|n_u,n_v\rangle\) to
\(|n_u\pm2,n_v\mp2\rangle\) within each total number, compute the
exact non-abelian energy levels for small total number, and place
them in the operator-level constraint of the nineteenth cycle, so
that the branch labels \(\tau_e\) of the physical sectors read the
gauge coupling.
\end{decision}

\begin{assessment}[V6 acquisition]
\label{ass:v20v6}
The first non-abelian order is now an integrable Hamiltonian flow on
the two colour amplitudes, with the quartic Hamiltonian (9.2),
conserved colour energy, and a one-dimensional exchange cycle between
exact turning points; the secular laws of V1--V4 are its linearization
and agree with it exactly at six rational data.
\end{assessment}

\begin{record}[V6 append record]
\label{rec:v20v6}
V6 is appended to the validated V5 whose installed SHA--256 was
\texttt{C5D692446573656A76151DBC7408B75668F53D36CDCFBF4CEF086E17F04B8AA8}.
After installation only V6 is the active numeric SM note; frozen
archives and unrelated notes are unchanged.
\end{record}


\section{V7: the colour interaction in the number basis, and the branch labels of the constraint}
\label{sec:v20v7}

V6 reduced the first non-abelian order to the Hamiltonian
\(H=\tfrac3{32}(2|u|^2|v|^2+\operatorname{Re}u^2\bar v^2)\) on the two
colour amplitudes.  This note quantizes it in the number basis of the
two colour oscillators of the seed and computes its exact spectrum
for small total number.  The pair-exchange term couples
\(|n_u,n_v\rangle\) to \(|n_u\pm2,n_v\mp2\rangle\), so within each total
number \(n\) the Hamiltonian is a pair of tridiagonal blocks labelled
by the parity of \(n_u\), whose characteristic polynomials have
integer coefficients and are computed exactly for \(n\leq8\).  The
block containing the pure-colour states \(|n,0\rangle,|0,n\rangle\)
always has a negative level, an attractive channel whose ground
state at \(n=2\) is the antisymmetric pair \((|2,0\rangle-|0,2\rangle)/\sqrt2\)
with energy \(-\tfrac3{32}\kappa e^2\), while the blocks in which both
colours are oddly occupied have only positive levels for even \(n\).
Placed in the operator-level constraint of the nineteenth cycle, the
levels shift the sector energies \(e\) and hence the mean-curvature
labels \(\tau_e\): the emergent geometry reads the gauge coupling
through the spectrum of the colour interaction.

\subsection{The quantized averaged Hamiltonian}

\begin{definition}[Colour interaction operator]
\label{def:v20v7-operator}
Let \(\hat a,\hat b\) be the annihilation operators of the two colour
modes, with the classical amplitudes \(u=\lambda\hat a\), \(v=\lambda\hat b\)
in the seed's normalization and \(\kappa=\lambda^4\).  The normal-ordered
quantization of (9.2) is

\[
 \widehat H_{\rm NA}=\frac3{32}\kappa e^2\Bigl[2\hat n_u\hat n_v+\tfrac12\bigl(\hat a^2\hat b^{\dagger2}+\hat a^{\dagger2}\hat b^2\bigr)\Bigr],
                                                                    \tag{10.1}
\]

which commutes with the total number \(\hat n=\hat n_u+\hat n_v\) and
with the parity of \(\hat n_u\), and whose expectation in the
coherent state \(|\alpha,\beta\rangle\) is the classical
\(\kappa e^2H(\alpha,\beta)\).
\end{definition}

\begin{theorem}[Exact colour levels for \(n\leq8\)]
\label{thm:v20v7-levels}
In units of \(\tfrac3{32}\kappa e^2\), on the basis \(|k,n-k\rangle\),
\(\widehat H_{\rm NA}\) has diagonal \(2k(n-k)\) and couplings
\(\tfrac12\sqrt{(k+1)(k+2)(n-k)(n-k-1)}\) between \(k\) and \(k+2\).
Its blocks have the characteristic polynomials

\[
\begin{array}{c|l|l}
n&\hbox{even }k&\hbox{odd }k\\ \hline
2&(x-1)(x+1)&x-2\\
3&x^2-4x-3&x^2-4x-3\\
4&x(x^2-8x-12)&(x-3)(x-9)\\
5&x^3-20x^2+68x+80&x^3-20x^2+68x+80\\
6&(x^2-10x-15)(x^2-22x-15)&(x-10)(x^2-28x+120)\\
7&x^4-56x^3+882x^2-3384x-5103&\hbox{same}\\
8&(x^2-24x-28)(x^3-56x^2+560x+896)&(x^2-34x+217)(x^2-54x+497)
\end{array}                                                            \tag{10.2}
\]

For odd \(n\) the two blocks are isospectral (colour swap
\(k\leftrightarrow n-k\)).  The even-\(k\) block has a strictly negative
lowest level for every \(2\leq n\leq8\): \(-1\), \(2-\sqrt7\),
\(4-2\sqrt7\), \(\approx-0.918\), \(5-2\sqrt{10}\), \(\approx-1.143\),
\(\approx-1.399\); the odd-\(k\) blocks of even \(n\) have only positive
levels, with lowest \(2,3,14-2\sqrt{19},17-6\sqrt2\) for
\(n=2,4,6,8\).  At \(n=4\) the even block has an exact zero level.
\end{theorem}

\begin{proof}
The characteristic polynomial of a symmetric tridiagonal block
depends only on the diagonal and the squared couplings, which are
integers, and is computed by the three-term recurrence in exact
arithmetic; the factorizations over \(\mathbb Q\), the exact root
counts by Sturm sequences, and the closed forms of the quadratic
levels are certified.
\end{proof}

\begin{corollary}[Attractive and repulsive colour channels]
\label{cor:v20v7-channels}
The pure-colour pair states \(|2,0\rangle\), \(|0,2\rangle\) are split
by the exchange into the antisymmetric ground state with energy
\(-\tfrac3{32}\kappa e^2\) and the symmetric state with
\(+\tfrac3{32}\kappa e^2\), while the mixed state \(|1,1\rangle\) has
\(+\tfrac3{16}\kappa e^2\).  The attraction persists in the even block
at every computed \(n\), with the lowest level of order one in these
units, whereas the mean interaction of the block grows like \(n^2\):
the attractive channel is a boundary effect of the pair exchange, not
a bulk one.
\end{corollary}

\begin{proof}
Read from (10.2); the trace of the even block is
\(\sum_{k\ \rm even}2k(n-k)\sim n^3/6\), so the mean level grows like
\(n^2\) while the lowest level stays bounded over the computed range.
\end{proof}

\subsection{The constraint reads the coupling}

\begin{proposition}[Non-abelian shift of the branch labels]
\label{prop:v20v7-labels}
In the operator-level averaged constraint of the nineteenth cycle,
\(\tfrac23\hat\tau^2-\widehat E-2\Lambda\), let \(\widehat E\) be
augmented by \(\widehat H_{\rm NA}\) on the two colour modes.  Since
\(\widehat H_{\rm NA}\) commutes with the free energy of the colour
modes, the sectors of fixed total colour number \(n\) split into the
eigenstates of (10.1), with labels

\[
 \tfrac23\tau^2=e_0+\tfrac3{32}\kappa e^2\,x_j(n)+2\Lambda,              \tag{10.3}
\]

\(x_j(n)\) the roots of (10.2).  For \(\Lambda=0\) the Hubble rate of
the sector, \(H^2=\tau^2/9\), is lowered by the attractive channel and
raised by the repulsive ones; the pure-colour pair at \(n=2\) has two
sectors whose squared Hubble rates differ by \(\tfrac1{16}\kappa e^2\).
\end{proposition}

\begin{proof}
The induced inner product of the nineteenth cycle's V52 and V62 is
defined on each spectral subspace of \(\widehat E\); commuting
operators refine the spectral decomposition without changing the
trichotomy of V62.  The numbers are read from (10.2).
\end{proof}

\subsection{Executable certificate}

The verifier

\[
 \texttt{scripts/verify\_sm\_v20\_v7\_colour\_levels.py}              \tag{10.4}
\]

builds the blocks, computes the characteristic polynomials by the
exact recurrence, factors them, counts the roots exactly, extracts
the closed forms of the lowest levels, and checks the coherent-state
expectation.  It emits

\[
 \texttt{artifacts/sm\_v20\_v7\_colour\_levels.json}.                 \tag{10.5}
\]

The hashes are

\[
\begin{array}{c|l}
\hbox{object}&\hbox{SHA--256}\\ \hline
\hbox{verifier}&
\texttt{730CACB345F29FA9ED3AAA3BDB6AC3F24C8BC456FFD09438188E99CDBF584C71}\\
\hbox{canonical payload}&
\texttt{70B92591FDAE99B423863E5C9D3DA70E19783B648790FEBCA25A9CA2947F96DA}.
\end{array}                                                     \tag{10.6}
\]

\begin{firewall}[Normalization, two modes, averaged Hamiltonian]
\label{fire:v20v7}
The conversion constant \(\kappa\) between the classical amplitudes
and the seed's oscillators is left abstract, so all levels are stated
in its units; the quantization is that of the averaged Hamiltonian,
not of the Yang--Mills field; only two colour modes and no ghosts,
gauge fixing, or Gauss-law constraint of the quantum theory are
included; the levels beyond \(n=8\) and the large-\(n\) behaviour of
the attractive channel are not established.
\end{firewall}

\begin{decision}[V8 acquisition]
\label{dec:v20v7-next}
V8 will obtain the gravitational counterpart of (9.2): the averaged
Hamiltonian of the two light polarizations mediated by the
second-order metric, whose resonant coefficients at order \(e^0\) were
found in V3 to vanish in phase and to give the shift \(\tfrac18\) in
quadrature, by fitting the engine's third-order resonant coefficients
at rational amplitudes and phases to a \(U(1)\)-invariant quartic
form and certifying its Hamiltonian structure; V9 will combine it
with (9.2) into the complete first-order slow dynamics.
\end{decision}

\begin{assessment}[V7 acquisition]
\label{ass:v20v7}
The first non-abelian interaction is quantized in the seed's number
basis with exact spectra for total colour number up to eight: pair
exchange within fixed total number, isospectral parity blocks for odd
number, an attractive channel with ground energy
\(-\tfrac3{32}\kappa e^2\) at \(n=2\), and repulsive mixed channels.
The operator-level constraint reads the coupling through these
levels in its branch labels.
\end{assessment}

\begin{record}[V7 append record]
\label{rec:v20v7}
V7 is appended to the validated V6 whose installed SHA--256 was
\texttt{24D069F2789835C8CB0BD851D9894B31C12C08D926BC88BAC7E69136FD2B3FBC}.
After installation only V7 is the active numeric SM note; frozen
archives and unrelated notes are unchanged.
\end{record}


\section{V8: the gravitational self-interaction of light is a rotation of its polarization by its own spin}
\label{sec:v20v8}

V6 gave the averaged non-abelian dynamics of the two colour waves.
This note obtains the gravitational counterpart at order \(e^0\): the
averaged interaction of the two light polarizations \(f\) (along \(x\))
and \(g\) (along \(y\)) mediated by the second-order metric.  Two
things had to be settled first.  The resonant coefficients of the
third-order fields are covariant under time translation only if the
free data of the second-order metric are chosen covariantly; the
datum of V3, which set the anisotropy \(b_2\) and its velocity to
zero at \(t=0\), is not, and with it the coefficients do not fit any
\(U(1)\)-covariant form.  With the covariant prescription, the
constraints fixing \(a_2\) and the pure particular solutions for
\(b_2,c_2\), the coefficients fit exactly and are Hamiltonian, and the
Hamiltonian is remarkably simple:
\(H_{\rm grav}=\tfrac18(\operatorname{Im}\bar uv)^2\), a function of the
circular-polarization density \(S=\operatorname{Im}\bar uv\) alone.  The
flow is \(\dot u=\tfrac S8v\), \(\dot v=-\tfrac S8u\): light rotates its
own polarization plane at the rate \(\tfrac18S\) in the slow time
\(\varepsilon^2t\), conserving its energy and its spin; linearly
polarized light has no gravitational self-shift at this order, and
circularly polarized light of unit amplitude has the frequency shift
\(\tfrac18\) found in V3.

\subsection{Covariant second-order data}

\begin{proposition}[Covariance of the resonant coefficients]
\label{prop:v20v8-covariance}
For first-order data \(f_1=A\cos z\cos t\), \(g_1=B\cos z\cos(t-\varphi)\)
the second-order constraints determine \(a_2\) completely (its
\(k\neq0\) part at all times by the Hamiltonian constraint, its
initial velocity by the momentum constraint, which requires
\(\partial_ta_2(0)=X(z)\) with \(X'=-\tfrac12T_{tz}(0)\neq0\) for
\(\sin2\varphi\neq0\)), while \(b_2\) and \(c_2\) carry free
homogeneous parts: free waves \(\cos kt,\sin kt\) for \(k\neq0\), and
constant and linear terms.  The prescription that \(b_2\) and \(c_2\)
are the pure particular solutions of their equations is invariant
under time translation of the data; the prescription
\(b_2(0)=\partial_tb_2(0)=0\) of V3 is not, since for
\(\sin2\varphi\neq0\) it adds the free linearly growing anisotropy
\(-\tfrac18B^2\sin2\varphi\,t\) to \(b_2\), which feeds the third-order
resonant coefficients through \(-2\partial_tb\,\partial_tf\).
\end{proposition}

\begin{proof}
The particular solution of a mode equation with source
\(\sum c_{mn}t^ne^{imt}\) is a sum of the same exponentials with
polynomial coefficients and contains no homogeneous term, and
translation in \(t\) maps such sums to such sums; the constraint data
are functions of the matter at \(t=0\) and transform with it.  The free
linear term is the value of \(\partial_tb_2(0)\) of the particular
solution, \(\tfrac18\sin2\varphi\), for \(A=B=1\).  The non-covariant
datum was certified in the construction of this note to give
resonant coefficients that fail every \(U(1)\)-covariant cubic fit
at \((A,B,\varphi)=(2,1,\arctan\tfrac43)\), while the covariant one
fits, as stated next.
\end{proof}

\subsection{The gravitational Hamiltonian}

\begin{theorem}[Averaged gravitational interaction of two polarizations]
\label{thm:v20v8-hamiltonian}
With the covariant prescription, the \(e^0\) third-order resonant
coefficients define \(\dot u=i\Sigma_u\), \(\dot v=i\Sigma_v\) in the
slow time \(\varepsilon^2t\), and

\[
 \Sigma_u=\frac1{16}\bigl(|v|^2u-v^2\bar u\bigr),\qquad
 \Sigma_v=\frac1{16}\bigl(|u|^2v-u^2\bar v\bigr),                        \tag{11.1}
\]

exactly, that is, the Hamiltonian flow \(\dot u=i\partial H/\partial\bar u\)
of

\[
 H_{\rm grav}=\frac1{16}\bigl(|u|^2|v|^2-\operatorname{Re}\bar u^2v^2\bigr)
 =\frac18\bigl(\operatorname{Im}\bar uv\bigr)^2=\frac18S^2.              \tag{11.2}
\]

Equivalently \(\dot u=\tfrac S8v\), \(\dot v=-\tfrac S8u\): a rotation of
the polarization pair at angular rate \(\tfrac18S\), with \(N=|u|^2+|v|^2\)
and \(S\) conserved.  In particular a single linear polarization has
no gravitational self-shift (\(p=0\)), two linear polarizations in
phase do not interact, and the circular state \(u=1\), \(v=-i\)
(\(S=-1\)) has \(\dot u=\tfrac i8u\), the frequency shift \(\tfrac18\) of
V3.
\end{theorem}

\begin{proof}
The \(U(1)\)-covariant cubic form \(\Sigma_u=p|u|^2u+q|v|^2u+rv^2\bar u\)
is fitted exactly from the engine's coefficients at
\((A,B,\varphi)=(1,0,0),(1,1,0),(1,1,\tfrac\pi2)\), giving
\(p=0\), \(q=\tfrac1{16}\), \(r=-\tfrac1{16}\), and verified at four
further rational points of amplitude and phase; \(\Sigma_v\) is fitted
and verified likewise with \(p'=p\), \(q'=q\), \(r'=\bar r\), which is the
Hamiltonian condition.  The identity
\(|u|^2|v|^2-\operatorname{Re}\bar u^2v^2=2(\operatorname{Im}\bar uv)^2\)
gives (11.2), and the rotation form is certified directly at all seven
points.  Constraints and rate law hold identically at orders two and
three in every run.
\end{proof}

\begin{corollary}[Gravitational versus non-abelian interaction]
\label{cor:v20v8-comparison}
In the variables \(N\), \(P=|u|^2-|v|^2\), \(S=\operatorname{Im}\bar uv\),
\(C=\operatorname{Re}\bar uv\), with \(|u|^2|v|^2=S^2+C^2=\tfrac14(N^2-P^2)\),

\[
 H_{\rm grav}=\frac18S^2,\qquad
 H_{\rm NA}=\frac3{32}\bigl(3|u|^2|v|^2-2S^2\bigr)=\frac9{128}(N^2-P^2)-\frac3{16}S^2.
                                                                    \tag{11.3}
\]

The gravitational term depends on the spin only and generates no
exchange (\(\dot P=0\)); the non-abelian term couples spin and exchange
with opposite sign in \(S^2\).  For circularly polarized light the two
interactions compete: at \(|u|=|v|=1\) in quadrature, \(H_{\rm grav}=\tfrac18\)
and \(H_{\rm NA}=\tfrac9{32}-\tfrac3{16}=\tfrac3{32}\), the values of V3
and V2.
\end{corollary}

\begin{proof}
\(\operatorname{Re}\bar u^2v^2=C^2-S^2\) and \(|u|^2|v|^2=C^2+S^2\);
insert into (9.2) and (11.2).
\end{proof}

\subsection{Executable certificate}

The verifier

\[
 \texttt{scripts/verify\_sm\_v20\_v8\_gravitational\_hamiltonian.py}  \tag{11.4}
\]

builds the second and third orders from the symbolically generated
system of V3 with the covariant prescription and the
momentum-constraint velocity, checks all residuals, fits and
verifies (11.1), checks the Hamiltonian condition, the V3 values, and
the rotation form at seven points.  It emits

\[
 \texttt{artifacts/sm\_v20\_v8\_gravitational\_hamiltonian.json}.     \tag{11.5}
\]

The hashes are

\[
\begin{array}{c|l}
\hbox{object}&\hbox{SHA--256}\\ \hline
\hbox{verifier}&
\texttt{218F142E40395579190B9A8D2FA40C9C5713D026F1CC12D62A1821BE722D6E91}\\
\hbox{canonical payload}&
\texttt{E85CC72D332A4C51FA94AA31FFDF0EBA617917E99849C316F7FA96314773D0E2}.
\end{array}                                                     \tag{11.6}
\]

\begin{firewall}[Averaging, prescription dependence, no background feedback]
\label{fire:v20v8}
The result is the leading averaged dynamics for two standing
polarizations of equal wavenumber on the torus; its validity time is
that of the averaging theorem.  The free second-order metric data are
fixed by a prescription, physically the absence of free gravitons
and of free homogeneous anisotropy at second order, and other
prescriptions change the coefficients as shown; the background
expansion (the secular \(t^2\) terms of \(a_2,c_2\)) is carried by the
particular solutions and enters the resonant coefficients only through
non-secular pieces at this order.  The result is classical.
\end{firewall}

\begin{decision}[V9 acquisition]
\label{dec:v20v8-next}
V9 will combine (9.2) and (11.2) into the complete first-order slow
dynamics \(H=\varepsilon^2H_{\rm grav}+e^2\kappa'H_{\rm NA}\) of the two
colour waves, which is a \(U(1)\)-invariant quartic on two modes and
hence integrable, determine its fixed points and their stability as a
function of the ratio \(e^2/\varepsilon^2\) of the couplings, and
quantize the combined Hamiltonian in the number basis as in V7.  V10
is the audit and the first ten-version sweep of the cycle.
\end{decision}

\begin{assessment}[V8 acquisition]
\label{ass:v20v8}
The gravitational self-interaction of light in the plane-symmetric
emergent cosmology is an exact Hamiltonian rotation of the
polarization pair by its own circular-polarization density,
\(H_{\rm grav}=\tfrac18(\operatorname{Im}\bar uv)^2\), with no
self-shift of linear polarizations and the shift \(\tfrac18\) of
circular ones; the covariance requirement on the second-order data
that makes the coefficients well defined is identified and the
non-covariant datum of V3 is shown to spoil it.
\end{assessment}

\begin{record}[V8 append record]
\label{rec:v20v8}
V8 is appended to the validated V7 whose installed SHA--256 was
\texttt{3D58963CEA65EF2B220CAF3CE3C2D5694E24CCB575C0F2E70C3A15C66E553DF5}.
After installation only V8 is the active numeric SM note; frozen
archives and unrelated notes are unchanged.
\end{record}


\section{V9: the polarization sphere: gravity and colour compete at the critical coupling \texorpdfstring{$e^2=2/3$}{e2 = 2/3}}
\label{sec:v20v9}

V6 and V8 gave the two averaged interactions of the colour waves,
the non-abelian one and the gravitational one, as quartic
Hamiltonians on the amplitudes \((u,v)\).  This note combines them.
Both are invariant under the common phase, so the slow dynamics lives
on the Poincar\'e sphere of the Stokes vector
\(s=(|u|^2-|v|^2,\,2\operatorname{Re}\bar uv,\,2\operatorname{Im}\bar uv)\),
\(|s|=N\), where the combined Hamiltonian is the quadratic form
\(H=\beta_2s_2^2+\beta_3s_3^2\) with \(\beta_2=\tfrac9{128}e^2\) from
colour alone and \(\beta_3=\tfrac1{32}+\tfrac3{128}e^2\) from gravity
and colour together, and the flow is the rigid-body-like rotation
\(\dot s=2\,s\times\nabla H\).  The six axis points, the linear
polarizations along the colour axes, the linear polarizations at
forty-five degrees, and the circular polarizations, are the fixed
points; their stability is decided by the ordering of \(0,\beta_2,\beta_3\),
which changes at \(e^2=\tfrac23\), where the two coefficients coincide,
the Hamiltonian becomes axially symmetric, and the colour exchange
freezes.  Below the critical coupling circular polarization is
stable and forty-five-degree linear polarization is the saddle; above
it the roles are exchanged.  The quantized combined Hamiltonian shows
the same critical coupling as a level crossing of the two-quantum
pair states.

\subsection{The Stokes form}

\begin{theorem}[Combined slow dynamics on the Poincar\'e sphere]
\label{thm:v20v9-stokes}
With \(H=H_{\rm grav}+e^2H_{\rm NA}\) of (11.2) and (9.2), in the slow
time \(\varepsilon^2t\),

\[
 H=\beta_2s_2^2+\beta_3s_3^2,\qquad \beta_2=\frac9{128}e^2,\qquad
 \beta_3=\frac1{32}+\frac3{128}e^2,                                     \tag{12.1}
\]

and the flow \(\dot u=i\partial H/\partial\bar u\), \(\dot v=i\partial H/\partial\bar v\)
is

\[
 \dot s=2\,s\times\nabla_sH=\bigl(4(\beta_3-\beta_2)s_2s_3,\ -4\beta_3s_1s_3,\ 4\beta_2s_1s_2\bigr),
                                                                    \tag{12.2}
\]

which conserves \(|s|=N\) and \(H\).  The six axis points \(\pm N\hat s_i\)
are fixed, and the linearization in the tangent plane at each has
zero trace and determinant

\[
 \det{}_1=\frac9{1024}e^2N^2(3e^2+4),\qquad
 \det{}_2=\frac9{512}e^2N^2(3e^2-2),\qquad
 \det{}_3=-\frac{N^2}{512}(3e^2-2)(3e^2+4),                             \tag{12.3}
\]

so that the colour-axis linear polarizations \(\pm N\hat s_1\) are
elliptic for every \(e\neq0\), the forty-five-degree linear
polarizations \(\pm N\hat s_2\) are elliptic for \(e^2>\tfrac23\) and
hyperbolic for \(e^2<\tfrac23\), and the circular polarizations
\(\pm N\hat s_3\) are elliptic for \(e^2<\tfrac23\) and hyperbolic for
\(e^2>\tfrac23\).  At \(e^2=\tfrac23\), \(\beta_2=\beta_3\), \(H\) is
symmetric about \(\hat s_1\), and \(s_1=P\) is conserved: the colour
exchange is frozen.
\end{theorem}

\begin{proof}
The identities \(|u|^2|v|^2=\tfrac14(s_2^2+s_3^2)\),
\(\operatorname{Re}\bar u^2v^2=\tfrac14(s_2^2-s_3^2)\) turn (9.2) and
(11.2) into (12.1); the flow of the Stokes components under the
complex equations is certified to equal (12.2) with the bracket
constant \(2\); the determinants, the vanishing traces, and the
critical value are certified symbolically, and \(\dot s_1=0\) at
\(e^2=\tfrac23\).  Zero trace makes the tangent eigenvalues
\(\pm\sqrt{-\det}\), imaginary for positive determinant.
\end{proof}

\begin{corollary}[The competition]
\label{cor:v20v9-competition}
Gravity alone (\(e=0\)) rotates the polarization plane about the
circular axis and has no exchange; colour alone (\(e\to\infty\)
relative to gravity) has the forty-five-degree states as maxima, the
circular states as saddles, and the exchange cycle of V6.  The
gravitational spin term raises the circular states until, at
\(e^2=\tfrac23\), they cross the forty-five-degree states; the exchange
cycle of V6 exists only for \(e^2>\tfrac23\) in the presence of gravity,
and for \(e^2<\tfrac23\) the trajectories encircle the circular states
instead, with a bounded and reversed exchange.
\end{corollary}

\begin{proof}
Read off the ordering of \(0,\beta_2,\beta_3\) in (12.1): the axis
with the intermediate coefficient is the saddle of a quadratic form on
the sphere.
\end{proof}

\subsection{The quantized combined Hamiltonian}

\begin{theorem}[Levels with both interactions]
\label{thm:v20v9-levels}
On the total-number-\(n\) subspace, with
\(\widehat H=\kappa[\tfrac18\widehat S^2+e^2\tfrac3{32}(2\hat n_u\hat n_v+\tfrac12(\hat a^2\hat b^{\dagger2}+\hat a^{\dagger2}\hat b^2))]\),
\(\widehat S=(\hat a^\dagger\hat b-\hat b^\dagger\hat a)/2i\), the
characteristic polynomials in units of \(\kappa\) are, with \(E=e^2\),

\[
\begin{aligned}
 n=2:&\quad (32x-3E)\bigl(16x-2-3E\bigr)\bigl(32x-4+3E\bigr),\\
 n=3:&\quad \bigl(1024x^2-(384E+320)x-27E^2+72E+9\bigr)^2,\\
 n=4:&\quad (8x-1)(32x-4-27E)(32x-16-9E)\bigl(256x^2-(192E+128)x-27E^2+60E\bigr),
\end{aligned}                                                          \tag{12.4}
\]

up to constant factors.  At \(n=2\) the levels are \(\tfrac3{32}E\),
\(\tfrac18+\tfrac3{16}E\), and \(\tfrac18-\tfrac3{32}E\): the symmetric
pair, the mixed state, and the antisymmetric pair, the last two
lifted by the gravitational spin energy \(\tfrac18\kappa\); the
antisymmetric and symmetric pair levels cross at \(E=\tfrac23\), the
classical critical coupling.  For \(n=3\) the parity blocks remain
isospectral, and for \(n=4\) the exact level \(\tfrac18\) and the zero
level of V7 (now the root of the last factor at \(E=0\)) persist.
\end{theorem}

\begin{proof}
The operators \(\hat n_u\hat n_v\), \(\hat a^2\hat b^{\dagger2}\), and
\(\hat a^\dagger\hat b\) preserve \(n\) and are written on the basis
\(|k,n-k\rangle\); \(\widehat H\) is certified Hermitian and its
characteristic polynomials are computed and factored exactly with
\(E\) symbolic.
\end{proof}

\subsection{Executable certificate}

The verifier

\[
 \texttt{scripts/verify\_sm\_v20\_v9\_polarization\_sphere.py}        \tag{12.5}
\]

certifies (12.1)--(12.3) symbolically, the invariants, the frozen
exchange at criticality, and (12.4).  It emits

\[
 \texttt{artifacts/sm\_v20\_v9\_polarization\_sphere.json}.           \tag{12.6}
\]

The hashes are

\[
\begin{array}{c|l}
\hbox{object}&\hbox{SHA--256}\\ \hline
\hbox{verifier}&
\texttt{D49593EF004DF9AB2A06EB1B1A0B3DB3517347F7980E4BEDD31411E9B0CCBA12}\\
\hbox{canonical payload}&
\texttt{91970B9D2AA09F212C18680DE2D4D4BC3B96E6FB524E6ABB26F77566D816F240}.
\end{array}                                                     \tag{12.7}
\]

\begin{firewall}[Two modes, leading averaging, normalization of the coupling]
\label{fire:v20v9}
The critical value \(e^2=\tfrac23\) refers to the coupling in the
normalization of V1, with \(8\pi\) absorbed into the fields, and to the
equal-wavenumber two-mode ansatz; it is a statement about the
leading averaged dynamics and its fixed points, not about the full
Yang--Mills--Einstein evolution.  The quantum levels carry the
abstract constant \(\kappa\) and the classical value \(\tfrac18\) of the
spin energy is used without an independent quantum derivation.
\end{firewall}

\begin{decision}[V10 acquisition]
\label{dec:v20v9-next}
V10 is the compulsory companion audit and the first ten-version
sweep of the cycle, which will fix the direction of V11--V19: whether
to extend the exact framework to unequal wavenumbers and three
colours, where the averaged dynamics ceases to be integrable, or to
turn to the fermionic sector.
\end{decision}

\begin{assessment}[V9 acquisition]
\label{ass:v20v9}
The complete first-order slow dynamics of the two colour waves is a
rotation on the Poincar\'e sphere generated by
\(\beta_2s_2^2+\beta_3s_3^2\); gravity and colour compete through the
spin term, the circular and forty-five-degree polarizations exchange
stability at \(e^2=\tfrac23\), where the colour exchange freezes, and
the quantized Hamiltonian shows the same critical coupling as a level
crossing of the pair states.
\end{assessment}

\begin{record}[V9 append record]
\label{rec:v20v9}
V9 is appended to the validated V8 whose installed SHA--256 was
\texttt{D7913BC8632486CC257A30E9841D5DD6B0C0E954CF8475DA42EDED70EB9CB33B}.
After installation only V9 is the active numeric SM note; frozen
archives and unrelated notes are unchanged.
\end{record}


\section{V10: companion audit and the first ten-version sweep of the twentieth cycle}
\label{sec:v20v10}

This is the required audit at the tenth version and the sweep of
V1--V9.  Every active companion note present in the shared folder was
inspected read-only.

\subsection{Companion snapshot}

\begin{record}[V10 companion snapshot]
\label{rec:v20v10-snapshot}
The files present were

\[
\begin{array}{c|r|r}
\hbox{file and SHA--256}&\hbox{bytes}&\hbox{lines}\\ \hline
\texttt{Renewal\_NS\_Stability\_Research\_Notes\_V26.tex}
&278283&5319\\
\multicolumn{3}{l}{
\texttt{82C887F8C2C69E4F28FAD7C3DC8DE5B9099CB5A4BF431EBA1FFF3E91935978AD}}
\\
\texttt{Renewal\_GRH\_Cofinal\_Visibility\_Attack\_V61.tex}
&591674&17536\\
\multicolumn{3}{l}{
\texttt{F8A1EDDCD98C5184FB88D501D0B17ABE6A923AFC03A87797BC6D9C818E527D41}}
\\
\texttt{Renewal\_Yang\_Mills\_Mass\_Gap\_Research\_Notes\_V35.tex}
&520002&12860\\
\multicolumn{3}{l}{
\texttt{6C014195DC1E1A52F7943D3840A669EB813C202A77B90B80F31E251AB1C1BB9D}}
\\
\texttt{Renewal\_YM\_Parallel\_Research\_Notes\_V5.tex}
&54174&1351\\
\multicolumn{3}{l}{
\texttt{306F1BB84CFA8A8D47EA569EC17E9545F9867C8D32B548EC9D9C444B4F3C0512}}.
\end{array}                                                     \tag{13.1}
\]

All four files are unchanged since the V5 audit.
\end{record}

\begin{decision}[V10 audit decision]
\label{dec:v20v10-audit}
No companion theorem, numerical constant, or executable artifact is
imported; nothing has changed in the companions since V5.
\end{decision}

\begin{proof}
Identity of the four hashes with those of (8.1).
\end{proof}

\subsection{Sweep of V1--V9}

\begin{proposition}[What V1--V9 established]
\label{prop:v20v10-sweep}
The first nine versions of the cycle opened the non-abelian sector of
the plane-symmetric emergent cosmology and closed its first order
completely:

\begin{enumerate}
\item the consistent \(SU(2)\) ansatz \(A^1=f\,dx\), \(A^2=g\,dy\) with
its exact Einstein--Yang--Mills system, the colour-magnetic energy
with negative longitudinal pressure, and the first non-abelian effects
in phase: effective mass \(\tfrac9{16}e^2\), a resonant third
harmonic, \(\Lambda_4\) shifted by \(-\tfrac3{16}e^2\) (V1);
\item the phase dependence, with non-abelian generation of a
gravitational wave at rate \(\tfrac1{48}e^2|\sin\varphi|\) (V2, V4);
\item the abelian fourth order with symbolically generated sources,
the exactness of the Bianchi--I reduction for electromagnetic
radiation, and the gravitational frequency shift \(\tfrac18\) of
circularly polarized light (V3);
\item the exact phase laws: colour energy exchange
\(\tfrac3{32}e^2\sin2\varphi\), effective masses
\(\tfrac3{16}+\tfrac3{32}\cos2\varphi\), \(\Lambda_4=-\tfrac3{16}e^2\cos^2\varphi\)
(V4);
\item the averaged colour dynamics as the integrable Hamiltonian flow
of \(\tfrac3{32}(2|u|^2|v|^2+\operatorname{Re}u^2\bar v^2)\) with its
exchange cycle (V6), and its quantization with exact colour levels for
total number up to eight, the attractive antisymmetric pair, and the
non-abelian shift of the constraint's branch labels (V7);
\item the gravitational self-interaction of light as the rotation of
its polarization by its own spin, \(H_{\rm grav}=\tfrac18(\operatorname{Im}\bar uv)^2\),
with the covariance requirement on second-order data (V8);
\item the combined dynamics on the Poincar\'e sphere and the critical
coupling \(e^2=\tfrac23\) at which gravity and colour exchange the
stability of circular and forty-five-degree polarizations, classically
and as a level crossing (V9).
\end{enumerate}
\end{proposition}

\begin{proof}
Items are the certified theorems of the cited sections.
\end{proof}

\begin{proposition}[Direction for V11--V19]
\label{prop:v20v10-direction}
The first order of the two-colour, equal-wavenumber sector is
complete and integrable; extending it to unequal wavenumbers or three
colours leaves integrability and enters the chaotic Yang--Mills
mechanics, where the exact tools yield only finite-order secular
coefficients.  The programme therefore turns to the remaining
Standard Model structures that plane symmetry can carry exactly.  V11
opens the Higgs sector: an \(SU(2)\) doublet in the vacuum
configuration coupled to the colour ansatz of V1, whose covariant
derivative gives the gauge waves a mass, so that the massive-scalar
law of the nineteenth cycle (null dust plus dust in the ratio
\(k^2:m^2\)) can be tested on massive vector waves and the emergent
geometry can read the Higgs vacuum through the gauge-boson mass.
V12--V14 treat the Higgs boson itself and its coupling to the colour
exchange; V15 is the audit; V16--V19 attempt the fermionic sector in
plane symmetry, where the diagonal-metric compatibility of a Dirac
field must first be settled.
\end{proposition}

\begin{proof}
This is a decision.  The non-integrability claim for three colours is
the classical result on Yang--Mills mechanics and is not used as a
theorem here.
\end{proof}

\begin{firewall}[Audit and sweep only]
\label{fire:v20v10}
No computation is added in V10.  The cycle's results are exact
statements about the first order of averaged dynamics and about
finite-order secular coefficients in plane symmetry; no continuum
limit, no quantum Yang--Mills field, and no Standard Model--Einstein
continuum are claimed.
\end{firewall}

\begin{decision}[V11 acquisition]
\label{dec:v20v10-next}
V11 derives the plane-symmetric Einstein--Yang--Mills--Higgs system
for the doublet \(\Phi=(0,v+h(z,t))/\sqrt2\) with the colour ansatz of
V1, certifies the gauge-boson mass \(m_A^2=e^2v^2/4\) in the wave
equations of \(f\) and \(g\), and computes the second-order dynamics of
a massive gauge standing wave, comparing its secular accelerations with
the massive-scalar law.
\end{decision}

\begin{assessment}[V10 acquisition]
\label{ass:v20v10}
The audit is current and the sweep records a closed first chapter:
the non-abelian sector of the plane-symmetric emergent cosmology at
first order, from the exact system to the integrable averaged
dynamics, its quantization, and its competition with gravity.  The
direction is the Higgs sector, then the fermions.
\end{assessment}

\begin{record}[V10 append record]
\label{rec:v20v10}
V10 is appended to the validated V9 whose installed SHA--256 was
\texttt{6BB908246C035306F46DCA18861DB73A4D98749B00CDE7A6D71C50836533C915}.
After installation only V10 is the active numeric SM note; the
companions, frozen archives, and all unrelated notes are unchanged.
The next compulsory companion audit is V15 and the next sweep V20.
\end{record}


\section{V11: the Higgs sector: the geometry reads the vacuum expectation value through the gauge-boson mass}
\label{sec:v20v11}

This note opens the Higgs sector in the plane-symmetric emergent
cosmology.  An \(SU(2)\) doublet in the unitary configuration
\(\Phi=(0,v+h(z,t))/\sqrt2\) is coupled to the colour ansatz of V1,
here with a single colour wave \(A^1=f\,dx\).  The exact
Einstein--Yang--Mills--Higgs system is derived symbolically from the
action, the diagonal ansatz is shown to remain consistent, the gauge
wave acquires the mass \(m_A^2=e^2v^2/4\) and the Higgs boson the
mass \(m_h^2=2\lambda v^2\), and the fourth-order dynamics of a
massive gauge standing wave is computed with symbolically generated
sources and all constraints verified.  The second-order secular
accelerations are exactly those of the massive scalar of the
nineteenth cycle with the same dispersion relation: the emergent
geometry reads the Higgs vacuum expectation value through the
gauge-boson mass, by the same law by which it read the mass of a
scalar.  The gauge wave also displaces the Higgs field, lowering the
local vacuum expectation value in proportion to its energy.

\subsection{The Einstein--Yang--Mills--Higgs system}

\begin{proposition}[Doublet stress and field equations]
\label{prop:v20v11-system}
With \(D_\mu\Phi=\partial_\mu\Phi-ieA^a_\mu\tfrac{\tau^a}2\Phi\),
\(V=\lambda(\Phi^\dagger\Phi-\tfrac12v^2)^2\), and \(A^1=f\,dx\),
\(A^2=A^3=0\), the doublet has
\(|D\Phi|^2=-\tfrac12h_t^2+\tfrac12e^{-2c}h_z^2+\tfrac18e^2(v+h)^2f^2e^{-2a-2b}\),
its stress \(T^\Phi=2\operatorname{Re}(D\Phi)^\dagger D\Phi-g(|D\Phi|^2+V)\)
has vanishing \(tx,ty,xy,zx,zy\) components, so the diagonal ansatz
is consistent, and

\[
 \rho_\Phi=T^\Phi_{tt}=\tfrac12h_t^2+\tfrac12e^{-2c}h_z^2+V+\tfrac18e^2(v+h)^2f^2e^{-2a-2b}.
                                                                    \tag{14.1}
\]

The Euler--Lagrange equations of
\(\sqrt{-g}\,(-\tfrac14F^2-|D\Phi|^2-V)\) are

\[
\begin{aligned}
 &f_{tt}+f_t(c_t-2b_t)-e^{-2c}\bigl(f_{zz}-f_z(2b_z+c_z)\bigr)+\tfrac14e^2(v+h)^2f=0,\\
 &h_{tt}+h_t(2a_t+c_t)-e^{-2c}\bigl(h_{zz}+h_z(2a_z-c_z)\bigr)
  +\lambda(v+h)\bigl((v+h)^2-v^2\bigr)+\tfrac14e^2(v+h)f^2e^{-2a-2b}=0,
\end{aligned}                                                          \tag{14.2}
\]

so that \(m_A^2=\tfrac14e^2v^2\) and \(m_h^2=2\lambda v^2\), and the
Einstein equations are \(\operatorname{Ric}=\Lambda g+T-\tfrac12(\operatorname{tr}T)g\)
with \(T=T^{\rm YM}+T^\Phi\).
\end{proposition}

\begin{proof}
Symbolic computation; the forms (14.2) are certified against the
Euler--Lagrange expressions, the off-diagonal stress components are
certified to vanish, and (14.1) is certified.
\end{proof}

\subsection{The massive gauge wave}

\begin{theorem}[Fourth-order dynamics of a massive gauge standing wave]
\label{thm:v20v11-wave}
Take \(e^2=3\), \(v=2\), \(\lambda=\tfrac58\), so \(m_A^2=3\),
\(m_h^2=5\), and the wave \(f_1=\cos z\cos2t\) (\(\omega^2=1+m_A^2\))
with no first-order Higgs wave, the time-symmetric conformal datum
for the metric, and the Higgs and third-order gauge fields taken as
the pure particular solutions of their equations.  Then, with the
perturbative system generated symbolically from (14.2) and the
Einstein equations, the Hamiltonian and momentum constraints and
the rate law hold identically at orders two, three, and four, and

\[
\begin{aligned}
 \Lambda_2&=-1,\qquad w_2=\tfrac1{32}\cos2z,\qquad
 \langle a_2\rangle_{\rm sec}=\tfrac3{128}-\tfrac5{16}t^2,\qquad
 \langle c_2\rangle_{\rm sec}=-\tfrac3{128}-\tfrac3{16}t^2,\\
 h_2&=-\frac3{40}-\frac1{24}\cos2z+\frac3{88}\cos4t+\frac3{56}\cos2z\cos4t,\\
 b_2&=\frac1{16}-\frac1{16}\cos4t-\frac1{16}\cos2z+\frac1{8}\cos2z\cos2t-\frac1{16}\cos2z\cos4t,
\end{aligned}                                                          \tag{14.3}
\]

with \(\Lambda_4=-\tfrac{4334587}{60712960}\), fourth-order secular
coefficients \(-\tfrac{373}{6144}t^4\) transverse and
\(-\tfrac{73}{6144}t^4\) longitudinal, and the frequency shift
\(\delta=\tfrac{37687}{197120}\) of the gauge wave from its
gravitational and Higgs-mediated self-interaction.
\end{theorem}

\begin{proof}
Certified by the generated system, with the massive mode solver
handling the resonant \((k,m)=(1,2)\) components by the polynomial
recursion.  The bounded fields \(b_2\), \(h_2\) are read from the
certified coefficients.
\end{proof}

\begin{corollary}[The massive-scalar law holds for massive vectors]
\label{cor:v20v11-law}
The second-order secular coefficients of (14.3), \(-\tfrac5{16}\) and
\(-\tfrac3{16}\), are the values \(-(\omega^2+1)/16\) and
\(-(\omega^2-1)/16\) of the nineteenth cycle's massive scalar at
\(\omega=2\), and \(\Lambda_2=-\omega^2/4\) likewise: the averaged
stress of the massive gauge wave is \(\rho=\omega^2/4\), \(p_z=\tfrac14\),
\(p_x=p_y=0\), null dust plus dust in the ratio \(k^2:m_A^2\).  The
emergent geometry therefore reads the Higgs vacuum expectation value
\(v\) through \(m_A^2=e^2v^2/4\) exactly as it reads a scalar mass,
and cannot distinguish, at this order and in the averaged stress, a
massive vector from a massive scalar of equal mass and energy.
\end{corollary}

\begin{proof}
The averages of the stress components of Proposition
\ref{prop:v20v11-system} on \(f_1\): the Yang--Mills part gives
\(\tfrac12(f_t^2+f_z^2)\) in \(\rho\) and \(T_{zz}\) and
\(\mp\tfrac12(f_t^2-f_z^2)\) transversally, the mass part
\(\tfrac18e^2v^2f^2=\tfrac12m_A^2f^2\) adds to \(\rho\), subtracts from
\(T_{zz}\), and contributes \(+\tfrac12m_A^2f^2\) to \(T_{xx}\) only;
the time averages with \(\omega^2=1+m_A^2\) give the stated values,
and the secular coefficients are certified.
\end{proof}

\begin{proposition}[Displacement of the vacuum]
\label{prop:v20v11-vev}
The gauge wave sources the Higgs field through
\(\tfrac14e^2(v+h)f^2\): at second order the mean displacement is
\(\langle h_2\rangle=-e^2v\langle f_1^2\rangle/(4m_h^2)=-\tfrac3{40}\)
with \(\langle f_1^2\rangle=\tfrac14\), so the local vacuum expectation
value is lowered by the gauge energy, with a static modulation
\(-\tfrac1{24}\cos2z\) and a forced oscillation at twice the wave
frequency.
\end{proposition}

\begin{proof}
The mean of the second-order Higgs equation
\(h_{2,tt}+m_h^2h_2=-\tfrac14e^2vf_1^2\) gives
\(\langle h_2\rangle=-e^2v\langle f_1^2\rangle/(4m_h^2)=-3\cdot2\cdot\tfrac14/(4\cdot5)=-\tfrac3{40}\);
certified.
\end{proof}

\subsection{Executable certificate}

The verifier

\[
 \texttt{scripts/verify\_sm\_v20\_v11\_higgs\_sector.py}              \tag{14.4}
\]

derives the system, certifies (14.1)--(14.2) and the ansatz
consistency, generates the perturbative sources, integrates orders one
to four, and checks constraints, rate law, (14.3), and the
massive-scalar law.  It emits

\[
 \texttt{artifacts/sm\_v20\_v11\_higgs\_sector.json}.                 \tag{14.5}
\]

The hashes are

\[
\begin{array}{c|l}
\hbox{object}&\hbox{SHA--256}\\ \hline
\hbox{verifier}&
\texttt{6B2A6A61082279F1549D07F1315C3AF92C3BE00692D43D51375AB62120F10421}\\
\hbox{canonical payload}&
\texttt{2BB2AA5446DA02916443436482967CB474074ABAE59690F8B838FB2F07407366}.
\end{array}                                                     \tag{14.6}
\]

\begin{firewall}[Unitary configuration, one colour, one set of couplings]
\label{fire:v20v11}
The doublet is taken in the unitary configuration with real \(h\), so
the Goldstone directions and the \(A^2,A^3\) couplings are absent;
one colour wave, one set of couplings, and the classical fields only;
the fourth-order numbers depend on \(m_h\) through the particular
solutions and are not claimed to have a closed form in the couplings.
\end{firewall}

\begin{decision}[V12 acquisition]
\label{dec:v20v11-next}
V12 will add the Higgs wave itself, \(h_1=\cos z\cos(\sqrt{1+m_h^2}\,t)\)
with \(m_h^2=3\) so that its frequency is integral, and the second
colour wave, to obtain the averaged interactions among the Higgs
boson and the two massive gauge bosons at first order, in the manner
of V6 and V8, and to ask whether the colour exchange survives the
gauge-boson mass.
\end{decision}

\begin{assessment}[V11 acquisition]
\label{ass:v20v11}
The Higgs sector is opened with the exact plane-symmetric
Einstein--Yang--Mills--Higgs system, the gauge and Higgs masses in
their field equations, and a fully verified fourth-order solution for
a massive gauge wave.  The massive-scalar law of the nineteenth cycle
holds for massive vectors: the geometry reads the Higgs vacuum
expectation value through the gauge-boson mass, and the gauge wave
lowers the local vacuum expectation value by
\(e^2v\langle f^2\rangle/(4m_h^2)\).
\end{assessment}

\begin{record}[V11 append record]
\label{rec:v20v11}
V11 is appended to the validated V10 whose installed SHA--256 was
\texttt{87AC5549BA81F391C3B3D1F8FE6376723B641F377FE9B8EDAACCB36D9A96D21F}.
After installation only V11 is the active numeric SM note; frozen
archives and unrelated notes are unchanged.
\end{record}


\section{V12: the Higgs boson and the massive gauge boson: an averaged Hamiltonian in flat space that gravity makes non-Hamiltonian}
\label{sec:v20v12}

V11 placed a single massive gauge wave in the Higgs vacuum.  This
note adds the Higgs wave and computes the averaged first-order
interaction of the two, at equal masses \(m_A^2=m_h^2=3\) so that
both oscillate at frequency two and every cubic resonance is present.
The third-order resonant coefficients are fitted exactly to the six
charge-one cubic monomials in the amplitudes \((u,w)\) of the gauge
and Higgs waves and verified at three further rational points, both
for the matter system on the fixed flat metric and for the full
Einstein--Yang--Mills--Higgs system with covariant second-order data.
In flat space the averaged flow is Hamiltonian, with a real quartic
Hamiltonian whose four coefficients are all negative: the Higgs and
gauge waves lower each other's frequencies and exchange energy
through the pair term \(\bar u^2w^2\).  With gravity every structural
condition of a Hamiltonian flow still holds except one: the cross
coefficients of \(|w|^2u\) in \(\dot u\) and of \(|u|^2w\) in \(\dot w\)
differ by exactly \(-\tfrac3{32}=-e^2/32\).  The expanding background
acts differently on the conformally coupled gauge wave and on the
non-conformal Higgs wave, and this action is not a function of the
two amplitudes alone: the background is a third slow variable, and
the averaged two-mode flow with gravity is not closed as a
Hamiltonian system.

\subsection{Setting and prescription}

\begin{proposition}[Covariant second-order data with matter sources]
\label{prop:v20v12-prescription}
For \(f_1=\operatorname{Re}(ue^{2it})\cos z\), \(h_1=\operatorname{Re}(we^{2it})\cos z\)
with \(e^2=3\), \(v=2\), \(\lambda=\tfrac38\), the second- and third-order
metric is fixed covariantly by: the \(k\neq0\) part of \(a\) from the
Hamiltonian constraint at \(t=0\) and its velocity from the momentum
constraint, \(\partial_ta^{(n)}(0)=X\), \(X'=-(\hbox{momentum source})/(\hbox{coefficient of }a_{tz})\);
the \(k=0\) part of \(a\) as the pure particular solution, without free
constant or linear term, since with a matter source in the transverse
equation the datum of vanishing mean Hubble rate is not
time-translation covariant; and \(b\), \(c\) as pure particular
solutions.  The second- and third-order matter fields are the pure
particular solutions of the massive mode equations, with the resonant
components handled by the polynomial recursion.  With this
prescription the constraints and the rate law hold identically at
orders two and three for every datum used.
\end{proposition}

\begin{proof}
Certified in every run.  The two non-covariant choices that were
excluded, the zero-velocity anisotropy of V3 and the zero mean Hubble
rate, were each found in the construction of V8 and of this note to
break the covariance of the resonant coefficients at data with
\(\sin2\varphi\neq0\).
\end{proof}

\subsection{The flat-space averaged Hamiltonian}

\begin{theorem}[Averaged Higgs--gauge interaction in flat space]
\label{thm:v20v12-flat}
On the fixed flat metric the resonant coefficients define
\(\dot u=i\Sigma_u\), \(\dot w=i\Sigma_w\) in the slow time
\(\varepsilon^2t\) with

\[
 \Sigma_u=\alpha_0|u|^2u+\gamma_0|w|^2u+\rho_0\bar uw^2,\qquad
 \Sigma_w=\beta_0|w|^2w+\gamma_0|u|^2w+\rho_0\bar wu^2,               \tag{15.1}
\]

\[
 \alpha_0=-\frac{1109}{11648},\quad\beta_0=-\frac{1881}{11648},\quad
 \gamma_0=-\frac{1031}{5824},\quad\rho_0=-\frac{591}{11648},          \tag{15.2}
\]

exactly, with no \(|u|^2\bar uw\) or \(|w|^2\bar uw\) terms; that is,
the Hamiltonian flow of

\[
 H_0=\frac{\alpha_0}2|u|^4+\frac{\beta_0}2|w|^4+\gamma_0|u|^2|w|^2+\rho_0\operatorname{Re}\bar u^2w^2.
                                                                    \tag{15.3}
\]

Both waves lower their own and each other's frequencies, and the pair
term \(\rho_0\operatorname{Re}\bar u^2w^2\) converts a gauge pair into a
Higgs pair and back.  On the Poincar\'e sphere of \((u,w)\), with
\(s=(|u|^2-|w|^2,2\operatorname{Re}\bar uw,2\operatorname{Im}\bar uw)\)
and \(|s|=N\),

\[
 H_0=\hbox{const}+\frac{\alpha_0+\beta_0}8s_1^2+\frac{\alpha_0-\beta_0}4Ns_1
 +\frac{\gamma_0+\rho_0}4s_2^2+\frac{\gamma_0-\rho_0}4s_3^2,             \tag{15.3a}
\]

a quadratic form plus a linear term along the exchange axis from the
unequal self-couplings, so the motion is integrable as in V9, with the
fixed points displaced from the axes by the linear term.
\end{theorem}

\begin{proof}
The six coefficients of each of \(\Sigma_u,\Sigma_w\) are determined
exactly from six rational data
\((A_f,A_h,\varphi)=(1,1,0),(1,1,\tfrac\pi2),(1,0,0),(0,1,0),(1,2,\arctan\tfrac43),(2,1,\arctan\tfrac{12}5)\)
and verified at three further points; the nine Hamiltonian conditions
(reality of \(\alpha,\beta,\gamma\), equality of the two cross
coefficients, conjugation of the pair coefficients, and the vanishing
of the \(\sigma,\tau\) families) are certified.  The Stokes form
follows from \(|u|^2|v|^2=\tfrac14(s_2^2+s_3^2)\),
\(\operatorname{Re}\bar u^2w^2=\tfrac14(s_2^2-s_3^2)\), and
\(|u|^4+|w|^4=\tfrac12(N^2+s_1^2)\).
\end{proof}

\subsection{With gravity}

\begin{theorem}[Gravitational correction and the loss of the Hamiltonian form]
\label{thm:v20v12-gravity}
With the Einstein equations and the covariant prescription, the
resonant coefficients keep the form (15.1) with

\[
 \alpha=\frac{6211}{46592},\quad\beta=-\frac{4157}{46592},\quad
 \gamma_u=-\frac{2395}{23296},\quad\gamma_w=-\frac{211}{23296},\quad
 \rho=-\frac{4275}{46592},                                            \tag{15.4}
\]

so that the gravitational corrections are
\(\Delta\alpha=\tfrac{117}{512}\), \(\Delta\beta=\tfrac{37}{512}\),
\(\Delta\gamma_u=\tfrac{19}{256}\), \(\Delta\gamma_w=\tfrac{43}{256}\),
\(\Delta\rho=-\tfrac{21}{512}\), all conditions of a Hamiltonian flow
hold except the equality of the cross coefficients, and

\[
 \gamma_u-\gamma_w=-\frac3{32}=-\frac{e^2}{32}.                         \tag{15.5}
\]

The averaged two-mode flow with gravity is therefore not the
Hamiltonian flow of any quartic function of \((u,w)\).
\end{theorem}

\begin{proof}
Same fits and verifications as in Theorem \ref{thm:v20v12-flat} on
the full system; the asymmetry (15.5) is certified exactly.  A
Hamiltonian flow \(\dot u=i\partial H/\partial\bar u\), \(\dot w=i\partial H/\partial\bar w\)
with \(H\) a real quartic has equal cross coefficients, since both are
the coefficient of \(|u|^2|w|^2\) in \(H\); rescaling the symplectic
weights of the two modes would scale the cross and the pair
coefficients by the same ratio, which the certified equality of the
pair coefficients excludes.
\end{proof}

\begin{proposition}[Origin of the asymmetry]
\label{prop:v20v12-origin}
The asymmetry is gravitational and absent in flat space.  The only
couplings of the two waves to the metric that distinguish them are
the transverse expansion \(2a_t\), which damps the Higgs wave and not
the conformally invariant gauge wave, and the anisotropy \(b_t\) and
the transverse conformal factor \(e^{-2a-2b}\) in the gauge mass term,
which act on the gauge wave and on the Higgs--gauge vertex; the
second-order metric that carries these couplings is sourced by both
waves and has secular parts that are not adiabatically eliminable.
The slow dynamics closes only on the enlarged space of the two
amplitudes and the background variables, as the nineteenth cycle's
Bianchi--I reduction already required for the metric alone.
\end{proposition}

\begin{proof}
The first sentence is Theorems \ref{thm:v20v12-flat} and
\ref{thm:v20v12-gravity}.  The list of asymmetric couplings is read
from (14.2); the last sentence is the interpretation, not a further
computation.
\end{proof}

\subsection{Executable certificate}

The verifier

\[
 \texttt{scripts/verify\_sm\_v20\_v12\_higgs\_gauge\_hamiltonian.py}  \tag{15.6}
\]

builds the perturbative system of V11 with the Higgs wave at first
order, runs orders one to three in flat space and with gravity at
nine rational data each, checks all residuals and constraints, fits
and verifies (15.1), and certifies (15.2), (15.4), (15.5) and the
Hamiltonian conditions.  It emits

\[
 \texttt{artifacts/sm\_v20\_v12\_higgs\_gauge\_hamiltonian.json}.     \tag{15.7}
\]

The hashes are

\[
\begin{array}{c|l}
\hbox{object}&\hbox{SHA--256}\\ \hline
\hbox{verifier}&
\texttt{E9C7E94FF645F00EF2622EBEFB4D30A6C91C9E1EEC46CD6544A1D48D638FD937}\\
\hbox{canonical payload}&
\texttt{0FC5988020C37137864794CCB672CA67E7CC8C6CA41117C394DFB5E0469C91DF}.
\end{array}                                                     \tag{15.8}
\]

\begin{firewall}[Equal masses, one coupling set, leading averaging, prescription]
\label{fire:v20v12}
The numbers (15.2), (15.4) hold for the chosen couplings and equal
masses; the identification \(-\tfrac3{32}=-e^2/32\) is a reading of one
value, not a proved dependence on the coupling.  The flat-space
system ignores the Einstein equations and is a comparison, not a
solution of the programme's equations.  The non-Hamiltonian
conclusion refers to the two-mode reduction with the stated covariant
prescription; an enlarged reduction including the background is not
constructed.
\end{firewall}

\begin{decision}[V13 acquisition]
\label{dec:v20v12-next}
V13 will vary the gauge coupling at fixed masses, \(e^2\in\{3,12,27\}\)
with \(v\) adjusted so that \(m_A^2=3\), to determine whether the
asymmetry is exactly \(-e^2/32\) and how the flat-space coefficients
depend on \(e^2\) and \(\lambda\), and will then construct the enlarged
slow system with the background Hubble rate as a third variable, in
which the flow should regain a Hamiltonian form.  V14 closes the
Higgs chapter and V15 is the audit.
\end{decision}

\begin{assessment}[V12 acquisition]
\label{ass:v20v12}
The averaged interaction of the Higgs boson and the massive gauge
boson is computed exactly: a negative-definite quartic Hamiltonian in
flat space, with mutual frequency lowering and pair conversion, and a
gravitational correction that breaks the Hamiltonian form by the
exact asymmetry \(-\tfrac3{32}\) between the two cross terms, because
the expanding background couples to the non-conformal Higgs and to
the conformal gauge wave differently.
\end{assessment}

\begin{record}[V12 append record]
\label{rec:v20v12}
V12 is appended to the validated V11 whose installed SHA--256 was
\texttt{28BA7FAE8F9CDAB2B1424721D0B09632A7854BFF7CA40432A892240C16E627B0}.
After installation only V12 is the active numeric SM note; frozen
archives and unrelated notes are unchanged.
\end{record}


\section{V13: the coupling dependence separates cleanly: matter scales with \texorpdfstring{$e^2$}{e2}, gravity does not}
\label{sec:v20v13}

V12 found the averaged Higgs--gauge interaction at one set of
couplings and read its gravitational asymmetry as \(-e^2/32\).  This
note varies the gauge coupling at fixed masses,
\((e^2,v,\lambda)=(3,2,\tfrac38),(12,1,\tfrac32),(27,\tfrac23,\tfrac{27}8)\),
all with \(m_A^2=m_h^2=3\), and repeats the flat-space and full
fits.  The result refutes the reading and replaces it by an exact
separation: the flat-space coefficients are proportional to \(e^2\),
so the whole matter interaction is \(e^2\) times a fixed quartic
Hamiltonian, while the gravitational corrections, including the
asymmetry \(-\tfrac3{32}\), are the same at all three couplings.  At
fixed masses the geometry's contribution to the slow dynamics of the
Higgs and gauge waves is independent of the gauge coupling: gravity
reads the masses and the amplitudes, not the coupling.

\subsection{Three couplings at fixed masses}

\begin{theorem}[Separation of the matter and gravitational parts]
\label{thm:v20v13-separation}
For the three coupling sets, with the covariant prescription and the
fits of V12 (six fitting data, two further verification data at each
coupling, flat and full), the resonant coefficients have the form
(15.1) at every coupling, are Hamiltonian in flat space, and satisfy:

\begin{enumerate}
\item the flat-space coefficients are exactly proportional to
\(e^2\):
\[
 (\alpha_0,\beta_0,\gamma_0,\rho_0)=\frac{e^2}{34944}\,(-1109,\,-1881,\,-2062,\,-591);  \tag{16.1}
\]
\item the gravitational corrections are identical at the three
couplings:
\[
 \Delta\alpha=\frac{117}{512},\quad\Delta\beta=\frac{37}{512},\quad
 \Delta\gamma_u=\frac{19}{256},\quad\Delta\gamma_w=\frac{43}{256},\quad
 \Delta\rho=-\frac{21}{512};                                            \tag{16.2}
\]
\item the asymmetry is \(\gamma_u-\gamma_w=-\tfrac3{32}\) at
\(e^2=3,12,27\), independent of the coupling; the reading \(-e^2/32\)
of V12 is refuted.
\end{enumerate}
\end{theorem}

\begin{proof}
Certified at the three couplings; the fits are exact and verified
out of sample, the proportionality is checked coefficient by
coefficient against the \(e^2=3\) values, the corrections are checked
equal, and the asymmetry equal to \(-\tfrac3{32}\).
\end{proof}

\begin{corollary}[Structure of the slow dynamics]
\label{cor:v20v13-structure}
At fixed masses the averaged first-order dynamics of the massive gauge
wave and the Higgs wave is

\[
 \dot u=i\bigl(e^2\,\partial_{\bar u}\widehat H+\Sigma^{\rm grav}_u\bigr),\qquad
 \dot w=i\bigl(e^2\,\partial_{\bar w}\widehat H+\Sigma^{\rm grav}_w\bigr),        \tag{16.3}
\]

with \(\widehat H\) the coupling-free quartic of (16.1) and
\(\Sigma^{\rm grav}\) the coupling-free gravitational part (16.2), which
is not a Hamiltonian vector field on \((u,w)\).  The competition
between matter and gravity in the Higgs sector is therefore governed
by \(e^2\) alone, as it was in the colour sector of V9, but here the
gravitational part is not of Hamiltonian form and no critical coupling
can be read from a stability exchange on a sphere.
\end{corollary}

\begin{proof}
Linearity of the third-order equations in their sources: the sources
of the resonant coefficients split into those proportional to the
coupling (the cubic and quartic vertices, whose strengths \(e^2v\),
\(e^2\), and \(\lambda\) all scale as \(e^2\) at fixed \(m_A,m_h\), their
second-order-perturbation products \((e^2v)^2/m_h^2\) scaling likewise)
and the metric couplings, which involve only the stresses of the
first-order waves and hence only the masses and amplitudes; Theorem
\ref{thm:v20v13-separation} certifies the consequence.
\end{proof}

\begin{proposition}[Why the proportionality is exact]
\label{prop:v20v13-why}
With \(m_A^2=e^2v^2/4\) and \(m_h^2=2\lambda v^2\) fixed, the vertices
of (14.2) are \(\tfrac14e^2(2vh+h^2)f\) and
\(\lambda(3vh^2+h^3)+\tfrac14e^2(vf^2+hf^2)\); the cubic vertex
strengths \(e^2v/2\), \(3\lambda v\), \(e^2v/4\) scale as \(1/v\propto e\)
and the quartic ones as \(e^2\), while the second-order fields they
source scale as \(1/v\) and enter the third order multiplied by a
cubic vertex, giving \(1/v^2\propto e^2\).  Every term of the
flat-space resonant coefficients is therefore proportional to
\(e^2\), and the gravitational terms carry no coupling at all.
\end{proposition}

\begin{proof}
Bookkeeping of the vertices under \(v\to v/s\), \(e^2\to s^2e^2\),
\(\lambda\to s^2\lambda\), which fixes both masses.
\end{proof}

\subsection{Executable certificate}

The verifier

\[
 \texttt{scripts/verify\_sm\_v20\_v13\_coupling\_variation.py}        \tag{16.4}
\]

rebuilds the perturbative system at each coupling, runs the flat and
full fits of V12 with verification data, and certifies (16.1)--(16.2)
and the constant asymmetry.  It emits

\[
 \texttt{artifacts/sm\_v20\_v13\_coupling\_variation.json}.           \tag{16.5}
\]

The hashes are

\[
\begin{array}{c|l}
\hbox{object}&\hbox{SHA--256}\\ \hline
\hbox{verifier}&
\texttt{873D5CE055F6424FA3CA4365AF47841BA9D8E055A0A202726103E3B0C92534A5}\\
\hbox{canonical payload}&
\texttt{9397B5F12CEAC20AC2E132AABD2F5DFC9FD1B8B51304EBFB029EF875A6F47D1C}.
\end{array}                                                     \tag{16.6}
\]

\begin{firewall}[Fixed masses, three couplings, equal masses]
\label{fire:v20v13}
The separation is certified at three couplings with \(m_A=m_h\); the
scaling argument of Proposition \ref{prop:v20v13-why} is exact for the
vertices but the mass dependence of the gravitational part is not
explored, and unequal masses, where the cubic resonances change, are
not treated.
\end{firewall}

\begin{decision}[V14 acquisition]
\label{dec:v20v13-next}
V14 closes the Higgs chapter by enlarging the slow system: it will
add the mean expansion rate and the homogeneous anisotropy as slow
variables, determine their coupling to the two amplitudes from the
certified secular structure, and test whether the enlarged system is
Hamiltonian with the asymmetry (16.2) reproduced; V15 is the audit.
\end{decision}

\begin{assessment}[V13 acquisition]
\label{ass:v20v13}
At fixed masses the averaged Higgs--gauge dynamics separates exactly
into a matter part proportional to \(e^2\), Hamiltonian, and a
coupling-independent gravitational part, non-Hamiltonian on the two
amplitudes with the constant asymmetry \(-\tfrac3{32}\); the V12
reading of the asymmetry as \(-e^2/32\) is corrected.
\end{assessment}

\begin{record}[V13 append record]
\label{rec:v20v13}
V13 is appended to the validated V12 whose installed SHA--256 was
\texttt{DA7CF5FFA5866E0B2ABCEBAB0E7912750A3987058F0B4A07253A87A7C71B18DB}.
After installation only V13 is the active numeric SM note; frozen
archives and unrelated notes are unchanged.
\end{record}


\section{V14: the mass laws of the gravitational part, and the close of the Higgs chapter}
\label{sec:v20v14}

V13 separated the averaged Higgs--gauge dynamics into a matter part
proportional to \(e^2\) and a coupling-independent gravitational part.
This note determines the gravitational part as a function of the
common mass.  At five equal masses, \(m_A^2=m_h^2=m^2=\omega^2-1\) with
\(\omega=2,\dots,6\), the flat and full fits of V12 are repeated; the
flat coefficients are Hamiltonian at every mass, and the five
gravitational corrections are exact rational functions of \(\omega\)
of the form \((p\omega^4+q\omega^2+r)/\omega^3\), read off three masses,
unique in a search over two-term and three-term monomial families
with twelve candidate denominators, and verified at two further
masses.  The asymmetry of V12--V13 is
\(-(\omega^2+2)/(32\omega)\).  In the massless limit the gravitational
self-shift of the gauge wave and the pair coefficient vanish, as they
did for two light polarizations in V8, while the Higgs self-shift and
the cross coefficients do not: the non-conformal coupling of the
scalar survives the massless limit.

\subsection{Five masses}

\begin{theorem}[Mass laws of the gravitational corrections]
\label{thm:v20v14-laws}
With \(v=2\), \(e^2=m^2\), \(\lambda=m^2/8\), so that
\(m_A^2=m_h^2=m^2\) and \(\omega=\sqrt{1+m^2}=2,3,4,5,6\), the
covariant flat and full fits of V12 hold at every mass with the
structure (15.1), the flat flow is Hamiltonian, and the gravitational
corrections are

\[
\begin{aligned}
 \Delta\alpha&=\frac{21\omega^4-27\omega^2+6}{128\,\omega^3},\qquad
 \Delta\beta=\frac{\omega^4+13\omega^2+6}{128\,\omega^3},\\
 \Delta\gamma_u&=\frac{3\omega^4-3\omega^2+2}{64\,\omega^3},\qquad
 \Delta\gamma_w=\frac{5\omega^4+\omega^2+2}{64\,\omega^3},\qquad
 \Delta\rho=\frac{-3\omega^4+\omega^2+2}{128\,\omega^3},
\end{aligned}                                                          \tag{17.1}
\]

so that

\[
 \gamma_u-\gamma_w=-\frac{\omega^2+2}{32\,\omega}=-\frac{m^2+3}{32\sqrt{1+m^2}}.
                                                                    \tag{17.2}
\]

The values at \(\omega=2\) are those of (16.2).
\end{theorem}

\begin{proof}
The fits are exact at each mass with two verification data; the
laws (17.1) were determined from \(\omega=2,3,4\) and are certified
to reproduce the independently computed values at \(\omega=5\)
(\(\tfrac{1557}{2000},\tfrac{239}{4000},\tfrac{901}{4000},\tfrac{197}{500},-\tfrac{231}{2000}\))
and \(\omega=6\)
(\(\tfrac{4375}{4608},\tfrac{295}{4608},\tfrac{1891}{6912},\tfrac{3259}{6912},-\tfrac{1925}{13824}\));
(17.2) is their difference.
\end{proof}

\begin{corollary}[Massless limit and large-mass limit]
\label{cor:v20v14-limits}
As \(\omega\to1\), \(\Delta\alpha\to0\), \(\Delta\rho\to0\),
\(\Delta\beta\to\tfrac5{32}\), \(\Delta\gamma_u\to\tfrac1{32}\),
\(\Delta\gamma_w\to\tfrac18\): the conformally coupled gauge wave loses
its gravitational self-shift and its pair coupling to the scalar,
consistent with the vanishing self-shift of a single light
polarization in V8, while the scalar keeps a self-shift and both
cross terms survive with the massless asymmetry \(-\tfrac3{32}\).  As
\(\omega\to\infty\) all corrections grow linearly in \(\omega\),
\(\Delta\alpha\sim\tfrac{21}{128}\omega\), \(\Delta\beta\sim\tfrac1{128}\omega\),
\(\Delta\gamma_u\sim\tfrac3{64}\omega\), \(\Delta\gamma_w\sim\tfrac5{64}\omega\),
\(\Delta\rho\sim-\tfrac3{128}\omega\), the gravitational frequency shifts
of heavy waves being proportional to their energy \(\omega^2/4\) per
unit amplitude divided by their frequency.
\end{corollary}

\begin{proof}
Limits of (17.1).
\end{proof}

\subsection{The Higgs chapter}

\begin{proposition}[What the Higgs chapter established]
\label{prop:v20v14-chapter}
In the plane-symmetric emergent cosmology with the unitary doublet
and one colour wave: the exact Einstein--Yang--Mills--Higgs system
and its masses (V11); the massive-scalar law for massive gauge waves,
so that the geometry reads the vacuum expectation value through
\(m_A\) (V11); the averaged Higgs--gauge interaction as a negative
quartic Hamiltonian in flat space and its non-Hamiltonian
gravitational correction (V12); the exact separation of the matter
part, proportional to \(e^2\), from the coupling-independent
gravitational part (V13); and the mass laws (17.1)--(17.2) of the
gravitational part (V14).  What the chapter did not do: the Goldstone
directions and the second colour, the enlarged slow system in which
the gravitational part would regain a Hamiltonian form, and any
quantum treatment of the Higgs field.
\end{proposition}

\begin{proof}
Record of V11--V14.
\end{proof}

\subsection{Executable certificate}

The verifier

\[
 \texttt{scripts/verify\_sm\_v20\_v14\_mass\_variation.py}            \tag{17.3}
\]

rebuilds the system at each mass, runs the flat and full fits with
verification data, and certifies the structure, the V13 values at
\(\omega=2\), and the laws (17.1)--(17.2) at all five masses.  It
emits

\[
 \texttt{artifacts/sm\_v20\_v14\_mass\_variation.json}.               \tag{17.4}
\]

The hashes are

\[
\begin{array}{c|l}
\hbox{object}&\hbox{SHA--256}\\ \hline
\hbox{verifier}&
\texttt{768206D79E8CB16F931D98FC0F0966B291B05FC0F9F73571BD7482B644B6C4B2}\\
\hbox{canonical payload}&
\texttt{F71DEE8AB6363D7384E9D0D23C49D050D67C5D9B93D210DE663E75816A93A0CF}.
\end{array}                                                     \tag{17.5}
\]

\begin{firewall}[Equal masses, integer frequencies, read-off laws]
\label{fire:v20v14}
The laws (17.1) are certified at five integer frequencies and are not
proved for general \(\omega\); equal masses are assumed throughout,
and unequal masses, where the pair resonance \(\bar u^2w^2\) is
detuned, are not treated; the massless limit of the corollary is a
limit of the formulas, not a computation at \(m=0\), where the Higgs
potential degenerates.
\end{firewall}

\begin{decision}[V15 acquisition]
\label{dec:v20v14-next}
V15 is the compulsory companion audit.  V16 opens the fermionic
sector: the Dirac equation on the plane-symmetric diagonal metric
with a tetrad adapted to it, the spinor stress tensor for a
\((z,t)\)-dependent spinor, and the question whether some spinor
polarization keeps the off-diagonal stress components zero, so that a
fermion wave can be placed in the same exact framework.
\end{decision}

\begin{assessment}[V14 acquisition]
\label{ass:v20v14}
The gravitational part of the averaged Higgs--gauge dynamics is
determined as exact rational functions of the common frequency,
verified out of sample at two masses, with the asymmetry
\(-(\omega^2+2)/(32\omega)\); the massless limit recovers the
vanishing self-shift of the conformal gauge wave and exposes the
surviving non-conformal scalar terms.  The Higgs chapter closes with
the vacuum expectation value, the coupling, and the masses each read
by the geometry through an exact law.
\end{assessment}

\begin{record}[V14 append record]
\label{rec:v20v14}
V14 is appended to the validated V13 whose installed SHA--256 was
\texttt{68098EBFB3A420A2F0F4FFFE6FD6D406A566111670B276BED4243D9A12FD3672}.
After installation only V14 is the active numeric SM note; frozen
archives and unrelated notes are unchanged.
\end{record}


\section{V15: compulsory companion audit before the fermionic sector}
\label{sec:v20v15}

This is the required audit at the fifteenth version.  Every active
companion note present in the shared folder was inspected read-only
before the fermionic sector announced in V14 begins.

\begin{record}[V15 companion snapshot]
\label{rec:v20v15-snapshot}
The files present were

\[
\begin{array}{c|r|r}
\hbox{file and SHA--256}&\hbox{bytes}&\hbox{lines}\\ \hline
\texttt{Renewal\_NS\_Stability\_Research\_Notes\_V26.tex}
&278283&5319\\
\multicolumn{3}{l}{
\texttt{82C887F8C2C69E4F28FAD7C3DC8DE5B9099CB5A4BF431EBA1FFF3E91935978AD}}
\\
\texttt{Renewal\_GRH\_Cofinal\_Visibility\_Attack\_V62.tex}
&602301&17779\\
\multicolumn{3}{l}{
\texttt{54CDB21330B0B7A062D350CDD091364395374203DFA3B1DF3676667F667C4A5F}}
\\
\texttt{Renewal\_Yang\_Mills\_Mass\_Gap\_Research\_Notes\_V36.tex}
&539561&13242\\
\multicolumn{3}{l}{
\texttt{57DE9862277623D9AEE2DE7870ED86627E5620DB18DB3C0E8ACF5F1EF16817A5}}
\\
\texttt{Renewal\_YM\_Parallel\_Research\_Notes\_V5.tex}
&54174&1351\\
\multicolumn{3}{l}{
\texttt{306F1BB84CFA8A8D47EA569EC17E9545F9867C8D32B548EC9D9C444B4F3C0512}}.
\end{array}                                                     \tag{18.1}
\]

Since V10 the GRH note advanced from V61 to V62 (the Appell
reformulation of the real-axis problem) and the Yang--Mills mass-gap
note from V35 to V36 (levelwise nondegeneracy, lift chains, the exact
rank-45 vertical covariance, and the reduction of global
nondegeneracy of the fixed point to a single unproved inequality);
the NS and YM parallel notes are unchanged.
\end{record}

\begin{decision}[V15 audit decision]
\label{dec:v20v15-audit}
No companion theorem, numerical constant, or executable artifact is
imported.  The GRH advance is a hypergeometric reformulation of a
polynomial inequality; the Yang--Mills advance is a Gaussian
nondegeneracy statement for a Euclidean blocking.  Neither concerns
the Lorentzian Einstein--matter systems of this cycle.  The V14
direction stands.
\end{decision}

\begin{proof}
The changed files are the GRH and Yang--Mills notes; their new
sections contain no constrained datum, no Einstein evolution, and no
averaged mode dynamics.
\end{proof}

\begin{assessment}[V15 acquisition]
\label{ass:v20v15}
The audit is current.  Since V10 the cycle has opened and closed the
Higgs chapter: the exact Einstein--Yang--Mills--Higgs system, the
massive-scalar law for massive gauge waves, the averaged Higgs--gauge
Hamiltonian in flat space with its non-Hamiltonian gravitational
correction, the exact separation of the matter part proportional to
\(e^2\) from the coupling-independent gravitational part, and the
mass laws of the latter.  No Standard Model--Einstein continuum is
claimed.
\end{assessment}

\begin{record}[V15 append record]
\label{rec:v20v15}
V15 is appended to the validated V14 whose installed SHA--256 was
\texttt{A93D49DC7E60B2F63226D1E78A7D6A6FC5EE9DBFFE240BF1308FFD5DCF7678EE}.
After installation only V15 is the active numeric SM note; the
companions, frozen archives, and all unrelated notes are unchanged.
The next compulsory companion audit is V20, which is also the second
ten-version sweep of the cycle.
\end{record}


\section{V16: the fermionic sector opens: a spin-polarized Dirac wave is exact null dust}
\label{sec:v20v16}

The last Standard Model structure absent from the plane-symmetric
framework is the fermion.  This note places the Dirac field on the
diagonal plane-symmetric metric with an adapted tetrad, fixes and
certifies the conventions by the on-shell conservation of the
stress tensor, and finds that a Dirac wave in one eigenspace of the
spin along the propagation axis is compatible with the diagonal
ansatz: its stress has no transverse pressure at all, it never
sources the gravitational polarization, and its only off-diagonal
component is proportional to that polarization, so the consistent
sector is \(b=0\).  For the massless standing wave in one spin sector
on a homogeneous background the result is exact and complete: the
Dirac equation is solved in closed form, the stress is homogeneous
null dust \(\rho=p_z=e^{-2a-2c}\) with vanishing transverse pressure
and vanishing flux, and the Einstein equations reduce, without
averaging and to all orders, to the Bianchi--I null-dust system that
the nineteenth cycle obtained for gravitational waves only by
averaging, and that this cycle found to fail for gravitons and to
hold through fourth order for light.  For the polarized massless
fermion the reduction is not an approximation but an identity.

\subsection{Conventions and the sector}

\begin{proposition}[Dirac field on the plane-symmetric metric]
\label{prop:v20v16-dirac}
With the tetrad \(e^A=(dt,e^{a+b}dx,e^{a-b}dy,e^cdz)\), the
mostly-plus gamma matrices \(\gamma^A=i\gamma^A_D\), the spin
connection \(\Gamma_\mu=\tfrac18\omega_{\mu AB}[\gamma^A,\gamma^B]\), and
\(\bar\psi=\psi^\dagger\gamma^0_D\), the stress tensor
\(T_{\mu\nu}=-\tfrac14[\bar\psi\gamma_{(\mu}\nabla_{\nu)}\psi-(\nabla_{(\mu}\bar\psi)\gamma_{\nu)}\psi]\)
is real and is conserved on shell.  For a spinor in the \(\Sigma^3=+1\)
eigenspace, \(\psi=(u,0,w,0)\) with \(u,w\) functions of \((z,t)\), the
Dirac equation reduces to

\[
\begin{aligned}
 &i\,u_t+i\bigl(a_t+\tfrac12c_t\bigr)u+i\,e^{-c}\bigl(w_z+a_zw\bigr)-mu=0,\\
 &-i\,w_t-i\bigl(a_t+\tfrac12c_t\bigr)w-i\,e^{-c}\bigl(u_z+a_zu\bigr)-mw=0,
\end{aligned}                                                          \tag{19.1}
\]

independent of \(b\), the eigenspace being preserved; its stress has
\(T_{tx}=T_{ty}=T_{xz}=T_{yz}=0\), \(T_{xx}=T_{yy}=0\), and

\[
 T_{xy}=-\tfrac12e^{2a-c}\bigl[(|u|^2+|w|^2)e^{c}b_t+(u\bar w+w\bar u)b_z\bigr],
                                                                    \tag{19.2}
\]

so that the diagonal ansatz is consistent exactly when \(b=0\), which
is preserved because the wave sources no polarization.
\end{proposition}

\begin{proof}
All statements are certified symbolically: the Clifford relations,
the reduction of the Dirac operator, the absence of \(b\), the
vanishing components, the reality of the surviving ones, the form of
\(T_{xy}\), and the two components of \(\nabla^\mu T_{\mu\nu}=0\) after
elimination of the time derivatives by (19.1).  The conservation test
fixed the index placement of the spin connection, whose naive form
failed it.
\end{proof}

\subsection{The exact massless standing wave}

\begin{theorem}[Polarized massless Dirac wave is exact null dust]
\label{thm:v20v16-nulldust}
Let \(a=a(t)\), \(c=c(t)\), \(b=0\), and

\[
 \psi=e^{-a-c/2}\,e^{-i\theta(t)}\,(\cos z,\,0,\,i\sin z,\,0),\qquad \dot\theta=e^{-c}.
                                                                    \tag{19.3}
\]

Then \(\psi\) solves the massless Dirac equation on the metric
\(-dt^2+e^{2a}(dx^2+dy^2)+e^{2c}dz^2\) exactly, its stress is

\[
 T_{tt}=e^{-2c}T_{zz}=\rho:=e^{-2a-2c},\qquad T_{xx}=T_{yy}=T_{tz}=0,\qquad\operatorname{tr}T=0,
                                                                    \tag{19.4}
\]

homogeneous and static up to the redshift factor, and the Einstein
equations \(\operatorname{Ric}=\Lambda g+T\) reduce exactly to

\[
 \ddot a+\dot a(2\dot a+\dot c)=\Lambda,\qquad
 \ddot c+\dot c(2\dot a+\dot c)=\Lambda+\rho,\qquad
 \dot a^2+2\dot a\dot c=\Lambda+\rho,                                 \tag{19.5}
\]

the Bianchi--I null-dust system (94.2) of the nineteenth cycle with
\(8\pi\rho\to\rho\) and \(\rho\propto e^{-2a-2c}\), the momentum
constraint holding identically.
\end{theorem}

\begin{proof}
Substitution of (19.3) into (19.1) with \(m=0\) gives zero identically
once \(\dot\theta=e^{-c}\); the stress components (19.4) and the
tracelessness are certified; the Ricci tensor of the homogeneous
metric and the forms (19.5) are certified against
\(\operatorname{Ric}-\Lambda g-T\), including \(T_{tz}=0\) and the
momentum equation.
\end{proof}

\begin{corollary}[Exact versus averaged reductions]
\label{cor:v20v16-comparison}
The system (19.5) governs the balanced torus of the nineteenth cycle
only after averaging over the wave phase and fails at fourth order for
the graviton (V92 there); it holds through fourth order for
electromagnetic radiation (V3); and it is an identity for the
polarized massless Dirac standing wave, whose stress carries neither
the second-harmonic breathing nor the inhomogeneous modulation of the
bosonic waves, because a spin-\(\tfrac12\) standing wave of definite
\(\Sigma^3\) has \(|u|^2+|w|^2=\cos^2z+\sin^2z\) constant.  In
particular the focusing theorem of the nineteenth cycle (its V92)
applies to the fermion cosmology with the exact lifetime of (19.5),
and the collapse bound \(t_*\leq4/\sqrt{-3\Lambda}\) of its V91 is
exact for it.
\end{corollary}

\begin{proof}
Theorem \ref{thm:v20v16-nulldust} and the cited results, whose
hypotheses (19.5) satisfies exactly.
\end{proof}

\begin{proposition}[What the geometry reads off the fermion]
\label{prop:v20v16-reading}
At the level of the exact solution the geometry reads the polarized
massless fermion as null dust along its propagation axis, like
gravitational or electromagnetic radiation on average, but with no
transverse pressure, no anisotropic stress, and no oscillation: the
fermionic standing wave is the one matter configuration in the
programme whose emergent cosmology is a closed system of ordinary
differential equations without any approximation.  Its spin couples
to the gravitational polarization through (19.2), so a fermion in a
polarized gravitational wave background is the first place where the
spin becomes visible, outside the consistent sector of this note.
\end{proposition}

\begin{proof}
Proposition \ref{prop:v20v16-dirac} and Theorem \ref{thm:v20v16-nulldust}.
\end{proof}

\subsection{Executable certificate}

The verifier

\[
 \texttt{scripts/verify\_sm\_v20\_v16\_dirac\_null\_dust.py}          \tag{19.6}
\]

builds the tetrad, spin connection, Dirac operator, and stress
tensor symbolically, certifies Proposition \ref{prop:v20v16-dirac}
including on-shell conservation, and certifies Theorem
\ref{thm:v20v16-nulldust}.  It emits

\[
 \texttt{artifacts/sm\_v20\_v16\_dirac\_null\_dust.json}.             \tag{19.7}
\]

The hashes are

\[
\begin{array}{c|l}
\hbox{object}&\hbox{SHA--256}\\ \hline
\hbox{verifier}&
\texttt{6AC781D784130C457923EE0E2E19CA08C404DED38A749E5D6748B70E6D6223ED}\\
\hbox{canonical payload}&
\texttt{8643A6855E4D740BF72D7580BDA991EAD713A7D239403D8BB1BCB0E92E54BE3B}.
\end{array}                                                     \tag{19.8}
\]

\begin{firewall}[Classical spinor, one sector, massless, homogeneous background]
\label{fire:v20v16}
The Dirac field is a classical commuting spinor, not a quantized or
Grassmann field, so its stress is a c-number and its physical
interpretation as fermionic matter is by analogy; only the
\(\Sigma^3=+1\) sector, the massless case, and the homogeneous
background are treated exactly; the mixed sector, the massive wave,
and inhomogeneous backgrounds are not.
\end{firewall}

\begin{decision}[V17 acquisition]
\label{dec:v20v16-next}
V17 treats the massive polarized standing wave: on the homogeneous
background the Dirac equation reduces to a two-level system in the
rescaled amplitudes \((\tilde u,\tilde w)\) with Hamiltonian matrix
\(\bigl(\begin{smallmatrix}m&e^{-c}\\ e^{-c}&-m\end{smallmatrix}\bigr)\),
whose Rabi oscillation between the balanced state and its partner
makes the stress inhomogeneous; V17 will derive the exact system and
its second-order secular structure and compare with the massive
scalar law.
\end{decision}

\begin{assessment}[V16 acquisition]
\label{ass:v20v16}
The fermionic sector is opened with certified conventions, the
consistency of a spin-polarized Dirac wave with the diagonal
plane-symmetric ansatz, and an exact theorem: the massless polarized
standing wave is homogeneous null dust and reduces the Einstein--Dirac
system, without approximation, to the Bianchi--I null-dust cosmology
of the programme.
\end{assessment}

\begin{record}[V16 append record]
\label{rec:v20v16}
V16 is appended to the validated V15 whose installed SHA--256 was
\texttt{B02C6B7C01AAB4605AFE7B24A559DACDE33C9BD0D25D76D03CAD3B2C31EC7308}.
After installation only V16 is the active numeric SM note; frozen
archives and unrelated notes are unchanged.
\end{record}


\section{V17: the massive polarized fermion wave: a static inhomogeneous stress with the universal equation of state}
\label{sec:v20v17}

V16 solved the massless polarized Dirac standing wave exactly.  This
note adds the mass.  In flat space the polarized standing wave of a
Dirac field of mass \(m\) is again exact, a two-component
positive-frequency state with amplitude ratio \(\omega-m\) between
its two nonzero components, and its stress is static but no longer
homogeneous: the energy density is modulated as
\(\cos^2z+(\omega-m)^2\sin^2z\) while the longitudinal pressure
\(\omega-m\) stays homogeneous, there is no transverse pressure, no
flux, and the trace is the mass term.  The averaged equation of
state is \(\langle p_z\rangle/\langle\rho\rangle=1/\omega^2\), the same
as for the massive scalar and the massive gauge wave: the emergent
geometry reads the fermion mass by the same dispersion law as every
other massive field.  Because the stress is static, the second-order
metric is obtained in closed form for all \(m\), with a static
inhomogeneous conformal factor and quadratic secular terms whose means
obey the massive law, and it is certified against the full Einstein
equations expanded to second order.

\subsection{The flat massive wave}

\begin{theorem}[Massive polarized standing wave]
\label{thm:v20v17-wave}
In flat space, in the \(\Sigma^3=+1\) sector, with \(\omega=\sqrt{1+m^2}\),

\[
 \psi=e^{-i\omega t}\bigl(\cos z,\ 0,\ i(\omega-m)\sin z,\ 0\bigr)         \tag{20.1}
\]

solves the Dirac equation of mass \(m\) exactly, and its stress is

\[
 \rho=\omega\bigl(\cos^2z+(\omega-m)^2\sin^2z\bigr),\qquad p_z=\omega-m,\qquad
 p_x=p_y=0,\qquad T_{tz}=0,\qquad \operatorname{tr}T=-m\bar\psi\psi,       \tag{20.2}
\]

with \(\langle\rho\rangle=\omega^2(\omega-m)\), so that
\(\langle p_z\rangle/\langle\rho\rangle=1/\omega^2\).
\end{theorem}

\begin{proof}
The Dirac equation (19.1) on the flat metric reduces for
\(\psi=e^{-i\omega t}(U\cos z,0,iW\sin z,0)\) to the two-level system
\(i\partial_t(U,W)=\bigl(\begin{smallmatrix}m&1\\1&-m\end{smallmatrix}\bigr)(U,W)\),
whose positive-frequency eigenvector is \((1,\omega-m)\); the stress
components and the trace identity are certified symbolically with
\(m\) symbolic, and the mean uses \(1+(\omega-m)^2=2\omega(\omega-m)\).
\end{proof}

\begin{corollary}[The massive law for the three spins]
\label{cor:v20v17-law}
The averaged stresses of the massive scalar (nineteenth cycle, V98),
the massive gauge wave (V11), and the massive polarized fermion
(20.2) all have \(p_x=p_y=0\) and \(\langle p_z\rangle/\langle\rho\rangle=k^2/\omega^2\)
at wavenumber \(k=1\): null dust plus dust in the ratio \(k^2:m^2\).
The geometry reads a mass, not a spin, from the equation of state.
\end{corollary}

\begin{proof}
Theorem \ref{thm:v20v17-wave} with the cited results.
\end{proof}

\subsection{The second-order metric in closed form}

\begin{theorem}[Closed-form second-order Einstein--Dirac solution]
\label{thm:v20v17-second}
With the time-symmetric conformally flat metric datum and the wave
(20.1) at amplitude \(\varepsilon\), the second-order metric is, for
every \(m\),

\[
\begin{aligned}
 \Lambda_2&=-\omega^2(\omega-m),\qquad
 \alpha(z)=\frac{m\,\omega(\omega-m)}8\cos2z,\qquad b_2=0,\\
 a_2&=\alpha(z)+\Bigl[\Lambda_2+\tfrac12\bigl(\langle\rho\rangle-p_z\bigr)\Bigr]\frac{t^2}2,\qquad
 c_2=\alpha(z)+\frac{p_z-\rho(z)}2\,\frac{t^2}2,
\end{aligned}                                                          \tag{20.3}
\]

and every component of \(\operatorname{Ric}-\Lambda g-T+\tfrac12(\operatorname{tr}T)g\)
vanishes at second order.  The mean secular accelerations obey

\[
 2\langle\ddot a_2\rangle=-\bigl(\langle\rho\rangle+p_z\bigr)=-(\omega-m)(\omega^2+1),\qquad
 2\langle\ddot c_2\rangle=-\bigl(\langle\rho\rangle-p_z\bigr)=-m^2(\omega-m),
                                                                    \tag{20.4}
\]

the massive law of the nineteenth cycle, and the longitudinal
secular term is inhomogeneous, \(\propto(p_z-\rho(z))\), as for the
graviton and the scalar and unlike the light and the massless fermion.
\end{theorem}

\begin{proof}
The stress is static, so the second-order equations have static
sources: the Hamiltonian constraint fixes \(\Lambda_2=-\langle\rho\rangle\)
and \(\alpha''=-\tfrac12(\rho-\langle\rho\rangle)\); the transverse
equation \(a_{2,tt}-a_{2,zz}=\Lambda_2+\tfrac12(\rho-p_z)\) then has
the homogeneous time coefficient of (20.3), the \(z\)-dependence of its
source being carried by \(\alpha''\); the longitudinal equation gives
\(c_{2,tt}=\tfrac12(p_z-\rho)\); \(b\) is unsourced and the momentum
constraint holds since \(T_{tz}=0\).  The six independent components
of the second-order Einstein equations are certified to vanish
symbolically, and the means (20.4) follow from (20.2).
\end{proof}

\begin{proposition}[Static versus oscillating stresses]
\label{prop:v20v17-static}
The second-order metric of a single-frequency polarized fermion wave
contains no oscillatory term at all, massless or massive, whereas the
bosonic waves of this cycle and the last produce second-harmonic
breathing and, for the graviton and the scalar, inhomogeneous secular
modulations tied to it.  The fermion's inhomogeneous secular term is
static in origin, \(\rho(z)\) being modulated by the mass through
\((\omega-m)^2\neq1\); in the massless limit it disappears and the
exact null-dust reduction of V16 is recovered, while in the
large-mass limit \(\omega-m\to0\) the wave's energy, pressure, and
back-reaction all vanish at fixed amplitude.
\end{proposition}

\begin{proof}
(20.2)--(20.3) and their limits; the statements about the bosons are
those of V3, V12, and the nineteenth cycle's V92 and V98.
\end{proof}

\subsection{Executable certificate}

The verifier

\[
 \texttt{scripts/verify\_sm\_v20\_v17\_massive\_fermion.py}           \tag{20.5}
\]

certifies the wave, the stress (20.2), the trace identity, the
equation of state, the closed forms (20.3) against the full
second-order Einstein equations, and the means (20.4), with \(m\)
symbolic.  It emits

\[
 \texttt{artifacts/sm\_v20\_v17\_massive\_fermion.json}.              \tag{20.6}
\]

The hashes are

\[
\begin{array}{c|l}
\hbox{object}&\hbox{SHA--256}\\ \hline
\hbox{verifier}&
\texttt{1EF0B0A350875E808555DC8DDD0128F1D9B060B92AF2DDDFB6FF64135D36B1F5}\\
\hbox{canonical payload}&
\texttt{4DF6EA24400D8487BED411C9F27B15F9D89212B0859233E4354AC9387ECA07AE}.
\end{array}                                                     \tag{20.7}
\]

\begin{firewall}[Classical spinor, second order, one sector]
\label{fire:v20v17}
The spinor is a classical commuting field; its negative-frequency
partner has negative energy, a pathology of the classical
Einstein--Dirac system that the quantum theory removes by normal
ordering and that this note does not address.  Only the second-order
metric is computed; the third-order back-reaction on the wave and the
mixed-sector spinors are not treated.
\end{firewall}

\begin{decision}[V18 acquisition]
\label{dec:v20v17-next}
V18 couples the fermion to the gauge field: an \(SU(2)\) doublet of
Dirac fields whose isospin-up component lies in the \(\Sigma^3=+1\)
sector and whose isospin-down component lies in the \(\Sigma^3=-1\)
sector, so that the gauge term \(-ie f\gamma^x\tau^1/2\) of the colour
ansatz maps the doublet into itself; V18 will derive the plane-symmetric
Einstein--Yang--Mills--Dirac system for this spin--isospin-locked
doublet, certify the consistency of the diagonal ansatz, and obtain
the fermion current that sources the colour wave.
\end{decision}

\begin{assessment}[V17 acquisition]
\label{ass:v20v17}
The massive polarized fermion wave is exact in flat space with a
static inhomogeneous energy density and homogeneous longitudinal
pressure, its equation of state is the universal \(k^2/\omega^2\) of
all massive fields of the programme, and its second-order metric is
closed form for every mass, certified against the full Einstein
equations, with the massive secular law and an inhomogeneous
longitudinal term of static origin.
\end{assessment}

\begin{record}[V17 append record]
\label{rec:v20v17}
V17 is appended to the validated V16 whose installed SHA--256 was
\texttt{2B59AFB2F239BFB095FD7A32087197ED0DFA12C81D7A89E87A52D0BD774E21CF}.
After installation only V17 is the active numeric SM note; frozen
archives and unrelated notes are unchanged.
\end{record}


\section{V18: the fermion meets the gauge field: a spin--isospin-locked doublet that keeps the plane-symmetric framework exact}
\label{sec:v20v18}

V16--V17 placed a single polarized Dirac field in the plane-symmetric
framework.  This note couples an \(SU(2)\) doublet of Dirac fields to
the colour wave \(A^1=f\,dx\) of V1.  The gauge term of the covariant
derivative is \(-ief\gamma^x\tau^1/2\): it swaps the isospin components
and, through \(\gamma^x\), flips the spin sector, so a doublet whose
isospin-up component has spin \(\Sigma^3=+1\) and whose isospin-down
component has spin \(\Sigma^3=-1\) is mapped into itself.  For this
locked doublet the four Dirac equations, the colour current, and the
fermion stress are derived and certified.  The gauge coupling
produces two new stress components: a transverse pressure along the
colour direction \(x\), which sources the gravitational polarization
\(b\), and an off-diagonal component \(T_{xy}=-\tfrac12efe^aS\) with
\(S=\operatorname{Im}(u_2\bar w_1+w_2\bar u_1)\), which would break the
diagonal metric ansatz.  The resolution is a further pairing: on the
subspace \(\psi_2=iP\psi_1\), where \(P\) maps the sector-plus spinor
\((u,0,w,0)\) to \((0,u,0,-w)\), the two isospin components' equations
coincide, so the subspace is invariant under the full gauged dynamics
with mass and metric, \(S\) vanishes identically, and every
\(b\)-dependent part of \(T_{xy}\) cancels between the two components.
On the paired subspace the Einstein--Yang--Mills--Dirac system is
exactly consistent with the diagonal ansatz, and the fermion sources
the gauge field and the gravitational polarization through the same
bilinear \(\operatorname{Im}(\bar w_1u_1)\).

\subsection{The locked doublet}

\begin{proposition}[Gauged Dirac equations and current]
\label{prop:v20v18-doublet}
For \(\Psi=(\psi_1,\psi_2)\) with \(\psi_1=(u_1,0,w_1,0)\), \(\psi_2=(0,u_2,0,w_2)\)
and the colour potential \(A^1=f\,dx\), the covariant derivative
\(D_\mu\Psi=\nabla_\mu\Psi-ieA^a_\mu\tfrac{\tau^a}2\Psi\) preserves the
sector structure, and the Dirac equation
\((\gamma^\mu D_\mu-m)\Psi=0\) reads

\[
\begin{aligned}
 &i\,u_{1,t}+i(a_t+\tfrac12c_t)u_1+i\,e^{-c}(w_{1,z}+a_zw_1)-mu_1+\tfrac12ef\,e^{-a-b}w_2=0,\\
 &-i\,w_{1,t}-i(a_t+\tfrac12c_t)w_1-i\,e^{-c}(u_{1,z}+a_zu_1)-mw_1-\tfrac12ef\,e^{-a-b}u_2=0,
\end{aligned}                                                          \tag{21.1}
\]

with the isospin-down equations obtained by the exchange
\(1\leftrightarrow2\).  The \(f\)-linear part of the Dirac Lagrangian is
the real current \(J=e\,e^{-a-b}\operatorname{Re}(\psi_1^\dagger\alpha^1\psi_2)\),
\(\alpha^1=\gamma^0_D\gamma^1_D\), and the Yang--Mills equation of \(f\)
becomes (3.3) with \(g=0\) and the source \(e^{2a+2b}J\) on the right.
\end{proposition}

\begin{proof}
Certified symbolically: the vanishing of the cross-sector rows of the
operator, the form (21.1), the reality and form of the current, and
the Euler--Lagrange equation of \(f\) from the total Lagrangian.
\end{proof}

\begin{theorem}[Stress of the locked doublet]
\label{thm:v20v18-stress}
The fermion stress has \(T_{tx}=T_{ty}=T_{xz}=T_{yz}=0\), \(T_{yy}=0\),

\[
 e^{-2a-2b}T_{xx}=-\tfrac12ef\,e^{-a-b}\bigl(u_1\bar w_2+u_2\bar w_1+w_1\bar u_2+w_2\bar u_1\bigr),
                                                                    \tag{21.2}
\]

and, at \(b=0\),

\[
 T_{xy}=-\tfrac12ef\,e^{a}\,S,\qquad S=\operatorname{Im}\bigl(u_2\bar w_1+w_2\bar u_1\bigr),
                                                                    \tag{21.3}
\]

while for \(b\neq0\) it carries in addition the terms of (19.2) for
each component, with opposite signs of the \(b_t\) part for the two
sectors.  The gauge coupling thus gives the fermion a transverse
pressure along the colour direction only, and an off-diagonal stress
\(T_{xy}\) that is present at \(b=0\).
\end{theorem}

\begin{proof}
Certified symbolically from the gauge-covariant Belinfante tensor.
\end{proof}

\subsection{The paired subspace}

\begin{theorem}[Exact consistency on the paired subspace]
\label{thm:v20v18-paired}
Let \(P(u,0,w,0)=(0,u,0,-w)\) and consider \(\psi_2=iP\psi_1\), that
is, \(u_2=iu_1\), \(w_2=-iw_1\).  Then:
\begin{enumerate}
\item the subspace is invariant under (21.1) and its partner with
mass, metric, and colour field arbitrary: the isospin-down equations
are \(i\) times the isospin-up ones;
\item \(S\equiv0\), and the full \(T_{xy}\), including its \(b_t\) and
\(b_z\) parts, vanishes identically for any \(b\);
\item the anisotropic stress is
\[
 e^{-2a-2b}T_{xx}-e^{-2a+2b}T_{yy}=2ef\,e^{-a-b}\operatorname{Im}(\bar w_1u_1),
                                                                    \tag{21.4}
\]
and the current is \(J=-2e\,e^{-a-b}\operatorname{Im}(\bar w_1u_1)\).
\end{enumerate}
Consequently the plane-symmetric Einstein--Yang--Mills--Dirac system
with the diagonal metric, the colour wave \(f\), and the paired locked
doublet is exactly consistent: the fermion sources the colour wave
and the gravitational polarization through the single bilinear
\(\operatorname{Im}(\bar w_1u_1)\), and the polarization \(b\), once
present, acts on the fermion only through the stress, never through
its equation of motion.
\end{theorem}

\begin{proof}
All three items are certified symbolically; the invariance is the
identity \(D_2=(iD_{1,0},-iD_{1,2})\) on the subspace.  The last
sentence restates that (21.1) contains no \(b\).
\end{proof}

\begin{proposition}[What the geometry reads off the gauged fermion]
\label{prop:v20v18-reading}
For the massless standing wave of V16 on the paired subspace,
\(\operatorname{Im}(\bar w_1u_1)=-\tfrac12\sin2z\): the fermion pair
sources a static colour field at wavenumber two and, through (21.4),
a gravitational polarization at wavenumber two in the presence of the
colour wave, while its own energy stays the homogeneous null dust of
V16.  The first gauge--fermion effects in the emergent cosmology are
therefore a colour condensate at twice the fermion wavenumber and an
anisotropic stress proportional to the product of the colour field
and the fermion pair density, both invisible in the averaged energy
and pressure.
\end{proposition}

\begin{proof}
Insert \(u_1=e^{-it}\cos z\), \(w_1=ie^{-it}\sin z\) into
\(\operatorname{Im}(\bar w_1u_1)\); the rest is Theorem
\ref{thm:v20v18-paired}.
\end{proof}

\subsection{Executable certificate}

The verifier

\[
 \texttt{scripts/verify\_sm\_v20\_v18\_fermion\_gauge.py}             \tag{21.5}
\]

derives the gauged Dirac operator, current, Yang--Mills source, and
stress for the locked doublet, and certifies Proposition
\ref{prop:v20v18-doublet} and Theorems \ref{thm:v20v18-stress} and
\ref{thm:v20v18-paired}.  It emits

\[
 \texttt{artifacts/sm\_v20\_v18\_fermion\_gauge.json}.                \tag{21.6}
\]

The hashes are

\[
\begin{array}{c|l}
\hbox{object}&\hbox{SHA--256}\\ \hline
\hbox{verifier}&
\texttt{45CBB228237BC641FC906657B209E656567CC660B9DF2184956C05BF0D19C5F8}\\
\hbox{canonical payload}&
\texttt{161B316B2CA7C2F2FFF851712B11777C39A5C94EF6EB82768E605FA6D63B989E}.
\end{array}                                                     \tag{21.7}
\]

\begin{firewall}[One colour, classical doublet, no dynamics computed]
\label{fire:v20v18}
Only the colour component \(A^1\) along \(x\) is present, and only the
paired locked doublet is shown consistent; the general locked doublet
has \(T_{xy}\neq0\) and leaves the diagonal framework.  The doublet is
classical; no second-order dynamics of the coupled system is computed
in this note, and the Higgs--fermion Yukawa coupling is absent.
\end{firewall}

\begin{decision}[V19 acquisition]
\label{dec:v20v18-next}
V19 computes the first dynamics of the exact fermion--gauge system on
the paired subspace: the static colour field and gravitational
polarization sourced at wavenumber two by the massless fermion pair,
their back-reaction on the fermion at the next order, and the
second-order metric, all with the exact engine; V20 is the audit and
the second ten-version sweep.
\end{decision}

\begin{assessment}[V18 acquisition]
\label{ass:v20v18}
The fermion is coupled to the gauge field within the exact framework:
a spin--isospin-locked doublet is mapped into itself by the colour
coupling, its stress acquires a colour-direction pressure and an
off-diagonal component, and on the invariant paired subspace the
off-diagonal component vanishes identically, leaving an exactly
consistent Einstein--Yang--Mills--Dirac plane-symmetric system in
which the fermion pair density sources both the colour wave and the
gravitational polarization.
\end{assessment}

\begin{record}[V18 append record]
\label{rec:v20v18}
V18 is appended to the validated V17 whose installed SHA--256 was
\texttt{3E1A7EC03F9FD5FC498249DB6C5CBD1A2C03E34470F3E76D65145B7AECF2DDC9}.
After installation only V18 is the active numeric SM note; frozen
archives and unrelated notes are unchanged.
\end{record}


\section{V19: the colour condensate of a fermion pair, its self-energy, and the uniform anisotropy it drives}
\label{sec:v20v19}

V18 established the exactly consistent Einstein--Yang--Mills--Dirac
plane-symmetric system on the paired locked doublet.  This note
computes its first dynamics for the massless standing fermion wave
of V16.  The fermion pair carries the colour current \(e\sin2z\),
which induces a static colour field \(f_2=\tfrac e4\sin2z\), a
condensate at twice the fermion wavenumber with energy density
\(\tfrac{e^2}8\cos^22z\).  The induced field acts back on the fermion
through the gauge term of (21.1): the resulting source is, at
wavenumber one, exactly proportional to the fermion mode itself, so
it is resonant and shifts the fermion frequency by \(-e^2\varepsilon^2/16\),
the colour self-energy of the pair, with no amplitude growth; the
wavenumber-three response is bounded.  At fourth order the
anisotropic stress of the fermion--colour coupling and that of the
colour field itself add to the exactly homogeneous value \(-e^2/8\):
the gravitational polarization is driven uniformly,
\(b_4=\tfrac{e^2}{16}t^2\), with no static modulation.  The emergent
geometry thus reads the fermion--gauge coupling as a uniform
anisotropic stress and the fermion reads it as a mass-like frequency
shift.

\subsection{The condensate}

\begin{theorem}[Induced colour field of the fermion pair]
\label{thm:v20v19-condensate}
For the paired doublet with \(\psi_1=\varepsilon e^{-it}(\cos z,0,i\sin z,0)\)
on the flat background, the colour current of Proposition
\ref{prop:v20v18-doublet} is \(J=\varepsilon^2e\sin2z\), the Yang--Mills
equation \(f_{tt}-f_{zz}=J\) has the static solution

\[
 f_2=\frac e4\sin2z,                                                    \tag{22.1}
\]

whose energy density and anisotropic stress are
\(\rho_f=\tfrac{e^2}8\cos^22z\), mean \(\tfrac{e^2}{16}\), and
\(\pi_f=-\tfrac{e^2}8\cos^22z\), all at order \(\varepsilon^4\).
\end{theorem}

\begin{proof}
\(\operatorname{Im}(\bar w_1u_1)=-\tfrac12\sin2z\) and (21.4); the
static solution and the stresses are certified.
\end{proof}

\subsection{Colour self-energy of the pair}

\begin{theorem}[Frequency shift of the fermion pair]
\label{thm:v20v19-shift}
The gauge source of the third-order fermion equations (21.1) with
\(f=\varepsilon^2f_2\), written in the form \(i\partial_t(u,w)=H(u,w)+\Sigma\),
has the wavenumber-one part

\[
 \Sigma^{(1)}=\sigma\,(\cos z,\,i\sin z)\,e^{-it},\qquad \sigma=-\frac{e^2}{16}\varepsilon^2,
                                                                    \tag{22.2}
\]

exactly along the mode, so the resonant response is
\(\psi_3=-i\sigma t\,(\cos z,0,i\sin z,0)e^{-it}\): the fermion
frequency becomes \(1+\sigma=1-e^2\varepsilon^2/16\) and the amplitude
does not grow.  The wavenumber-three part is the bounded response
\(e^{-it}(\tfrac{e^2}{64}\cos3z,0,-i\tfrac{e^2}{64}\sin3z,0)\varepsilon^3\).
\end{theorem}

\begin{proof}
The sources are \(S_u=\tfrac12ef_2w_2=\tfrac{e^2}8\sin2z\sin z\,e^{-it}\)
and \(S_w=-\tfrac12ef_2u_2=-\tfrac{ie^2}8\sin2z\cos z\,e^{-it}\) on the
paired subspace, giving \(\Sigma=(-S_u,S_w)\); its wavenumber-one
components are \(-\tfrac{e^2}{16}(\cos z,i\sin z)e^{-it}\), certified;
the secular solution is certified to satisfy
\(i\partial_t\psi_3-H\psi_3=\Sigma^{(1)}\), and the wavenumber-three
coefficients are solved exactly.
\end{proof}

\begin{corollary}[A mass from the coupling]
\label{cor:v20v19-mass}
By the massive dispersion law \(\omega^2=1+m^2\) of V17, the shift
(22.2) is that of an effective mass
\(m_{\rm eff}^2=2\sigma=-e^2\varepsilon^2/8\), negative: the colour
condensate of the pair lowers the pair's frequency below the
massless value, the opposite sign to the colour-induced masses of
the gauge waves themselves (V1), where the other colour's field
raised the frequency.
\end{corollary}

\begin{proof}
\(\omega=1+\sigma\) with \(\omega^2\approx1+2\sigma\); the comparison
is with Theorem \ref{thm:v20v1-mass}.
\end{proof}

\subsection{The uniform anisotropy}

\begin{theorem}[Fourth-order gravitational polarization]
\label{thm:v20v19-anisotropy}
The fourth-order anisotropic stress of the system is

\[
 \pi_4=\tfrac12\bigl(e^{-2a-2b}T_{xx}-e^{-2a+2b}T_{yy}\bigr)_4
 =ef_2\operatorname{Im}(\bar w_1u_1)+\pi_f=-\frac{e^2}8\sin^22z-\frac{e^2}8\cos^22z=-\frac{e^2}8,
                                                                    \tag{22.3}
\]

exactly homogeneous, so the polarization equation \(b_{tt}-b_{zz}+\pi=0\)
gives \(b_4=\tfrac{e^2}{16}t^2\) with no static \(z\)-dependent part:
the fermion--colour system drives a uniformly growing anisotropy of
the transverse metric, the \(x\) direction (the colour direction)
expanding relative to \(y\).
\end{theorem}

\begin{proof}
The fermion part of (21.4) with (22.1) and the colour part
\(\pi_f\) of Theorem \ref{thm:v20v19-condensate} are certified to sum
to \(-e^2/8\); the zero-mean static part of \(b_4\) is certified to
vanish.
\end{proof}

\begin{proposition}[Reading of the fermion--gauge coupling]
\label{prop:v20v19-reading}
The emergent geometry reads the coupling of a fermion pair to the
colour field at fourth order as a uniform anisotropic stress
\(-e^2\varepsilon^4/8\) with no wavenumber content, distinct from every
bosonic signature of this cycle, which carried modulations at twice
the wavenumber; the fermion itself reads the coupling as a negative
effective mass squared \(-e^2\varepsilon^2/8\).  Both effects are exact
consequences of the colour condensate (22.1).
\end{proposition}

\begin{proof}
Theorems \ref{thm:v20v19-shift} and \ref{thm:v20v19-anisotropy}.
\end{proof}

\subsection{Executable certificate}

The verifier

\[
 \texttt{scripts/verify\_sm\_v20\_v19\_colour\_condensate.py}         \tag{22.4}
\]

certifies the current, the condensate and its stresses, the resonant
form of the third-order source, the secular solution and its
frequency shift, the wavenumber-three response, and (22.3).  It emits

\[
 \texttt{artifacts/sm\_v20\_v19\_colour\_condensate.json}.            \tag{22.5}
\]

The hashes are

\[
\begin{array}{c|l}
\hbox{object}&\hbox{SHA--256}\\ \hline
\hbox{verifier}&
\texttt{97EC2CCFA0E68D3D4282314C6F105E9543D171C68B3B58DE11C01FE202CD807B}\\
\hbox{canonical payload}&
\texttt{641EFC26D3C59C9BD6566D682161842B1ED794F3AD985148F64E153BF3B2F110}.
\end{array}                                                     \tag{22.6}
\]

\begin{firewall}[Flat background, lowest orders, classical fields]
\label{fire:v20v19}
The condensate, the frequency shift, and the anisotropy are computed
on the flat background to the lowest order at which each appears;
the fermion's null-dust metric (V16) and its redshift are not
combined with them, the colour field's back-reaction on the metric at
fourth order is only its anisotropic part, and the fields are
classical.
\end{firewall}

\begin{decision}[V20 acquisition]
\label{dec:v20v19-next}
V20 is the compulsory companion audit and the second ten-version
sweep of the cycle, closing the Higgs and fermion chapters and fixing
the direction of V21--V29.
\end{decision}

\begin{assessment}[V19 acquisition]
\label{ass:v20v19}
The first dynamics of the exact fermion--gauge system is computed:
a colour condensate at twice the fermion wavenumber, the colour
self-energy that lowers the pair's frequency by \(e^2\varepsilon^2/16\),
and a uniform fourth-order anisotropic stress \(-e^2/8\) that drives
the gravitational polarization as \(\tfrac{e^2}{16}t^2\).
\end{assessment}

\begin{record}[V19 append record]
\label{rec:v20v19}
V19 is appended to the validated V18 whose installed SHA--256 was
\texttt{181374F97996126692E92958461E6E797F294619A2962AE837198D5C678B7160}.
After installation only V19 is the active numeric SM note; frozen
archives and unrelated notes are unchanged.
\end{record}


\section{V20: companion audit and the second ten-version sweep of the twentieth cycle}
\label{sec:v20v20}

This is the required audit at the twentieth version and the sweep of
V10--V19.  Every active companion note present in the shared folder
was inspected read-only.

\subsection{Companion snapshot}

\begin{record}[V20 companion snapshot]
\label{rec:v20v20-snapshot}
The files present were

\[
\begin{array}{c|r|r}
\hbox{file and SHA--256}&\hbox{bytes}&\hbox{lines}\\ \hline
\texttt{Renewal\_NS\_Stability\_Research\_Notes\_V26.tex}
&278283&5319\\
\multicolumn{3}{l}{
\texttt{82C887F8C2C69E4F28FAD7C3DC8DE5B9099CB5A4BF431EBA1FFF3E91935978AD}}
\\
\texttt{Renewal\_GRH\_Cofinal\_Visibility\_Attack\_V67.tex}
&640849&18652\\
\multicolumn{3}{l}{
\texttt{4A050D2FE5D6B29E11A878BD2EBC071A1FB1227F2636E98D7844D86C5F80AA58}}
\\
\texttt{Renewal\_Yang\_Mills\_Mass\_Gap\_Research\_Notes\_V38.tex}
&572590&13875\\
\multicolumn{3}{l}{
\texttt{CDC39872DC193FADBBB04B923A0370825E8EE9C1A3B08E498B2D4516F67801DE}}
\\
\texttt{Renewal\_YM\_Parallel\_Research\_Notes\_V5.tex}
&54174&1351\\
\multicolumn{3}{l}{
\texttt{306F1BB84CFA8A8D47EA569EC17E9545F9867C8D32B548EC9D9C444B4F3C0512}}.
\end{array}                                                     \tag{23.1}
\]

Since V15 the GRH note advanced from V62 to V67 (the corrected
real-axis criterion with a Gaussian representation of the Appell
polynomial, a signed-weight moment representation, its Fourier
symbol, and the equal-height and center-dependent symbols as
products of shifted parabolic operators) and the Yang--Mills mass-gap
note from V36 to V38 (the block-average rule as the upstream
Gaussian tower, harmonic extensions, the shell decomposition of the
free propagator, and the level-\(j\) shell functional); the NS and YM
parallel notes are unchanged.
\end{record}

\begin{decision}[V20 audit decision]
\label{dec:v20v20-audit}
No companion theorem, numerical constant, or executable artifact is
imported.  The GRH advances concern representations of a real-axis
positivity criterion for Appell polynomials; the Yang--Mills advances
concern a Euclidean Gaussian blocking tower.  Neither has an object
in common with the Lorentzian Einstein--matter systems of this cycle.
\end{decision}

\begin{proof}
The changed files are the GRH and Yang--Mills notes; their new
sections contain no constrained datum, no plane-symmetric field
system, and no averaged mode dynamics.
\end{proof}

\subsection{Sweep of V10--V19}

\begin{proposition}[What V10--V19 established]
\label{prop:v20v20-sweep}
The second block of the cycle placed the remaining Standard Model
structures in the exact plane-symmetric framework:

\begin{enumerate}
\item the Higgs sector: the exact Einstein--Yang--Mills--Higgs system
with \(m_A^2=e^2v^2/4\), \(m_h^2=2\lambda v^2\); the massive-scalar
law for massive gauge waves, so that the geometry reads the vacuum
expectation value through \(m_A\); the gauge wave's lowering of the
local vacuum expectation value (V11);
\item the averaged Higgs--gauge interaction: a negative-definite
quartic Hamiltonian in flat space and a gravitational correction that
breaks the Hamiltonian form of the two-mode reduction by the cross
asymmetry, exactly \(-(\omega^2+2)/(32\omega)\) at equal masses (V12,
V14);
\item the exact separation at fixed masses of the matter part,
proportional to \(e^2\), from the coupling-independent gravitational
part, and the mass laws of the latter as rational functions of the
frequency verified out of sample (V13, V14);
\item the fermionic sector: certified Dirac conventions on the
plane-symmetric metric; a spin-polarized Dirac wave keeps the
diagonal ansatz with \(b=0\), has no transverse pressure and sources
no polarization; the massless polarized standing wave is exact null
dust and reduces the Einstein--Dirac system without approximation to
the Bianchi--I null-dust cosmology (V16);
\item the massive polarized wave: static inhomogeneous energy
density, homogeneous longitudinal pressure, the universal equation of
state \(k^2/\omega^2\), and a closed-form second-order metric certified
against the full equations for every mass (V17);
\item the fermion--gauge coupling: a spin--isospin-locked doublet is
mapped into itself by the colour wave; on the invariant paired
subspace \(\psi_2=iP\psi_1\) the off-diagonal stress vanishes
identically and the Einstein--Yang--Mills--Dirac system is exactly
consistent, the fermion pair sourcing the colour wave and the
gravitational polarization through \(\operatorname{Im}(\bar w_1u_1)\)
(V18);
\item the colour condensate \(\tfrac e4\sin2z\) of a fermion pair, its
self-energy lowering the pair's frequency by \(e^2\varepsilon^2/16\),
and the uniform fourth-order anisotropic stress \(-e^2/8\) driving the
polarization as \(\tfrac{e^2}{16}t^2\) (V19).
\end{enumerate}
\end{proposition}

\begin{proof}
Items are the certified theorems of the cited sections.
\end{proof}

\begin{proposition}[Direction for V21--V29]
\label{prop:v20v20-direction}
The framework now carries, exactly and in plane symmetry, the
gravitational, gauge, Higgs, and fermionic fields, each with its
first effects computed and each read by the emergent geometry
through an exact law: masses through the equation of state, the vacuum
expectation value through the gauge mass, the gauge coupling through
the colour-magnetic flux and the colour levels, the relative phases
and spins through polarization sources and exchange rates.  What is
missing to make the plane-symmetric system a toy Standard Model is
the Yukawa coupling of the fermion doublet to the Higgs doublet, which
gives the fermion its mass from the vacuum expectation value and
couples the fermion wave to the Higgs wave.  V21--V23 add the Yukawa
term to the locked doublet, derive the plane-symmetric
Einstein--Yang--Mills--Higgs--Dirac system, certify its consistency on
the paired subspace, and compute its first effects; V24 assembles the
full system; V25 is the audit; V26--V29 return to the quantum reading
of the combined interactions in the number basis and to the
non-integrable three-field averaged dynamics, before the V30 sweep.
\end{proposition}

\begin{proof}
This is a decision.
\end{proof}

\begin{firewall}[Audit and sweep only]
\label{fire:v20v20}
No computation is added in V20.  The cycle's results are exact
statements about plane-symmetric classical fields, their first-order
averaged dynamics, and finite-order secular coefficients; no
continuum limit, no quantum field theory of any sector, and no
Standard Model--Einstein continuum are claimed.
\end{firewall}

\begin{decision}[V21 acquisition]
\label{dec:v20v20-next}
V21 derives the Yukawa coupling \(-y(\bar\Psi\Phi\psi_R+\hbox{h.c.})\)
in the plane-symmetric framework for the locked doublet and the
unitary Higgs, identifies the fermion mass \(m_f=yv/\sqrt2\) in the
gauged Dirac equations of V18, certifies the consistency of the
paired subspace with the Yukawa term, and computes the first
fermion--Higgs effect: the shift of the fermion frequency by the Higgs
wave.
\end{decision}

\begin{assessment}[V20 acquisition]
\label{ass:v20v20}
The audit is current and the sweep records the completion of the
matter sectors in the exact plane-symmetric framework: Higgs, massive
gauge, and fermion fields, their equations of state, their averaged
interactions with each other and with gravity, and the first
fermion--gauge dynamics.  The direction is the Yukawa coupling and
the assembly of the toy Standard Model, then its quantum reading.
\end{assessment}

\begin{record}[V20 append record]
\label{rec:v20v20}
V20 is appended to the validated V19 whose installed SHA--256 was
\texttt{C3DADE1E3CFDAE19D8AF21BC606138F1684A409357B63DBDC2386BF3F7B32368}.
After installation only V20 is the active numeric SM note; the
companions, frozen archives, and all unrelated notes are unchanged.
The next compulsory companion audit is V25 and the next sweep V30.
\end{record}


\section{V21: the Yukawa coupling of the locked doublet, the fermion mass from the vacuum, and the Yukawa self-energy of a fermion pair}
\label{sec:v20v21}

V18--V19 coupled the locked fermion doublet to the colour wave; V11
placed the unitary Higgs \(\Phi=(0,v+h)/\sqrt2\) in the framework.
This note adds the Yukawa coupling.  For the doublet to remain locked
and paired the two isospin components must receive the same mass, so
the coupling is the isospin-degenerate one,
\(-y(\bar\Psi\Phi\psi_R+\bar\Psi\tilde\Phi\psi_R+\hbox{h.c.})\) with
equal up and down couplings, which in unitary gauge is the Dirac mass
term \(-\tfrac{y}{\sqrt2}(v+h)\bar\Psi\Psi\).  The fermion mass of
V17 is thereby identified as \(m=yv/\sqrt2\), the vacuum expectation
value read through the Yukawa coupling, and the Higgs wave modulates
it as \(m_f=m(1+h/v)\).  Three exact statements follow.  The paired
subspace of V18 is invariant under the gauged Dirac equations with an
arbitrary field-dependent mass \(m_f(z,t)\), so the exactly consistent
system of V18 survives the Yukawa term.  The Higgs equation acquires
the source \(-\tfrac{y}{\sqrt2}\bar\Psi\Psi\), which on the paired
subspace is \(-\tfrac{2m}v(|u_1|^2-|w_1|^2)\).  For the massive
standing wave of V17 this source is static, the Higgs is displaced
statically at wavenumbers zero and two, the pair lowers the local
vacuum expectation value, and the induced mass modulation acts back
on the pair resonantly: its frequency is shifted by the Yukawa
self-energy \(\sigma_Y\varepsilon^2\) with \(\sigma_Y\) in closed form,
negative for every mass, so the Yukawa self-interaction, like the
colour self-energy of V19, lowers the pair's frequency.

\subsection{The Yukawa term and the paired subspace}

\begin{proposition}[Yukawa coupling of the locked doublet]
\label{prop:v20v21-yukawa}
With the unitary Higgs of V11 the isospin-degenerate Yukawa coupling is
\(\mathcal L_Y=-\tfrac{y}{\sqrt2}(v+h)\bar\Psi\Psi\).  It enters the
gauged Dirac equations (21.1) of V18 through the replacement
\(m\to m_f=\tfrac{y}{\sqrt2}(v+h)\) and the Higgs equation of V11
through the source \(-\tfrac y{\sqrt2}\bar\Psi\Psi\):

\[
 h_{tt}-h_{zz}+\lambda(v+h)\bigl((v+h)^2-v^2\bigr)+\frac y{\sqrt2}\bar\Psi\Psi=0
                                                                    \tag{24.1}
\]

on the flat background.  The fermion mass of V17 is \(m=yv/\sqrt2\).
\end{proposition}

\begin{proof}
The Higgs Euler--Lagrange equation with the external bilinear is
certified from the Lagrangian; the Dirac equation with the
field-dependent mass is the Euler--Lagrange equation of \(\bar\Psi\).
\end{proof}

\begin{theorem}[The paired subspace survives the Yukawa term]
\label{thm:v20v21-paired}
For an arbitrary real function \(m_f(z,t)\) in place of \(m\), the
gauged Dirac equations of the locked doublet on the paired subspace
\(\psi_2=iP\psi_1\) satisfy \(D_2=(iD_{1,0},-iD_{1,2})\): the subspace
is invariant.  On it the scalar bilinear is

\[
 \bar\Psi\Psi=2\bigl(|u_1|^2-|w_1|^2\bigr),                           \tag{24.2}
\]

so the exactly consistent Einstein--Yang--Mills--Dirac system of
Theorem \ref{thm:v20v18-paired} extends to the
Einstein--Yang--Mills--Higgs--Dirac system with the fermion sourcing
the Higgs through (24.2), the colour wave and the polarization through
\(\operatorname{Im}(\bar w_1u_1)\), and the Higgs acting on the fermion
only through its mass.
\end{theorem}

\begin{proof}
The operator is rebuilt from the gamma matrices of V16 with the gauge
term of V18 and the mass \(m_f(z,t)\); the locking, the form of the
sector-plus equations, the invariance identity, and (24.2) are
certified symbolically.  The metric terms of (21.1) are identical for
the two isospin components and were certified in V18, so the identity
holds with the metric present.
\end{proof}

\subsection{Static Higgs displacement by a massive fermion pair}

\begin{theorem}[Higgs displacement]
\label{thm:v20v21-displacement}
For the massive standing wave of V17 on the paired subspace,
\(\psi_1=\varepsilon e^{-i\omega t}(\cos z,0,iW\sin z,0)\),
\(\omega=\sqrt{1+m^2}\), \(W=\omega-m\), the bilinear is static,

\[
 \bar\Psi\Psi=2\varepsilon^2\bigl(\cos^2z-W^2\sin^2z\bigr)
 =2\varepsilon^2W\bigl(m+\omega\cos2z\bigr),                          \tag{24.3}
\]

by the identities \(1-W^2=2mW\), \(1+W^2=2\omega W\), and the
second-order Higgs response is the static

\[
 h_2=-\frac{2mW}{v}\Bigl[\frac m{m_h^2}+\frac{\omega\cos2z}{4+m_h^2}\Bigr],
                                                                    \tag{24.4}
\]

with mean \(\langle h_2\rangle=-2m^2W/(vm_h^2)<0\): the fermion pair
lowers the local vacuum expectation value, and with it its own mass by
\(\langle\delta m\rangle=-2m^3W\varepsilon^2/(v^2m_h^2)\).
\end{theorem}

\begin{proof}
The Fourier coefficients, the identities, the solution of
\(-h_2''+m_h^2h_2=-\tfrac mv\bar\Psi\Psi\), and the mean are certified.
\end{proof}

\subsection{Yukawa self-energy}

\begin{theorem}[Frequency shift by the induced Higgs field]
\label{thm:v20v21-shift}
The mass modulation \(\delta m=\tfrac mvh_2\) enters the free massive
system \(i\partial_t\psi=H_0\psi\), \(H_0(u,w)=(-iw_z+mu,-iu_z-mw)\), as
the perturbation \(\delta m\operatorname{diag}(1,-1)\).  \(H_0\) is
Hermitian for the periodic inner product, and the mode-parallel part
of the perturbation is the resonant source
\(\sigma_Y\psi_0\) with

\[
 \sigma_Y=\frac{\langle\psi_0|\delta m\operatorname{diag}(1,-1)|\psi_0\rangle}{\langle\psi_0|\psi_0\rangle}
 =-\frac{m^2}{v^2}\,W\Bigl[\frac{2m^2}{\omega m_h^2}+\frac{\omega}{4+m_h^2}\Bigr],
                                                                    \tag{24.5}
\]

whose secular solution \(\psi_3=-i\sigma_Yt\,\psi_0\) shifts the
frequency to \(\omega+\sigma_Y\varepsilon^2\) with no amplitude growth.
\(\sigma_Y<0\) for every \(m>0\): the Yukawa self-energy of the pair
lowers its frequency, as the colour self-energy \(-e^2/16\) of V19
did.  At \(m=1\), \(m_h=2\), \(v=1\) it equals
\(-\tfrac34+\tfrac{3\sqrt2}8\).
\end{theorem}

\begin{proof}
The Hermiticity of \(H_0\), the free solution \(\psi_0\), the closed
form (24.5), and the secular solution
\(i\partial_t\psi_3-H_0\psi_3=\sigma_Y\psi_0\) are certified.  The sign
follows from \(W>0\).
\end{proof}

\begin{corollary}[Reading of the Yukawa coupling]
\label{cor:v20v21-reading}
The emergent framework reads the Yukawa coupling twice: through the
fermion mass \(m=yv/\sqrt2\), which the geometry sees as the equation
of state \(1/\omega^2\) of V17, and through the self-energy (24.5),
quadratic in \(y\) at fixed \(m_h\), which the fermion sees as a
lowering of its frequency and the Higgs sees as a static depression of
the vacuum expectation value.  In the massless limit \(m\to0\) at fixed
\(m/v\) the self-energy tends to \(-(m^2/v^2)/(4+m_h^2)\), entirely from
the wavenumber-two Higgs response; for \(m\to\infty\) it grows as
\(-(m^2/v^2)\bigl[m_h^{-2}+\tfrac12(4+m_h^2)^{-1}\bigr]\).
\end{corollary}

\begin{proof}
Insert \(W\to1\), \(\omega\to1\) and \(W\sim1/(2m)\), \(\omega\sim m\)
into (24.5).
\end{proof}

\subsection{Executable certificate}

The verifier

\[
 \texttt{scripts/verify\_sm\_v20\_v21\_yukawa.py}                     \tag{24.6}
\]

certifies Proposition \ref{prop:v20v21-yukawa}, Theorem
\ref{thm:v20v21-paired}, the closed forms (24.3)--(24.5), the
Hermiticity of \(H_0\), and the secular solution.  It emits

\[
 \texttt{artifacts/sm\_v20\_v21\_yukawa.json}.                        \tag{24.7}
\]

The hashes are

\[
\begin{array}{c|l}
\hbox{object}&\hbox{SHA--256}\\ \hline
\hbox{verifier}&
\texttt{B75891227CB41AAC32DCB8A8D746827F74BD0C4FE58134B480751D8D819EF70A}\\
\hbox{canonical payload}&
\texttt{D93FD8DA5E3AFC5DCD9FB7DD922906A49320E40629CAE5FED429024EB16C426D}.
\end{array}                                                     \tag{24.8}
\]

\begin{firewall}[Degenerate Yukawa, flat background, lowest order]
\label{fire:v20v21}
Only the isospin-degenerate coupling \(y_u=y_d\) is treated; an
isospin-split mass breaks the pairing identity and its consistency with
the diagonal ansatz is not established here.  The Higgs displacement
and the self-energy are computed on the flat background at the lowest
order at which each appears, without the colour wave, without the
fermion's null-dust metric, and with classical fields.
\end{firewall}

\begin{decision}[V22 acquisition]
\label{dec:v20v21-next}
V22 examines the isospin-split Yukawa coupling \(y_u\neq y_d\): the
mass-split locked doublet, the fate of the paired subspace, the
off-diagonal stress it generates, and whether a modified pairing
restores the diagonal ansatz; then V23 computes the fermion--Higgs
exchange dynamics (a Higgs wave and a fermion pair, averaged), before
the V24 assembly.
\end{decision}

\begin{assessment}[V21 acquisition]
\label{ass:v20v21}
The Yukawa coupling is placed in the exact framework: the fermion
mass is the vacuum expectation value read through \(y\), the paired
locked doublet is invariant with a field-dependent mass, the Higgs
equation acquires the fermion source (24.2), and the first
fermion--Higgs effect is computed exactly: a static depression of the
vacuum expectation value by a fermion pair and the negative Yukawa
self-energy (24.5).
\end{assessment}

\begin{record}[V21 append record]
\label{rec:v20v21}
V21 is appended to the validated V20 whose installed SHA--256 was
\texttt{840CB99179E4E8C3B2EDCD9D735187D3599556A2FF8E4F67BB50488976DD57F1}.
After installation only V21 is the active numeric SM note; frozen
archives and unrelated notes are unchanged.
\end{record}


\section{V22: the mass-split doublet, the cross polarization of the gravitational wave, and beat-frequency radiation}
\label{sec:v20v22}

V21 required the isospin-degenerate Yukawa coupling for the paired
subspace to survive.  This note treats the split coupling
\(y_u\neq y_d\), that is, masses \(m_1\neq m_2\) for the two
components of the locked doublet.  The pairing identity of V18 then
fails by a defect proportional to \(m_2-m_1\), so the off-diagonal
stress \(T_{xy}=-\tfrac12efS\) no longer vanishes.  Without the colour
wave nothing is lost: the two components are independent massive
polarized waves of V17 and the diagonal ansatz holds.  With the colour
wave present the split standing waves beat: the bilinears
\(S\) and the colour current oscillate at the beat frequency
\(\Delta=\omega_2-\omega_1\), so the fermion pair drives the second,
cross polarization of the plane gravitational wave, absent from every
diagonal system of this series, and drives colour radiation at
wavenumber two.  The linearized cross polarization \(d=g_{xy}\) is
certified to decouple from the diagonal metric at linear order and to
obey \(d_{tt}-d_{zz}=2T_{xy}\).  Both radiations are resonant at the
mass splitting \(\omega_2-\omega_1=2\), where the cross-polarized
gravitational wave and the colour wave grow linearly in time with
certified rates; the cross polarization is also resonant at
\(\omega_2-\omega_1=4\).

\subsection{The defect and the free split doublet}

\begin{theorem}[Pairing defect and the stress of a locked doublet]
\label{thm:v20v22-defect}
With masses \(m_1,m_2\) in the sector-plus and sector-minus equations
of V18 on the flat background, the pairing \(u_2=iu_1\), \(w_2=-iw_1\)
gives

\[
 D_{2,1}-iD_{1,0}=i(m_1-m_2)u_1,\qquad D_{2,3}+iD_{1,2}=i(m_2-m_1)w_1,
                                                                    \tag{25.1}
\]

so the paired subspace is invariant if and only if \(m_1=m_2\).  For an
arbitrary locked doublet the flat gauge-covariant stress has
\(T_{tx}=T_{ty}=T_{xz}=T_{yz}=T_{yy}=0\) and

\[
 T_{xy}=-\tfrac12ef\,S,\qquad S=\operatorname{Im}(u_2\bar w_1+w_2\bar u_1).
                                                                    \tag{25.2}
\]

Consequently at \(f=0\) the split doublet is two independent massive
polarized waves, each obeying the laws of V17 with its own mass, and
the diagonal ansatz holds.
\end{theorem}

\begin{proof}
The operator is built from the gamma matrices of V16 with the gauge
term of V18; (25.1) and (25.2) are certified symbolically.
\end{proof}

\subsection{Beating bilinears}

\begin{theorem}[Split standing waves]
\label{thm:v20v22-beat}
Let \(\omega_i=\sqrt{1+m_i^2}\), \(W_i=\omega_i-m_i\), and
\(\psi_1=e^{-i\omega_1t}(\cos z,iW_1\sin z)\),
\(\psi_2=e^{-i\omega_2t}(i\cos z,W_2\sin z)\), each certified to solve
its free massive equations.  Then, with \(\Delta=\omega_2-\omega_1\),

\[
 S=-\frac{W_1+W_2}2\sin2z\sin\Delta t,\qquad
 J=e\,\frac{W_1+W_2}2\sin2z\cos\Delta t,\qquad
 \pi=\tfrac12T_{xx}=-\frac{ef}4(W_1+W_2)\sin2z\cos\Delta t,
                                                                    \tag{25.3}
\]

reducing at \(m_1=m_2=0\) to \(J=e\sin2z\) of V19.  The mass-split pair
therefore carries a colour current, an anisotropic stress, and an
off-diagonal stress all at wavenumber two and beat frequency
\(\Delta\); the off-diagonal one vanishes identically only at
\(\Delta=0\).
\end{theorem}

\begin{proof}
Certified symbolically from (25.2), the current of V18, and the flat
stress.
\end{proof}

\subsection{The cross polarization}

\begin{theorem}[Linearized cross polarization]
\label{thm:v20v22-cross}
For the metric \(-dt^2+dx^2+dy^2+2d(z,t)\,dx\,dy+dz^2\) the Einstein
tensor has \(G_{xy}=\tfrac12(d_{tt}-d_{zz})+O(d^2)\) and every other
component \(O(d^2)\).  Hence the cross polarization sourced by a
locked doublet in a colour field obeys, at linear order,

\[
 d_{tt}-d_{zz}=2T_{xy}=-efS.                                          \tag{25.4}
\]
\end{theorem}

\begin{proof}
The Einstein tensor is computed exactly and expanded to first order in
\(d\); the vanishing of the other components and the form of
\(G_{xy}\) are certified.
\end{proof}

\begin{theorem}[Resonant beat-frequency radiation]
\label{thm:v20v22-resonance}
With the colour standing wave \(f=\varepsilon_c\cos z\cos t\) of V1 and
fermion amplitude \(\varepsilon\), the source of (25.4) is

\[
 A_0(\sin z+\sin3z)\bigl(\sin(\Delta+1)t+\sin(\Delta-1)t\bigr),\qquad
 A_0=\frac{e\varepsilon_c\varepsilon^2(W_1+W_2)}8.                     \tag{25.5}
\]

For generic \(\Delta\) the response is bounded,
\(d=A_0\sum_{k=1,3}\sum_{\Omega=\Delta\pm1}\sin kz\sin\Omega t/(k^2-\Omega^2)\).
At \(\Delta=2\) both wavenumbers are resonant and the secular part is

\[
 d=-\frac{A_0}2\,t\sin z\cos t-\frac{A_0}6\,t\sin3z\cos3t,             \tag{25.6}
\]

while at \(\Delta=4\) only \(-\tfrac{A_0}6t\sin3z\cos3t\) is secular.
The colour current of (25.3) sources \(f_{tt}-f_{zz}=J\) with the
bounded response \(A_J\sin2z\cos\Delta t/(4-\Delta^2)\),
\(A_J=e\varepsilon^2(W_1+W_2)/2\), and at \(\Delta=2\) the secular
\(\tfrac{A_J}4t\sin2z\sin2t\).  The resonant splitting is
\(m_2^2=m_1^2+4\omega_1+4\), for instance \(m_1=0\), \(m_2=2\sqrt2\).
\end{theorem}

\begin{proof}
The decomposition (25.5), the generic particular solutions, the
secular solutions at \(\Delta=2\) and \(\Delta=4\), and the splitting
condition are certified.
\end{proof}

\begin{proposition}[Reading of the mass splitting]
\label{prop:v20v22-reading}
The emergent geometry distinguishes a mass-degenerate from a
mass-split doublet through the polarization content of the
gravitational wave: the degenerate pair drives only the plus
polarization \(b\) and the colour condensate is static (V19), while
the split pair in a colour field drives the cross polarization \(d\)
and radiates colour at the beat frequency \(\omega_2-\omega_1\), with
linear growth exactly at the splittings \(2\) and \(4\) in units of the
wavenumber.  The mass difference of the doublet is thereby readable
from the frequency of the cross-polarized wave.
\end{proposition}

\begin{proof}
Theorems \ref{thm:v20v22-beat}--\ref{thm:v20v22-resonance} together
with Theorem \ref{thm:v20v18-paired}.
\end{proof}

\subsection{Executable certificate}

The verifier

\[
 \texttt{scripts/verify\_sm\_v20\_v22\_split\_doublet.py}             \tag{25.7}
\]

certifies Theorems \ref{thm:v20v22-defect}--\ref{thm:v20v22-resonance}.
It emits

\[
 \texttt{artifacts/sm\_v20\_v22\_split\_doublet.json}.                \tag{25.8}
\]

The hashes are

\[
\begin{array}{c|l}
\hbox{object}&\hbox{SHA--256}\\ \hline
\hbox{verifier}&
\texttt{6BA5316CD8E1DCBD46E4F8AA041A6AA823926F44FF392E5ACA85467E29C11EEE}\\
\hbox{canonical payload}&
\texttt{6C983DB3412E1190F3F553C2D6943E75C5E21726267089D2BF333533DF61241E}.
\end{array}                                                     \tag{25.9}
\]

\begin{firewall}[Linear cross polarization, external colour wave, flat background]
\label{fire:v20v22}
The cross polarization is treated at linear order on the flat
background and the colour wave is external and unperturbed by the
fermions; the nonlinear two-polarization plane metric, the
back-reaction of the radiated colour wave on the pair, and the
fermions' own null-dust metric are not combined.  The Higgs field
enters only through the masses.
\end{firewall}

\begin{decision}[V23 acquisition]
\label{dec:v20v22-next}
V23 computes the fermion--Higgs exchange dynamics: a Higgs standing
wave and a degenerate fermion pair, the resonant projections of the
mass modulation by the Higgs wave onto the pair, the pair's source
in the Higgs equation, and the averaged two-mode dynamics with its
conserved quantities, completing the first effects of the Yukawa
coupling before the V24 assembly of the full plane system.
\end{decision}

\begin{assessment}[V22 acquisition]
\label{ass:v20v22}
The mass-split doublet is placed in the framework: the pairing
defect (25.1), the universal stress (25.2), the beating bilinears
(25.3), the linearized cross polarization law (25.4), and the resonant
beat-frequency radiation of cross-polarized gravitational waves and
colour at the splittings \(2\) and \(4\), with certified secular rates.
\end{assessment}

\begin{record}[V22 append record]
\label{rec:v20v22}
V22 is appended to the validated V21 whose installed SHA--256 was
\texttt{00A6756579FD884034339D93BD41159EC12065B6FAD0216E6A21206C02291CD5}.
After installation only V22 is the active numeric SM note; frozen
archives and unrelated notes are unchanged.
\end{record}


\section{V23: fermion--Higgs exchange: an exact two-mode reduction, its invariants, and the classical Higgs--pair resonance}
\label{sec:v20v23}

V21 computed the static Higgs displacement by a fermion pair and the
Yukawa self-energy.  This note computes the dynamical exchange
between the pair and a homogeneous Higgs oscillation \(h(t)\), the
zero-wavenumber Higgs mode.  The key structural fact is that a
homogeneous mass modulation keeps the two-dimensional wavenumber-one
subspace \(u\propto\cos z\), \(w\propto\sin z\) of the polarized Dirac
system exactly invariant; it is spanned by the pair mode
\(\psi_0\) of V17 at frequency \(\omega\) and its negative-frequency
partner \(\psi_-\) at \(-\omega\).  The fermion dynamics is therefore an
exact two-level system driven by the Higgs, and the zero-wavenumber
Higgs equation is sourced by the mean pair bilinear, which is a
quadratic form in the two amplitudes.  The coupled system conserves
the fermion norm and a total energy exactly.  The Higgs drives the
transition \(\psi_0\to\psi_-\) resonantly exactly when
\(m_h=2\omega=2\sqrt{1+m^2}\), the classical form of the threshold for
the Higgs to convert into a fermion--antifermion pair; the averaged
dynamics at resonance is the three-wave system with its Manley--Rowe
invariants, and with an undepleted Higgs the pair mode is converted
fully into its partner in the time \(\pi/(2gH_0)\) with
\(g=m/(2\omega v)=y/(2\sqrt2\,\omega)\).

\subsection{The exact two-mode reduction}

\begin{theorem}[Invariant subspace and projected dynamics]
\label{thm:v20v23-subspace}
Let \(\psi_0=(\cos z,iW\sin z)\) and \(\psi_-=(W\cos z,-i\sin z)\),
\(W=\omega-m\).  Both solve the free massive system of V17, with
frequencies \(\omega\) and \(-\omega\), are orthogonal, and have the
norm \(2\pi\omega W\).  Under the mass modulation
\(m\to m+\tfrac mvh(t)\) with \(h\) independent of \(z\), the subspace
\(\psi=a_0(t)\psi_0+a_-(t)\psi_-\) is exactly invariant and

\[
 i\dot a=\Bigl[\omega\sigma_3+\frac mvh(t)B_0\Bigr]a,\qquad
 B_0=\frac1\omega\begin{pmatrix}m&1\\1&-m\end{pmatrix}.               \tag{26.1}
\]
\end{theorem}

\begin{proof}
The free solutions, the inner products, the projections, and the
vanishing of the remainder of \(H\psi\) outside the subspace are
certified.
\end{proof}

\begin{theorem}[Mean bilinear and the coupled system]
\label{thm:v20v23-coupled}
On the subspace the mean of the pair bilinear is
\(\langle|u|^2-|w|^2\rangle_z=\omega W\,a^\dagger B_0a\), so with the
paired doublet of amplitude \(\varepsilon\) the zero-wavenumber Higgs
equation of (24.1) is

\[
 \ddot h+m_h^2h=-\frac{2m\varepsilon^2\omega W}v\,a^\dagger B_0a.       \tag{26.2}
\]

The system (26.1)--(26.2) conserves exactly

\[
 N=|a_0|^2+|a_-|^2,\qquad
 E=\tfrac12\dot h^2+\tfrac12m_h^2h^2+2\varepsilon^2\omega W\,a^\dagger\Bigl[\omega\sigma_3+\frac mvhB_0\Bigr]a.
                                                                    \tag{26.3}
\]
\end{theorem}

\begin{proof}
The mean, and the vanishing of \(\dot N\) and \(\dot E\) on the
equations of motion, are certified symbolically.
\end{proof}

\subsection{The resonance}

\begin{theorem}[Three-wave system at \(m_h=2\omega\)]
\label{thm:v20v23-resonance}
Write \(a_0=c_0e^{-i\omega t}\), \(a_-=c_-e^{i\omega t}\),
\(h=\tfrac12(He^{-im_ht}+\bar He^{im_ht})\).  The Higgs couples \(c_0\)
to \(c_-\) without net frequency exactly when \(m_h=2\omega\), and
keeping the resonant terms with \(H\) slowly varying gives

\[
 i\dot c_0=gHc_-,\qquad i\dot c_-=g\bar Hc_0,\qquad i\dot H=\kappa c_0\bar c_-,\qquad
 g=\frac m{2\omega v}=\frac y{2\sqrt2\,\omega},\quad \kappa=\frac{m\varepsilon^2W}{\omega v},
                                                                    \tag{26.4}
\]

with the exact invariants \(N\), \(|H|^2-2\varepsilon^2W|c_-|^2\), and
\(K=g(H\bar c_0c_-+\bar Hc_0\bar c_-)\).  For an undepleted real
Higgs amplitude \(H_0\) the solution is \(c_0=\cos gH_0t\),
\(c_-=-i\sin gH_0t\): the pair mode is converted entirely into its
negative-frequency partner in the time \(\pi/(2gH_0)\), and back.
\end{theorem}

\begin{proof}
The resonant parts of (26.1) and of the \(e^{-2i\omega t}\) component
of (26.2), the three invariants, the ratio \(\kappa/g=2\varepsilon^2W\),
and the Rabi solution are certified.
\end{proof}

\begin{proposition}[Reading of the exchange]
\label{prop:v20v23-reading}
The condition \(m_h=2\sqrt{1+m^2}>2m\) is the classical
plane-symmetric form of the threshold \(m_h>2m_f\) for the Higgs to
produce a fermion pair: the negative-frequency partner \(\psi_-\) is
the antiparticle mode, and the conversion \(\psi_0\leftrightarrow\psi_-\)
at fixed norm is the classical shadow of pair creation and
annihilation, stimulated only, since both the pure pair and the pure
Higgs are fixed points of (26.4).  The Manley--Rowe invariant is the
energy balance \(2\omega^2|H|^2+2\varepsilon^2\omega^2W(|c_0|^2-|c_-|^2)\):
each unit of norm transferred to the partner mode releases the energy
\(2\omega=m_h\) into the Higgs.  The rate \(g\) is linear in the Yukawa
coupling and inversely proportional to the fermion energy.
\end{proposition}

\begin{proof}
The threshold and the invariants are Theorem
\ref{thm:v20v23-resonance}; the fixed points are read off (26.4).
\end{proof}

\subsection{Executable certificate}

The verifier

\[
 \texttt{scripts/verify\_sm\_v20\_v23\_higgs\_fermion\_exchange.py}    \tag{26.5}
\]

certifies Theorems \ref{thm:v20v23-subspace}--\ref{thm:v20v23-resonance}.
It emits

\[
 \texttt{artifacts/sm\_v20\_v23\_higgs\_fermion\_exchange.json}.       \tag{26.6}
\]

The hashes are

\[
\begin{array}{c|l}
\hbox{object}&\hbox{SHA--256}\\ \hline
\hbox{verifier}&
\texttt{1C384895EE970110424E6746CA199BBF3551772968DABFFA971C1CB1531BDE85}\\
\hbox{canonical payload}&
\texttt{889BC183B24B950DF3F57CC777FC699193893E7EFC4140093C31B618B38B6D79}.
\end{array}                                                     \tag{26.7}
\]

\begin{firewall}[Zero-wavenumber Higgs, flat background, classical]
\label{fire:v20v23}
Only the homogeneous Higgs mode is retained; the wavenumber-two
Higgs response of V21 feeds back on the pair at the next order and is
dropped, as are the Higgs self-interaction, the gauge field, and
gravity.  The three-wave system is the rotating-wave truncation at
exact resonance; the exact system is (26.1)--(26.2).  The fields are
classical: the conversion is stimulated, and no spontaneous decay is
described.
\end{firewall}

\begin{decision}[V24 acquisition]
\label{dec:v20v23-next}
V24 assembles the full plane-symmetric toy Standard Model system:
the diagonal metric \((a,b,c)\) of V91, the colour wave, the unitary
Higgs, and the paired locked doublet with the Yukawa mass, with every
equation and source collected from V1, V11, V16, V18, and V21, and
certifies the consistency of the assembled system (Bianchi identity
on shell) as a single symbolic object.
\end{decision}

\begin{assessment}[V23 acquisition]
\label{ass:v20v23}
The fermion--Higgs exchange is reduced exactly to a two-level fermion
driven by the homogeneous Higgs, with the exact invariants (26.3); the
resonance at \(m_h=2\omega\) is the classical pair threshold, the
averaged dynamics is the three-wave system (26.4) with its
Manley--Rowe energy balance, and the conversion rate is
\(g=y/(2\sqrt2\,\omega)\).
\end{assessment}

\begin{record}[V23 append record]
\label{rec:v20v23}
V23 is appended to the validated V22 whose installed SHA--256 was
\texttt{050E3D7F0398933B5C62D0791A1612CA5752E78C690F4F5F99D18E4BE501CE07}.
After installation only V23 is the active numeric SM note; frozen
archives and unrelated notes are unchanged.
\end{record}


\section{V24: the assembled plane-symmetric toy Standard Model system, its exact consistency, and a corrigendum to V19}
\label{sec:v20v24}

This note assembles into a single symbolic object the fields placed
one by one in the exact framework: the diagonal metric \((a,b,c)\) of
V91, the colour wave \(A^1=f\,dx\) of V1, the unitary Higgs
\(\Phi=(0,v+h)/\sqrt2\) of V11, and the paired locked doublet
\(\psi_1=(u_1,0,w_1,0)\), \(\psi_2=iP\psi_1\) of V18 with the Yukawa
mass \(M=y(v+h)/\sqrt2\) of V21.  Every matter equation is derived
from the total Lagrangian, the total stress is shown to respect the
diagonal ansatz, and its covariant conservation is certified on
shell, for arbitrary metric functions, using the matter equations
only: the assembled Einstein--Yang--Mills--Higgs--Dirac system is
exactly consistent.  The five Einstein equations are written in the
\((a,b,c)\) form of (3.2) with the general sources, and the sources
are given explicitly.  In the course of the assembly the anisotropic
stress of the fermion--colour system of V19 was recomputed from the
total stress in the single convention of (3.1), and the value
\(-e^2/8\) of (22.3) is found to have combined its two parts in
opposite conventions; the corrected fourth-order polarization source
is \(-\tfrac{e^2}8\cos4z\), of mean zero, and the uniform growth
claimed in V19 is withdrawn.

\subsection{The assembled system}

\begin{theorem}[Matter equations]
\label{thm:v20v24-matter}
From the total Lagrangian on the metric
\(-dt^2+e^{2a+2b}dx^2+e^{2a-2b}dy^2+e^{2c}dz^2\), with
\(W_\pm=e^{-2a\pm2b}\), the matter equations of the assembled system
are

\[
\begin{aligned}
 &f_{tt}+f_t(c_t-2b_t)-e^{-2c}\bigl(f_{zz}-f_z(2b_z+c_z)\bigr)+\tfrac14e^2(v+h)^2f
 =e^{2a+2b}J,\qquad J=-2e\,e^{-a-b}\operatorname{Im}(\bar w_1u_1),\\
 &h_{tt}+h_t(2a_t+c_t)-e^{-2c}\bigl(h_{zz}+h_z(2a_z-c_z)\bigr)
 +\lambda(v+h)\bigl((v+h)^2-v^2\bigr)+\tfrac14e^2(v+h)f^2W_-
 +\sqrt2\,y\bigl(|u_1|^2-|w_1|^2\bigr)=0,
\end{aligned}                                                          \tag{27.1}
\]

together with the Dirac equations (21.1) with \(m\to M=y(v+h)/\sqrt2\)
and the partner components \(u_2=iu_1\), \(w_2=-iw_1\); the
sector-minus equations are \((iD_{1,0},-iD_{1,2})\) on the curved
metric with the field-dependent mass.
\end{theorem}

\begin{proof}
The Euler--Lagrange equations of \(f\) and \(h\) are computed from
the bosonic Lagrangian of V11 plus the algebraic \(f\)- and
\(h\)-dependence of the fermion Lagrangian, certified to be the
current coupling \(-2e\,e^{-a-b}f\operatorname{Im}(\bar w_1u_1)\) and the
Yukawa source \(-\tfrac y{\sqrt2}\bar\Psi\Psi\); the forms (27.1), the
locking, the Dirac form, and the pairing identity are certified.
\end{proof}

\begin{theorem}[Total stress, ansatz, conservation]
\label{thm:v20v24-stress}
Let \(T=T^{YM}+T^{\Phi}+T^{\Psi}\) with the stresses of V1, V11, and
V18.  On the paired subspace
\(T_{tx}=T_{ty}=T_{xy}=T_{xz}=T_{yz}=0\), so the diagonal ansatz is
exactly consistent, and \(\nabla^\mu T_{\mu\nu}=0\) for every \(\nu\)
on shell, using only the matter equations, for arbitrary \(a,b,c\).
Hence the Einstein equations \(G_{\mu\nu}+\Lambda g_{\mu\nu}=T_{\mu\nu}\)
are compatible with the matter dynamics by the Bianchi identity.
\end{theorem}

\begin{proof}
The vanishing of the five off-diagonal components and the on-shell
divergence (the fermion time derivatives from the Dirac equations,
\(f_{tt}\) and \(h_{tt}\) from (27.1)) are certified symbolically.
\end{proof}

\begin{theorem}[Einstein equations with general sources]
\label{thm:v20v24-einstein}
With \(\rho=T_{tt}\), \(q=T_{tz}\), \(p_z=e^{-2c}T_{zz}\),
\(p_\perp=\tfrac12(W_-T_{xx}+W_+T_{yy})\), and, in the convention of
(3.1), \(\pi=\tfrac12(W_+T_{yy}-W_-T_{xx})\), the Einstein equations are

\[
\begin{aligned}
 &a_t^2+2a_tc_t-b_t^2-e^{-2c}\bigl(3a_z^2-2a_zc_z+2a_{zz}+b_z^2\bigr)-\rho=\Lambda,\\
 &a_{tz}+a_ta_z-a_zc_t+b_tb_z+\tfrac12q=0,\\
 &b_{tt}+b_t(2a_t+c_t)-e^{-2c}\bigl(b_{zz}+b_z(2a_z-c_z)\bigr)+\pi=0,\\
 &a_{tt}+2a_t^2+a_tc_t-e^{-2c}\bigl(a_{zz}+2a_z^2-a_zc_z\bigr)=\Lambda+\tfrac12(\rho-p_z),\\
 &c_{tt}+c_t^2+2a_tc_t-e^{-2c}\bigl(2a_z^2-2a_zc_z+2a_{zz}+2b_z^2\bigr)=\Lambda+\tfrac12(\rho+p_z)-p_\perp,
\end{aligned}                                                          \tag{27.2}
\]

each being the stated combination of
\(\operatorname{Ric}-\Lambda g-T+\tfrac12g\operatorname{tr}T\), and on the
paired subspace the sources are

\[
\begin{aligned}
 \rho&=\tfrac12W_-(f_t^2+e^{-2c}f_z^2)+\tfrac12h_t^2+\tfrac12e^{-2c}h_z^2+V(h)+\tfrac18e^2(v+h)^2f^2W_-
 -2\operatorname{Im}(\bar u_1\partial_tu_1+\bar w_1\partial_tw_1),\\
 p_z&=\tfrac12W_-(f_t^2+e^{-2c}f_z^2)+\tfrac12h_t^2+\tfrac12e^{-2c}h_z^2-V(h)-\tfrac18e^2(v+h)^2f^2W_-
 +2e^{-c}\operatorname{Im}(\bar u_1\partial_zw_1+\bar w_1\partial_zu_1),\\
 q&=W_-f_tf_z+h_th_z+e^{c}\operatorname{Im}(\bar u_1\partial_tw_1+\bar w_1\partial_tu_1)
 -\operatorname{Im}(\bar u_1\partial_zu_1+\bar w_1\partial_zw_1),\\
 p_\perp&=\tfrac12h_t^2-\tfrac12e^{-2c}h_z^2-V(h)+ef\,e^{-a-b}\operatorname{Im}(\bar w_1u_1),\\
 \pi&=\tfrac12W_-(f_t^2-e^{-2c}f_z^2)-\tfrac18e^2(v+h)^2f^2W_--ef\,e^{-a-b}\operatorname{Im}(\bar w_1u_1),
\end{aligned}                                                          \tag{27.3}
\]

with \(V(h)=\lambda\bigl((v+h)^2-v^2\bigr)^2/4\).
\end{theorem}

\begin{proof}
The five identities and the forms of the bosonic energy density and
of \(\pi\) are certified; the remaining forms are the certified
components of the total stress rewritten with
\(\operatorname{Im}\).
\end{proof}

\begin{proposition}[What the assembled system contains]
\label{prop:v20v24-content}
The system (27.1)--(27.3) with (21.1) is the exact plane-symmetric
toy Standard Model of this cycle: one colour of an \(SU(2)\) gauge
field with the Higgs-generated mass \(e^2(v+h)^2/4\), the Higgs with
mass \(2\lambda v^2\) and the gauge and Yukawa sources, a fermion pair
with the Yukawa mass and the colour coupling, and the two
gravitational degrees of freedom \(a,b\) with the lapse-like \(c\),
all sourced by the five combinations (27.3).  Every result of
V1--V23 is a statement about a subsystem or a truncation of it.
\end{proposition}

\begin{proof}
By inspection of the equations and of the cited sections.
\end{proof}

\subsection{Corrigendum to V19}

\begin{proposition}[Corrected fourth-order polarization source of V19]
\label{prop:v20v24-corrigendum}
For the massless fermion pair of V16 with its colour condensate
\(f=\varepsilon^2\tfrac e4\sin2z\) on the flat background, the
anisotropic stress in the convention of (3.1), which is the one
entering the polarization equation of (27.2), is

\[
 \pi_4=\tfrac12\bigl(f_t^2-f_z^2\bigr)_4-ef_2\operatorname{Im}(\bar w_1u_1)
 =-\frac{e^2}8\cos^22z+\frac{e^2}8\sin^22z=-\frac{e^2}8\cos4z,        \tag{27.4}
\]

of mean zero.  Equation (22.3) combined the fermion part with the
sign of \(\tfrac12(W_-T_{xx}-W_+T_{yy})\) and the colour part with the
opposite convention of (3.1); its value \(-e^2/8\) and the uniform
growth \(b_4=\tfrac{e^2}{16}t^2\) of Theorem \ref{thm:v20v19-anisotropy}
are withdrawn, as is the uniform-anisotropy reading of Proposition
\ref{prop:v20v19-reading} and item 7 of Proposition
\ref{prop:v20v20-sweep} in that respect.  The corrected fourth-order
polarization is static,

\[
 b_4=\frac{e^2}{128}\cos4z,                                            \tag{27.5}
\]

a modulation at four times the fermion wavenumber with no secular
part.  The condensate, its energy, the colour self-energy
\(-e^2/16\), and the wavenumber-three response of V19 are unaffected.
\end{proposition}

\begin{proof}
(27.4), its vanishing mean, and (27.5) as the static solution of the
polarization equation are certified in the V24 verifier from the
total stress of Theorem \ref{thm:v20v24-stress}.  The \(\pi\) of
(25.3) is \(\tfrac12(T_{xx}-T_{yy})\), the negative of the convention
of (3.1), and was not used in a polarization equation.
\end{proof}

\subsection{Executable certificate}

The verifier

\[
 \texttt{scripts/verify\_sm\_v20\_v24\_assembled\_system.py}          \tag{27.6}
\]

certifies Theorems \ref{thm:v20v24-matter}--\ref{thm:v20v24-einstein}
and Proposition \ref{prop:v20v24-corrigendum}.  It emits

\[
 \texttt{artifacts/sm\_v20\_v24\_assembled\_system.json}.             \tag{27.7}
\]

The hashes are

\[
\begin{array}{c|l}
\hbox{object}&\hbox{SHA--256}\\ \hline
\hbox{verifier}&
\texttt{354444C4C200CE72CB0037D5B664F340DA4F63AF4E12F5C88DD9934FD73E1825}\\
\hbox{canonical payload}&
\texttt{4ADBAD50CA25EAD725D42E8492BFE6ECC76D20696526BFD8664461B28FAF8A8C}.
\end{array}                                                     \tag{27.8}
\]

\begin{firewall}[One colour, paired doublet, classical]
\label{fire:v20v24}
The assembled system carries one colour component along \(x\), the
unitary-gauge Higgs, and the paired locked doublet with degenerate
Yukawa coupling; the second colour of V1, the split doublet and the
cross polarization of V22, and the general locked doublet are outside
it.  All fields are classical; no continuum limit and no quantum
field theory is claimed.
\end{firewall}

\begin{decision}[V25 acquisition]
\label{dec:v20v24-next}
V25 is the compulsory companion audit.  V26--V29 then return to the
quantum reading of the combined interactions in the number basis,
beginning with the three-wave system (26.4), whose quantization is
the Jaynes--Cummings-type model of Higgs--pair conversion, and with
the non-integrable three-field averaged dynamics.
\end{decision}

\begin{assessment}[V24 acquisition]
\label{ass:v20v24}
The plane-symmetric toy Standard Model is assembled and certified
exactly consistent: matter equations (27.1) from the total
Lagrangian, diagonal total stress on the paired subspace, on-shell
covariant conservation for arbitrary metric functions, the Einstein
equations (27.2) with the explicit sources (27.3), and the
corrigendum (27.4)--(27.5) to V19.
\end{assessment}

\begin{record}[V24 append record]
\label{rec:v20v24}
V24 is appended to the validated V23 whose installed SHA--256 was
\texttt{A925736E892B79BD074229945663F27BE3732B92C8341A759E10749063163DCE}.
After installation only V24 is the active numeric SM note; frozen
archives and unrelated notes are unchanged.  The next compulsory
companion audit is V25 and the next sweep V30.
\end{record}


\section{V25: companion audit, and a convention ledger after the corrigendum}
\label{sec:v20v25}

This is the required audit at the twenty-fifth version.  Every
active companion note present in the shared folder was inspected
read-only.  Since V24 corrected a sign in V19, the audit also records
a self-audit of the cycle's anisotropic-stress conventions.

\subsection{Companion snapshot}

\begin{record}[V25 companion snapshot]
\label{rec:v20v25-snapshot}
The active companion files present were

\[
\begin{array}{c|r|r}
\hbox{file and SHA--256}&\hbox{bytes}&\hbox{lines}\\ \hline
\texttt{Renewal\_NS\_Stability\_Research\_Notes\_V26.tex}
&278283&5319\\
\multicolumn{3}{l}{
\texttt{82C887F8C2C69E4F28FAD7C3DC8DE5B9099CB5A4BF431EBA1FFF3E91935978AD}}
\\
\texttt{Renewal\_GRH\_Cofinal\_Visibility\_Attack\_V67.tex}
&640849&18652\\
\multicolumn{3}{l}{
\texttt{4A050D2FE5D6B29E11A878BD2EBC071A1FB1227F2636E98D7844D86C5F80AA58}}
\\
\texttt{Renewal\_Yang\_Mills\_Mass\_Gap\_Research\_Notes\_V41.tex}
&601237&14452\\
\multicolumn{3}{l}{
\texttt{2EEC02EB12520BAA0651AB07A171970809451FB0DEB83B596B084312F9EE4961}}
\\
\texttt{Renewal\_YM\_Parallel\_Research\_Notes\_V5.tex}
&54174&1351\\
\multicolumn{3}{l}{
\texttt{306F1BB84CFA8A8D47EA569EC17E9545F9867C8D32B548EC9D9C444B4F3C0512}}.
\end{array}                                                     \tag{28.1}
\]

Since V20 the NS, GRH, and YM parallel notes are unchanged by hash;
the Yang--Mills mass-gap note advanced from V38 to V41 (the
level-zero shell response of the block-average rule, its own audit
and sweep, and the tree-gauge representative of the tower's shells).
\end{record}

\begin{decision}[V25 audit decision]
\label{dec:v20v25-audit}
No companion theorem, numerical constant, or executable artifact is
imported.  The Yang--Mills advances concern gauge representatives of
a Euclidean blocking tower and have no object in common with the
Lorentzian assembled system (27.1)--(27.3).
\end{decision}

\begin{proof}
The only changed file is the Yang--Mills note; its new sections
contain no plane-symmetric field system, no constrained datum, and
no averaged mode dynamics.
\end{proof}

\subsection{Self-audit of conventions}

\begin{record}[Anisotropic-stress convention ledger]
\label{rec:v20v25-conventions}
The polarization equation of (3.2) and (27.2) carries
\(\pi=\tfrac12(W_+T_{yy}-W_-T_{xx})\).  In this convention the
certified parts are: colour, \(+\tfrac12W_-(f_t^2-e^{-2c}f_z^2)\)
for \(A^1=f\,dx\) (3.1); Higgs gauge mass,
\(-\tfrac18e^2(v+h)^2f^2W_-\) (27.3); paired fermion doublet,
\(-ef\,e^{-a-b}\operatorname{Im}(\bar w_1u_1)\) (27.3), the negative of
the quantity \(\tfrac12(W_-T_{xx}-W_+T_{yy})\) displayed in (21.4),
(22.3), and (25.3).  Every polarization computed by the symbolic
source generator of V3 (V2, V4, V8, V9, V12--V14) took its sources
from the exact Einstein system and is unaffected; the only
hand-combined polarization source of the cycle was (22.3), corrected
in (27.4)--(27.5).  V22's (25.4) was derived directly from the
linearized Einstein tensor and is unaffected.
\end{record}

\begin{firewall}[Audit only]
\label{fire:v20v25}
No computation is added in V25 beyond the ledger, which restates
certified results of V24.
\end{firewall}

\begin{decision}[V26 acquisition]
\label{dec:v20v25-next}
V26 gives the quantum reading of the Higgs--pair conversion: the
three-wave system (26.4) in symmetric canonical form, its
quantization as a Jaynes--Cummings-type model with the Higgs mode as
the oscillator and the pair mode and its partner as the two levels,
the exact dressed spectrum, the vacuum Rabi splitting, the
spontaneous conversion absent classically, and the certified
agreement of the large-number limit with the classical Rabi rate.
\end{decision}

\begin{assessment}[V25 acquisition]
\label{ass:v20v25}
The audit is current, nothing is imported, and the convention ledger
fixes the sign of every anisotropic-stress contribution of the cycle
in the single convention of the polarization equation.
\end{assessment}

\begin{record}[V25 append record]
\label{rec:v20v25}
V25 is appended to the validated V24 whose installed SHA--256 was
\texttt{F5C7AA0DB0EF023EDF75E139DB5260687C9B70288F0811C385085CCC15077AF6}.
After installation only V25 is the active numeric SM note; the
companions, frozen archives, and all unrelated notes are unchanged.
The next compulsory companion audit is V30, together with the sweep
of V20--V29.
\end{record}


\section{V26: the quantum reading of Higgs--pair conversion: a Jaynes--Cummings model, its dressed spectrum, and the spontaneous channel}
\label{sec:v20v26}

V23 reduced the resonant Higgs--pair exchange to the three-wave
system (26.4).  This note reads it quantum mechanically in the number
basis.  After the rescaling \(H=\sqrt{\kappa/g}\,\eta\) the system is
the symmetric three-wave system with the Hamiltonian
\(K=G(\eta\bar c_0c_-+\bar\eta c_0\bar c_-)\), \(G=\sqrt{\kappa g}\),
and its canonical quantization is the Jaynes--Cummings model: the
Higgs mode is the oscillator, the pair mode and its
negative-frequency partner are the two levels, and the interaction
\(G(a^\dagger\sigma_-+a\sigma_+)\) converts the pair mode into its
partner while emitting one Higgs quantum.  The Heisenberg equations
reproduce (26.4) under the mean-field replacement, the spectrum is
exactly the ladder of doublets \(\pm G\sqrt{n+1}\), the large-number
Rabi frequency \(G\sqrt n\) agrees with the classical rate \(gH_0\) of
V23, and the model has what the classical system lacks: from the
state with one pair mode and no Higgs quantum, the conversion
proceeds spontaneously as the vacuum Rabi oscillation
\(P=\sin^2Gt\).  This is the quantum shadow of \(f\bar f\to h\) and
\(h\to f\bar f\) in the plane-symmetric framework, with the rate
\(G=y\varepsilon\sqrt W/(2\omega)\) fixed by the Yukawa coupling.

\subsection{Canonical form}

\begin{proposition}[Symmetric three-wave system]
\label{prop:v20v26-canonical}
With \(H=\sqrt{\kappa/g}\,\eta\), (26.4) becomes
\(i\dot c_0=G\eta c_-\), \(i\dot c_-=G\bar\eta c_0\),
\(i\dot\eta=Gc_0\bar c_-\), \(G=\sqrt{\kappa g}\), which are Hamilton's
equations \(i\dot x=\partial K/\partial\bar x\) of

\[
 K=G\bigl(\eta\bar c_0c_-+\bar\eta c_0\bar c_-\bigr),\qquad
 G=\sqrt{\kappa g}=\frac{m\varepsilon\sqrt W}{\sqrt2\,\omega v}=\frac{y\varepsilon\sqrt W}{2\omega}.
                                                                    \tag{29.1}
\]
\end{proposition}

\begin{proof}
The rescaling and Hamilton's equations are certified.
\end{proof}

\subsection{The Jaynes--Cummings model}

\begin{theorem}[Quantization and exact spectrum]
\label{thm:v20v26-jc}
Let \(a\) be the Higgs mode operator and let \(|0\rangle\), \(|-\rangle\)
denote the pair mode and its partner occupied, with
\(\sigma_-=|-\rangle\langle0|\).  The quantization of (29.1) at exact
resonance is

\[
 H_{JC}=G\bigl(a^\dagger\sigma_-+a\sigma_+\bigr),                      \tag{29.2}
\]

which conserves the excitation number \(a^\dagger a+|0\rangle\langle0|\),
has the Heisenberg equations \(i\dot a=G\sigma_-\),
\(i\dot\sigma_-=-Ga\sigma_z\) whose mean-field form with
\(a\to\eta\), \(\sigma_-\to\bar c_-c_0\), \(\sigma_z\to|c_0|^2-|c_-|^2\)
is exactly (29.1), and has the spectrum

\[
 E^\pm_n=\pm G\sqrt{n+1}\ \hbox{on}\ \{|0,n\rangle,|-,n+1\rangle\},\qquad
 E=0\ \hbox{on}\ |-,0\rangle.                                          \tag{29.3}
\]

From \(|0,0\rangle\) the state \(|-,1\rangle\) is reached with
probability \(\sin^2Gt\); for large \(n\) the Rabi frequency
\(G\sqrt n\) equals the classical \(gH_0\) with \(H_0^2=(\kappa/g)n\).
\end{theorem}

\begin{proof}
The conservation law, the Heisenberg equations and their mean-field
form, the spectrum in a truncated Fock basis and in closed form on
each doublet, the vacuum Rabi oscillation, and the classical limit
are certified.
\end{proof}

\begin{proposition}[Reading]
\label{prop:v20v26-reading}
The Jaynes--Cummings ladder is the number-basis form of the energy
balance of V23: each doublet exchanges one Higgs quantum of energy
\(m_h=2\omega\) against the conversion of the pair mode into its
partner, which in the sea reinterpretation is the annihilation of a
fermion--antifermion pair into a Higgs quantum and its inverse.  The
vacuum Rabi splitting \(2G\) is the strong-coupling signature of the
Yukawa interaction at the threshold \(m_h=2\sqrt{1+m^2}\); the
spontaneous channel \(\sin^2Gt\) has no classical counterpart, the
pure pair being a fixed point of (26.4).
\end{proposition}

\begin{proof}
Theorems \ref{thm:v20v23-resonance} and \ref{thm:v20v26-jc}.
\end{proof}

\subsection{Executable certificate}

The verifier

\[
 \texttt{scripts/verify\_sm\_v20\_v26\_jaynes\_cummings.py}           \tag{29.4}
\]

certifies Proposition \ref{prop:v20v26-canonical} and Theorem
\ref{thm:v20v26-jc}.  It emits

\[
 \texttt{artifacts/sm\_v20\_v26\_jaynes\_cummings.json}.              \tag{29.5}
\]

The hashes are

\[
\begin{array}{c|l}
\hbox{object}&\hbox{SHA--256}\\ \hline
\hbox{verifier}&
\texttt{42843F7FDAD67B1155A42A03F24B6F3E157E3E1DAE185568AC8710D29C2931B1}\\
\hbox{canonical payload}&
\texttt{086E890F1E77CEE40945316EAE36A24F04C8A29C1F350006FD6A100C14553D90}.
\end{array}                                                     \tag{29.6}
\]

\begin{firewall}[Model quantization, resonant truncation]
\label{fire:v20v26}
The quantization is that of the rotating-wave three-wave model, not
of the field theory: only the homogeneous Higgs mode and the two
wavenumber-one fermion modes are quantized, at exact resonance, and
the one-quantum normalization of the classical amplitude
\(\varepsilon\) per unit transverse area is not fixed, so \(G\) is
given in terms of \(\varepsilon\).  Detuning, the colour and Yukawa
self-energies, and all other modes are absent.
\end{firewall}

\begin{decision}[V27 acquisition]
\label{dec:v20v26-next}
V27 adds the detuning: the self-energies \(-e^2/16\) of V19 and
\(\sigma_Y\) of V21 shift the pair frequency and hence the resonance
condition, and the classical three-wave system and the
Jaynes--Cummings model are solved exactly with detuning, giving the
shifted resonance line, the generalized Rabi frequency, and the
maximal conversion, all in closed form.
\end{decision}

\begin{assessment}[V26 acquisition]
\label{ass:v20v26}
The Higgs--pair conversion is read in the number basis as a
Jaynes--Cummings model with rate \(G=y\varepsilon\sqrt W/(2\omega)\),
exact spectrum (29.3), classical limit certified, and a spontaneous
conversion channel absent from the classical dynamics.
\end{assessment}

\begin{record}[V26 append record]
\label{rec:v20v26}
V26 is appended to the validated V25 whose installed SHA--256 was
\texttt{683ECB50D3491A92DB35627AA478F78F5042880070BE8AAD3E1A96D85AEB164B}.
After installation only V26 is the active numeric SM note; frozen
archives and unrelated notes are unchanged.
\end{record}


\section{V27: detuning of the Higgs--pair resonance by the pair's own colour and Yukawa self-energies}
\label{sec:v20v27}

V23 and V26 treated the exchange at exact resonance \(m_h=2\omega\).
The pair, however, shifts its own levels: by the colour condensate
(V19) and by the static Higgs displacement (V21).  This note computes
the shifts of both levels of the two-level fermion, the pair mode
\(\psi_0\) and its partner \(\psi_-\), in the same background fields,
and finds that the colour self-energy shifts them oppositely while
the Yukawa self-energy shifts them by different amounts whose
difference is exactly the mass shift produced by the depression of
the vacuum expectation value.  The gap \(\omega_0-\omega_-\) therefore
moves below \(2\omega\) by a certified amount, the resonance is
detuned, and the detuned dynamics is solved exactly: the classical
three-wave system keeps three invariants, the undepleted-Higgs
conversion obeys the generalized Rabi formula, and the detuned
Jaynes--Cummings doublets are \(-\delta/2\pm\sqrt{\delta^2/4+G^2(n+1)}\)
with maximal conversion \(G^2(n+1)/(G^2(n+1)+\delta^2/4)\).

\subsection{Self-energies of both levels}

\begin{theorem}[Colour and Yukawa shifts of the pair mode and its partner]
\label{thm:v20v27-shifts}
The colour source of V19 is, in Schr\"odinger form, the operator
\(V_c=-\tfrac12ef\,\sigma_y\) on \((u,w)\).  The massive pair mode
carries the current \(J=eW\sin2z\) and induces the condensate
\(f_2=\tfrac14eW\sin2z\); the first-order shifts of \(\psi_0\) and
\(\psi_-\) in it are

\[
 \sigma_c^{\pm}=\mp\frac{e^2W}{16\omega},                              \tag{30.1}
\]

reducing to \(-e^2/16\) of V19 at \(m=0\).  In the Higgs displacement
(24.4) the shifts are

\[
 \sigma_Y^+=-\frac{m^2W}{v^2}\Bigl[\frac{2m^2}{\omega m_h^2}+\frac{\omega}{4+m_h^2}\Bigr],\qquad
 \sigma_Y^-=\frac{m^2W}{v^2}\Bigl[\frac{2m^2}{\omega m_h^2}-\frac{\omega}{4+m_h^2}\Bigr],
                                                                    \tag{30.2}
\]

so that the wavenumber-two part of \(h_2\) shifts both levels equally
and only the vacuum depression detunes:
\(\sigma_Y^+-\sigma_Y^-=-4m^4W/(\omega v^2m_h^2)=2(m/\omega)\langle\delta m\rangle\),
with \(\langle\delta m\rangle=(m/v)\langle h_2\rangle\) of Theorem
\ref{thm:v20v21-displacement}.  The gap of the two-level fermion is

\[
 \omega_0-\omega_-=2\omega+\varepsilon^2\Delta,\qquad
 \Delta=-\frac{e^2W}{8\omega}-\frac{4m^4W}{\omega v^2m_h^2}<0.          \tag{30.3}
\]
\end{theorem}

\begin{proof}
The operator form, the current and condensate, the four projections,
the gap identity with the mass shift, and (30.3) are certified.
\end{proof}

\subsection{Detuned dynamics}

\begin{theorem}[Detuned three-wave system and generalized Rabi formula]
\label{thm:v20v27-classical}
Let \(\delta=2\omega+\varepsilon^2\Delta-m_h\).  The resonant reduction of
V23 becomes \(i\dot c_0=gHc_-e^{i\delta t}\),
\(i\dot c_-=g\bar Hc_0e^{-i\delta t}\), \(i\dot H=\kappa c_0\bar c_-e^{-i\delta t}\),
which in the frame \(c_-=de^{-i\delta t}\) is the autonomous system
\(i\dot c_0=gHd\), \(i\dot d=-\delta d+g\bar Hc_0\), \(i\dot H=\kappa c_0\bar d\)
with the exact invariants

\[
 N=|c_0|^2+|d|^2,\qquad |H|^2-\frac\kappa g|d|^2,\qquad
 K_\delta=g(H\bar c_0d+\bar Hc_0\bar d)-\delta|d|^2.                    \tag{30.4}
\]

For an undepleted real Higgs amplitude \(H_0\), with \(\Omega=gH_0\),

\[
 |d(t)|^2=\frac{\Omega^2}{\Omega^2+\delta^2/4}\sin^2\Bigl(\sqrt{\Omega^2+\delta^2/4}\,t\Bigr),
                                                                    \tag{30.5}
\]

the two-level eigenvalues being \(-\delta/2\pm\sqrt{\Omega^2+\delta^2/4}\).
\end{theorem}

\begin{proof}
The frame change, the three invariants, the eigenvalues, and (30.5)
from the exact propagator are certified.
\end{proof}

\begin{theorem}[Detuned Jaynes--Cummings ladder]
\label{thm:v20v27-jc}
With the detuning \(\delta\) the doublet \(\{|0,n\rangle,|-,n+1\rangle\}\)
of (29.2) has the eigenvalues \(-\delta/2\pm\sqrt{\delta^2/4+G^2(n+1)}\)
relative to the excitation-number ladder, and the conversion from
\(|0,n\rangle\) reaches at most

\[
 P_{\max}=\frac{G^2(n+1)}{G^2(n+1)+\delta^2/4}.                       \tag{30.6}
\]
\end{theorem}

\begin{proof}
Certified from the doublet block.
\end{proof}

\begin{proposition}[Reading]
\label{prop:v20v27-reading}
The resonance line of the Higgs--pair conversion sits at
\(m_h=2\omega+\varepsilon^2\Delta\), below the bare threshold, shifted by
the colour self-energy and by the vacuum depression, both
proportional to the pair intensity: the line moves with the number
of pairs.  The width of the resonance is set by \(G\sqrt{n+1}\), so
that at fixed \(m_h\) the conversion is complete only when the
intensity-dependent detuning is small compared to the Rabi
splitting, a nonlinear condition coupling the gauge coupling, the
Yukawa coupling and the Higgs self-coupling through (30.3).
\end{proposition}

\begin{proof}
Theorems \ref{thm:v20v27-shifts}--\ref{thm:v20v27-jc}.
\end{proof}

\subsection{Executable certificate}

The verifier

\[
 \texttt{scripts/verify\_sm\_v20\_v27\_detuned\_exchange.py}          \tag{30.7}
\]

certifies Theorems \ref{thm:v20v27-shifts}--\ref{thm:v20v27-jc}.  It
emits

\[
 \texttt{artifacts/sm\_v20\_v27\_detuned\_exchange.json}.             \tag{30.8}
\]

The hashes are

\[
\begin{array}{c|l}
\hbox{object}&\hbox{SHA--256}\\ \hline
\hbox{verifier}&
\texttt{61F90C1B333C09A98EC9C96D0CC8BAEC73921EDDF7BE04DE9C6E27412BB10356}\\
\hbox{canonical payload}&
\texttt{51C01C4AA57672A147C5C08B70486413A5363E51791BF814D09D40CB66B023E9}.
\end{array}                                                     \tag{30.9}
\]

\begin{firewall}[Frozen backgrounds, first order]
\label{fire:v20v27}
The condensate and the Higgs displacement are those induced by the
pair mode alone and are held fixed while the levels are shifted at
first order; the self-consistent change of the backgrounds during
conversion, the partner's own condensate, and the level shifts from
the gravitational sector are not included.
\end{firewall}

\begin{decision}[V28 acquisition]
\label{dec:v20v27-next}
V28 gives the quantum reading of the colour channel of V22: at the
splitting \(\omega_2-\omega_1=2\) the mass-split doublet radiates the
wavenumber-two colour wave resonantly; the corresponding three-wave
system, its Jaynes--Cummings quantization with the colour mode as the
oscillator, and its rate in terms of the gauge coupling and the
masses are derived and certified, giving the colour analogue of
(29.3).
\end{decision}

\begin{assessment}[V27 acquisition]
\label{ass:v20v27}
The level shifts of both fermion modes by the colour condensate and
the Higgs displacement are computed exactly, the gap shift (30.3) is
identified as the vacuum-depression mass shift plus the colour
self-energy, and the detuned exchange is solved exactly classically
(30.4)--(30.5) and quantum mechanically (30.6).
\end{assessment}

\begin{record}[V27 append record]
\label{rec:v20v27}
V27 is appended to the validated V26 whose installed SHA--256 was
\texttt{A060CE41E7E2C359AA873108A2A3A84422A49F834D880042B54447DA7D566B5B}.
After installation only V27 is the active numeric SM note; frozen
archives and unrelated notes are unchanged.
\end{record}


\section{V28: the colour channel of the mass-split doublet: heavy member to light member plus a colour quantum}
\label{sec:v20v28}

V22 found that the mass-split locked doublet in a colour field
radiates colour resonantly at the splitting \(\omega_2-\omega_1=2\).
This note derives the resonant three-wave system behind that
radiation, with the colour mass \(m_A\) of V11 included, and reads it
quantum mechanically as in V26.  The gauged Dirac equations of the
locked doublet are put in Schr\"odinger form, in which the colour
field couples the two isospin components through
\(-\tfrac12ef\,\sigma_x\); the two standing waves of V22, normalized,
are the eigenmodes; the resonance is \(\omega_2=\omega_1+\omega_A\) with
\(\omega_A=\sqrt{4+m_A^2}\), and the rotating-wave projections give the
symmetric coupling \(g_c\) between the heavy and the light member
mediated by the colour amplitude, while the colour current gives the
colour amplitude's own equation.  The result is again a three-wave
system with Manley--Rowe invariants and an energy balance
\(\omega_2=\omega_1+\omega_A\), whose quantization is a Jaynes--Cummings
model in which the heavy member decays into the light member and one
colour quantum: the plane-symmetric shadow of a heavy fermion
decaying to its isospin partner and a gauge boson, with the rate
\(G_c\) fixed by the gauge coupling and the two masses.

\subsection{The three-wave system}

\begin{theorem}[Schr\"odinger form and resonant reduction]
\label{thm:v20v28-threewave}
On the flat background the gauged Dirac equations of the locked
doublet are \(i\partial_t\Psi_1=H_1\Psi_1-\tfrac12ef\sigma_x\Psi_2\),
\(i\partial_t\Psi_2=H_2\Psi_2-\tfrac12ef\sigma_x\Psi_1\), with
\(H_1(u,w)=(-iw_z+m_1u,-iu_z-m_1w)\),
\(H_2(u,w)=(iw_z+m_2u,iu_z-m_2w)\), and the unit-normalized modes
\(\varphi_1=(\cos z,iW_1\sin z)/\sqrt{2\pi\omega_1W_1}\),
\(\varphi_2=(i\cos z,W_2\sin z)/\sqrt{2\pi\omega_2W_2}\) are eigenmodes
with eigenvalues \(\omega_1,\omega_2\).  For the colour wave
\(f=\tfrac12(Fe^{-i\omega_At}+\bar Fe^{i\omega_At})\sin2z\),
\(\omega_A^2=4+m_A^2\), and \(\Psi_j=\varepsilon c_j\varphi_je^{-i\omega_jt}\),
the resonance \(\omega_2=\omega_1+\omega_A\) gives the rotating-wave system

\[
 i\dot c_1=-g_c\bar Fc_2,\qquad i\dot c_2=-g_cFc_1,\qquad i\dot F=-\kappa_c\bar c_1c_2,
                                                                    \tag{31.1}
\]

\[
 g_c=\frac{e(W_1+W_2)}{16\sqrt{\omega_1W_1\omega_2W_2}},\qquad
 \kappa_c=\frac{e\varepsilon^2(W_1+W_2)}{8\pi\omega_A\sqrt{\omega_1W_1\omega_2W_2}},
                                                                    \tag{31.2}
\]

the colour equation being \(f_{tt}-f_{zz}+m_A^2f=J\) with
\(J=e\operatorname{Re}(\bar u_1w_2+\bar w_1u_2)\), which is purely
\(\sin2z\).  The system conserves \(|c_1|^2+|c_2|^2\),
\(|F|^2+(\kappa_c/g_c)|c_2|^2\), \(K=-g_c(\bar F\bar c_1c_2+F c_1\bar c_2)\),
and the energy \(\omega_1|c_1|^2+\omega_2|c_2|^2+\omega_A(g_c/\kappa_c)|F|^2\).
\end{theorem}

\begin{proof}
The Schr\"odinger form is certified against the rows of V18; the
eigenmodes, normalizations, resonant projections, the current's
wavenumber content, the colour amplitude equation, the four
invariants, and the Hamiltonian form after \(F=\sqrt{\kappa_c/g_c}\,\eta\)
are certified.
\end{proof}

\subsection{Quantum reading}

\begin{theorem}[Colour Jaynes--Cummings model]
\label{thm:v20v28-jc}
The quantization of (31.1) is \(H=-G_c(a^\dagger\sigma_-+a\sigma_+)\)
with \(\sigma_-=|1\rangle\langle2|\), the heavy member decaying into the
light member and one colour quantum, with

\[
 G_c=\sqrt{\kappa_cg_c}=\frac{e\varepsilon(W_1+W_2)}{8\sqrt{2\pi\,\omega_A\omega_1W_1\omega_2W_2}},
                                                                    \tag{31.3}
\]

the ladder \(\pm G_c\sqrt{n+1}\) on \(\{|2,n\rangle,|1,n+1\rangle\}\), and
the spontaneous decay \(P_{|1,1\rangle}(t)=\sin^2G_ct\) from
\(|2,0\rangle\).  At \(m_A=0\) the resonance is \(\omega_2-\omega_1=2\)
of V22, \(m_2^2=m_1^2+4\omega_1+4\).
\end{theorem}

\begin{proof}
Certified as in V26.
\end{proof}

\begin{proposition}[Reading]
\label{prop:v20v28-reading}
The two Jaynes--Cummings models of the cycle are the two decay
channels of the toy Standard Model read in the number basis: the
Yukawa channel (29.2), pair into Higgs, with rate
\(G=y\varepsilon\sqrt W/(2\omega)\), and the gauge channel (31.3), heavy
member into light member and colour quantum, with rate \(G_c\)
proportional to \(e\).  Both are kinematically gated by exact
threshold conditions, \(m_h=2\omega\) and \(\omega_2=\omega_1+\omega_A\),
in which the geometry-independent frequencies of the plane modes
play the role of the masses.  The emergent geometry enters only
through the modes' frequencies and the transverse area implicit in
\(\varepsilon\).
\end{proposition}

\begin{proof}
Theorems \ref{thm:v20v26-jc} and \ref{thm:v20v28-jc}.
\end{proof}

\subsection{Executable certificate}

The verifier

\[
 \texttt{scripts/verify\_sm\_v20\_v28\_colour\_channel.py}            \tag{31.4}
\]

certifies Theorems \ref{thm:v20v28-threewave} and
\ref{thm:v20v28-jc}.  It emits

\[
 \texttt{artifacts/sm\_v20\_v28\_colour\_channel.json}.               \tag{31.5}
\]

The hashes are

\[
\begin{array}{c|l}
\hbox{object}&\hbox{SHA--256}\\ \hline
\hbox{verifier}&
\texttt{7949E09C5AEA943BFA3D5FDABBE1E2ED432CB0300F9A60511AA549C19637CE5D}\\
\hbox{canonical payload}&
\texttt{4328D4DFB43D402F0BE81B0DC6C5A2D0A343F599ACEA94B1B03502E105C5A723}.
\end{array}                                                     \tag{31.6}
\]

\begin{firewall}[Rotating-wave model, flat, one colour mode]
\label{fire:v20v28}
Only the wavenumber-two colour mode and the two wavenumber-one
fermion modes are retained at exact resonance on the flat
background; the off-diagonal stress and the cross polarization of
V22 driven by the same beat, the Higgs, and the self-energies of V27
are not included, and the one-quantum normalization of
\(\varepsilon\) is not fixed.  Identities with nested radicals are
certified by exact simplification where it succeeds and otherwise by
high-precision evaluation at random rational points, as the verifier
records.
\end{firewall}

\begin{decision}[V29 acquisition]
\label{dec:v20v28-next}
V29 completes the classical side of both channels: the depleted
three-wave dynamics is reduced by its invariants to a single cubic
quadrature, the full conversion is shown to be periodic with the
period a complete elliptic integral, and the maximal conversion and
the period are given in closed form, before the V30 audit and sweep.
\end{decision}

\begin{assessment}[V28 acquisition]
\label{ass:v20v28}
The colour channel of the split doublet is reduced exactly to the
three-wave system (31.1)--(31.2) with its four invariants and
quantized as a Jaynes--Cummings model with the rate (31.3): the toy
Standard Model now has both its decay channels in the number basis.
\end{assessment}

\begin{record}[V28 append record]
\label{rec:v20v28}
V28 is appended to the validated V27 whose installed SHA--256 was
\texttt{82236A1FCC33640205255718F0478B03DB95AC457A97686E45B1AE33EC5F23E9}.
After installation only V28 is the active numeric SM note; frozen
archives and unrelated notes are unchanged.
\end{record}


\section{V29: the depleted conversion in both channels: a cubic quadrature, complete periodic conversion, and the elliptic period}
\label{sec:v20v29}

V23 and V28 gave the rotating-wave three-wave systems of the Yukawa
and colour channels, and V26--V28 their quantum ladders.  This note
closes their classical dynamics beyond the undepleted-boson
approximation.  In the symmetric form the three invariants reduce
the motion of the converted occupation \(P=|c_-|^2\) to a single cubic
quadrature.  From a pure pair mode in the presence of a boson
amplitude \(\eta_0\), the conversion is complete and periodic: \(P\)
rises to one and returns, with the closed form
\((1-k^2)\operatorname{sn}^2/\operatorname{dn}^2\) and the period
\(2K(k)/(G\sqrt{1+\eta_0^2})\), a complete elliptic integral of the
first kind, which tends to the Rabi period as the boson amplitude
grows.  The pure pair without boson and the pure boson are fixed
points: the classical dynamics has no spontaneous channel, whereas
the quantum ladders of V26 and V28 do.

\subsection{The quadrature}

\begin{theorem}[Cubic quadrature of the three-wave system]
\label{thm:v20v29-quadrature}
For \(i\dot c_0=G\eta c_-\), \(i\dot c_-=G\bar\eta c_0\),
\(i\dot\eta=Gc_0\bar c_-\), with \(N=|c_0|^2+|c_-|^2\),
\(M=|\eta|^2-|c_-|^2\), \(K=2G\operatorname{Re}(\eta\bar c_0c_-)\), the
occupation \(P=|c_-|^2\) satisfies \(\dot P=-2G\operatorname{Im}(\eta\bar c_0c_-)\)
and

\[
 \dot P^2=4G^2P(N-P)(M+P)-K^2.                                        \tag{32.1}
\]

The pure pair (\(c_-=\eta=0\)) and the pure boson (\(c_0=c_-=0\)) are
fixed points.
\end{theorem}

\begin{proof}
The invariants, \(\dot P\), the identity (32.1) on the equations of
motion, and the fixed points are certified symbolically.
\end{proof}

\subsection{Complete conversion}

\begin{theorem}[Elliptic solution and period]
\label{thm:v20v29-elliptic}
From \(c_0=1\), \(c_-=0\), \(\eta=\eta_0>0\), so that \(N=1\),
\(M=\eta_0^2\), \(K=0\), the solution is

\[
 P(t)=(1-k^2)\frac{\operatorname{sn}^2(u,k)}{\operatorname{dn}^2(u,k)},\qquad
 u=G\sqrt{1+\eta_0^2}\,t,\qquad k^2=\frac1{1+\eta_0^2},                \tag{32.2}
\]

which reaches \(P=1\), complete conversion of the pair mode into its
partner with \(|\eta|^2=1+\eta_0^2\), at half the period

\[
 T=\frac{2K(k)}{G\sqrt{1+\eta_0^2}}=\frac1G\int_0^1\frac{dP}{\sqrt{P(1-P)(P+\eta_0^2)}},
                                                                    \tag{32.3}
\]

and \(TG\eta_0\to\pi\) as \(\eta_0\to\infty\), the Rabi period of V23.
The same statements hold for the colour channel (31.1) with
\(G=G_c\).
\end{theorem}

\begin{proof}
(32.2) is certified against a thirty-digit integration of the
three-wave system at two parameter sets and five times each, with
complete conversion at \(T/2\) and return at \(T\) to \(10^{-18}\); the
integral (32.3) is certified against \(K(k)\) after the substitution
\(P=\sin^2\theta\); the Rabi limit is certified symbolically.  The sign
of \(G\) is immaterial.
\end{proof}

\begin{proposition}[Reading]
\label{prop:v20v29-reading}
Classically, a pair mode in a boson field is converted completely
and periodically, with the period lengthening from the Rabi value
as the boson is depleted, and no conversion at all without a seed
boson; quantum mechanically (V26, V28) the seed is supplied by the
vacuum and the conversion proceeds at rate \(G\).  The two readings
agree in the large-number limit, as certified in V26, and differ
precisely in the spontaneous channel, which is the only place in
this cycle where the number basis adds a process to the exact
classical plane-symmetric dynamics.
\end{proposition}

\begin{proof}
Theorems \ref{thm:v20v29-elliptic}, \ref{thm:v20v26-jc}, and
\ref{thm:v20v28-jc}.
\end{proof}

\subsection{Executable certificate}

The verifier

\[
 \texttt{scripts/verify\_sm\_v20\_v29\_depleted\_conversion.py}       \tag{32.4}
\]

certifies Theorems \ref{thm:v20v29-quadrature} and
\ref{thm:v20v29-elliptic}.  It emits

\[
 \texttt{artifacts/sm\_v20\_v29\_depleted\_conversion.json}.          \tag{32.5}
\]

The hashes are

\[
\begin{array}{c|l}
\hbox{object}&\hbox{SHA--256}\\ \hline
\hbox{verifier}&
\texttt{F54880D49908CC4DE241AF897857D071C68A8D1D2D6F5DBA5122855D8ACB6E5F}\\
\hbox{canonical payload}&
\texttt{5A7890C928370E16835E34B797549528ECFE64B7523D2F50524C3CB272021DF1}.
\end{array}                                                     \tag{32.6}
\]

\begin{firewall}[Resonant model, numerical elliptic certificate]
\label{fire:v20v29}
The results concern the rotating-wave three-wave models at exact
resonance; the elliptic closed form and period are certified
numerically at thirty digits rather than by a symbolic identity for
the Jacobi functions, which the symbolic engine lacks.  Detuning
(V27), self-consistent backgrounds, and gravity are absent.
\end{firewall}

\begin{decision}[V30 acquisition]
\label{dec:v20v29-next}
V30 is the compulsory companion audit and the third ten-version
sweep of the cycle, recording what V20--V29 established, the
corrigendum to V19, and the direction of V31--V39.
\end{decision}

\begin{assessment}[V29 acquisition]
\label{ass:v20v29}
The classical conversion dynamics of both channels is closed: the
cubic quadrature (32.1), the elliptic solution (32.2), and the period
(32.3), with the Rabi limit and the absence of a classical
spontaneous channel.
\end{assessment}

\begin{record}[V29 append record]
\label{rec:v20v29}
V29 is appended to the validated V28 whose installed SHA--256 was
\texttt{FABFFCA733C715FA2C8A914C3D230D75B05FDD9346DAB98551322DA62E849C82}.
After installation only V29 is the active numeric SM note; frozen
archives and unrelated notes are unchanged.
\end{record}


\section{V30: companion audit and the third ten-version sweep of the twentieth cycle}
\label{sec:v20v30}

This is the required audit at the thirtieth version and the sweep of
V20--V29.  Every active companion note present in the shared folder
was inspected read-only.

\subsection{Companion snapshot}

\begin{record}[V30 companion snapshot]
\label{rec:v20v30-snapshot}
The active companion files present were

\[
\begin{array}{c|r|r}
\hbox{file and SHA--256}&\hbox{bytes}&\hbox{lines}\\ \hline
\texttt{Renewal\_NS\_Stability\_Research\_Notes\_V26.tex}
&278283&5319\\
\multicolumn{3}{l}{
\texttt{82C887F8C2C69E4F28FAD7C3DC8DE5B9099CB5A4BF431EBA1FFF3E91935978AD}}
\\
\texttt{Renewal\_GRH\_Cofinal\_Visibility\_Attack\_V69.tex}
&674175&19330\\
\multicolumn{3}{l}{
\texttt{1FC4490F5816480A88309D1DA0E33FCCEBE7BB4BA698D9BE7B147A02A1D6632F}}
\\
\texttt{Renewal\_Yang\_Mills\_Mass\_Gap\_Research\_Notes\_V41.tex}
&601237&14452\\
\multicolumn{3}{l}{
\texttt{2EEC02EB12520BAA0651AB07A171970809451FB0DEB83B596B084312F9EE4961}}
\\
\texttt{Renewal\_YM\_Parallel\_Research\_Notes\_V5.tex}
&54174&1351\\
\multicolumn{3}{l}{
\texttt{306F1BB84CFA8A8D47EA569EC17E9545F9867C8D32B548EC9D9C444B4F3C0512}}.
\end{array}                                                     \tag{33.1}
\]

Since V25 the GRH note advanced from V67 to V69 (the leading form
and collision hierarchy of a real-axis determinant, a one-pair closed
form, the pairing form of a heat image, and a bridge ledger); the NS,
Yang--Mills, and YM parallel notes are unchanged by hash.
\end{record}

\begin{decision}[V30 audit decision]
\label{dec:v20v30-audit}
No companion theorem, numerical constant, or executable artifact is
imported.  The GRH advances concern determinants and heat images of
Appell-type polynomials and have no object in common with the
plane-symmetric field systems, resonant reductions, or number-basis
models of this cycle.
\end{decision}

\begin{proof}
The only changed file is the GRH note; its new sections contain no
constrained datum, no field system, and no three-wave dynamics.
\end{proof}

\subsection{Sweep of V20--V29}

\begin{proposition}[What V20--V29 established]
\label{prop:v20v30-sweep}
The third block of the cycle completed the toy Standard Model and
read its interactions in the number basis:

\begin{enumerate}
\item the Yukawa coupling: the fermion mass as the vacuum expectation
value read through \(y\), the paired doublet invariant with a
field-dependent mass, the static depression of the vacuum expectation
value by a fermion pair, and the negative Yukawa self-energy in closed
form (V21);
\item the mass-split doublet: the pairing defect, the universal stress
\(T_{xy}=-\tfrac12efS\), the beating bilinears at \(\omega_2-\omega_1\),
the linearized cross polarization \(d_{tt}-d_{zz}=2T_{xy}\), and
resonant cross-polarized gravitational and colour radiation at the
splittings \(2\) and \(4\) (V22);
\item fermion--Higgs exchange: the exact two-mode reduction of the
polarized Dirac field under a homogeneous Higgs, its exact norm and
energy invariants, the resonance \(m_h=2\omega\) as the classical pair
threshold, and the three-wave system with its Manley--Rowe balance
(V23);
\item the assembled Einstein--Yang--Mills--Higgs--Dirac plane system:
matter equations from the total Lagrangian, diagonal total stress on
the paired subspace, on-shell covariant conservation for arbitrary
metric functions, the Einstein equations with the explicit sources
(27.3), and the corrigendum to V19: the fourth-order polarization
source of the fermion--colour system is \(-\tfrac{e^2}8\cos4z\), of
mean zero, with the static response \(\tfrac{e^2}{128}\cos4z\) and no
secular growth (V24);
\item the audit with the anisotropic-stress convention ledger (V25);
\item the quantum reading of both decay channels as Jaynes--Cummings
models, pair into Higgs with rate \(G=y\varepsilon\sqrt W/(2\omega)\)
(V26) and heavy member into light member plus colour quantum with rate
\(G_c\) (V28), with exact ladders, vacuum Rabi oscillations, and
certified classical limits;
\item detuning by the pair's own colour and Yukawa self-energies, the
gap shift identified as the vacuum-depression mass shift plus the
colour self-energy, and the detuned dynamics solved exactly (V27);
\item the depleted classical conversion: the cubic quadrature, the
complete periodic conversion in Jacobi functions, the elliptic
period, and the absence of a classical spontaneous channel (V29).
\end{enumerate}
\end{proposition}

\begin{proof}
Items are the certified theorems of the cited sections.
\end{proof}

\begin{proposition}[Direction for V31--V39]
\label{prop:v20v30-direction}
The cycle now holds an exact plane-symmetric toy Standard Model
coupled to gravity, with its two decay channels understood mode by
mode, classically and in the number basis.  What separates this from
a continuum is that every result concerns a finite set of resonant
modes.  The upstream step is therefore the passage from resonant
triads to many modes: V31--V33 extend the three-wave reductions to
all resonant triads of the plane system, with the frequency-matching
conditions as exact selection rules and the couplings as certified
functions of the wavenumbers; V34 derives, under the random-phase
hypothesis, the wave-kinetic equation of the decay channels, the
classical Boltzmann equation whose collision term is built from the
certified rates, which is the first continuum object of the cycle;
V35 is the audit; V36--V38 read the conversions gravitationally, as
the change of the equation of state seen by the averaged Einstein
equations of V16--V17 and V23, and extend the linear cross
polarization of V22 to the exact two-polarization plane metric; V39
prepares the V40 sweep.
\end{proposition}

\begin{proof}
This is a decision.
\end{proof}

\begin{firewall}[Audit and sweep only]
\label{fire:v20v30}
No computation is added in V30.  The cycle's results are exact
statements about plane-symmetric classical fields, their resonant
reductions, and model quantizations of those reductions; no
continuum limit, no quantum field theory of any sector, and no
Standard Model--Einstein continuum are claimed.
\end{firewall}

\begin{decision}[V31 acquisition]
\label{dec:v20v30-next}
V31 enumerates the resonant triads of the Yukawa channel for
fermion modes of general wavenumber \(k\) and Higgs modes of
wavenumber \(q\): the selection rules \(k_0-k_-=q\) and
\(\omega_{k_0}+\omega_{k_-}=\Omega_q\), the mode functions, and the
coupling \(g(k_0,k_-;q)\) certified from the projections, with the
homogeneous case \(q=0\) reducing to (26.4).
\end{decision}

\begin{assessment}[V30 acquisition]
\label{ass:v20v30}
The audit is current and the sweep records the completion of the
toy Standard Model in the exact plane framework, its two decay
channels in both readings, the corrigendum to V19, and the direction
toward the wave-kinetic continuum.
\end{assessment}

\begin{record}[V30 append record]
\label{rec:v20v30}
V30 is appended to the validated V29 whose installed SHA--256 was
\texttt{EE8BCE49E9F0C8C3ED05B34E1A12B32DC62C13FCE6A6CE703D62FB5FDBABD282}.
After installation only V30 is the active numeric SM note; the
companions, frozen archives, and all unrelated notes are unchanged.
The next compulsory companion audit is V35 and the next sweep V40.
\end{record}


\section{V31: the resonant triads of the Yukawa channel for general wavenumbers: selection rules and couplings}
\label{sec:v20v31}

V23--V29 treated the Yukawa channel for the wavenumber-one pair and
the homogeneous Higgs.  This note extends it to all wavenumbers.
The cosine family of polarized Dirac modes on the circle, positive
and negative frequency, is written down for every wavenumber \(k\)
and certified; a Higgs mode \(\cos qz\) of frequency
\(\Omega_q=\sqrt{q^2+m_h^2}\) couples the positive mode \(k_0\) to the
negative mode \(k_-\) only if \(q=|k_0-k_-|\) or \(q=k_0+k_-\), the
wavenumber selection rule, and resonantly only if
\(\omega_{k_0}+\omega_{k_-}=\Omega_q\), the frequency selection rule; the
resonant three-wave coupling is a certified closed function of the
triad.  These are the exact plane-symmetric forms of momentum and
energy conservation in the decay of a Higgs quantum into a
fermion--antifermion pair, and they define the resonance manifold
on which the wave-kinetic description of V34 will be built.

\subsection{Modes and selection rules}

\begin{theorem}[Cosine-family modes]
\label{thm:v20v31-modes}
For every integer \(k\ge1\), with \(\omega_k=\sqrt{k^2+m^2}\) and
\(W_k=\omega_k-m\),

\[
 \psi_k^+=\Bigl(\cos kz,\ i\frac{W_k}k\sin kz\Bigr),\qquad
 \psi_k^-=\Bigl(\frac{W_k}k\cos kz,\ -i\sin kz\Bigr)                   \tag{34.1}
\]

are eigenmodes of \(H_0\) with eigenvalues \(\omega_k\) and \(-\omega_k\),
orthogonal, of norm \(2\pi\omega_kW_k/k^2\); \(k=1\) gives the modes of
V23.
\end{theorem}

\begin{proof}
Certified for symbolic integer \(k\) (eigenmodes) and for
\(k=1,2,3\) (norms and orthogonality).
\end{proof}

\begin{theorem}[Selection rule and matrix element]
\label{thm:v20v31-selection}
The matrix element
\(\langle\psi_{k_-}^-|\cos qz\operatorname{diag}(1,-1)|\psi_{k_0}^+\rangle\)
vanishes unless \(q=|k_0-k_-|\) or \(q=k_0+k_-\), and then equals

\[
 \frac\pi2\Bigl(\frac{W_{k_-}}{k_-}+s\,\frac{W_{k_0}}{k_0}\Bigr),\qquad
 s=+1\ (q=|k_0-k_-|),\quad s=-1\ (q=k_0+k_-),                         \tag{34.2}
\]

for \(q\ge1\), and \(2\pi W_k/k\) for \(q=0\), \(k_0=k_-=k\).
\end{theorem}

\begin{proof}
The product-to-sum identities for
\(\cos qz\cos k_-z\cos k_0z\) and \(\cos qz\sin k_-z\sin k_0z\) are
certified for symbolic integers, giving the rule; the values are
certified on seven allowed triads, six forbidden ones, and the
\(q=0\) cases \(k=1,2,3\).
\end{proof}

\subsection{The triad three-wave system}

\begin{theorem}[Resonant couplings]
\label{thm:v20v31-coupling}
With unit-normalized modes \(\varphi_{k_0}^+\), \(\varphi_{k_-}^-\), the
Higgs mode \(h_q=\tfrac12(He^{-i\Omega_qt}+\bar He^{i\Omega_qt})\cos qz\), and
\(\Psi=\varepsilon(c_0\varphi_{k_0}^+e^{-i\omega_{k_0}t}+c_-\varphi_{k_-}^-e^{i\omega_{k_-}t})\),
the rotating-wave reduction at \(\omega_{k_0}+\omega_{k_-}=\Omega_q\) is
\(i\dot c_0=gHc_-\), \(i\dot c_-=g\bar Hc_0\), \(i\dot H=\kappa c_0\bar c_-\)
with

\[
 g(k_0,k_-;q)=\frac m{8v}\,\frac{W_{k_-}k_0+s\,W_{k_0}k_-}{\sqrt{\omega_{k_0}W_{k_0}\omega_{k_-}W_{k_-}}}\ (q\ge1),\qquad
 g(k,k;0)=\frac{mk}{2v\omega_k},                                        \tag{34.3}
\]

\[
 \frac\kappa g=\frac{4\varepsilon^2}{\pi\Omega_q}\ (q\ge1),\qquad
 \frac\kappa g=\frac{2\varepsilon^2}{\pi\Omega_0}\ (q=0).                \tag{34.4}
\]

At \(q=0\), \(k=1\) this is (26.4): \(g=m/(2\omega v)\), and
\(\kappa/g\) times the mode norm \(2\pi\omega W\) is the \(2\varepsilon^2W\)
of V23.  The invariants and the elliptic dynamics of V23 and V29
apply verbatim to every triad.
\end{theorem}

\begin{proof}
The projections of the exact Dirac equation with the mass
modulation and of the exact Higgs equation with the pair source are
certified on the triads \((2,1;1)\), \((3,1;2)\), \((2,1;3)\),
\((1,1;0)\), \((2,2;0)\), against (34.3)--(34.4); the V23
consistency is certified.
\end{proof}

\begin{proposition}[Resonance manifold]
\label{prop:v20v31-manifold}
Each allowed triad is resonant on the surface
\(m_h^2=(\omega_{k_0}+\omega_{k_-})^2-q^2\) in the \((m,m_h)\) plane, which
lies on or above the threshold \(m_h=2m\), with equality exactly for
the sum triads \(q=2k\), \(k_0=k_-=k\), where the two fermions move
together.  In the sea
reinterpretation the triad is the decay \(h(q)\to f(k_0)\bar f(-k_-)\)
(for \(s=+1\)) or \(h(q)\to f(k_0)\bar f(k_-)\) (for \(s=-1\)), the
standing modes carrying both signs of the wavenumber: wavenumber and
frequency selection are the exact plane-symmetric forms of momentum
and energy conservation, and the coupling (34.3) is the amplitude
of the decay.
\end{proposition}

\begin{proof}
Theorems \ref{thm:v20v31-selection}--\ref{thm:v20v31-coupling}; the
threshold is certified symbolically for the marginal triad and at
sample masses for the others.
\end{proof}

\subsection{Executable certificate}

The verifier

\[
 \texttt{scripts/verify\_sm\_v20\_v31\_yukawa\_triads.py}             \tag{34.5}
\]

certifies Theorems \ref{thm:v20v31-modes}--\ref{thm:v20v31-coupling}
and the threshold statement, the marginal triad \((1,1;2)\)
symbolically.  It emits

\[
 \texttt{artifacts/sm\_v20\_v31\_yukawa\_triads.json}.                \tag{34.6}
\]

The hashes are

\[
\begin{array}{c|l}
\hbox{object}&\hbox{SHA--256}\\ \hline
\hbox{verifier}&
\texttt{697D8534FE2CEC58A685A845623617E34B6242AA05E4FFA1E2F5915914B39582}\\
\hbox{canonical payload}&
\texttt{8787EC323C8986B68FFF8E991F43A45131892C84F1A79E95804FC5D1DB7D4236}.
\end{array}                                                     \tag{34.7}
\]

\begin{firewall}[Cosine family, one triad at a time]
\label{fire:v20v31}
Only the cosine family of fermion modes and the cosine Higgs modes
are treated, so that the families do not mix; each triad is reduced
in isolation at exact resonance, and the general formulas (34.3)--(34.4)
are certified on concrete triads while the selection rule is
certified symbolically.  Nothing is said yet about many triads
acting together.
\end{firewall}

\begin{decision}[V32 acquisition]
\label{dec:v20v31-next}
V32 does the same for the colour channel: the sector-minus modes of
general wavenumber, the colour modes \(\sin qz\), the selection rules
for the heavy-to-light transition, and the coupling
\(g_c(k_1,k_2;q)\), reducing to (31.2) at \((1,1;2)\).
\end{decision}

\begin{assessment}[V31 acquisition]
\label{ass:v20v31}
The Yukawa channel is extended to all wavenumbers: the modes (34.1),
the selection rule (34.2), the couplings (34.3)--(34.4), and the
resonance manifold, each certified.
\end{assessment}

\begin{record}[V31 append record]
\label{rec:v20v31}
V31 is appended to the validated V30 whose installed SHA--256 was
\texttt{463F18E74798F85C0428C9F22CCA8A075E9291241CE2BEB7A64A3B916D3C7151}.
After installation only V31 is the active numeric SM note; frozen
archives and unrelated notes are unchanged.
\end{record}


\section{V32: the resonant triads of the colour channel for general wavenumbers}
\label{sec:v20v32}

V31 extended the Yukawa channel to all wavenumbers; this note does
the same for the colour channel of V28.  The light member of the
locked doublet lives in the sector-plus family of V31 with mass
\(m_1\); the heavy member lives in the sector-minus family, whose
modes for every wavenumber are written down and certified; a colour
mode \(\sin qz\) of frequency \(\omega_{A,q}=\sqrt{q^2+m_A^2}\) couples
the heavy mode \(k_2\) to the light mode \(k_1\) only if
\(q=k_1+k_2\) or \(q=|k_1-k_2|\), and resonantly only if
\(\omega_{2,k_2}=\omega_{1,k_1}+\omega_{A,q}\).  The resonant coupling is
a certified closed function of the triad, with the colour-side
coefficient fixed by the current, and the whole reduces to (31.2)
at \((1,1;2)\).  Together with V31 this gives the complete list of
resonant triads of the toy Standard Model's two decay channels.

\subsection{Modes and selection rule}

\begin{theorem}[Sector-minus modes and the selection rule]
\label{thm:v20v32-modes}
For the heavy member with \(H_2(u,w)=(iw_z+m_2u,iu_z-m_2w)\) and every
integer \(k\ge1\), \(\varphi_{2,k}=(i\cos kz,(W_{2k}/k)\sin kz)\) with
\(\omega_{2k}=\sqrt{k^2+m_2^2}\), \(W_{2k}=\omega_{2k}-m_2\), is an eigenmode
of eigenvalue \(\omega_{2k}\) and norm \(2\pi\omega_{2k}W_{2k}/k^2\); the
light member's modes are \(\psi_k^+\) of (34.1) with \(m_1\).  The
matrix element \(\langle\varphi_{1,k_1}|\sin qz\,\sigma_x|\varphi_{2,k_2}\rangle\)
equals

\[
 \frac\pi2\Bigl(\frac{W_{2k_2}}{k_2}+\frac{W_{1k_1}}{k_1}\Bigr)\ (q=k_1+k_2),\quad
 \frac\pi2\Bigl(\frac{W_{2k_2}}{k_2}-\frac{W_{1k_1}}{k_1}\Bigr)\ (q=k_2-k_1),\quad
 \frac\pi2\Bigl(\frac{W_{1k_1}}{k_1}-\frac{W_{2k_2}}{k_2}\Bigr)\ (q=k_1-k_2),
                                                                    \tag{35.1}
\]

and vanishes otherwise.
\end{theorem}

\begin{proof}
The eigenmodes are certified for symbolic \(k\), the norms for
\(k=1,2,3\), the product-to-sum identities for symbolic integers,
and (35.1) on eight allowed and seven forbidden triads.
\end{proof}

\subsection{The triad three-wave system}

\begin{theorem}[Resonant colour couplings]
\label{thm:v20v32-coupling}
With unit-normalized modes, the colour mode
\(f=\tfrac12(Fe^{-i\omega_{A,q}t}+\bar Fe^{i\omega_{A,q}t})\sin qz\), and
\(\Psi_j=\varepsilon c_j\varphi_{j,k_j}e^{-i\omega_{j,k_j}t}\), the
rotating-wave reduction at \(\omega_{2,k_2}=\omega_{1,k_1}+\omega_{A,q}\) is
\(i\dot c_1=-g_c\bar Fc_2\), \(i\dot c_2=-g_cFc_1\),
\(i\dot F=-\kappa_c\bar c_1c_2\) with

\[
 g_c(k_1,k_2;q)=\frac e{16}\,\frac{\mathcal N}{\sqrt{\omega_{1k_1}W_{1k_1}\omega_{2k_2}W_{2k_2}}},\qquad
 \mathcal N=\begin{cases}W_{2k_2}k_1+W_{1k_1}k_2&(q=k_1+k_2)\\ W_{2k_2}k_1-W_{1k_1}k_2&(q=k_2-k_1)\\ W_{1k_1}k_2-W_{2k_2}k_1&(q=k_1-k_2),\end{cases}
                                                                    \tag{35.2}
\]

\[
 \frac{\kappa_c}{g_c}=\frac{2\varepsilon^2}{\pi\omega_{A,q}},              \tag{35.3}
\]

the colour equation being \(f_{tt}-f_{zz}+m_A^2f=J\) with the current
of V28 projected on \(\sin qz\).  At \((1,1;2)\) this is (31.2).  The
invariants of V28 and the elliptic dynamics of V29 apply to every
triad, and each triad is resonant on the surface
\(m_A^2=(\omega_{2,k_2}-\omega_{1,k_1})^2-q^2\).
\end{theorem}

\begin{proof}
Certified on the triads \((1,1;2)\), \((1,2;3)\), \((1,2;1)\),
\((2,1;1)\), \((3,1;2)\) from the exact Dirac and colour equations;
the V28 consistency is certified.
\end{proof}

\begin{proposition}[The complete resonant triad list]
\label{prop:v20v32-list}
The toy Standard Model of V24 has, in the cosine and sine families
retained, exactly the resonant triads of Theorems
\ref{thm:v20v31-coupling} and \ref{thm:v20v32-coupling}: Higgs into
pair and heavy into light plus colour, each with a wavenumber
selection rule (three or two branches), a frequency condition, and
a coupling given in closed form by the masses and wavenumbers.
Every other cubic interaction of the assembled system, in
particular the Higgs--colour and fermion--gravity couplings, either
lacks a resonant branch in these families or enters at higher order
in the amplitudes.
\end{proposition}

\begin{proof}
The cubic vertices of (27.1) and (21.1) are the Yukawa term, the
colour--fermion term, the Higgs--colour term \(e^2vhf^2\) and the
gravitational couplings; the first two are the certified channels,
the third is quadratic in \(f\) and enters the averaged dynamics of
V12 without a decay branch, and the gravitational couplings are of
higher order in the plane framework (V8, V12).
\end{proof}

\subsection{Executable certificate}

The verifier

\[
 \texttt{scripts/verify\_sm\_v20\_v32\_colour\_triads.py}             \tag{35.4}
\]

certifies Theorems \ref{thm:v20v32-modes} and
\ref{thm:v20v32-coupling}.  It emits

\[
 \texttt{artifacts/sm\_v20\_v32\_colour\_triads.json}.                \tag{35.5}
\]

The hashes are

\[
\begin{array}{c|l}
\hbox{object}&\hbox{SHA--256}\\ \hline
\hbox{verifier}&
\texttt{89B1957B0ACB43CB9EA66CB294FADADB994381ABE511C975706F32758BBEAB51}\\
\hbox{canonical payload}&
\texttt{5D91F7CD00F43C58A2B74B032D4DD2900BB6099E6073C05478ADCBA73AFDA189}.
\end{array}                                                     \tag{35.6}
\]

\begin{firewall}[Families and isolation]
\label{fire:v20v32}
As in V31, one family of fermion modes and the sine colour modes are
retained, each triad is reduced in isolation at exact resonance, and
the general formulas are certified on concrete triads.  The
statement of Proposition \ref{prop:v20v32-list} about the other
vertices is a reading of the equations, not a certified exclusion.
\end{firewall}

\begin{decision}[V33 acquisition]
\label{dec:v20v32-next}
V33 assembles the many-triad system: for a finite set of fermion
and boson modes, the coupled rotating-wave equations of all resonant
triads sharing modes, their Hamiltonian, the conserved excitation
numbers, and the failure of integrability when two triads share a
mode, exhibited on the smallest such system.
\end{decision}

\begin{assessment}[V32 acquisition]
\label{ass:v20v32}
The colour channel is extended to all wavenumbers: the heavy-member
modes, the three-branch selection rule (35.1), and the couplings
(35.2)--(35.3), each certified, completing the resonant triad list
of the toy Standard Model.
\end{assessment}

\begin{record}[V32 append record]
\label{rec:v20v32}
V32 is appended to the validated V31 whose installed SHA--256 was
\texttt{0D34D381F7E8B175E40A0F760570B1D3ED910A55D6C962A5019D3831BEE93106}.
After installation only V32 is the active numeric SM note; frozen
archives and unrelated notes are unchanged.
\end{record}


\section{V33: the many-triad system and the smallest shared-mode system: symmetries, invariants, and what integrability would need}
\label{sec:v20v33}

V31--V32 listed the resonant triads of the two decay channels one at
a time.  This note assembles them: for any finite set of resonant
triads the rotating-wave equations are Hamilton's equations of a
single cubic Hamiltonian, the sum of the triad Hamiltonians in the
symmetric normalization, and modes shared between triads couple the
triads.  The smallest shared-mode system, one pair mode converted by
two Higgs modes into two partner modes, is analysed exactly: it has
a \(U(1)^3\) symmetry, exactly three quadratic invariants, the cubic
Hamiltonian, and, certified by finite linear algebra, no further
phase-invariant polynomial invariant of degree up to four.  Its
reduced phase space has two degrees of freedom with the Hamiltonian
as one invariant, so Liouville integrability would require a second
independent invariant, which does not exist among low-degree
polynomials.  Each triad alone is the integrable system of V29,
recovered on an invariant subspace.

\subsection{The many-triad Hamiltonian}

\begin{theorem}[Many-triad rotating-wave system]
\label{thm:v20v33-hamiltonian}
For resonant triads \(\tau\) with pump mode \(p_\tau\), daughter modes
\(d_\tau,d'_\tau\), and symmetric couplings \(G_\tau\), the rotating-wave
equations are \(i\dot x=\partial K/\partial\bar x\) with

\[
 K=\sum_\tau G_\tau\bigl(x_{p_\tau}\bar x_{d_\tau}x_{d'_\tau}+\hbox{c.c.}\bigr),
                                                                    \tag{36.1}
\]

where a mode shared by several triads appears in each of their
terms.  For the pair mode \(c_0\) converted by two Higgs modes
\(\eta_a,\eta_b\) into partners \(c_a,c_b\),

\[
 K=G_a(\eta_a\bar c_0c_a+\hbox{c.c.})+G_b(\eta_b\bar c_0c_b+\hbox{c.c.}),
                                                                    \tag{36.2}
\]

giving \(i\dot c_0=G_a\eta_ac_a+G_b\eta_bc_b\), \(i\dot c_a=G_a\bar\eta_ac_0\),
\(i\dot\eta_a=G_ac_0\bar c_a\), and likewise for \(b\).
\end{theorem}

\begin{proof}
Certified on (36.2); the general case is the same computation term
by term.
\end{proof}

\subsection{Invariants of the shared-mode system}

\begin{theorem}[Symmetry and invariants]
\label{thm:v20v33-invariants}
The system (36.2) is invariant under \(U(1)^3\) with charges
\(c_0:(1,0,0)\), \(c_a:(0,1,0)\), \(c_b:(0,0,1)\), \(\eta_a:(1,-1,0)\),
\(\eta_b:(1,0,-1)\).  The space of quadratic polynomial invariants is
exactly three-dimensional, spanned by

\[
 N=|c_0|^2+|c_a|^2+|c_b|^2,\qquad M_a=|\eta_a|^2-|c_a|^2,\qquad M_b=|\eta_b|^2-|c_b|^2,
                                                                    \tag{36.3}
\]

and the phase-invariant polynomial invariants of degree at most four
form the eleven-dimensional space spanned by \(1\), \(N\), \(M_a\),
\(M_b\), \(K\), and the six products of pairs of (36.3): there is no
new invariant of degree at most four.  The subspace \(c_b=\eta_b=0\)
is invariant and carries the single-triad dynamics of V29.
\end{theorem}

\begin{proof}
The symmetry is certified.  The invariant spaces are the null spaces
of the linear conditions \(\dot Q=0\) on the coefficients of a general
quadratic form, respectively of a general phase-invariant polynomial
of degree at most four, computed exactly at \(G_a=1\), \(G_b=2\); the
eleven known invariants are certified to be independent.  The
invariant subspace is certified.
\end{proof}

\begin{proposition}[What integrability would need]
\label{prop:v20v33-reduced}
The ten-dimensional phase space of (36.2), reduced by the three
symmetries at fixed charges, is four-dimensional with \(K\) as its
Hamiltonian; Liouville integrability of the reduced system would
require one further independent invariant.  None exists among
phase-invariant polynomials of degree at most four; the single-triad
subsystems are integrable (V29).  The many-triad dynamics of the toy
Standard Model is therefore, already for two triads sharing the
pair mode, outside the class solved by the invariants that solve
each triad.
\end{proposition}

\begin{proof}
Dimension count from Theorem \ref{thm:v20v33-invariants}; the
exclusion is the certified linear algebra.
\end{proof}

\subsection{Executable certificate}

The verifier

\[
 \texttt{scripts/verify\_sm\_v20\_v33\_shared\_mode.py}               \tag{36.4}
\]

certifies Theorems \ref{thm:v20v33-hamiltonian} and
\ref{thm:v20v33-invariants}.  It emits

\[
 \texttt{artifacts/sm\_v20\_v33\_shared\_mode.json}.                  \tag{36.5}
\]

The hashes are

\[
\begin{array}{c|l}
\hbox{object}&\hbox{SHA--256}\\ \hline
\hbox{verifier}&
\texttt{40B06C03B6A81723095BEDF76DF546CFD5673F4C0C2545DB5427B7354ECF4754}\\
\hbox{canonical payload}&
\texttt{EE26AB49E1E5A873200C9C583670039EB8F209AB7D659E951E5C19926CE3D800}.
\end{array}                                                     \tag{36.6}
\]

\begin{firewall}[Polynomial exclusion only]
\label{fire:v20v33}
Non-integrability is not proved: only the absence of polynomial
invariants of degree at most four, at one generic pair of couplings,
is certified.  The many-triad Hamiltonian is the rotating-wave
truncation at exact resonance of every triad; detunings and
non-resonant terms are dropped.
\end{firewall}

\begin{decision}[V34 acquisition]
\label{dec:v20v33-next}
V34 derives the wave-kinetic equation of the decay channels: the
exact second-order solution of a detuned triad from random-phase
data, its phase average, the resonance integral that turns the
coherent \(t^2\) growth into a rate, the collision term
\(2\pi G^2\rho(0)[n_\eta n_--n_0(n_-+n_\eta)]\), its conservation laws,
the \(H\)-theorem, and the Rayleigh--Jeans equilibrium.
\end{decision}

\begin{assessment}[V33 acquisition]
\label{ass:v20v33}
The many-triad rotating-wave system is a single cubic Hamiltonian
(36.1); the smallest shared-mode system has exactly the invariants
(36.3) and \(K\) among polynomials of degree at most four, so each
triad's integrability does not extend to shared modes.
\end{assessment}

\begin{record}[V33 append record]
\label{rec:v20v33}
V33 is appended to the validated V32 whose installed SHA--256 was
\texttt{C420167A73CC18C3B56563AED20356AC4E34C174B6F5A8695F768ACB5310B171}.
After installation only V33 is the active numeric SM note; frozen
archives and unrelated notes are unchanged.
\end{record}


\section{V34: the wave-kinetic equation of the decay channels: collision term, conservation laws, \texorpdfstring{$H$}{H}-theorem, and Rayleigh--Jeans equilibrium}
\label{sec:v20v34}

V31--V33 established the resonant triads and their coupled
Hamiltonian.  This note derives the first continuum object of the
cycle: the wave-kinetic equation that governs the mode occupations
of the two decay channels when many triads with a continuum of
detunings act on random-phase data.  The derivation is exact at
second order in the coupling: the detuned triad is solved to
second order from constant data, the phase average kills every
monomial with a net phase, and what survives is the collision term
\(C=n_\eta n_--n_0(n_-+n_\eta)\) multiplied by the kernel
\((1-\cos\delta t)/\delta^2\).  The same collision term appears exactly
as the second time derivative of the occupation along the resonant
flow.  Integrating the kernel over a continuum of detunings gives
\(\pi|t|\), turning the coherent \(t^2\) growth into a rate: the
kinetic equation \(\dot n_0=2\pi G^2\rho(0)\,C\).  It conserves the
two Manley--Rowe numbers and the energy on the resonance, its
entropy production is a certified non-negative square, and its
equilibria are the Rayleigh--Jeans distributions with additive
chemical potentials.  This is the classical statistical mechanics
of the toy Standard Model's decays, with rates fixed by the
certified couplings of V31--V32.

\subsection{Second-order solution and phase average}

\begin{theorem}[Random-phase second-order occupations]
\label{thm:v20v34-second-order}
For the detuned triad \(i\dot c_0=G\eta c_-e^{-i\delta t}\),
\(i\dot c_-=G\bar\eta c_0e^{i\delta t}\), \(i\dot\eta=Gc_0\bar c_-e^{i\delta t}\)
with constant data of moduli \(\sqrt{n_0},\sqrt{n_-},\sqrt{n_\eta}\) and
independent uniformly distributed phases,

\[
 \langle n_0(t)\rangle=n_0+2G^2C\,\frac{1-\cos\delta t}{\delta^2}+O(G^3),\qquad
 \langle n_-(t)\rangle=n_--2G^2C\,\frac{1-\cos\delta t}{\delta^2},\qquad
 \langle n_\eta(t)\rangle=n_\eta-2G^2C\,\frac{1-\cos\delta t}{\delta^2},
                                                                    \tag{37.1}
\]

\[
 C=n_\eta n_--n_0(n_-+n_\eta).                                        \tag{37.2}
\]

Along the exact resonant flow, for any data,
\(\ddot n_0=2G^2[|\eta|^2|c_-|^2-|c_0|^2(|c_-|^2+|\eta|^2)]\) identically,
whose value at \(t=0\) is the \(\delta\to0\) limit \(G^2Ct^2\) of (37.1).
\end{theorem}

\begin{proof}
The second-order solution is computed by term-wise exact
integration; the phase average keeps the monomials with equal
powers of each amplitude and its conjugate; (37.1) is certified as
an exact polynomial identity, and the second-derivative identity is
certified on the equations of motion.
\end{proof}

\subsection{The kinetic equation}

\begin{theorem}[Resonance integral and rate]
\label{thm:v20v34-rate}
\(\int_{-\infty}^\infty(1-\cos\delta t)\,\delta^{-2}\,d\delta=\pi|t|\).
Hence for a family of triads with detunings of smooth density
\(\rho(\delta)\) about resonance, the phase-averaged occupations obey,
for \(1/\rho\)-scale times short compared with the depletion time,

\[
 \dot n_0=RC,\qquad \dot n_-=\dot n_\eta=-RC,\qquad R=2\pi G^2\rho(0),
                                                                    \tag{37.3}
\]

the wave-kinetic equation of the channel.
\end{theorem}

\begin{proof}
The integral is certified through the half-angle identity, the
change of variables, and the Dirichlet integral
\(\int\sin^2u\,u^{-2}du=\pi\); (37.3) is (37.1) summed over the family.
\end{proof}

\begin{theorem}[Conservation, \(H\)-theorem, equilibrium]
\label{thm:v20v34-H}
The system (37.3) conserves \(N=n_0+n_-\), \(M=n_\eta-n_-\), and, on the
resonance \(\omega_0=\omega_-+\Omega\), the energy
\(\omega_0n_0+\omega_-n_-+\Omega n_\eta\).  The entropy
\(S=\ln n_0+\ln n_-+\ln n_\eta\) satisfies

\[
 \dot S=R\,\frac{C^2}{n_0n_-n_\eta}\ \ge0,                             \tag{37.4}
\]

with equality if and only if \(1/n_0=1/n_-+1/n_\eta\), which holds for
the Rayleigh--Jeans distributions \(n=T/(\omega-\mu)\) with the
frequencies on the resonance and additive chemical potentials
\(\mu_0=\mu_-+\mu_\eta\).
\end{theorem}

\begin{proof}
All statements are certified symbolically.
\end{proof}

\begin{proposition}[Reading]
\label{prop:v20v34-reading}
The wave-kinetic equation is the classical Boltzmann equation of
the toy Standard Model's decays: the collision term (37.2) has the
stimulated structure \(n_\eta n_--n_0n_--n_0n_\eta\), the classical
limit of the quantum \((n_\eta+1)(n_-+1)n_0-n_\eta n_-(n_0+1)\)-type term
whose spontaneous piece the Jaynes--Cummings ladders of V26 and V28
supply; its \(H\)-theorem drives the occupations to the
Rayleigh--Jeans equilibrium, the classical thermal state.  The rate
\(R\) is quadratic in the Yukawa or gauge coupling through \(G\) of
(34.3) or (35.2) and proportional to the density of resonant
triads, the plane-symmetric form of the phase-space density in a
decay rate.
\end{proposition}

\begin{proof}
Theorems \ref{thm:v20v34-rate}--\ref{thm:v20v34-H} and the ladders
of V26 and V28.
\end{proof}

\subsection{Executable certificate}

The verifier

\[
 \texttt{scripts/verify\_sm\_v20\_v34\_wave\_kinetic.py}              \tag{37.5}
\]

certifies Theorems \ref{thm:v20v34-second-order}--\ref{thm:v20v34-H}.
It emits

\[
 \texttt{artifacts/sm\_v20\_v34\_wave\_kinetic.json}.                 \tag{37.6}
\]

The hashes are

\[
\begin{array}{c|l}
\hbox{object}&\hbox{SHA--256}\\ \hline
\hbox{verifier}&
\texttt{CC31A144DB70C4716685D181E23D955985BABD5F45DEAEAEB002404146771C0D}\\
\hbox{canonical payload}&
\texttt{0F9C9AAF2C0391F10776E619CCFA42928DD9910CC78FA34D87BDDFE16CD74444}.
\end{array}                                                     \tag{37.7}
\]

\begin{firewall}[Second order, random phases, continuum of detunings assumed]
\label{fire:v20v34}
The kinetic equation is derived at second order in the coupling
from random-phase data and requires a continuum of detunings with
smooth density; in the discrete plane framework the detunings are
discrete, so \(\rho\) is an assumption about the many-mode limit, not
a certified object.  The closure beyond second order, the
quantum collision term, and the coupling to the geometry are not
treated.
\end{firewall}

\begin{decision}[V35 acquisition]
\label{dec:v20v34-next}
V35 is the compulsory companion audit.  V36 then reads the
conversions gravitationally: the averaged Einstein equations of
V16--V17 and V23 with the occupations of the kinetic equation as
sources, the change of the equation of state as pairs convert into
Higgs quanta, and the Rayleigh--Jeans equilibrium as a radiation
universe of V96.
\end{decision}

\begin{assessment}[V34 acquisition]
\label{ass:v20v34}
The wave-kinetic equation (37.3) of the decay channels is derived
exactly at second order with its collision term (37.2), certified
conservation laws, the \(H\)-theorem (37.4), and the Rayleigh--Jeans
equilibrium: the first continuum object of the cycle.
\end{assessment}

\begin{record}[V34 append record]
\label{rec:v20v34}
V34 is appended to the validated V33 whose installed SHA--256 was
\texttt{BC8B4CB9ED17A7224B86F65631311456CEDE52F1CA9A376968E996D48950EC95}.
After installation only V34 is the active numeric SM note; frozen
archives and unrelated notes are unchanged.  The next compulsory
companion audit is V35 and the next sweep V40.
\end{record}


\section{V35: companion audit}
\label{sec:v20v35}

This is the required audit at the thirty-fifth version.  Every
active companion note present in the shared folder was inspected
read-only.

\subsection{Companion snapshot}

\begin{record}[V35 companion snapshot]
\label{rec:v20v35-snapshot}
The active companion files present were

\[
\begin{array}{c|r|r}
\hbox{file and SHA--256}&\hbox{bytes}&\hbox{lines}\\ \hline
\texttt{Renewal\_NS\_Stability\_Research\_Notes\_V26.tex}
&278283&5319\\
\multicolumn{3}{l}{
\texttt{82C887F8C2C69E4F28FAD7C3DC8DE5B9099CB5A4BF431EBA1FFF3E91935978AD}}
\\
\texttt{Renewal\_GRH\_Cofinal\_Visibility\_Attack\_V81.tex}
&777669&21553\\
\multicolumn{3}{l}{
\texttt{5B2A0570E0C894178887EB4BBF58C6B2B48F87EDB8E906C25B3BB8CE270D2986}}
\\
\texttt{Renewal\_Yang\_Mills\_Mass\_Gap\_Research\_Notes\_V41.tex}
&601237&14452\\
\multicolumn{3}{l}{
\texttt{2EEC02EB12520BAA0651AB07A171970809451FB0DEB83B596B084312F9EE4961}}
\\
\texttt{Renewal\_YM\_Parallel\_Research\_Notes\_V5.tex}
&54174&1351\\
\multicolumn{3}{l}{
\texttt{306F1BB84CFA8A8D47EA569EC17E9545F9867C8D32B548EC9D9C444B4F3C0512}}.
\end{array}                                                     \tag{38.1}
\]

Since V30 the GRH note advanced from V69 to V81 (a growing filter
and its Euler positivity, last-zero bounds, conjugate-pair factors,
the bump form and its Gaussian limit as an Appell symbol, escape
theorems for stacked pairs, the Gamma-ladder form of a block, and a
finite-\(M\) transfer problem); the NS, Yang--Mills, and YM parallel
notes are unchanged by hash.
\end{record}

\begin{decision}[V35 audit decision]
\label{dec:v20v35-audit}
No companion theorem, numerical constant, or executable artifact is
imported.  The GRH advances concern zeros of filtered polynomials
under a heat flow and have no object in common with the resonant
triads, the kinetic equation, or the plane-symmetric field systems
of this cycle.
\end{decision}

\begin{proof}
The only changed file is the GRH note; its new sections contain no
field system, no three-wave dynamics, and no kinetic equation.
\end{proof}

\begin{firewall}[Audit only]
\label{fire:v20v35}
No computation is added in V35.
\end{firewall}

\begin{decision}[V36 acquisition]
\label{dec:v20v35-next}
V36 reads the conversions gravitationally: the averaged energy
density and longitudinal and transverse pressures of every mode of
V31 (fermion) and of the Higgs modes, the mixture's equation of
state as a function of the occupations, its change under the
kinetic equation at fixed energy, and the exact power-law Bianchi--I
response to an anisotropic fluid with \(p_\perp=0\) and
\(p_z=w\rho\), with the null-dust and dust limits.
\end{decision}

\begin{assessment}[V35 acquisition]
\label{ass:v20v35}
The audit is current and nothing is imported.
\end{assessment}

\begin{record}[V35 append record]
\label{rec:v20v35}
V35 is appended to the validated V34 whose installed SHA--256 was
\texttt{914B522906F666E2C8AB2E805FDC6771B2923496E1D564BCDE26DEB91A024397}.
After installation only V35 is the active numeric SM note; the
companions, frozen archives, and all unrelated notes are unchanged.
The next compulsory companion audit is V40, together with the sweep
of V30--V39.
\end{record}


\section{V36: the gravitational reading of the conversions: equations of state of the modes, the kinetic mixture, and the Bianchi--I response}
\label{sec:v20v36}

V34 gave the kinetic equation of the decay channels.  This note
reads it gravitationally.  Every plane mode carries, after
averaging, an energy density, a longitudinal pressure and no
transverse pressure, with the universal equation of state
\(k^2/\omega_k^2\) for fermion modes and \(q^2/\Omega_q^2\) for Higgs
modes; a kinetic mixture of occupations therefore has an energy
density and a longitudinal pressure that are linear in the
occupations, and the collision term of V34 conserves the former and
changes the latter by a certified amount: conversions change the
equation of state at fixed energy.  The geometry's response to an
anisotropic fluid with \(p_\perp=0\) and \(p_z=w\rho\) is computed
exactly in the power-law class of Bianchi--I metrics, with the dust
limit the isotropic Einstein--de Sitter law and the \(w\to1\) limit
the vacuum Kasner metric, so that a null dust cannot be carried by
the power law, in agreement with the different null-dust metric of
V16.  The emergent geometry thus reads the decays as a drift of
\(w\) and responds by changing its expansion exponents.

\subsection{Equations of state of the modes}

\begin{theorem}[Averaged stresses]
\label{thm:v20v36-eos}
On the flat background, with the flat forms of (27.3): the fermion
pair mode \(\psi_k^+\) of (34.1) has, per unit norm on the circle,
\(\rho=\omega_k\), \(p_z=k^2/\omega_k\), \(p_\perp=0\); the Higgs mode
\(h=A\cos qz\cos\Omega_qt\) with the quadratic potential has
\(\langle\rho\rangle=c_qA^2\Omega_q^2/2\), \(\langle p_z\rangle=c_qA^2q^2/2\),
\(\langle p_\perp\rangle=0\), \(c_0=1\), \(c_{q\ge1}=\tfrac12\).  Hence

\[
 \frac{p_z}\rho=\frac{k^2}{\omega_k^2}\ \hbox{(fermion)},\qquad
 \frac{p_z}\rho=\frac{q^2}{\Omega_q^2}\ \hbox{(Higgs)},\qquad p_\perp=0.
                                                                    \tag{39.1}
\]
\end{theorem}

\begin{proof}
Certified for \(k=1,2,3\) and \(q=0,1,2\) by \(z\)- and \(t\)-averages.
\end{proof}

\begin{theorem}[Kinetic mixture]
\label{thm:v20v36-mixture}
For occupations \(n_i\) of modes with energies \(\omega_i\) per quantum
and wavenumbers \(k_i\), \(\rho=\sum n_i\omega_i\) and
\(p_z=\sum n_ik_i^2/\omega_i\).  Under the kinetic equation (37.3) on the
resonance \(\omega_0=\omega_-+\Omega\),

\[
 \dot\rho=0,\qquad \dot p_z=RC\Bigl(\frac{k_0^2}{\omega_0}-\frac{k_-^2}{\omega_-}-\frac{q^2}\Omega\Bigr).
                                                                    \tag{39.2}
\]

For the homogeneous-Higgs channel \(q=0\) the pressure of the
converted pair is lost: the mixture softens toward dust as pairs
become Higgs quanta at rest.
\end{theorem}

\begin{proof}
Certified symbolically.
\end{proof}

\subsection{The Bianchi--I response}

\begin{theorem}[Power-law response to an anisotropic fluid]
\label{thm:v20v36-kasner}
For \(p_\perp=0\), \(p_z=w\rho\), \(\Lambda=0\), the Einstein equations
(27.2) with \(b=0\) and homogeneous sources have the unique
power-law solution with nonzero density

\[
 a=\alpha\ln t,\quad c=\gamma\ln t,\quad \rho=\frac{\rho_0}{t^2},\qquad
 \alpha=\frac{2(1-w)}{w^2+3},\quad \gamma=\frac{2(1+w)}{w^2+3},\quad
 \rho_0=\frac{4(1-w)(w+3)}{(w^2+3)^2},                                   \tag{39.3}
\]

and the conservation law \(\dot\rho+\rho(2\dot a+\dot c)+w\rho\dot c=0\)
follows.  At \(w=0\) it is the isotropic Einstein--de Sitter law
\(\alpha=\gamma=\tfrac23\), \(\rho_0=\tfrac43\); as \(w\to1\) it
degenerates to the vacuum Kasner metric \(\alpha=0\), \(\gamma=1\),
\(\rho_0=0\), which satisfies the null-dust equations of V16 with
zero density: a null dust is not carried by the \(t^{-2}\) power law.
\end{theorem}

\begin{proof}
The three equations are solved exactly for \((\alpha,\gamma,\rho_0)\);
uniqueness among solutions with \(\rho_0\neq0\), the conservation
law, and both limits are certified.
\end{proof}

\begin{proposition}[Reading]
\label{prop:v20v36-reading}
The emergent geometry reads the decay channels through a single
number, the longitudinal equation of state \(w=p_z/\rho\) of the
kinetic mixture, which the collision term drifts according to (39.2)
at fixed energy; the expansion exponents respond through (39.3),
the transverse expansion \(\alpha\) slowing and the longitudinal
\(\gamma\) quickening as \(w\) grows, and the mixture in
Rayleigh--Jeans equilibrium fixes \(w\) by the occupations
\(T/(\omega-\mu)\) of the modes retained.  This is the first place in
the cycle where a statistical state of the matter, rather than a
single mode, determines the geometry.
\end{proposition}

\begin{proof}
Theorems \ref{thm:v20v36-mixture} and \ref{thm:v20v36-kasner} with
Theorem \ref{thm:v20v34-H}.
\end{proof}

\subsection{Executable certificate}

The verifier

\[
 \texttt{scripts/verify\_sm\_v20\_v36\_kinetic\_geometry.py}          \tag{39.4}
\]

certifies Theorems \ref{thm:v20v36-eos}--\ref{thm:v20v36-kasner}.
It emits

\[
 \texttt{artifacts/sm\_v20\_v36\_kinetic\_geometry.json}.             \tag{39.5}
\]

The hashes are

\[
\begin{array}{c|l}
\hbox{object}&\hbox{SHA--256}\\ \hline
\hbox{verifier}&
\texttt{174F10EFDE8CD25D284E2C59B62207C9EE6B15EE0EC47319F2F443251508EF1C}\\
\hbox{canonical payload}&
\texttt{3FE24585F45A04B04BF6E4659D30887B1D0179F5619130E2E4447EC57D4A3C42}.
\end{array}                                                     \tag{39.6}
\]

\begin{firewall}[Instantaneous sources, quadratic Higgs potential, power-law class]
\label{fire:v20v36}
The mixture's stresses are those of flat-space modes taken at an
instant; the redshift of the modes by the expansion, the
self-consistent evolution of the occupations in the expanding
metric, and the cubic and quartic Higgs terms are not included.  The
Bianchi--I response is the exact power-law class only, with
constant \(w\); a drifting \(w\) is read adiabatically.
\end{firewall}

\begin{decision}[V37 acquisition]
\label{dec:v20v36-next}
V37 removes the first limitation: the kinetic occupations in the
expanding Bianchi--I metric, the redshift of the fermion and Higgs
modes, the resulting drift of the resonance conditions, and the
combined system of kinetic equation and Einstein equations for a
single triad in the homogeneous averaged metric.
\end{decision}

\begin{assessment}[V36 acquisition]
\label{ass:v20v36}
The conversions are read gravitationally: the equations of state
(39.1) of all modes, the kinetic mixture's drift (39.2) at fixed
energy, and the exact power-law Bianchi--I response (39.3) with its
dust and Kasner limits.
\end{assessment}

\begin{record}[V36 append record]
\label{rec:v20v36}
V36 is appended to the validated V35 whose installed SHA--256 was
\texttt{9FA5A4318379EAD1C365A85D32771703E0087EB893F7423848CDF7D68BA6389D}.
After installation only V36 is the active numeric SM note; frozen
archives and unrelated notes are unchanged.
\end{record}


\section{V37: modes in the expanding metric: exact fermion two-level reduction, redshift, resonance drift, and Landau--Zener conversion}
\label{sec:v20v37}

V36 read the kinetic mixture with instantaneous flat-space
stresses.  This note places the modes in the expanding homogeneous
metric \(-dt^2+e^{2a}(dx^2+dy^2)+e^{2c}dz^2\) with \(a,c\) functions of
time.  The polarized Dirac mode of comoving wavenumber \(k\) reduces
exactly, with the conformal prefactor \(e^{-a-c/2}\) of V16, to a
two-level system whose off-diagonal element is the physical
wavenumber \(ke^{-c}\): the comoving norm is exactly conserved, the
instantaneous frequency is \(\sqrt{k^2e^{-2c}+m^2}\), and the only
coupling between the positive- and negative-frequency branches is
the non-adiabatic term \(m\dot\kappa/(2\omega^2)\), which vanishes for
massless fermions and for homogeneous modes: gravitational pair
creation in the plane framework is the Landau--Zener physics of this
two-level system.  Higgs modes redshift in the WKB sense with a
certified second-order residual.  The frequencies of a triad
redshift at different rates, and the drift of the detuning per unit
longitudinal expansion is exactly the pressure combination of V36:
a conversion that changes the equation of state also moves off
resonance.  A triad swept through resonance converts by the
Landau--Zener law, certified numerically.

\subsection{Fermion and Higgs modes in the expanding metric}

\begin{theorem}[Exact two-level reduction]
\label{thm:v20v37-twolevel}
With the curved Dirac operator of V16 on the homogeneous metric,
\(u=e^{-a-c/2}U(t)\cos kz\), \(w=e^{-a-c/2}iV(t)\sin kz\) satisfy the
Dirac equation if and only if

\[
 i\frac d{dt}\begin{pmatrix}U\\V\end{pmatrix}
 =\begin{pmatrix}m&\kappa\\ \kappa&-m\end{pmatrix}\begin{pmatrix}U\\V\end{pmatrix},\qquad
 \kappa=ke^{-c(t)},                                                    \tag{40.1}
\]

so \(|U|^2+|V|^2\) is exactly conserved.  The instantaneous
eigenvalues are \(\pm\omega\), \(\omega=\sqrt{\kappa^2+m^2}\), and between
the instantaneous eigenvectors \(|\pm\rangle\)

\[
 \langle-|\tfrac d{dt}|+\rangle=\frac{m\dot\kappa}{2\omega^2},\qquad
 \langle+|\tfrac d{dt}|+\rangle=0,                                     \tag{40.2}
\]

which vanishes for \(m=0\) and for \(k=0\).
\end{theorem}

\begin{proof}
The reduction is certified from the spin-connection Dirac operator;
the eigenvectors, (40.2), and its vanishing are certified.
\end{proof}

\begin{theorem}[Higgs WKB]
\label{thm:v20v37-wkb}
The Higgs mode \(\cos qz\) obeys
\(\ddot h+(2\dot a+\dot c)\dot h+(q^2e^{-2c}+m_h^2)h=0\), and
\(h=e^{-(2a+c)/2}\Omega^{-1/2}e^{-i\int\Omega dt}\), \(\Omega^2=q^2e^{-2c}+m_h^2\),
leaves the residual
\(\bigl[-\tfrac12(2a+c)\ddot{}-\tfrac14(2a+c)\dot{}^2+\tfrac34\dot\Omega^2/\Omega^2-\tfrac12\ddot\Omega/\Omega\bigr]h\),
of second order in the rates.
\end{theorem}

\begin{proof}
Certified symbolically.
\end{proof}

\subsection{Resonance drift and sweep}

\begin{theorem}[Drift of the detuning]
\label{thm:v20v37-drift}
With all frequencies redshifting, \(\omega_k=\sqrt{k^2e^{-2c}+m^2}\) and
\(\Omega_q=\sqrt{q^2e^{-2c}+m_h^2}\), the detuning
\(\delta=\Omega_q-\omega_{k_0}-\omega_{k_-}\) of a Yukawa triad obeys

\[
 \frac{d\delta}{dc}=e^{-2c}\Bigl(\frac{k_0^2}{\omega_{k_0}}+\frac{k_-^2}{\omega_{k_-}}-\frac{q^2}{\Omega_q}\Bigr),
                                                                    \tag{40.3}
\]

the pressure combination of (39.2): a conversion is stationary under
expansion exactly when it is pressure-neutral.
\end{theorem}

\begin{proof}
Certified symbolically.
\end{proof}

\begin{theorem}[Landau--Zener sweep]
\label{thm:v20v37-lz}
For the undepleted-boson two-level problem with linearly swept
detuning, \(i\dot c_0=\Omega_Rc_-\), \(i\dot c_-=\Omega_Rc_0-vtc_-\), the
survival of the pair mode across the resonance is
\(\exp(-2\pi\Omega_R^2/v)\).
\end{theorem}

\begin{proof}
Certified numerically at twenty digits for two parameter sets as
the adiabatic-to-adiabatic survival over \([-80,80]\), which
converges as \(T^{-3}\), to \(2\times10^{-6}\).
\end{proof}

\begin{proposition}[Reading]
\label{prop:v20v37-reading}
In the expanding plane universe every fermion mode is a two-level
system (40.1) driven by the longitudinal scale factor; the geometry
creates pairs only through the massive, inhomogeneous coupling
(40.2), and the massless fermions of V16 are exactly conformal.  A
decay channel that softens the equation of state drifts off
resonance at the rate (40.3), so its total yield per expansion
e-fold is a Landau--Zener yield governed by \(\Omega_R^2/v\), with
\(v\) proportional to the pressure lost in the conversion: the
expansion self-limits the conversions that change it.
\end{proposition}

\begin{proof}
Theorems \ref{thm:v20v37-twolevel}--\ref{thm:v20v37-lz}.
\end{proof}

\subsection{Executable certificate}

The verifier

\[
 \texttt{scripts/verify\_sm\_v20\_v37\_expanding\_modes.py}           \tag{40.4}
\]

certifies Theorems \ref{thm:v20v37-twolevel}--\ref{thm:v20v37-lz}.
It emits

\[
 \texttt{artifacts/sm\_v20\_v37\_expanding\_modes.json}.              \tag{40.5}
\]

The hashes are

\[
\begin{array}{c|l}
\hbox{object}&\hbox{SHA--256}\\ \hline
\hbox{verifier}&
\texttt{40CCEC55C67F491FD01C9EFF870C72CB0C43FE7A56F2B7C4D309BC098722ABA2}\\
\hbox{canonical payload}&
\texttt{08CBE0AB87BE4CB7F1D103DE2C5243B934E302FABA4BCEFDCBD4EDA2A62BBEE1}.
\end{array}                                                     \tag{40.6}
\]

\begin{firewall}[Homogeneous metric, prescribed expansion, numerical Landau--Zener]
\label{fire:v20v37}
The metric is homogeneous and prescribed; the back-reaction of the
modes on \(a,c\) and the self-consistent kinetic--Einstein system are
not solved.  The Landau--Zener law is certified numerically, not
symbolically, and only for the undepleted-boson two-level
truncation.  Identities with nested radicals are certified at random
rational points where exact simplification fails.
\end{firewall}

\begin{decision}[V38 acquisition]
\label{dec:v20v37-next}
V38 extends the linear cross polarization of V22 to the exact
two-polarization plane metric: the transverse metric as a unimodular
matrix \(e^{2a}M(P,Q)\), the Einstein equations for \((a,c,P,Q)\),
the identification of the two polarization equations as the
hyperbolic-plane sigma model on the \((t,z)\) base weighted by the
transverse volume, the reduction to (27.2) at \(Q=0\), and the
matter source of the cross polarization.
\end{decision}

\begin{assessment}[V37 acquisition]
\label{ass:v20v37}
Modes in the expanding metric are reduced exactly (40.1), their
branch coupling (40.2) identifies gravitational pair creation with
Landau--Zener transitions, the Higgs redshift is certified in WKB
form, the detuning drift (40.3) is the pressure combination of V36,
and the sweep yield is the Landau--Zener law.
\end{assessment}

\begin{record}[V37 append record]
\label{rec:v20v37}
V37 is appended to the validated V36 whose installed SHA--256 was
\texttt{5F256233C2065F77D38117CF6731EA369637034A1469942EFEC23D29C039F281}.
After installation only V37 is the active numeric SM note; frozen
archives and unrelated notes are unchanged.
\end{record}


\section{V38: the exact two-polarization plane metric: the polarizations as a hyperbolic-plane sigma model}
\label{sec:v20v38}

V22 treated the cross polarization of the gravitational wave only
at linear order on the flat background.  This note gives the exact
framework.  The transverse metric is written as \(e^{2a}M\) with the
unimodular matrix \(M(P,Q)\), so that \(Q=0\), \(P=2b\) is the diagonal
metric of the whole series and \(g_{xy}=e^{2a+P}Q\) is the cross
polarization.  The Einstein tensor is computed exactly for
\((a,c,P,Q)\) and its traceless transverse part is certified to be an
invertible combination, with coefficients depending only on
\((P,Q)\), of two wave equations \(\mathrm{Eq}_P\) and \(\mathrm{Eq}_Q\):
the plus and cross polarizations obey the equations of a sigma model
whose target is the hyperbolic plane \(dP^2+e^{2P}dQ^2\), on the
\((t,z)\) base, weighted by the transverse volume \(e^{2a+c}\) and
the longitudinal metric.  At \(Q=0\) the whole system reduces to the
vacuum forms of (27.2), and on the flat background the linear
\(Q\)-equation is that of V22.  The gravitational polarizations of
the plane framework are therefore an exactly known nonlinear
system, the classical chiral-field form of plane gravitational
waves, into which the toy Standard Model's traceless transverse
stress enters as a source.

\subsection{The metric and the polarization equations}

\begin{theorem}[Two-polarization metric and traceless transverse Einstein equations]
\label{thm:v20v38-metric}
Let \(g=\operatorname{diag}(-1,e^{2a}M,e^{2c})\) with

\[
 M=\begin{pmatrix}e^{P}&e^{P}Q\\ e^{P}Q&e^{P}Q^2+e^{-P}\end{pmatrix},\qquad \det M=1,\qquad \sqrt{-g}=e^{2a+c}.
                                                                    \tag{41.1}
\]

With \(G^i{}_j=e^{-2a}(M^{-1}G_{\rm tr})^i{}_j\) the mixed transverse
Einstein tensor, \(X=\tfrac12(G^x{}_x-G^y{}_y)\), \(Y=G^x{}_y\), and

\[
\begin{aligned}
 \mathrm{Eq}_P&=P_{tt}+P_t(2a_t+c_t)-e^{-2c}\bigl(P_{zz}+P_z(2a_z-c_z)\bigr)-e^{2P}\bigl(Q_t^2-e^{-2c}Q_z^2\bigr),\\
 \mathrm{Eq}_Q&=Q_{tt}+Q_t(2a_t+c_t)-e^{-2c}\bigl(Q_{zz}+Q_z(2a_z-c_z)\bigr)+2\bigl(P_tQ_t-e^{-2c}P_zQ_z\bigr),
\end{aligned}                                                          \tag{41.2}
\]

one has identically

\[
 X=\tfrac12\mathrm{Eq}_P-\tfrac12e^{2P}Q\,\mathrm{Eq}_Q,\qquad
 Y=Q\,\mathrm{Eq}_P+\tfrac12\bigl(1-e^{2P}Q^2\bigr)\mathrm{Eq}_Q,        \tag{41.3}
\]

with determinant \(\tfrac14(1+e^{2P}Q^2)>0\).  Hence the traceless
transverse Einstein equations are equivalent to
\(\mathrm{Eq}_P=S_P\), \(\mathrm{Eq}_Q=S_Q\), where \((S_P,S_Q)\) is the
inverse combination of the mixed traceless stress.
\end{theorem}

\begin{proof}
The determinant, the volume, and (41.3) as an exact identity in
\((a,c,P,Q)\) are certified; the coefficients are read off the second
time derivatives and certified to depend only on \((P,Q)\).
\end{proof}

\begin{theorem}[Sigma-model form]
\label{thm:v20v38-sigma}
\(\mathrm{Eq}_P\) and \(\mathrm{Eq}_Q\) are the Euler--Lagrange
equations of

\[
 L=e^{2a+c}\Bigl[-\bigl(P_t^2+e^{2P}Q_t^2\bigr)+e^{-2c}\bigl(P_z^2+e^{2P}Q_z^2\bigr)\Bigr],
                                                                    \tag{41.4}
\]

the sigma model with target metric \(dP^2+e^{2P}dQ^2\), the
hyperbolic plane of curvature \(-1\), on the \((t,z)\) base with the
weight \(e^{2a+c}\) and the longitudinal factor \(e^{-2c}\).
\end{theorem}

\begin{proof}
Certified by comparing the Euler--Lagrange equations with (41.2) up
to the normalizations \(2e^{2a+c}\) and \(2e^{2a+c+2P}\).
\end{proof}

\subsection{Reductions}

\begin{theorem}[Diagonal and linear limits]
\label{thm:v20v38-limits}
At \(Q=0\), \(P=2b\): \(\tfrac12\mathrm{Eq}_P\) is the polarization
equation of (27.2) with \(\pi=0\), \(G_{xy}=0\), and the constraints
and the \(a\)- and \(c\)-equations reduce to the vacuum forms of
(27.2).  On the flat background \(a=c=P=0\), \(\mathrm{Eq}_Q=Q_{tt}-Q_{zz}\)
and \(g_{xy}=Q\), the linear cross polarization of V22.
\end{theorem}

\begin{proof}
Certified symbolically.
\end{proof}

\begin{proposition}[Reading]
\label{prop:v20v38-reading}
The two polarizations of the plane gravitational wave form a
nonlinear but exactly structured system: geodesic motion in the
hyperbolic plane driven by the base.  The plus polarization alone is
a free wave in the diagonal framework; the cross polarization
couples to it through the curvature of the target, and at second
order a cross wave sources the plus polarization through
\(-e^{2P}(Q_t^2-e^{-2c}Q_z^2)\).  The mass-split doublet of V22, which
drives \(Q\) at the beat frequency, therefore also drives \(P\) at
twice the beat frequency through the target curvature, and the
degenerate doublet, which drives only \(P\), leaves \(Q\) free.
\end{proposition}

\begin{proof}
Theorems \ref{thm:v20v38-metric}--\ref{thm:v20v38-limits}; the
second-order statement is read off (41.2).
\end{proof}

\subsection{Executable certificate}

The verifier

\[
 \texttt{scripts/verify\_sm\_v20\_v38\_two\_polarizations.py}         \tag{41.5}
\]

certifies Theorems \ref{thm:v20v38-metric}--\ref{thm:v20v38-limits}.
It emits

\[
 \texttt{artifacts/sm\_v20\_v38\_two\_polarizations.json}.            \tag{41.6}
\]

The hashes are

\[
\begin{array}{c|l}
\hbox{object}&\hbox{SHA--256}\\ \hline
\hbox{verifier}&
\texttt{58B5EE8A4A7BFA4D5506A141E1259E4202173E9F270C49D6F390B4ED1D5484D9}\\
\hbox{canonical payload}&
\texttt{1ADF8302C93E5AE3BECEEEFF4BD58DFE8A9272A796C06AE382CEA6B421FBB2D3}.
\end{array}                                                     \tag{41.7}
\]

\begin{firewall}[Vacuum polarization equations; sources not yet computed]
\label{fire:v20v38}
The matter sources \((S_P,S_Q)\) are defined by (41.3) but not
computed for any field in this note; the Dirac field in the
non-diagonal transverse metric needs a tetrad adapted to \(M\), which
is not constructed.  The remaining Einstein equations with \(Q\neq0\)
are not displayed.
\end{firewall}

\begin{decision}[V39 acquisition]
\label{dec:v20v38-next}
V39 computes the sources: the traceless transverse stress of the
colour wave and of the Higgs in the two-polarization metric, the
resulting \((S_P,S_Q)\), the reduction to \(-2\pi\) of (3.1) at
\(Q=0\), and the general form of the \(Q\)-source in terms of the
mixed stress, with the split doublet's \(T_{xy}\) of (25.2) as the
linear seed; then V40 is the audit and sweep of V30--V39.
\end{decision}

\begin{assessment}[V38 acquisition]
\label{ass:v20v38}
The exact two-polarization plane metric (41.1) is placed in the
framework, its polarization equations (41.2) are certified to be the
traceless transverse Einstein equations (41.3) and the
hyperbolic-plane sigma model (41.4), with the diagonal and linear
limits of the series recovered.
\end{assessment}

\begin{record}[V38 append record]
\label{rec:v20v38}
V38 is appended to the validated V37 whose installed SHA--256 was
\texttt{D4E8D9124B3BBA09492C03E31AE8DE138F04F4B80BD154A488CF067C1F6EA1D9}.
After installation only V38 is the active numeric SM note; frozen
archives and unrelated notes are unchanged.
\end{record}


\section{V39: the matter sources of the two gravitational polarizations}
\label{sec:v20v39}

V38 identified the two polarizations with a hyperbolic-plane sigma
model whose traceless transverse Einstein equations are
\(\mathrm{Eq}_P=S_P\), \(\mathrm{Eq}_Q=S_Q\).  This note computes the
sources.  The source map inverts (41.3): \((S_P,S_Q)\) is a rational
function of \((P,Q)\) applied to the mixed traceless stress
\((X_T,Y_T)\).  For the colour wave the sources are certified in closed
form: the colour anisotropy \(\Pi_f\) drives the plus polarization
with the factor \(1-e^{2P}Q^2\) and the cross polarization with the
factor \(2Q\), so a single colour cannot seed a cross polarization
but amplifies an existing one.  The Higgs has an isotropic
transverse stress and sources neither polarization.  For a general
transverse stress the sources at \(Q=0\) are \(-2\pi\) of (3.1) and
\(2e^{-2a-P}T_{xy}\), and on the flat background the linear
\(Q\)-source is \(2T_{xy}\), the law (25.4) of V22: the split
doublet's off-diagonal stress is the only seed of the cross
polarization in the toy Standard Model, and once seeded the cross
polarization is amplified by the colour anisotropy and fed back to
the plus polarization through the target curvature.

\subsection{The source map and the fields}

\begin{theorem}[Sources of the polarization equations]
\label{thm:v20v39-map}
With \(T^i{}_j=e^{-2a}(M^{-1}T_{\rm tr})^i{}_j\), \(X_T=\tfrac12(T^x{}_x-T^y{}_y)\),
\(Y_T=T^x{}_y\), the Einstein equations give \(\mathrm{Eq}_P=S_P\),
\(\mathrm{Eq}_Q=S_Q\) with

\[
 S_P=\frac{2\bigl[(1-e^{2P}Q^2)X_T+e^{2P}QY_T\bigr]}{1+e^{2P}Q^2},\qquad
 S_Q=\frac{2\bigl[Y_T-2QX_T\bigr]}{1+e^{2P}Q^2}.                         \tag{42.1}
\]

For the colour wave \(A^1=f\,dx\), with
\(\Pi_f=e^{-2a-P}\bigl(f_t^2-e^{-2c}f_z^2\bigr)\),

\[
 S_P=-\bigl(1-e^{2P}Q^2\bigr)\Pi_f,\qquad S_Q=2Q\,\Pi_f,               \tag{42.2}
\]

reducing at \(Q=0\), \(P=2b\) to \(S_P=-2\pi\) with \(\pi\) of (3.1) and
\(S_Q=0\).  For the Higgs with the quadratic potential
\(X_T=Y_T=0\) identically.
\end{theorem}

\begin{proof}
The inverse map, the colour stress in the two-polarization metric,
(42.2), the \(Q=0\) reduction, and the vanishing of the Higgs
sources are certified symbolically.
\end{proof}

\begin{theorem}[General transverse stress]
\label{thm:v20v39-general}
For arbitrary \(T_{xx},T_{xy},T_{yy}\), at \(Q=0\)

\[
 S_P=e^{-2a}\bigl(e^{-P}T_{xx}-e^{P}T_{yy}\bigr)=W_-T_{xx}-W_+T_{yy}=-2\pi,\qquad
 S_Q=2e^{-2a-P}T_{xy},                                                 \tag{42.3}
\]

and on the flat background \(S_Q=2T_{xy}+O(Q)\), which is (25.4).
\end{theorem}

\begin{proof}
Certified symbolically.
\end{proof}

\begin{proposition}[Seeding and amplification of the cross polarization]
\label{prop:v20v39-reading}
In the toy Standard Model of V24 extended by the split doublet of
V22, the only off-diagonal transverse stress is
\(T_{xy}=-\tfrac12efS\) of (25.2), so the cross polarization is seeded
only by a mass-split doublet in a colour field, at the beat
frequency; once present it is amplified by the colour anisotropy
through \(2Q\Pi_f\) and feeds the plus polarization through
\(-e^{2P}(Q_t^2-e^{-2c}Q_z^2)\) of (41.2) and through the factor
\(1-e^{2P}Q^2\) in (42.2).  The Higgs is polarization-blind.  The
polarization content of the plane gravitational wave thus records
the isospin structure of the matter: degenerate doublets and
scalars leave the wave linearly polarized.
\end{proposition}

\begin{proof}
Theorems \ref{thm:v20v39-map}--\ref{thm:v20v39-general} with (25.2)
and (41.2).
\end{proof}

\subsection{Executable certificate}

The verifier

\[
 \texttt{scripts/verify\_sm\_v20\_v39\_polarization\_sources.py}      \tag{42.4}
\]

certifies Theorems \ref{thm:v20v39-map} and \ref{thm:v20v39-general}.
It emits

\[
 \texttt{artifacts/sm\_v20\_v39\_polarization\_sources.json}.         \tag{42.5}
\]

The hashes are

\[
\begin{array}{c|l}
\hbox{object}&\hbox{SHA--256}\\ \hline
\hbox{verifier}&
\texttt{7560A45AA4237D45AC87785575F92150ED4A1A0551805C66F53CA1E8D2632BE3}\\
\hbox{canonical payload}&
\texttt{6CBF287D1F7EFA1809809ABA4A01D35E5835270D8FAC829C255D9913918707DF}.
\end{array}                                                     \tag{42.6}
\]

\begin{firewall}[Bosonic stresses only in the two-polarization metric]
\label{fire:v20v39}
The colour and Higgs stresses are computed exactly in the
two-polarization metric; the fermion stress there requires the
adapted tetrad and is used only through its flat and \(Q=0\) forms.
The remaining Einstein equations with \(Q\neq0\) and their sources
are not displayed.
\end{firewall}

\begin{decision}[V40 acquisition]
\label{dec:v20v39-next}
V40 is the compulsory companion audit and the fourth ten-version
sweep of the cycle, recording V30--V39 and fixing the direction of
V41--V49.
\end{decision}

\begin{assessment}[V39 acquisition]
\label{ass:v20v39}
The sources (42.1)--(42.3) of the two polarizations are certified:
the colour wave drives both once a cross polarization exists, the
Higgs drives neither, and the split doublet's off-diagonal stress is
the only seed of the cross polarization.
\end{assessment}

\begin{record}[V39 append record]
\label{rec:v20v39}
V39 is appended to the validated V38 whose installed SHA--256 was
\texttt{DC610B795C756042DE4C20B1EB21A30F83B71DCAE0A7760565D5EFE4388DDA53}.
After installation only V39 is the active numeric SM note; frozen
archives and unrelated notes are unchanged.
\end{record}


\section{V40: companion audit and the fourth ten-version sweep of the twentieth cycle}
\label{sec:v20v40}

This is the required audit at the fortieth version and the sweep of
V30--V39.  Every active companion note present in the shared folder
was inspected read-only.

\subsection{Companion snapshot}

\begin{record}[V40 companion snapshot]
\label{rec:v20v40-snapshot}
The active companion files present were

\[
\begin{array}{c|r|r}
\hbox{file and SHA--256}&\hbox{bytes}&\hbox{lines}\\ \hline
\texttt{Renewal\_NS\_Stability\_Research\_Notes\_V26.tex}
&278283&5319\\
\multicolumn{3}{l}{
\texttt{82C887F8C2C69E4F28FAD7C3DC8DE5B9099CB5A4BF431EBA1FFF3E91935978AD}}
\\
\texttt{Renewal\_GRH\_Cofinal\_Visibility\_Attack\_V85.tex}
&816931&22378\\
\multicolumn{3}{l}{
\texttt{7FF60585A8BFD46822FF1015F0BFFE449DC4123449F2E67A3AD3CEE3BE40ADEF}}
\\
\texttt{Renewal\_Yang\_Mills\_Mass\_Gap\_Research\_Notes\_V41.tex}
&601237&14452\\
\multicolumn{3}{l}{
\texttt{2EEC02EB12520BAA0651AB07A171970809451FB0DEB83B596B084312F9EE4961}}
\\
\texttt{Renewal\_YM\_Parallel\_Research\_Notes\_V5.tex}
&54174&1351\\
\multicolumn{3}{l}{
\texttt{306F1BB84CFA8A8D47EA569EC17E9545F9867C8D32B548EC9D9C444B4F3C0512}}.
\end{array}                                                     \tag{43.1}
\]

Since V35 the GRH note advanced from V81 to V85 (a saddle-point
form of an escape estimate uniform in a size parameter, exact escape
certificates to finite sizes, a verdict on a filter family, and a
real-factor zero law); the NS, Yang--Mills, and YM parallel notes are
unchanged by hash.
\end{record}

\begin{decision}[V40 audit decision]
\label{dec:v20v40-audit}
No companion theorem, numerical constant, or executable artifact is
imported.  The GRH advances concern zero-escape estimates for
filtered polynomials and have no object in common with the kinetic,
gravitational, or polarization results of this block.
\end{decision}

\begin{proof}
The only changed file is the GRH note; its new sections contain no
field system, no kinetic equation, and no plane metric.
\end{proof}

\subsection{Sweep of V30--V39}

\begin{proposition}[What V30--V39 established]
\label{prop:v20v40-sweep}
The fourth block of the cycle passed from single resonant modes to
the continuum objects and the exact geometry:

\begin{enumerate}
\item the resonant triads of both decay channels for all
wavenumbers: modes, selection rules with two or three branches,
frequency conditions, couplings in closed form, and the resonance
surfaces with the marginal equal-wavenumber triads (V31, V32);
\item the many-triad Hamiltonian and the smallest shared-mode
system: its \(U(1)^3\) symmetry, exactly three quadratic invariants,
and no further polynomial invariant of degree at most four, so that
the integrability of single triads does not extend to shared modes
(V33);
\item the wave-kinetic equation of the decay channels: second-order
random-phase solution, the exact second-derivative identity, the
resonance integral, the collision term
\(n_\eta n_--n_0(n_-+n_\eta)\), its conservation laws, the
\(H\)-theorem, and the Rayleigh--Jeans equilibrium (V34);
\item the audit (V35);
\item the gravitational reading of the conversions: the equations
of state \(k^2/\omega^2\) and \(q^2/\Omega^2\) of all modes, the kinetic
mixture's pressure drift at fixed energy, and the exact power-law
Bianchi--I response with its dust and Kasner limits (V36);
\item modes in the expanding metric: the exact fermion two-level
reduction with the physical wavenumber \(ke^{-c}\), the branch
coupling \(m\dot\kappa/(2\omega^2)\) as gravitational pair creation,
the Higgs WKB residual, the detuning drift equal to the pressure
combination, and the Landau--Zener sweep law (V37);
\item the exact two-polarization plane metric: the polarizations as
the hyperbolic-plane sigma model, the traceless transverse Einstein
equations as an invertible combination of the two wave equations,
and the diagonal and linear reductions (V38);
\item the matter sources of the polarizations: the colour wave's
\(-(1-e^{2P}Q^2)\Pi_f\) and \(2Q\Pi_f\), the polarization-blind Higgs,
the general \(Q=0\) sources \(-2\pi\) and \(2e^{-2a-P}T_{xy}\), and the
split doublet as the only seed of the cross polarization (V39).
\end{enumerate}
\end{proposition}

\begin{proof}
Items are the certified theorems of the cited sections.
\end{proof}

\begin{proposition}[Direction for V41--V49]
\label{prop:v20v40-direction}
The kinetic equation of V34 contains the continuum in one number,
the density \(\rho(0)\) of resonant triads, which in the discrete
plane framework is an assumption.  The upstream step is to compute
it: V41 rewrites the modes, couplings and triads on a circle of
length \(L\), with the certified scalings of the couplings in \(L\);
V42 takes \(L\to\infty\) and obtains the density of resonant triads
and the \(L\)-independent rate \(2\pi G^2\rho(0)\) of the Yukawa
channel as a closed function of \(y\), \(m\), \(m_h\), the
plane-symmetric decay width of the Higgs into a fermion pair derived
from the exact mode dynamics; V43 does the same for the colour
channel; V44 supplies the quantum spontaneous rate from the
Jaynes--Cummings ladders with the same density; V45 is the audit;
V46--V49 close the kinetic--Einstein system: the Boltzmann equation
with redshift in the Bianchi--I metric, its thermal state, and the
self-consistent expansion law, before the V50 sweep.
\end{proposition}

\begin{proof}
This is a decision.
\end{proof}

\begin{firewall}[Audit and sweep only]
\label{fire:v20v40}
No computation is added in V40.  The cycle's results are exact
statements about plane-symmetric classical fields, their resonant
reductions and kinetic limits under stated hypotheses, and model
quantizations of those reductions; no continuum limit in the sense
of a field theory, no quantum field theory of any sector, and no
Standard Model--Einstein continuum are claimed.
\end{firewall}

\begin{decision}[V41 acquisition]
\label{dec:v20v40-next}
V41 places the Yukawa channel on the circle of length \(L\):
wavenumbers \(2\pi n/L\), unit-normalized modes, the couplings
\(g\) and \(\kappa\) of (34.3)--(34.4) with their certified
\(L\)-scalings, and the resonant triads as integer solutions of the
frequency condition at finite \(L\).
\end{decision}

\begin{assessment}[V40 acquisition]
\label{ass:v20v40}
The audit is current and the sweep records the passage from single
modes to kinetic and geometric continuum objects: the resonant
triad lists, the wave-kinetic equation with its \(H\)-theorem, the
gravitational reading of conversions, the expanding-metric modes,
and the exact two-polarization geometry with its sources.  The
direction is the infinite-volume density of triads and the decay
widths.
\end{assessment}

\begin{record}[V40 append record]
\label{rec:v20v40}
V40 is appended to the validated V39 whose installed SHA--256 was
\texttt{071B55878A53B3F1E1F1A4BF1F9CD8426BE9F51D0247AC7D442A0873147269AB}.
After installation only V40 is the active numeric SM note; the
companions, frozen archives, and all unrelated notes are unchanged.
The next compulsory companion audit is V45 and the next sweep V50.
\end{record}


\section{V41: the Yukawa channel on a circle of length \texorpdfstring{$L$}{L}: scalings, exact finite-size resonances, and the density of triads}
\label{sec:v20v41}

V34 left the density of resonant triads as an assumption.  This
note begins its computation by placing the Yukawa channel on a
circle of length \(L\).  The modes, matrix elements and couplings of
V31 are recomputed with wavenumbers \(2\pi n/L\): the fermion-side
coupling is independent of \(L\) at fixed physical wavenumbers, the
boson-side coupling scales as \(1/L\), so the symmetric coupling
scales as \(L^{-1/2}\), the standard volume law.  Exact resonances
at finite \(L\) are exhibited: the marginal family at \(m_h=2m\) for
every \(L\), and, for massless fermions, the difference branch
resonant exactly when \(m_hL/4\pi\) is the square root of a product
of two integers.  The detuning as a function of the partner
wavenumber has, on both branches, the derivative equal to the
relative velocity of the boson and the partner, so the number of
triads per unit detuning is \(L/(2\pi|v_q-v_{\rm phys}|)\): the
one-dimensional phase-space factor emerges from the mode counting.

\subsection{Modes and couplings at finite \texorpdfstring{$L$}{L}}

\begin{theorem}[Scalings]
\label{thm:v20v41-scalings}
On \([0,L]\) with \(k_n=2\pi n/L\), \(q_j=2\pi j/L\), the modes of (34.1)
have norm \(L\omega_nW_n/k_n^2\), the matrix element of
\(\cos q_jz\operatorname{diag}(1,-1)\) between \(\psi_{n_-}^-\) and
\(\psi_{n_0}^+\) is \(\tfrac L4(W_-/k_-+sW_0/k_0)\) for \(j=|n_0-n_-|\)
(\(s=+1\)) or \(j=n_0+n_-\) (\(s=-1\)), \(LW/k\) for \(j=0\), and zero
otherwise; with unit-normalized modes the rotating-wave couplings
are

\[
 g=\frac m{8v}\,\frac{W_-k_0+sW_0k_-}{\sqrt{\omega_0W_0\omega_-W_-}}\ (j\ge1),\qquad
 g=\frac{mk}{2v\omega_k}\ (j=0),\qquad
 \frac\kappa g=\frac{8\varepsilon^2}{L\Omega_j}\ (j\ge1),\quad \frac\kappa g=\frac{4\varepsilon^2}{L\Omega_0}\ (j=0),
                                                                    \tag{44.1}
\]

so \(g\) is \(L\)-independent at fixed physical wavenumbers and
\(G=\sqrt{\kappa g}\propto L^{-1/2}\); \(L=2\pi\) reproduces
(34.3)--(34.4).
\end{theorem}

\begin{proof}
Certified on concrete triads through the substitution
\(z=L\theta/2\pi\).
\end{proof}

\subsection{Exact resonances and the triad density}

\begin{theorem}[Finite-size resonances]
\label{thm:v20v41-resonances}
The marginal family \(j=2n\), \(n_0=n_-=n\) is resonant if and only if
\(m_h=2m\), for every \(L\).  For massless fermions the difference
branch \(j=n_0-n_-\) is resonant if and only if

\[
 m_h=\frac{4\pi\sqrt{n_0n_-}}L.                                       \tag{44.2}
\]
\end{theorem}

\begin{proof}
Certified from \(\Omega_j^2-(\omega_{n_0}+\omega_{n_-})^2\).
\end{proof}

\begin{theorem}[Density of triads]
\label{thm:v20v41-density}
At fixed \(k_0\), on the branch \(q=k_0-sk_-\), the detuning
\(\delta(k_-)=\Omega_q-\omega_{k_0}-\omega_{k_-}\) satisfies

\[
 \frac{d\delta}{dk_-}=-sv_q-v_-,\qquad
 \Bigl|\frac{d\delta}{dk_-}\Bigr|=|v_q-v_{\rm phys}|,\qquad
 v_q=\frac q{\Omega_q},\quad v_{\rm phys}=-s\frac{k_-}{\omega_{k_-}},        \tag{44.3}
\]

so the number of resonant triads per unit detuning is
\(\rho=L/(2\pi|v_q-v_{\rm phys}|)\), with \(v_{\rm phys}\) the physical
velocity of the partner mode, of momentum \(-sk_-\).
\end{theorem}

\begin{proof}
Certified on both branches.
\end{proof}

\begin{proposition}[Reading]
\label{prop:v20v41-reading}
The two branches of the selection rule are the two kinematic
regimes of a two-body decay in one dimension: partner momentum
opposite to (\(s=+1\)) or along (\(s=-1\)) the fermion's, and the
density (44.3) is the inverse relative velocity, the one-dimensional
phase-space factor.  Since \(G^2\propto1/L\) and \(\rho\propto L\), the
kinetic rate \(2\pi G^2\rho\) of V34 has an \(L\)-independent limit,
computed in V42.
\end{proposition}

\begin{proof}
Theorems \ref{thm:v20v41-scalings} and \ref{thm:v20v41-density}.
\end{proof}

\subsection{Executable certificate}

The verifier

\[
 \texttt{scripts/verify\_sm\_v20\_v41\_circle\_length.py}             \tag{44.4}
\]

certifies Theorems \ref{thm:v20v41-scalings}--\ref{thm:v20v41-density}.
It emits

\[
 \texttt{artifacts/sm\_v20\_v41\_circle\_length.json}.                \tag{44.5}
\]

The hashes are

\[
\begin{array}{c|l}
\hbox{object}&\hbox{SHA--256}\\ \hline
\hbox{verifier}&
\texttt{02CF364802400927C00DE5AEAB00ECAB85A692489F680CDC40EAEB388C26960D}\\
\hbox{canonical payload}&
\texttt{1632AE0C5CD006A768DA1B94F2A0458290F3124B748F6E73F0D7DBB243379661}.
\end{array}                                                     \tag{44.6}
\]

\begin{firewall}[Cosine family, concrete triads, no limit taken]
\label{fire:v20v41}
As in V31 the general formulas are certified on concrete triads and
the selection rule symbolically; the limit \(L\to\infty\) is not
taken here, and the one-quantum normalization of \(\varepsilon\) is
not fixed.
\end{firewall}

\begin{decision}[V42 acquisition]
\label{dec:v20v41-next}
V42 takes \(L\to\infty\): the rate \(2\pi G^2\rho(0)\) as an
\(L\)-independent closed function of \(y\), \(m\), \(m_h\) and the
boson momentum on both branches, its rest-frame value
\(m^2k_*/(v^2m_h^2)\) with \(k_*=\sqrt{m_h^2/4-m^2}\), its threshold
behaviour, and the test of Lorentz covariance: whether the rate of a
moving boson equals the rest-frame rate times \(m_h/\Omega_q\).
\end{decision}

\begin{assessment}[V41 acquisition]
\label{ass:v20v41}
The Yukawa channel at finite \(L\) is certified: the scalings (44.1),
the exact resonances (44.2), and the triad density (44.3) with the
inverse relative velocity as the one-dimensional phase-space factor.
\end{assessment}

\begin{record}[V41 append record]
\label{rec:v20v41}
V41 is appended to the validated V40 whose installed SHA--256 was
\texttt{8832837386A87DB3DBD3B3531589272A4D02C4F0BBC7804B107CD066E7B4277B}.
After installation only V41 is the active numeric SM note; frozen
archives and unrelated notes are unchanged.
\end{record}



\section{V42: fixed-parent phase space, the order of limits, and a normalized Yukawa benchmark}
\label{sec:v20v42}

The latest source is V41 of the twentieth cycle, not V31 of the
nineteenth cycle last present in this task.  Its V40 sweep and V41
couplings set the present direction.  We retain their results and go
upstream at one essential distinction: a density computed with one
daughter momentum fixed is not the phase space for decay of a parent
with fixed momentum.  V41 also explicitly leaves one-quantum
normalization open.  Accordingly this section proves the fixed-parent
continuum density and a normalized benchmark, without identifying the
classical cosine amplitudes with quantum external states.

\subsection{Which quantity is held fixed}

Set \(M=m_h>2m>0\), \(\Omega=\sqrt{M^2+q^2}\),
\(E(p)=\sqrt{m^2+p^2}\), and fix the signed parent momentum \(q\).
Define
\[
 D_q(p)=\Omega-E(p)-E(q-p),\qquad
 k_*=\sqrt{M^2/4-m^2},\qquad
 \beta=q/\Omega,\quad \gamma=\Omega/M .
\]
The two daughters are labeled (a fermion and an antifermion); there
is no identical-particle factorial in this convention.

\begin{theorem}[Fixed-parent roots and Jacobian]
\label{thm:v20v42-roots}
There are exactly two energy-conserving roots:
\[
 p_\pm=\frac q2\pm\frac{\Omega}{M}k_*,
 \qquad E_\pm=\frac{\Omega}{2}\pm\frac qM k_*.
\]
The other daughter's momentum and energy at \(p_\pm\) are \(p_\mp,E_\mp\).
Moreover
\[
 D'_q(p)=-\frac p{E(p)}+\frac{q-p}{E(q-p)},\qquad
 E_+E_-\,|D'_q(p_\pm)|=M k_*.
\]
\end{theorem}
\begin{proof}
Boost the rest-frame on-shell pair
\((M/2,k_*)\), \((M/2,-k_*)\) by velocity \(\beta\).
This gives the displayed energies and momenta, all energies being
positive.  Conversely any on-shell pair summing to \((\Omega,q)\)
boosts back to a zero-total-momentum pair, whose energy constraint
forces momenta \(k_*\) and \(-k_*\).  Thus these are all the roots.
Differentiation gives \(D'_q\).  Finally direct expansion gives
\[
 p_+E_--p_-E_+=M k_*,
\]
which proves the Jacobian identity and simplicity of the roots.
\end{proof}

\begin{remark}[The scope of V41's density]
\label{rem:v20v42-density}
V41 differentiated with \(k_0\) fixed and the parent momentum varying
as \(q=k_0-sk_-\).  Its derivative \(-s v_q-v_-\) is correct for that
question.  The fixed-parent derivative above instead uses the relative
velocity of the two daughters.  Substituting V41's density directly
into a claimed parent decay width changes the constrained integration
and is not justified.  A standing cosine parent also contains both
signed momenta; it is not a single traveling parent eigenstate.
\end{remark}

\subsection{A rigorous density limit with a specified order}

Let \(\eta\) be a nonnegative smooth compactly supported function of
integral one, \(\eta_h(s)=h^{-1}\eta(s/h)\), \(h>0\), and let
\(\varphi\) be a smooth compactly supported function of daughter
momentum.  Take circle momenta \(p_n=2\pi n/L\).

\begin{theorem}[Sequential large-volume and energy-resolution limits]
\label{thm:v20v42-limit}
For the strict above-threshold parameters just specified,
\[
 \lim_{h\downarrow0}\lim_{L\to\infty}
 \frac{2\pi}{L}\sum_{n\in\mathbb Z}
 \varphi(p_n)\eta_h(D_q(p_n))
 =
 \sum_{\nu\in\{+,-\}}\frac{\varphi(p_\nu)}{|D'_q(p_\nu)|}.
\]
The same result holds along allowed parent momenta
\(q_L=2\pi j_L/L\to q\), replacing \(q\) inside the sum by \(q_L\).
\end{theorem}
\begin{proof}
For fixed \(h\), the integrand is continuous and compactly supported.
The mesh Riemann sums converge to
\(\int \varphi(p)\eta_h(D_q(p))\,dp\).  If \(q_L\to q\),
the corresponding integrands converge uniformly on the common
compact support, which leaves the same limit.
Choose disjoint neighborhoods of the two simple roots.  On the
remaining compact support \(D_q\) is bounded away from zero, so
for small enough \(h\) it contributes nothing.  On each root
neighborhood use the smooth coordinate \(s=D_q(p)\).
The approximate identity then converges to the value of
\(\varphi(p)/|D'_q(p)|\) at the root.  Adding the two terms proves
the assertion.
\end{proof}

Thus the leading density of daughter states per unit detuning is
\[
 \rho_q(0)=\frac{L}{2\pi}
 \left(\frac1{|D'_q(p_+)|}+\frac1{|D'_q(p_-)|}\right),
\]
understood as this smoothed asymptotic density, not a count of exact
finite-circle resonances.  No interchange of limits has been proved:
at finite \(L\), generic masses have no exact resonances, whereas
a delta peak at an exact resonance becomes singular as \(h\to0\).
The theorem is not a long-time weak-coupling kinetic limit.

\subsection{Explicitly normalized one-dimensional phase space}

Define, with metric signature \(+,-\) for this subsection,
\[
 \Phi_2(q)=\int_{\mathbb R^2}
 \frac{dp\,dr}{4E(p)E(r)}
 \delta(q-p-r)\delta(\Omega-E(p)-E(r)).
\]
This is the \(1+1\)-dimensional convention in which the factors
\((2\pi)^2\) of the conserving delta functions cancel the two
one-particle \(2\pi\) factors.

\begin{theorem}[Invariant phase space and a conditional width]
\label{thm:v20v42-width}
One has
\[
 \Phi_2(q)=\frac1{2Mk_*}.
\]
If a canonically normalized decay amplitude has the same squared
value \(A^2\) at the two labeled roots, define its leading-order
golden-rule width by
\[
 \Gamma(q)=\frac{A^2}{2\Omega}\Phi_2(q).
\]
Then
\[
 \Gamma(q)=\frac{A^2}{4\Omega M k_*}
 =\frac M\Omega\Gamma(0).
\]
\end{theorem}
\begin{proof}
Integrate first over \(r=q-p\).  The simple-root delta identity
gives two contributions, each equal to
\(1/(4E_+E_-|D'_q(p_\pm)|)=1/(4Mk_*)\).
This proves the phase-space formula.  The width statements follow
by multiplication with the specified flux factor.
\end{proof}

The invariant-phase-space and initial-state flux conventions agree
with the general particle-decay convention described by the
\href{https://pdg.lbl.gov/2025/reviews/rpp2025-rev-kinematics.pdf}
{Particle Data Group, Kinematics, section 49.4}.
The one-dimensional evaluation and all constants above have been
derived here.  A plane reduction of a four-dimensional field theory
does not inherit this normalization without a transverse-volume and
state-space identification.

\subsection{A canonical two-component Dirac--scalar benchmark}

\begin{theorem}[One Dirac species with scalar Yukawa coupling]
\label{thm:v20v42-yukawa}
Consider an independently normalized \(1+1\)-dimensional free-state
Dirac--scalar model with interaction
\(-yH\overline\Psi\Psi\), and adopt the completeness conventions
\[
 u(p)\overline u(p)=\slashed p+m,\qquad
 v(r)\overline v(r)=\slashed r-m,\qquad
 \operatorname{tr}I=2,\quad
 \operatorname{tr}(\gamma^\mu\gamma^\nu)=2\eta^{\mu\nu}.
\]
At tree level the summed squared amplitude is
\[
 A^2=y^2(M^2-4m^2)=4y^2k_*^2,
\]
and its golden-rule width is
\[
 \Gamma_{\rm can}(q)=\frac{y^2 k_*}{M\Omega},
 \qquad \Gamma_{\rm can}(0)=\frac{y^2 k_*}{M^2}.
\]
\end{theorem}
\begin{proof}
The amplitude is \(y\overline u(p)v(r)\).  Multiplication by its
complex conjugate, followed by the completeness identities, gives
\[
 A^2=y^2\operatorname{tr}
 ((\slashed p+m)(\slashed r-m))
 =2y^2(p\cdot r-m^2).
\]
Since \(p^2=r^2=m^2\) and \((p+r)^2=M^2\), this is
\(y^2(M^2-4m^2)\).  The preceding theorem gives the stated widths.
\end{proof}

There is one species and no additional spin, colour, or flavour
multiplicity here.  If \(y=m/v\) is chosen, the rest value is exactly
the expression proposed at the end of V41.  This is a benchmark
under explicit free-state conventions; it is not a proof that
V41's classical triad quantization has these conventions.
For fixed \(m>0\), writing \(M=2m+\epsilon\) shows
\(k_*=\sqrt{m\epsilon+\epsilon^2/4}\).  Thus this benchmark width
vanishes as a constant times \(\sqrt{\epsilon}\) at threshold,
despite the phase-space divergence \(1/k_*\).
The simple-root proof excludes \(M=2m\) itself.

\subsection{The finite-box normalization that must be matched}

\begin{proposition}[Necessary box coupling for the benchmark]
\label{prop:v20v42-box}
Under the canonical conventions above, the squared matrix element
between unit-normalized circle states is
\[
 |G_L(q;p)|^2=\frac{A^2}{8\Omega E(p)E(q-p)L}
\]
on shell.  This coupling, summed over both labeled daughter roots,
has the continuum golden-rule coefficient in
Theorem \ref{thm:v20v42-width}.
\end{proposition}
\begin{proof}
Each of the three external fields contributes
\((2EL)^{-1/2}\) and the spatial vertex integral contributes \(L\)
with a Kronecker momentum constraint.  Squaring gives the displayed
factor.  Equivalently it defines the box normalization explicitly.
Combining the Riemann-sum density \(L/(2\pi)\) with the factor \(2\pi\)
of the golden rule gives
\[
 2\pi\frac L{2\pi}\sum_{\nu=\pm}
 \frac{A^2}{8\Omega E_+E_-L|D'_q(p_\nu)|}
 =\frac{A^2}{4\Omega Mk_*}.
\]
For a rigorous version apply Theorem \ref{thm:v20v42-limit} to a
smooth compact cutoff equal to one near both roots and the smooth
coupling weight.  The integral localizes to those roots as
the broadening is removed.
\end{proof}

In contrast, V41's homogeneous standing-wave formula gives
\[
 L G_{\rm old}^2
 =\frac{4\varepsilon^2y^2 k_*^2}{M^3},
\]
whereas the canonical traveling-state formula at either root is
\[
 L |G_{\rm can}|^2=\frac{2y^2 k_*^2}{M^3}.
\]
For example, counting only positive \(k_*\) in the old standing-wave
sector compensates a factor of two from the two canonical roots
when \(\varepsilon=1\).  This agreement of a total rest-frame
coefficient does not determine the canonical transformation, the
vacuum, the spontaneous term, or the moving-parent multiplicities.
It is precisely the remaining normalization problem, rather than a
reason to set \(\varepsilon=1\) silently.

\subsection{Executable checks and claim boundary}

The exact-rational verifier
\[
 \texttt{scripts/verify\_sm\_v20\_v42\_fixed\_parent\_phase\_space.py}
\]
checks eight boosted mass-shell examples, both root identities, the
Jacobian, the box golden-rule factor, and time dilation.  It emits
\[
 \texttt{artifacts/sm\_v20\_v42\_fixed\_parent\_phase\_space.json}.
\]
Its SHA--256 is
\[
\texttt{8CD7B050AD23B6AC3F242F1508E5080A882F0CBFCDF699BAEEDAD8A61536EC13},
\]
and the payload SHA--256 is
\[
\texttt{008C219BE2C5C8B18489F84EE7526B03816A65FAA7FCD2C76F0EF89E35828F79}.
\]
These finite checks supplement the proofs; they are not proofs by
sampling of identities over all parameters.

\begin{firewall}[The remaining normalization and continuum gates]
\label{fire:v20v42}
The density limit is sequential and distributional, and the width is
a leading-order benchmark in an explicitly normalized
one-dimensional model.  No nonperturbative decay law, kinetic
limit, four-dimensional Standard Model width, interacting quantum
measure, or SM--Einstein continuum follows.  In particular the
classical cosine family and the canonical traveling-state model
have not been proved unitarily equivalent as interacting systems.
\end{firewall}

\begin{decision}[V43: normalization before another channel]
\label{dec:v20v42-next}
Before applying this calculation to the colour channel, derive the
canonical traveling-wave and standing-wave transformations in a
finite box, including both parity sectors, the Dirac anticommutators,
the scalar zero mode, and the interaction matrix elements.
Determine whether the retained cosine family is a closed quantum
subspace and how its classical \(\varepsilon\) enters normalized
occupations.  If it is not closed, identify the omitted states and
repair the truncation.  This upstream step replaces the planned
immediate transfer of an unnormalized rate to another channel.
\end{decision}

\begin{assessment}[V42 acquisition]
\label{ass:v20v42}
The parent-decay Jacobian, a precise large-circle density limit, and
the normalized two-body Yukawa benchmark now have analytic proofs.
The benchmark has the proposed rest value and the correct
time-dilation factor.  Its application to the existing reduced
model is explicitly conditional on the still-missing quantum
normalization and state-counting identification.
\end{assessment}

\begin{record}[V42 append record]
\label{rec:v20v42}
The predecessor V41 has SHA--256
\texttt{8ACA6F9361691EF65AF81A82F4356BCC99C01E4F323060F65AB3D559545055B4}.
Its text is retained byte for byte before the former final
document terminator.  Only its immediate active numeric version is
replaced on validated installation; archives and companion notes
are preserved.  The next compulsory companion audit is V45.
\end{record}



\section{V43: canonical box normalization and the missing standing-wave pair}
\label{sec:v20v43}

V42 left the identification of classical standing-wave amplitudes with
canonical quantum states open.  We now construct the canonical finite-box
model explicitly.  A single homogeneous scalar quantum already disproves
closure of the retained fermion parity family: its Yukawa vertex produces
two orthogonal pair channels, of equal weight.  This is a state-space
obstruction, not a choice of a numerical normalization constant.

\subsection{Canonical fields and the scalar zero mode}

Fix a circle of length \(L\), masses \(m,M>0\), and periodic momenta
\(p_n=2\pi n/L\).  Work at time zero, initially with finite mode cutoffs
so every displayed matrix element between finite-particle states is
well defined.  Let
\[
 \{c_p,c_r^\dagger\}=\{d_p,d_r^\dagger\}=\delta_{pr},
 \qquad [a_q,a_s^\dagger]=\delta_{qs},
\]
with other fermion anticommutators zero.  All pair states below have the
fixed creation order \(c^\dagger d^\dagger\).  Take
\[
 \alpha=\begin{pmatrix}0&1\\1&0\end{pmatrix},\quad
 B=\begin{pmatrix}1&0\\0&-1\end{pmatrix},\quad
 E_p=\sqrt{p^2+m^2},\quad \Omega_q=\sqrt{q^2+M^2},
\]
\[
 A_p=\sqrt{\frac{E_p+m}{2E_p}},\qquad
 B_p=\frac{p}{\sqrt{2E_p(E_p+m)}},\qquad
 U(p)=\binom{A_p}{B_p},\quad V(p)=\binom{B_p}{A_p}.
\]
The symbol \(B_p\) is a scalar component, distinct from the matrix \(B\).
The canonically normalized expansions are
\[
 \Psi(z)=\frac1{\sqrt L}\sum_p
 \left(c_pU(p)e^{ipz}+d_p^\dagger V(p)e^{-ipz}\right),
\]
\[
 H(z)=\sum_q\frac{a_qe^{iqz}+a_q^\dagger e^{-iqz}}
 {\sqrt{2\Omega_q L}}.
\]

\begin{lemma}[Free normalization and completeness]
\label{lem:v20v43-free}
The spinors have unit norm, satisfy
\[
 (\alpha p+Bm)U(p)=E_pU(p),\qquad
 (\alpha(-p)+Bm)V(p)=-E_pV(p),
\]
and \(U(p)U(p)^\dagger+V(-p)V(-p)^\dagger=I\).
Consequently the field equal-time CAR is the identity matrix times
the periodic delta function (or its Fourier-cutoff projection).
The scalar homogeneous coordinate is
\[
 H_0=\frac{a_0+a_0^\dagger}{\sqrt{2ML}},
\]
with one oscillator, not a cosine/sine pair.
\end{lemma}
\begin{proof}
Use \(A_p^2+B_p^2=1\), \(2A_pB_p=p/E_p\),
and \(A_p^2-B_p^2=m/E_p\) to multiply the displayed matrices.
Since \(V(-p)=(-B_p,A_p)^{\mathsf T}\), its projector is the complement
of that of \(U(p)\).  Substitution into the field CAR then gives
\(L^{-1}\sum_p e^{ip(z-w)}I\).  The scalar zero-mode expression is
the \(q=0\) term of its expansion.
\end{proof}

For \(q>0\) define \(a_C=(a_q+a_{-q})/\sqrt2\) and
\(a_S=(a_q-a_{-q})/(i\sqrt2)\).  Then the annihilation part of the
\(\{q,-q\}\) contribution to \(H\) is
\[
 \frac{a_C\cos qz-a_S\sin qz}{\sqrt{\Omega_q L}}.
\]
This derives the distinct homogeneous and nonhomogeneous factors.

\subsection{The traveling-state Yukawa matrix element}

Define the normal-ordered canonical interaction
\[
 H_I=y\int_0^L H(z):\Psi^\dagger(z)B\Psi(z):\,dz.
\]
For normalized states
\(|q\rangle=a_q^\dagger|0\rangle\) and
\(|p,r\rangle=c_p^\dagger d_r^\dagger|0\rangle\),
the part annihilating the parent and creating a pair is
\[
 H_{I,\mathrm{dec}}=
 \sum_{q=p+r}\frac{y\,U(p)^\dagger B V(r)}
 {\sqrt{2\Omega_q L}}\,a_qc_p^\dagger d_r^\dagger.
\]

\begin{theorem}[Exact box amplitude]
\label{thm:v20v43-vertex}
For the indicated external states,
\[
 \langle p,r|H_I|q\rangle
 =\delta_{q,p+r}\frac{y(A_pB_r-B_pA_r)}{\sqrt{2\Omega_q L}}.
\]
Upon putting \(u(p)=\sqrt{2E_p}U(p)\) and
\(v(r)=\sqrt{2E_r}V(r)\), this becomes
\[
 \frac{\mathcal M(p,r)}{\sqrt{8\Omega_qE_pE_rL}},
 \qquad \mathcal M=y\overline u(p)v(r),
\]
exactly the canonical box convention used in V42.
\end{theorem}
\begin{proof}
The spatial exponent is \(e^{i(q-p-r)z}\); integration supplies
\(L\delta_{q,p+r}\).  The two Dirac fields contribute \(L^{-1}\),
and the scalar contributes \((2\Omega_q L)^{-1/2}\).
The unit-normalized CAR states produce no extra degeneracy or norm
factor.  Finally \(\overline u=u^\dagger B\), giving the second formula.
\end{proof}

\subsection{Both standing-wave fermion sectors}

For each \(p>0\) and \(b=c,d\), set
\[
 b_C=\frac{b_p+b_{-p}}{\sqrt2},\qquad
 b_S=\frac{b_p-b_{-p}}{i\sqrt2}.
\]

\begin{proposition}[Unitary change of basis, not deletion of modes]
\label{prop:v20v43-car}
The \(C,S\) operators satisfy exactly the same CAR as the original
two momentum modes.  Their four particle--antiparticle pair states
\(c_X^\dagger d_Y^\dagger|0\rangle\), \(X,Y=C,S\), are orthonormal.
The transformation preserves the free Hamiltonian and total occupation
on each \(\{p,-p\}\) pair.
\end{proposition}
\begin{proof}
The coefficient matrix has orthonormal rows, so conjugating the
identity CAR matrix leaves it the identity.  Using the CAR twice gives
the inner product of pair states as \(\delta_{XX'}\delta_{YY'}\).
Both traveling modes have energy \(E_p\), so their free Hamiltonian
is \(E_p(b_p^\dagger b_p+b_{-p}^\dagger b_{-p})\), which is
unchanged by this unitary transformation.
\end{proof}

\begin{theorem}[Homogeneous scalar creates two standing-wave pairs]
\label{thm:v20v43-leak}
Fix \(p>0\).  The contribution of the two daughter configurations
\((p,-p)\) and \((-p,p)\) to \(H_{I,\mathrm{dec}}\) is
\[
 -h_p a_0(c_p^\dagger d_{-p}^\dagger-c_{-p}^\dagger d_p^\dagger),
 \qquad h_p=\frac{yp}{E_p\sqrt{2ML}}.
\]
In standing variables,
\[
 c_p^\dagger d_{-p}^\dagger-c_{-p}^\dagger d_p^\dagger
 =i(c_C^\dagger d_S^\dagger-c_S^\dagger d_C^\dagger).
\]
Thus the pair-production image of a scalar quantum has equal squared
amplitude \(h_p^2\) in the two orthogonal channels \(CS\) and \(SC\).
Retaining only the \(CS\) channel omits half of its squared norm.
\end{theorem}
\begin{proof}
At \(r=-p\),
\(A_pB_{-p}-B_pA_{-p}=-2A_pB_p=-p/E_p\);
the reversed pair has the opposite sign.  This proves the first
identity by Theorem \ref{thm:v20v43-vertex}.
Insert
\(b_p^\dagger=(b_C^\dagger-i b_S^\dagger)/\sqrt2\) and
\(b_{-p}^\dagger=(b_C^\dagger+i b_S^\dagger)/\sqrt2\)
without changing the fermion creation order.  Expanding gives the
second identity.  Orthonormality from Proposition
\ref{prop:v20v43-car} proves the norm assertion.
\end{proof}

This directly tests the old spinor family.  The coefficient of \(c_C\)
in the field is proportional to
\[
 \binom{A_p\cos pz}{iB_p\sin pz},
\]
and the coefficient of \(d_S^\dagger\), up to an overall phase, is
proportional to
\[
 \binom{B_p\cos pz}{-iA_p\sin pz}.
\]
Since \(B_p/A_p=(E_p-m)/p\), these are precisely the two
component profiles used in the earlier cosine-family construction.
That family contains \(c_C,d_S\); its complementary \(c_S,d_C\)
pair is nonzero in the same scalar vertex.

\begin{corollary}[Nonclosure and the minimal repair of this vertex]
\label{cor:v20v43-nonclosure}
Let \(\mathcal P\) project onto a Fock truncation containing the parent
state \(a_0^\dagger|0\rangle\) and the \(CS\) daughter pair but excluding
the \(SC\) pair.  For \(yp\neq0\),
\[
 (I-\mathcal P)H_I\mathcal P\neq0.
\]
Adding the \(SC\) pair repairs this particular rest-parent vertex.
Keeping both standing sectors for each signed momentum pair restores
the full finite-cutoff canonical Yukawa matrix elements.
\end{corollary}
\begin{proof}
The omitted component in Theorem \ref{thm:v20v43-leak} has coefficient
\(ih_p\neq0\).  This proves nonclosure even if other terms of \(H_I\)
are retained, since they cannot cancel this same specified matrix
element.  With both sectors present the transformation is unitary
and no matrix elements are dropped.
\end{proof}

Adding a single pair does not close the full interacting Fock space:
other parent momenta and higher-particle sectors remain.  Moreover
a finite Fourier cutoff is itself a projected interaction, not an
invariant subspace of the uncut local field theory.

\subsection{What becomes of the classical amplitude}

\begin{proposition}[Occupation normalization is not a field rescaling]
\label{prop:v20v43-epsilon}
For any canonical fermion mode \(b\), \(N_b=b^\dagger b\) is a projection,
so every positive normalized state has
\(0\leq\langle N_b\rangle\leq1\).  Replacing \(b\) by
\(\varepsilon b\) changes its CAR to
\(\{\varepsilon b,\overline\varepsilon b^\dagger\}=|\varepsilon|^2\);
it preserves canonical normalization only if \(|\varepsilon|=1\).
In a fermion-parity-invariant state, \(\langle\Psi\rangle=0\).
\end{proposition}
\begin{proof}
The CAR gives \(N_b^2=N_b\), hence the occupation bounds.
The rescaled anticommutator follows by bilinearity.
Fermion parity sends \(\Psi\) to \(-\Psi\); invariance of the state
therefore forces its expectation to equal its negative.
\end{proof}

Consequently an earlier commuting classical spinor amplitude
\(\varepsilon\) is not fixed by declaring it one quantum.  A comparison
must specify bilinear expectations or occupation matrices, Pauli
bounds, and the prepared many-body state.  The previous classical
resonant equations remain classical statements; their rate constants
cannot alone determine spontaneous fermionic production.

\subsection{Verification and next dependency}

The verifier
\[
 \texttt{scripts/verify\_sm\_v20\_v43\_parity\_channels.py}
\]
checks the four pair states' exact Gram matrix, the standing/traveling
identity, the equal retained and omitted weights, and a rational
mass-shell normalization example.  The script SHA--256 is
\[
\texttt{CAE593FC59B884F3DEACAB685074C2A3DD917A4C1CE9B612DE0B925189B91A8C}.
\]
The artifact
\(\texttt{artifacts/sm\_v20\_v43\_parity\_channels.json}\)
has SHA--256
\[
\texttt{7B29D37B064E1636539D275165EA85A158B721E13A8FC7E5630D6308F9DABB8E}.
\]

\begin{firewall}[Canonical finite-box model only]
\label{fire:v20v43}
The canonical state normalization and missing pair are proved for
one Dirac species and one scalar with a finite-mode regulator.
No four-dimensional chiral Standard Model identification, ultraviolet
limit, interacting vacuum, or gravitational coupling is constructed.
The pair norm calculation is a vertex statement, not an exponential
decay law or a Markovian kinetic theorem.
\end{firewall}

\begin{decision}[V44 acquisition]
\label{dec:v20v43-next}
Use the repaired canonical pair space to derive the exact finite-triad
occupation commutators, the spontaneous source, and the Pauli blocking
factors.  Compare them with the classical wave-kinetic collision term
of V34 and isolate the assumptions needed to replace correlated
expectations by occupation products.  A classical random-phase closure
must not be carried over to CAR variables without this derivation.
\end{decision}

\begin{assessment}[V43 acquisition]
\label{ass:v20v43}
The normalized finite-box Yukawa vertex now realizes V42's benchmark.
The old standing-wave family omits an equal-strength orthogonal pair
already for a homogeneous scalar, so its quantum closure fails.
Restoring both standing sectors repairs the canonical finite-mode
vertex; occupation statistics and a controlled kinetic limit are the
next distinct dependencies.
\end{assessment}

\begin{record}[V43 append record]
\label{rec:v20v43}
The predecessor V42 has SHA--256
\texttt{942554C4C4ABB6EBE7529CA33CBE73B16595331BD31D463FE7EA5FD45CBA1C11}.
Its complete prefix before the last document terminator is preserved.
Only the immediate active SM predecessor is replaced after validation.
The next companion audit is V45.
\end{record}



\section{V44: exact Yukawa occupation algebra, Pauli blocking, and the closure obstruction}
\label{sec:v20v44}

V43 repaired the missing standing-wave pair in the canonical box vertex.
The next issue is statistical: canonical fermion occupations obey CAR,
whereas V34 used commuting classical amplitudes.  We do not repeat the
two-state ladder diagonalization of V26.  Instead we derive the exact
occupation hierarchy, exhibit states with equal one-mode occupations
but different initial conversion curvature, and identify the extra
coherence generated when the two repaired channels share a scalar.

\subsection{Exact single-channel identities}

Let \(a\) be a scalar oscillator and \(c,d\) two CAR modes.  Define
\[
 N=a^\dagger a,\quad F=c^\dagger c,\quad D=d^\dagger d,\quad
 P^\dagger=c^\dagger d^\dagger,\quad P=dc,\quad X=aP^\dagger.
\]
The resonant-vertex Hamiltonian with a possible detuning is
\[
 H=H_0+gX+\overline gX^\dagger,\qquad
 H_0=\Omega N+E_cF+E_dD,\qquad
 \Delta=E_c+E_d-\Omega.
\]
Here \(g\) is the canonical finite-box coupling of V43; its phase is
allowed to be complex.  These identities hold on the finite-particle
core.  For evolution one may work in each finite invariant block fixed
by \(N+F\) and \(N+D\); the block Hamiltonian is a finite self-adjoint
matrix, so the differentiations below have no domain ambiguity.

\begin{theorem}[Occupation, coherence, and blocking operators]
\label{thm:v20v44-algebra}
Put
\[
 B=N(1-F)(1-D)-(N+1)FD=N-NF-ND-FD.
\]
Then
\[
 [X^\dagger,X]=B,\qquad
 [F,X]=[D,X]=X,\qquad [N,X]=-X.
\]
For a state evolved by \(H\), write
\(n=\langle N\rangle\), \(f=\langle F\rangle\),
\(d=\langle D\rangle\), \(z=g\langle X\rangle\).
Its exact equations are
\[
 \dot f=\dot d=-\dot n=2\operatorname{Im}z,\qquad
 \dot z=i\Delta z+i|g|^2\langle B\rangle,
\]
\[
 \ddot f=2\Delta\operatorname{Re}z+2|g|^2\langle B\rangle.
\]
The operators \(N+F,N+D,F-D\) are conserved.
\end{theorem}
\begin{proof}
The CAR gives
\[
 PP^\dagger=(1-F)(1-D),\qquad P^\dagger P=FD.
\]
Since the oscillator commutes with the fermions,
\(X^\dagger X=N(1-F)(1-D)\) and
\(XX^\dagger=(N+1)FD\).  Their difference is \(B\); expanding
cancels both triple-occupation terms and yields its second form.
The number commutators follow by counting creations and annihilations.
Use \(\dot O=i[H,O]\), \([H_0,X]=\Delta X\), and the displayed
commutators to obtain the equations.  The conserved sums commute
with each term of \(H\).
\end{proof}

For \(|n,0,0\rangle\), \(B=n\): a scalar can decay into an empty
pair without an initially occupied fermion mode.  For
\(|n,1,1\rangle\), \(B=-(n+1)\): inverse pair annihilation includes
the scalar vacuum contribution.  For either singly occupied pair
\(|n,1,0\rangle\) or \(|n,0,1\rangle\), both transition operators
vanish.  These are exact Pauli blocking statements.
In particular a number-diagonal initial state has \(z(0)=0\),
\(\dot f(0)=0\), and
\[
 f(t)=f(0)+|g|^2\langle B(0)\rangle t^2+O(t^3).
\]
This is short-time unitary curvature, not a constant irreversible rate.

\subsection{An explicit correction to the classical quantum reading}

\begin{record}[Correction of the Yukawa interpretation after V34]
\label{rec:v20v44-correction}
V34's wave-kinetic equations concern classical commuting amplitudes.
Its reading in Proposition \ref{prop:v20v34-reading} suggests a
quantum collision term with bosonic enhancement factors on all
three modes.  That suggestion must not be applied to the canonical
fermion--antifermion Yukawa channel of V43.  Its exact operator
factor is \(B\) above.  Under an additional independent-occupation
ansatz it becomes
\[
 B_{\rm fac}=n(1-f)(1-d)-(n+1)fd
            =n(1-f-d)-fd.
\]
The fermionic factors are \(1-f\) and \(1-d\), not \(1+f\) and
\(1+d\).  The classical identities of V34 are not altered by this
append-only correction; their promotion to a quantum Yukawa closure
is withdrawn pending an appropriate limit theorem.
\end{record}

\subsection{Same occupations do not imply the same evolution}

For the commuting number observables define
\(\operatorname{Cov}(A,B)=\langle AB\rangle-\langle A\rangle
\langle B\rangle\).

\begin{theorem}[Exact error of occupation factorization]
\label{thm:v20v44-correlations}
One has
\[
 \langle B\rangle-B_{\rm fac}
 =-\operatorname{Cov}(N,F)-\operatorname{Cov}(N,D)
  -\operatorname{Cov}(F,D).
\]
Even number-diagonal states with identical \(n,f,d\) can have distinct
\(\langle B\rangle\).
\end{theorem}
\begin{proof}
Insert expectations into \(B=N-NF-ND-FD\) and subtract the
factorized expression.  For a concrete counterexample fix \(N=1\)
and take the two equally weighted mixtures of fermion configurations
\[
 \rho_{\rm corr}:\quad 00,\ 11,\qquad
 \rho_{\rm anti}:\quad 10,\ 01.
\]
Both have \(n=1,f=d=1/2\) and zero initial \(z\).  Their exact
blocking expectations are respectively
\[
 \langle B\rangle_{\rm corr}=(1-2)/2=-1/2,\qquad
 \langle B\rangle_{\rm anti}=0.
\]
The occupation-only expression gives \(-1/4\) for both.
Thus Theorem \ref{thm:v20v44-algebra} gives different second
derivatives for equal one-mode data.
\end{proof}

This counterexample rules out an exact universal autonomous equation
for \(n,f,d\) alone.  Randomizing phases does not eliminate number
correlations, since both counterexample states already have no
off-diagonal number-basis coherence.

\subsection{The repaired channels also generate pair coherence}

Let \(i=1,\ldots,r\) denote disjoint fermion pairs coupled to one
oscillator.  Set \(P_i^\dagger=c_i^\dagger d_i^\dagger\),
\(X_i=aP_i^\dagger\), and
\[
 H_I=\sum_i(g_iX_i+\overline g_iX_i^\dagger).
\]
Disjointness holds for the two standing-wave pairs \(CS,SC\) found
in V43: they use distinct particle and antiparticle modes.

\begin{theorem}[Cross-channel hierarchy]
\label{thm:v20v44-cross}
For \(j\neq i\),
\[
 [X_j^\dagger,X_i]=-P_i^\dagger P_j.
\]
Consequently
\[
 \frac{d}{dt}\langle X_i\rangle
 =i\Delta_i\langle X_i\rangle+
 i\overline g_i\langle B_i\rangle
 -i\sum_{j\neq i}\overline g_j\langle P_i^\dagger P_j\rangle,
\]
where \(B_i=N(1-F_i)(1-D_i)-(N+1)F_iD_i\).
Thus repairing the state space does not yield independent triads
merely by adding their scalar occupation rates.
\end{theorem}
\begin{proof}
Even fermion pair monomials on disjoint modes commute.  Therefore
\[
 X_j^\dagger X_i=NP_i^\dagger P_j,\qquad
 X_iX_j^\dagger=(N+1)P_i^\dagger P_j.
\]
Subtract and use the single-channel commutator for \(j=i\).
The terms \([X_j,X_i]\) vanish, so the Heisenberg equation gives
the result.
\end{proof}

At a fixed momentum magnitude the two channels may instead be
organized into their coupled linear combination and its orthogonal
combination.  This simplifies the one-excitation block but does not
justify discarding correlations in general many-particle states.

\subsection{Conditional kinetic equation and its statistical consistency}

\begin{proposition}[Entropy production of the proposed CAR closure]
\label{prop:v20v44-entropy}
Suppose, as an additional model assumption, that a Markovian closure
has a nonnegative coefficient \(R\) and equations
\[
 \dot f=\dot d=-\dot n=R(A-C),\quad
 A=n(1-f)(1-d),\quad C=(n+1)fd.
\]
For \(n>0\), \(0<f,d<1\), its Bose--Fermi entropy
\[
 S=(n+1)\log(n+1)-n\log n
   -f\log f-(1-f)\log(1-f)
   -d\log d-(1-d)\log(1-d)
\]
satisfies
\[
 \dot S=R(A-C)\log(A/C)\geq0.
\]
On resonance, the Bose--Einstein and Fermi--Dirac occupations
at a common inverse temperature obey \(A=C\) when
\(\mu_H=\mu_c+\mu_d\).
\end{proposition}
\begin{proof}
Differentiate \(S\), use the assumed occupation equations, and collect
the logarithms:
\[
 -\log((n+1)/n)+\log((1-f)/f)+\log((1-d)/d)=\log(A/C).
\]
The logarithm is increasing, proving the inequality.
For the stated equilibrium distributions,
\(n/(n+1)=e^{-\beta(\Omega-\mu_H)}\) and
\((1-f)/f=e^{\beta(E_c-\mu_c)}\), with the analogous formula for \(d\).
Their product is one on resonance with the chemical-potential balance.
\end{proof}

This is a property of an explicitly assumed closure.  It is not a
derivation of that closure from unitary dynamics.  A proof of a
quantum kinetic limit must control number covariances and
cross-channel coherences, justify the weak-coupling and large-volume
time scales, and derive the coefficient with V42's order-of-limits
and normalization conventions.  Replacing these steps by an
entropy calculation would be circular.

\subsection{Executable certificate}

The verifier
\[
 \texttt{scripts/verify\_sm\_v20\_v44\_car\_kinetics.py}
\]
checks the exact number and pair commutators on 24 oscillator--fermion
basis states, using an unnormalized oscillator basis with no upper
cutoff, and checks the equal-marginal counterexample.  The script
SHA--256 is
\[
\texttt{F2DA140764AB45DC267BCC9FC5B0BF54C80C29E85183EF22F11F0BAE6B4267C3}.
\]
The artifact
\(\texttt{artifacts/sm\_v20\_v44\_car\_kinetics.json}\)
has SHA--256
\[
\texttt{178F39AF45E4EB86E0623290DF62C4F29B552E0263B3DF5A081C7BCD039B818B}.
\]
The proofs establish the general identities; the finite checks verify
signs and a counterexample rather than asserting a continuum theorem.

\begin{firewall}[No controlled quantum kinetic limit yet]
\label{fire:v20v44}
The operator equations and nonclosure examples are exact within the
canonical resonant-vertex model.  The entropy statement is conditional
on a Markovian occupation closure.  Neither establishes that closure,
a physical spontaneous decay semigroup, the full untruncated Yukawa
dynamics, a quantum Standard Model sector, or an Einstein continuum.
\end{firewall}

\begin{decision}[V45 audit, then V46 controlled memory calculation]
\label{dec:v20v44-next}
Perform the compulsory companion audit at V45.  Then go upstream to
the one-parent, one-pair continuum block, which has an exact scalar
memory equation.  Derive its spectral density from the canonically
normalized vertex and distinguish a finite-time memory evolution
from a Markovian exponential approximation.  This is a tractable
place to test the missing limit estimates before adopting a
many-occupation kinetic closure.
\end{decision}

\begin{assessment}[V44 acquisition]
\label{ass:v20v44}
Pauli blocking, the spontaneous terms, conserved occupations, and
cross-pair coherences now follow from exact CAR identities.
Equal occupations provably do not determine conversion curvature.
The classical-to-quantum inference after V34 is corrected, and the
appropriate Bose--Fermi entropy is proved only for an explicitly
conditional closure.
\end{assessment}

\begin{record}[V44 append record]
\label{rec:v20v44}
The predecessor V43 has SHA--256
\texttt{E83CD7BEB9ADF90F21B1817995668C7E17F9313841257234C3E98ED7892493DE}.
Its full prefix before the final document terminator is preserved.
Only that active SM predecessor is replaced after validation.
The next version V45 is the required companion audit.
\end{record}



\section{V45: compulsory companion audit before the quantum memory problem}
\label{sec:v20v45}

The active SM predecessor is V44.  The latest main Yang--Mills,
Navier--Stokes, and parallel Yang--Mills notes were inspected read-only.
The exact snapshots are:
\begin{itemize}
\item \texttt{Renewal\_Yang\_Mills\_Mass\_Gap\_Research\_Notes\_V43.tex},
619362 bytes, SHA--256
\texttt{7E8245FBDAECB0AF0B9E19DDB7FC87D0A614D14701FE45FDA8CFD6476DA1E06E}.
\item \texttt{Renewal\_NS\_Stability\_Research\_Notes\_V26.tex},
278283 bytes, SHA--256
\texttt{82C887F8C2C69E4F28FAD7C3DC8DE5B9099CB5A4BF431EBA1FFF3E91935978AD}.
\item \texttt{Renewal\_YM\_Parallel\_Research\_Notes\_V5.tex},
54174 bytes, SHA--256
\texttt{306F1BB84CFA8A8D47EA569EC17E9545F9867C8D32B548EC9D9C444B4F3C0512}.
\end{itemize}

Main YM has advanced from the V41 snapshot recorded at SM V40.
Its latest result bounds a conditioned free low-alias covariance
uniformly near the toron, by representing the constrained space as a
graph over the high modes instead of subtracting two divergent terms.
It explicitly leaves shell derivative and higher-level uniform bounds
open.  NS remains at its Oseen-kernel rank-one refinement, and parallel
YM retains its obstruction to volume-uniform tails of extensive global
topological charge.  The latter two sources are unchanged by hash.

\begin{decision}[Audit decision]
\label{dec:v20v45-audit}
No companion estimate is imported into a Yukawa spectral density or
kinetic limit.  Their variables and hypotheses do not establish a
decay memory estimate.  Main YM's useful methodological lesson is to
control a combined exact expression before separating divergent terms.
V46 therefore derives the exact one-parent resolvent and memory kernel
first, and checks ultraviolet convergence before proposing an
exponential approximation.
\end{decision}

\begin{firewall}[Audit only]
\label{fire:v20v45}
The companion results do not supply an interacting SM state, a
quantum decay theorem, or a gravitational continuum.  This audit is
not evidence of closure of any of those gates.
\end{firewall}

\begin{assessment}[V45 acquisition]
\label{ass:v20v45}
The required audit is current.  V44's statistical nonclosure remains
the relevant obstruction; the next substantive step is the exact
one-excitation memory problem.
\end{assessment}

\begin{record}[V45 append record]
\label{rec:v20v45}
The predecessor V44 has SHA--256
\texttt{3B31ACFCB9791BBE78F62879DE2DD4E1EBA0070C0C451AC528699A6C7A5143A2}.
Its prefix is preserved and only its active numeric file is replaced.
Companions and frozen archives are unchanged.  The next compulsory
companion audit and ten-version sweep are V50.
\end{record}



\section{V46: exact one-parent memory, the Yukawa spectral density, and an ultraviolet obstruction}
\label{sec:v20v46}

The corrected CAR equations of V44 do not close on occupation means.
Before pursuing a many-mode kinetic approximation we solve the
one-parent, one-pair block of the canonical resonant vertex exactly
at the level of a memory equation.  The parent is at rest throughout
this section.  Both signed daughter momenta, or equivalently both
standing pair channels of V43, are retained.  The resulting spectral
density reproduces V42 on shell but exposes a separate ultraviolet
divergence away from the energy shell.

\subsection{The precise reduced Hamiltonian}

Fix \(m>0\), \(M>2m\), and real \(y\neq0\).  Let
\(E_p=\sqrt{p^2+m^2}\), \(\omega(p)=2E_p\), and choose an even smooth
cutoff \(0\leq\chi_\Lambda\leq1\), equal to one for \(|p|\leq\Lambda\)
and zero for \(|p|\geq2\Lambda\).  Define
\[
 v_\Lambda(p)=
 \frac{yp\,\chi_\Lambda(p)}{E_p\sqrt{4\pi M}},\qquad
 \mathcal H=\mathbb C\oplus L^2(\mathbb R,dp).
\]
The scalar component is the parent amplitude; the second component
labels the ordered particle--antiparticle pair \((p,-p)\).
For a real bare parent energy \(\mu_\Lambda\), set
\[
 H_\Lambda\binom{A}{b}=
 \binom{\mu_\Lambda A+\langle v_\Lambda,b\rangle}
       {\omega b+v_\Lambda A}.
\]
Our inner product is conjugate-linear in the first entry.
The canonical cutoff benchmark has \(\mu_\Lambda=M\); retaining the
symbol \(\mu_\Lambda\) permits a later energy counterterm.

\begin{lemma}[Self-adjoint cutoff dynamics and normalization]
\label{lem:v20v46-H}
The operator \(H_\Lambda\) is self-adjoint on
\(\mathbb C\oplus D(\omega)\).  Its coupling normalization is the
large-circle limit of the V43 rest-parent vertex.
\end{lemma}
\begin{proof}
Since \(v_\Lambda\in L^2\), the off-diagonal operator is bounded,
self-adjoint, and has norm \(\|v_\Lambda\|_2\).
Adding it to \(\operatorname{diag}(\mu_\Lambda,\omega)\) preserves
self-adjointness on that domain.
V43 gives the traveling-pair coefficient, up to an irrelevant sign,
\[
 G_L(p)=\frac{yp}{E_p\sqrt{2ML}}.
\]
Replacing the state sum by \(L\,dp/(2\pi)\) multiplies its squared
coefficient by \(L/(2\pi)\), giving
\(|v_\Lambda(p)|^2=y^2p^2\chi_\Lambda(p)^2/(4\pi M E_p^2)\).
\end{proof}

This is the restriction of the pair-creation/annihilation vertex
to its one-excitation block.  It is not asserted to be an invariant
subspace of the full normal-ordered Yukawa interaction, which also
contains other creation and scattering terms.

\subsection{Exact memory and resolvent}

\begin{theorem}[Volterra equation without a Markov approximation]
\label{thm:v20v46-memory}
Start from the normalized parent \((A(0),b(0))=(1,0)\).
Then
\[
 b(t,p)=-i v_\Lambda(p)\int_0^t e^{-i\omega(p)(t-s)}A(s)\,ds,
\]
\[
 \dot A(t)=-i\mu_\Lambda A(t)
 -\int_0^t K_\Lambda(t-s)A(s)\,ds,\qquad
 K_\Lambda(u)=\int_{\mathbb R}|v_\Lambda(p)|^2e^{-i\omega(p)u}\,dp.
\]
Furthermore, for \(z\notin\mathbb R\),
\[
 \langle e_0,(z-H_\Lambda)^{-1}e_0\rangle
 =\frac1{z-\mu_\Lambda-\Sigma_\Lambda(z)},\qquad
 \Sigma_\Lambda(z)=\int_{\mathbb R}
 \frac{|v_\Lambda(p)|^2}{z-\omega(p)}\,dp,
\]
where \(e_0=(1,0)\).  The evolution obeys
\(|A(t)|^2+\|b(t)\|_2^2=1\).
\end{theorem}
\begin{proof}
Variation of constants in \(i\dot b=\omega b+v_\Lambda A\)
gives the first formula.  Substitution into
\(i\dot A=\mu_\Lambda A+\langle v_\Lambda,b\rangle\) gives the
Volterra equation, with Fubini justified by the finite \(L^2\) norm.
For the resolvent equation \((z-H_\Lambda)(A,b)=e_0\), solve
\(b=(z-\omega)^{-1}v_\Lambda A\) and substitute into its scalar
component.  Self-adjointness ensures invertibility for nonreal \(z\).
The final identity is conservation of the norm under unitary evolution.
\end{proof}

For the interaction-picture amplitude \(a(t)=e^{i\mu_\Lambda t}A(t)\),
the kernel is \(e^{i\mu_\Lambda u}K_\Lambda(u)\).  The memory
integral remains exact; replacing \(a(s)\) by \(a(t)\) and extending
the upper limit to infinity are additional operations requiring
estimates.

\subsection{The full spectral density, not only its on-shell value}

\begin{theorem}[Rest-parent pair spectral density]
\label{thm:v20v46-density}
Write
\(K_\Lambda(u)=\int_{2m}^{\infty}J_\Lambda(w)e^{-iwu}\,dw\).
With \(k(w)=\sqrt{w^2/4-m^2}\),
\[
 J_\Lambda(w)=
 \frac{y^2}{4\pi M}
 \sqrt{1-\frac{4m^2}{w^2}}\,\chi_\Lambda(k(w))^2,
 \qquad w>2m.
\]
If the cutoff is one at \(k_*\), its on-shell value gives
\[
 2\pi J_\Lambda(M)=\frac{y^2 k_*}{M^2},
 \qquad k_*=\sqrt{M^2/4-m^2},
\]
the V42 golden-rule benchmark.
\end{theorem}
\begin{proof}
The map \(p\mapsto w=2E_p\) has two inverse branches
\(\pm k(w)\) and derivative magnitude \(2|p|/E_p\).
Adding their weights gives
\[
 J_\Lambda(w)=
 2\frac{y^2 k(w)^2\chi_\Lambda(k(w))^2}
 {4\pi M E_{k(w)}^2}
 \frac{E_{k(w)}}{2k(w)}
 =\frac{y^2 k(w)}{4\pi M E_{k(w)}}\chi_\Lambda(k(w))^2.
\]
Since \(E_{k(w)}=w/2\), this is the stated expression.
At \(w=M\) the same identity gives the benchmark.
\end{proof}

Writing \(c=y^2/(4\pi M)\), the uncut density satisfies
\[
 J(w)=c\sqrt{1-4m^2/w^2},\qquad
 J(2m+\epsilon)\sim c\sqrt{\epsilon/m},\qquad
 J(w)=c+O(w^{-2})\quad(w\to\infty).
\]
The vanishing threshold density and the nonzero high-energy limit
are distinct features.  An on-shell finite width does not imply
convergence of the off-shell memory problem.

\subsection{Short-time survival and failure of an exact exponential}

\begin{proposition}[Cutoff quadratic survival law]
\label{prop:v20v46-short}
For every fixed cutoff,
\[
 |A(t)|^2=1-\|v_\Lambda\|_2^2t^2+o(t^2).
\]
It cannot equal \(e^{-\Gamma t}\) for all \(t\geq0\) with
\(\Gamma>0\).
\end{proposition}
\begin{proof}
The parent lies in \(D(H_\Lambda^2)\), since \(\omega v_\Lambda\)
is square integrable.  Its first two energy moments are
\(\mu_\Lambda\) and \(\mu_\Lambda^2+\|v_\Lambda\|_2^2\).
Taylor expansion of the survival amplitude, followed by squaring,
gives the result.  Its probability derivative at zero is zero,
whereas that of \(e^{-\Gamma t}\) is \(-\Gamma\).
\end{proof}

\subsection{Ultraviolet obstruction and a convergent subtraction}

\begin{theorem}[Unsubtracted divergence]
\label{thm:v20v46-UV}
The uncut form factor is not in \(L^2(dp)\).  At any fixed real
\(z_0<2m\),
\[
 \Sigma_\Lambda(z_0)\longrightarrow-\infty .
\]
For the cutoff family specified above its divergence is logarithmic
with leading coefficient \(-c\); the growing cutoff energy is
asymptotic to a constant times \(\Lambda\).
Thus bounded-perturbation self-adjointness at fixed cutoff does not
by itself define an uncut block with a fixed parent energy.
\end{theorem}
\begin{proof}
Since \(|v(p)|^2\to y^2/(4\pi M)>0\), its integral diverges.
Equivalently \(J(w)\to c>0\).  For large \(w\),
\[
 \frac{J(w)}{z_0-w}=-\frac c w+O(w^{-2}).
\]
The cutoff is one up to energy \(2\sqrt{\Lambda^2+m^2}\) and
zero above \(2\sqrt{4\Lambda^2+m^2}\).  Integration of \(-c/w\)
up to the first bound gives the logarithm; the transition region
has a bounded contribution since its energy ratio stays bounded.
The remaining integrable term cannot cancel the divergence.
\end{proof}

\begin{theorem}[One-subtracted self-energy]
\label{thm:v20v46-subtraction}
Fix \(z_0<2m\).  For \(z\in\mathbb C\setminus[2m,\infty)\),
\[
 \Sigma_\Lambda(z)-\Sigma_\Lambda(z_0)\ \longrightarrow\
 S(z):=\int_{2m}^{\infty}J(w)
 \frac{z_0-z}{(z-w)(z_0-w)}\,dw .
\]
The convergence is locally uniform on that slit plane.
Choosing formally
\(\mu_\Lambda=\mu_R-\Sigma_\Lambda(z_0)\) makes the scalar
resolvent denominator converge locally to \(z-\mu_R-S(z)\).
\end{theorem}
\begin{proof}
Subtract the two fractions before integrating:
\[
 \frac1{z-w}-\frac1{z_0-w}
 =\frac{z_0-z}{(z-w)(z_0-w)}.
\]
On each compact subset of the slit plane the integrand has an
integrable bound at threshold and is \(O(w^{-2})\) uniformly at
infinity.  Because \(0\leq J_\Lambda\leq J\) and
\(J_\Lambda(w)\to J(w)\) at every fixed \(w>2m\), dominated
convergence, with that uniform bound, proves the assertion.
The denominator identity is then direct substitution.
\end{proof}

Convergence of this one matrix element's denominator is not yet
norm-resolvent convergence of the full Hamiltonian.  The latter
requires control of all block entries and proof that the limit
is the resolvent of a densely defined self-adjoint operator.
The counterterm also changes the bare parent frequency, so a
physical mass identification must be imposed and proved rather
than inferred from the letter \(M\).

\subsection{Verification, scope, and next step}

The exact-rational verifier
\[
 \texttt{scripts/verify\_sm\_v20\_v46\_memory\_density.py}
\]
checks the density normalization at two rational mass-shell points
and the sign of the subtraction identity.  The script SHA--256 is
\[
\texttt{E757B438640425E235E5B6DB14A8E1D36694C539B3C061C1C2005B3375AB1E86}.
\]
The artifact
\(\texttt{artifacts/sm\_v20\_v46\_memory\_density.json}\)
has SHA--256
\[
\texttt{0AA3000B6EB6AC9C80037B6775C1702AEBADCF95A16D30325968D5B3AD015F18}.
\]

\begin{firewall}[Projected quantum block, no exponential or field-theory limit]
\label{fire:v20v46}
The cutoff Hamiltonian and memory equation concern a resonant
one-parent block in one spatial dimension.  The unsubtracted
ultraviolet limit fails in the bounded-form-factor construction.
The convergent subtraction is not yet an operator construction,
and no Markov limit, many-body Yukawa theory, chiral gauge theory,
or SM--Einstein continuum has been proved.
\end{firewall}

\begin{decision}[V47 acquisition]
\label{dec:v20v46-next}
Analyze the entire subtracted block resolvent, including its
off-diagonal vectors, and prove or disprove norm-resolvent
convergence to a self-adjoint renormalized one-parent Hamiltonian.
Only after that operator exists should a resonance pole or
weak-coupling exponential law be used.  This goes upstream of
the proposed kinetic--Einstein closure rather than building it
on an uncut Hamiltonian that has not been defined.
\end{decision}

\begin{assessment}[V46 acquisition]
\label{ass:v20v46}
The exact memory equation and its explicitly normalized spectral
density now replace an assumed instantaneous rate in this block.
The on-shell benchmark survives, but the full self-energy has
a logarithmic ultraviolet divergence.  A single subtraction
converges analytically and identifies the next operator-level
construction problem.
\end{assessment}

\begin{record}[V46 append record]
\label{rec:v20v46}
The predecessor V45 has SHA--256
\texttt{FA264CA71E5ED7CB8651E25216027F821D9D948F414F671650779D42682A448F}.
Its full prefix is preserved.  Only the immediate active SM
predecessor is replaced; companions and archives are untouched.
The next compulsory audit and sweep are V50.
\end{record}



\section{V47: norm-resolvent construction of the renormalized one-parent Hamiltonian}
\label{sec:v20v47}

V46 proved convergence of a subtracted scalar self-energy but deliberately
left the operator construction open.  Here all four resolvent blocks
converge in operator norm.  The limit is injective with dense range,
which rules out loss of the parent degree of freedom and yields a
densely defined self-adjoint Hamiltonian.  The result concerns precisely
the one-parent resonant block, not the full Yukawa field Hamiltonian.

\subsection{Resolvent-weighted form factors}

Use the Hilbert space and notation of V46:
\[
 \mathcal H=\mathbb C\oplus L^2(\mathbb R,dp),\quad
 \omega(p)=2\sqrt{p^2+m^2},\quad
 v(p)=\frac{yp}{\sqrt{p^2+m^2}\sqrt{4\pi M}}.
\]
Fix a real subtraction point \(z_0<2m\), a real parameter \(\mu_R\),
and the V46 cutoff \(v_\Lambda=\chi_\Lambda v\).  Set
\[
 \mu_\Lambda=\mu_R-\Sigma_\Lambda(z_0),\qquad
 H_\Lambda=
 \begin{pmatrix}\mu_\Lambda&\langle v_\Lambda,\cdot\rangle\\
                 v_\Lambda&\omega\end{pmatrix}.
\]
Let \(r(z)=(z-\omega)^{-1}\) for \(z\notin\mathbb R\), and put
\[
 h_z=r(z)v,\qquad h_{\Lambda,z}=r(z)v_\Lambda,\qquad
 \ell_z(f)=\langle h_{\bar z},f\rangle,\qquad
 \ell_{\Lambda,z}(f)=\langle h_{\Lambda,\bar z},f\rangle.
\]
The complex conjugation in the subscript is essential: the inner
product is conjugate-linear in its first argument.

\begin{lemma}[Square-integrable weighted coupling]
\label{lem:v20v47-weight}
For each nonreal \(z\), \(h_z\in L^2\), and
\(h_{\Lambda,z}\to h_z\) in \(L^2\), locally uniformly in \(z\)
off the real axis.  The corresponding row functionals converge
in operator norm.
\end{lemma}
\begin{proof}
At large \(|p|\), \(v(p)\) is bounded and
\(|z-\omega(p)|^{-1}=O(|p|^{-1})\); its square is integrable in one
dimension.  On bounded momentum intervals the nonzero imaginary part
of \(z\) bounds the denominator away from zero.  These estimates
are uniform on compact subsets off the real axis.  Since
\(\chi_\Lambda=1\) for \(|p|\leq\Lambda\) and
\(|\chi_\Lambda|\leq1\), the norm of the difference is bounded by
the integrable tail, uniformly on those compact subsets.  The
functional norm equals the norm of its representing vector.
\end{proof}

\subsection{The denominator cannot vanish away from the real axis}

Define
\[
 S_\Lambda(z)=\Sigma_\Lambda(z)-\Sigma_\Lambda(z_0),\qquad
 D_\Lambda(z)=z-\mu_R-S_\Lambda(z),
\]
\[
 S(z)=\int_{2m}^\infty J(w)
 \left(\frac1{z-w}-\frac1{z_0-w}\right)\,dw,\qquad
 D(z)=z-\mu_R-S(z).
\]
V46 proves \(S_\Lambda\to S\) locally uniformly on the slit plane.

\begin{lemma}[Imaginary-part bound]
\label{lem:v20v47-denominator}
For nonreal \(z\),
\[
 \operatorname{Im}D(z)=\operatorname{Im}z
 \left(1+\int_{2m}^{\infty}\frac{J(w)}{|z-w|^2}\,dw\right).
\]
The analogous formula holds with cutoffs.  In particular
\[
 |D(z)|,\ |D_\Lambda(z)|\geq|\operatorname{Im}z|.
\]
\end{lemma}
\begin{proof}
The subtraction term at \(z_0\) is real.  The imaginary part of
\((z-w)^{-1}\) is \(-\operatorname{Im}z/|z-w|^2\).
This imaginary-part integral is absolutely convergent even though
the unsubtracted real part diverges.  Substitution into \(D\)
gives the formula and its lower bound.
\end{proof}

\subsection{Full norm-resolvent limit}

\begin{theorem}[Convergence of all resolvent blocks]
\label{thm:v20v47-limit}
For every nonreal \(z\),
\[
 (z-H_\Lambda)^{-1}\longrightarrow R(z)
 \quad\hbox{in operator norm},
\]
where, for \((a,f)\in\mathbb C\oplus L^2\),
\[
 \begin{split}
 A&=D(z)^{-1}\bigl(a+\ell_z(f)\bigr),\\
 R(z)(a,f)&=(A,\ r(z)f+h_z A).
 \end{split}
\]
The convergence is locally uniform off the real axis.  Equivalently,
\[
 R(z)=
 \begin{pmatrix}
 D(z)^{-1}&D(z)^{-1}\ell_z\\
 h_zD(z)^{-1}&r(z)+h_zD(z)^{-1}\ell_z
 \end{pmatrix}.
\]
\end{theorem}
\begin{proof}
Solve the lower row of \((z-H_\Lambda)(A,b)=(a,f)\):
\[
 b=r(z)f+h_{\Lambda,z}A.
\]
Substitution into the upper row gives
\(D_\Lambda(z)A=a+\ell_{\Lambda,z}(f)\), since the bare energy
was chosen with the displayed counterterm.
This gives the same four-block formula at finite cutoff.
Lemma \ref{lem:v20v47-weight} gives convergence of both row and
column factors.  Local uniform convergence of \(S_\Lambda\) and
Lemma \ref{lem:v20v47-denominator} give convergence of
\(D_\Lambda^{-1}\).  Products of these convergent bounded factors
converge in operator norm; the free block \(r(z)\) does not change.
\end{proof}

\subsection{Why the limit is an actual self-adjoint resolvent}

\begin{theorem}[Renormalized operator exists and retains the parent]
\label{thm:v20v47-operator}
There is a unique densely defined self-adjoint operator \(H_R\) on
\(\mathcal H\) with
\[
 R(z)=(z-H_R)^{-1},\qquad z\notin\mathbb R.
\]
Thus \(H_\Lambda\to H_R\) in norm-resolvent sense.
\end{theorem}
\begin{proof}
Norm limits of the cutoff resolvents give
\[
 R(z)-R(w)=(w-z)R(z)R(w),\qquad
 R(z)^*=R(\bar z),\qquad
 \|R(z)\|\leq|\operatorname{Im}z|^{-1}.
\]
If \(R(z)(a,f)=0\), the scalar output \(A\) is zero.
The lower output is then \(r(z)f=0\).  Multiplication by
\((z-\omega)^{-1}\) is injective, so \(f=0\), and the top equation
gives \(a=0\).  Thus \(R(z)\) is injective.
Since \(R(z)^*=R(\bar z)\) is also injective, its range is dense.

The resolvent identity shows all ranges coincide: explicitly
\(R(z)=R(w)[I+(w-z)R(z)]\), and conversely with \(z,w\)
interchanged.  On this common dense range define
\[
 H_R=z-R(z)^{-1}.
\]
The same identity shows independence of \(z\); the inverse of an
injective bounded operator on its range is closed, so \(H_R\)
is closed.  To check symmetry, take \(u=R(z)x\) and
\(v=R(\bar z)y\).  Then
\[
 \langle H_Ru,v\rangle
 =\bar z\langle u,v\rangle-\langle x,v\rangle,\quad
 \langle u,H_Rv\rangle
 =\bar z\langle u,v\rangle-\langle u,y\rangle.
\]
The last terms coincide by \(R(z)^*=R(\bar z)\).
Both \(z-H_R\) and \(\bar z-H_R\) map the common domain onto
the whole Hilbert space.  A densely defined symmetric operator
with these two surjectivity properties is self-adjoint:
the orthogonal complements of the ranges are the two deficiency
spaces, so both vanish.  The defining identity gives its resolvent;
uniqueness follows since one resolvent determines its operator.
\end{proof}

\subsection{The domain contains a singular dressing}

\begin{proposition}[Explicit domain and absence of the bare parent]
\label{prop:v20v47-domain}
For any fixed nonreal \(z\),
\[
 D(H_R)=\{(A,b): A\in\mathbb C,\ b-h_zA\in D(\omega)\}.
\]
In particular the bare parent \(e_0=(1,0)\) does not belong to
\(D(H_R)\) when \(y\neq0\).
\end{proposition}
\begin{proof}
The range formula in Theorem \ref{thm:v20v47-limit} gives
\(b-h_zA=r(z)f\in D(\omega)\).
Conversely, given such \(A,b\), set
\(f=(z-\omega)(b-h_zA)\in L^2\) and
\(a=D(z)A-\ell_z(f)\).  The same formula gives
\(R(z)(a,f)=(A,b)\), proving the domain equality.
Finally
\[
 \omega h_z=\frac{\omega v}{z-\omega}=-v+z h_z.
\]
The first term is not square-integrable, while the second is;
hence \(h_z\notin D(\omega)\).  For \(e_0\) the required
condition would be \(-h_z\in D(\omega)\), which fails.
\end{proof}

The domain is independent of the chosen nonreal \(z\), as also follows
directly from the common-range argument.  It allows a correlated
parent/pair dressing, rather than the uncut bounded off-diagonal
operator of V46.  The pure parent remains a valid Hilbert-space
initial state and evolves unitarily under \(H_R\), but its second
spectral moment is infinite.  Consequently the fixed-cutoff
quadratic survival coefficient in V46 has no justified finite
cutoff-removal limit.  This domain fact does not by itself prove
or disprove a long-time exponential approximation.

\subsection{Independence of regulator profile and subtraction convention}

\begin{corollary}[Regulator and reference-point comparison]
\label{cor:v20v47-reference}
All cutoff profiles satisfying V46's support and boundedness
conditions give the same \(H_R\) at fixed \((z_0,\mu_R)\).
If \(z_1<2m\) is used instead, the same operator is obtained with
\[
 \mu_R^{(1)}=\mu_R+S(z_1).
\]
\end{corollary}
\begin{proof}
The norm-limit formula depends only on \(v\), \(S\), and the
renormalized parameter.  Those limits are independent of the
cutoff shape by the integrable tail estimates.
With subtraction point \(z_1\), the subtracted function is
\(S_1(z)=S(z)-S(z_1)\).  Therefore
\[
 z-\mu_R^{(1)}-S_1(z)=z-\mu_R-S(z),
\]
and all other resolvent blocks are identical.
\end{proof}

This is independence within the declared one-parameter subtraction
scheme.  A measured mass or resonance energy still requires an
additional renormalization condition; it is not automatically equal
to the parameter \(M\) used in the free vertex normalization.

\subsection{Exact algebra check and scope}

The verifier
\[
 \texttt{scripts/verify\_sm\_v20\_v47\_full\_resolvent.py}
\]
checks the full Schur inverse, the resolvent and adjoint identities,
and the subtraction-point shift on an exact rational finite-bath
example at nonreal spectral parameters.  It supplements, and does
not replace, the analytic norm estimates above.
The script SHA--256 is
\[
\texttt{9359B6D15242BF568408AC6C2035576590DA8A944AAEA37FE51F7DFD457F3C79}.
\]
The artifact
\(\texttt{artifacts/sm\_v20\_v47\_full\_resolvent.json}\)
has SHA--256
\[
\texttt{66F86B80A42B5B895F117E8662555B75A48290D074896AE6B66E6F7DDB9F4217}.
\]

\begin{firewall}[A renormalized block is not a renormalized field theory]
\label{fire:v20v47}
Theorem \ref{thm:v20v47-operator} removes the ultraviolet cutoff
only for the explicitly projected one-parent resonant Hamiltonian
in one spatial dimension.  It constructs no full Yukawa Fock-space
Hamiltonian, chiral Standard Model sector, positive Euclidean
physical measure, or coupled Einstein continuum.  Self-adjointness
alone gives neither a resonance lifetime nor a kinetic limit.
\end{firewall}

\begin{decision}[V48 acquisition]
\label{dec:v20v47-next}
Determine the renormalized parent's spectral measure from the
resolvent boundary values.  Identify any bound-state pole below
the pair threshold and its residue, and determine the continuous
density above threshold.  This decides whether the parent can
decay completely before an exponential approximation is proposed.
\end{decision}

\begin{assessment}[V47 acquisition]
\label{ass:v20v47}
One subtraction now yields a complete norm-resolvent construction,
not just a convergent scalar denominator.  Injectivity and dense
range prove that the limit is a self-adjoint operator on the
original Hilbert space.  Its dressed domain excludes the bare
parent from the operator domain, making the cutoff short-time
Taylor law nonuniform.  Spectral analysis is the next dependency.
\end{assessment}

\begin{record}[V47 append record]
\label{rec:v20v47}
The predecessor V46 has SHA--256
\texttt{01956D7BE8328C2211BA17408334507D82321F5CEFC91A7C5963DE4EB18C043B}.
Its complete prefix is retained.  Only the immediate active SM
predecessor is replaced on validation.  The next companion audit
and ten-version sweep remain V50.
\end{record}



\section{V48: the parent spectral measure and the bound-state obstruction to decay}
\label{sec:v20v48}

V47 constructed the renormalized one-parent Hamiltonian.  Its existence
does not imply decay: a point mass in the parent's spectral measure
would survive forever.  We now classify that measure completely for
the rest-parent block.  The result distinguishes an isolated bound
state, a critical threshold without an atom, and a purely continuous
parent spectrum.  These are statements about this projected model.

\subsection{Threshold function and the critical parameter}

Write \(a=2m>0\), \(c=y^2/(4\pi M)>0\), and retain the real subtraction
point \(z_0<a\).  V46--V47 give
\[
 J(w)=c\sqrt{1-a^2/w^2},\quad w>a,\qquad
 S(z)=\int_a^\infty J(w)
 \left(\frac1{z-w}-\frac1{z_0-w}\right)\,dw,
\]
\[
 D(z)=z-\mu_R-S(z),\qquad
 F(z)=\langle e_0,(z-H_R)^{-1}e_0\rangle=D(z)^{-1}.
\]
Let \(\nu\) be the probability spectral measure of the normalized
parent \(e_0=(1,0)\), so \(F(z)=\int(z-x)^{-1}\,d\nu(x)\).

\begin{lemma}[Finite threshold value and monotonicity]
\label{lem:v20v48-threshold}
The real limit \(S(a):=\lim_{x\uparrow a}S(x)\) is finite.
For \(x<a\),
\[
 D'(x)=1+\int_a^\infty\frac{J(w)}{(x-w)^2}\,dw>1.
\]
Furthermore \(D(x)\to-\infty\) as \(x\to-\infty\).  Define
\[
 \mu_{\rm crit}=a-S(a).
\]
Then \(D(a-)=\mu_{\rm crit}-\mu_R\).
\end{lemma}
\begin{proof}
Near threshold, \(J(a+u)=c\sqrt{2/a}\sqrt u+O(u^{3/2})\).
At \(x=a\) the singular part of the integrand is therefore
of order \(u^{-1/2}\), which is integrable.  At infinity the
subtracted integrand is \(O(w^{-2})\).  Dominated convergence
on the threshold and tail pieces proves the finite limit.
Differentiating on compact intervals below \(a\) gives
\(S'(x)=-\int J(w)/(x-w)^2\,dw\), hence the formula.
For \(x<z_0\), each subtracted integrand is nonnegative, so
\(S(x)\geq0\) and \(D(x)\leq x-\mu_R\to-\infty\).
\end{proof}

\begin{theorem}[Unique subthreshold pole and its weight]
\label{thm:v20v48-bound}
There is exactly one root \(x_b<a\) of \(D\) if
\(\mu_R<\mu_{\rm crit}\), and no such root if
\(\mu_R\geq\mu_{\rm crit}\).  When it exists, it contributes
the parent spectral mass
\[
 Z_b=\frac1{D'(x_b)}
 =\left(1+\int_a^\infty\frac{J(w)}{(x_b-w)^2}\,dw\right)^{-1},
 \qquad 0<Z_b<1.
\]
\end{theorem}
\begin{proof}
Strict monotonicity, the negative limit at minus infinity, and
the threshold value prove existence and uniqueness by the
intermediate value theorem.  If the threshold value is zero,
strict monotonicity still makes \(D(x)<0\) at every \(x<a\).
At a root the derivative is finite and strictly positive, so
\(F(z)\) has a simple pole with residue \(1/D'(x_b)\).
For a Stieltjes transform of a self-adjoint spectral measure,
this residue is its point mass.  Equivalently the full resolvent
formula of V47 gives the eigenvector proportional to
\((1,h_{x_b})\), whose squared norm is \(D'(x_b)\);
its normalized overlap squared with \(e_0\) is \(Z_b\).
\end{proof}

\subsection{Continuous spectrum seen by the parent}

Define the real principal-value function, with the subtraction kept
inside the integral,
\[
 s(x)=\operatorname{PV}\int_a^\infty J(w)
 \left(\frac1{x-w}-\frac1{z_0-w}\right)\,dw,\qquad x>a.
\]

\begin{theorem}[Absolutely continuous parent density above threshold]
\label{thm:v20v48-ac}
On \((a,\infty)\), the parent spectral measure has density
\[
 \rho(x)=
 \frac{J(x)}
 {(x-\mu_R-s(x))^2+\pi^2J(x)^2}.
\]
There are no parent atoms or singular-continuous components on
this open interval.
\end{theorem}
\begin{proof}
On any compact subinterval of \((a,\infty)\), \(J\) is smooth.
Split the integral into a local smooth piece and a remote
subtracted tail.  In the local piece subtract the numerator's
value at \(w=x\); the resulting quotient is continuous.
The remaining elementary Cauchy integral and the uniformly
convergent tail give the boundary value
\[
 S(x+i0)=s(x)-i\pi J(x)
\]
locally uniformly in \(x\).  Thus
\[
 F(x+i0)=\frac1{x-\mu_R-s(x)+i\pi J(x)}.
\]
The strictly positive \(J(x)\) bounds the denominator away
from zero on each compact subinterval, including for sufficiently
small positive imaginary parts.  Consequently \(F\) has a
locally bounded continuous boundary value.  Stieltjes inversion,
\(\rho=-\pi^{-1}\operatorname{Im}F(x+i0)\), gives the formula.
More explicitly the Poisson-kernel approximants to \(\nu\)
converge uniformly on compact subintervals to this continuous
density.  Testing against compactly supported continuous functions
identifies the entire restricted measure with \(\rho(x)\,dx\),
excluding an additional singular measure there.
\end{proof}

\subsection{The critical threshold is not a spectral atom}

\begin{theorem}[No parent point mass at \(a\)]
\label{thm:v20v48-noatom}
For every \(\mu_R\), including \(\mu_R=\mu_{\rm crit}\),
\(\nu(\{a\})=0\).
\end{theorem}
\begin{proof}
V47's imaginary-part identity at \(z=a+i\eta\), \(\eta>0\),
gives
\[
 \operatorname{Im}D(a+i\eta)
 =\eta+\eta\int_0^\infty
 \frac{J(a+u)}{u^2+\eta^2}\,du.
\]
There are \(u_0,C>0\) with \(J(a+u)\geq C\sqrt u\)
for \(0<u<u_0\).  For sufficiently small \(\eta\), restrict
the integral to \(\eta<u<2\eta\).  Substitution \(u=\eta t\)
gives a lower bound \(C'\sqrt\eta\), with \(C'>0\).
Therefore
\[
 |F(a+i\eta)|\leq(C'\sqrt\eta)^{-1},\qquad
 \eta|F(a+i\eta)|\longrightarrow0.
\]
For a finite spectral measure,
\(-\eta\operatorname{Im}F(a+i\eta)\to\nu(\{a\})\):
this follows directly by dominated convergence in
\(\int \eta^2/((a-x)^2+\eta^2)\,d\nu(x)\).
The bound proves the claim independently of whether \(D(a-)\)
vanishes.
\end{proof}

\begin{corollary}[Complete parent measure and survival limit]
\label{cor:v20v48-survival}
The parent measure is exactly
\[
 d\nu(x)=\rho(x)\mathbf1_{(a,\infty)}\,dx
 +\begin{cases}Z_b\delta_{x_b},&\mu_R<\mu_{\rm crit},\\
 0,&\mu_R\geq\mu_{\rm crit}.
 \end{cases}
\]
In particular \(\rho\in L^1\), with integral \(1-Z_b\) or \(1\).
The parent survival amplitude obeys
\[
 \langle e_0,e^{-itH_R}e_0\rangle
 =\begin{cases}
 Z_b e^{-itx_b}+o(1),&\mu_R<\mu_{\rm crit},\\
 o(1),&\mu_R\geq\mu_{\rm crit},
 \end{cases}\qquad t\to+\infty.
\]
Its survival probability therefore tends to \(Z_b^2\) or zero.
\end{corollary}
\begin{proof}
Below \(a\), the resolvent is analytic except at the root
classified above, so the measure there consists solely of
that pole if present.  On the open continuum the preceding
theorem excludes all other components; at the one remaining
point \(a\), there is no atom.  These sets exhaust the real
line, giving the complete measure.  Normalization follows
from \(\|e_0\|=1\) and the spectral theorem.
The Fourier transform of its \(L^1\) density tends to zero
by the Riemann--Lebesgue lemma.  The atom gives the stated
oscillatory term, and squaring gives the probability limit.
\end{proof}

Complete decay here means only that the overlap with the original
parent vector tends to zero.  It does not by itself establish
asymptotic completeness, a many-body particle interpretation, or
an exponential decay law.

\subsection{An explicit critical energy in one subtraction convention}

\begin{proposition}[Closed threshold integral at \(z_0=0\)]
\label{prop:v20v48-explicit}
Since \(a>0\), subtraction at \(z_0=0\) is allowed.  In this
convention,
\[
 S(a)=-c(1+\pi/2),\qquad
 \mu_{\rm crit}=2m+\frac{y^2}{4\pi M}(1+\pi/2).
\]
\end{proposition}
\begin{proof}
Keeping the subtraction combined,
\[
 S(a)=c\int_a^\infty\sqrt{1-a^2/w^2}
 \left(\frac1{a-w}+\frac1w\right)\,dw.
\]
The substitution \(x=a/w\) reduces this to
\[
 -c\int_0^1\frac{\sqrt{1-x^2}}{1-x}\,dx.
\]
Putting \(x=\cos\theta\) gives
\(-c\int_0^{\pi/2}(1+\cos\theta)\,d\theta=-c(1+\pi/2)\).
\end{proof}

This value refers to the specified subtraction convention, not
to a physical mass observable.  V47's reference-point change
shifts both \(\mu_R\) and \(\mu_{\rm crit}\) by the same amount;
the inequality deciding the bound state is invariant.

\subsection{Verification and next step}

The symbolic verifier
\[
 \texttt{scripts/verify\_sm\_v20\_v48\_spectral\_measure.py}
\]
checks the threshold integral, the Stieltjes-density sign, and
the \(\sqrt\eta\) scaling used in the no-atom proof.  Its SHA--256 is
\[
\texttt{E2E999CAA0C0DFF763060A0B7EFB35D9F506B5979AE2627BFC30B9BE8F7337EB}.
\]
The artifact
\(\texttt{artifacts/sm\_v20\_v48\_spectral\_measure.json}\)
has SHA--256
\[
\texttt{2E8E61655EF97E94CA520A9762FF22EF62ED247655B3B25F2BA9CA43203963F0}.
\]
These algebraic checks support the analytic proofs and do not
substitute for spectral inversion or the limiting arguments.

\begin{firewall}[Projected spectral theorem only]
\label{fire:v20v48}
The classification concerns the parent's spectral measure in the
renormalized rest-parent resonant block.  It proves no exponential
decay rate, weak-coupling kinetic theorem, unprojected Yukawa
Hamiltonian, Standard Model quantum state, or Einstein continuum.
\end{firewall}

\begin{decision}[V49 acquisition]
\label{dec:v20v48-next}
Derive threshold asymptotics of the continuous parent density
in the noncritical and critical regimes, then determine what
long-time survival tails follow with controlled Fourier
remainders.  The exact threshold tail must be understood before
an intermediate-time exponential approximation is claimed.
V50 remains the companion audit and ten-version sweep.
\end{decision}

\begin{assessment}[V48 acquisition]
\label{ass:v20v48}
The parent spectral measure is now explicit up to a real
principal-value integral.  A sharp renormalized-energy criterion
decides the unique bound-state contribution.  The critical
threshold has no atom, and the continuous component gives a
rigorous long-time survival limit.  This settles complete versus
incomplete parent decay within the constructed block.
\end{assessment}

\begin{record}[V48 append record]
\label{rec:v20v48}
The predecessor V47 has SHA--256
\texttt{FACB112191C8F19EED87C65E1786E0F19FE71125444F515A50D8ACE78B864FD6}.
Its complete prefix is retained.  Only the immediate active
SM predecessor is replaced after validation.
\end{record}



\section{V49: threshold density and controlled algebraic survival tails}
\label{sec:v20v49}

V48 determined whether a bound state prevents complete parent decay.
This section sharpens its continuous-spectrum result to explicit
long-time asymptotics.  The critical threshold and the ordinary
threshold have different powers.  The Fourier error is controlled
by integrable derivatives, not inferred from a fitted curve.

\subsection{Local boundary regularity}

Retain \(a=2m>0\), \(c=y^2/(4\pi M)>0\), and set
\[
 j=c\sqrt{2/a},\qquad
 \delta=\mu_{\rm crit}-\mu_R,\qquad
 A(x)=x-\mu_R-s(x).
\]
Then \(J(a+u)=j\sqrt u(1+O(u))\).

\begin{lemma}[Analytic real part at the threshold]
\label{lem:v20v49-analytic}
For sufficiently small \(u\geq0\),
\[
 J(a+u)=\sqrt u\,j(u),\qquad
 A(a+u)=\delta+u b(u),
\]
where \(j(u),b(u)\) are real analytic and \(j(0)=j>0\).
\end{lemma}
\begin{proof}
The formula
\[
 J(a+u)=c\frac{\sqrt{u(2a+u)}}{a+u}
\]
proves the first assertion.  For the real part of the Cauchy
transform, split the energy integral at \(a+\epsilon\) with
\(\epsilon>0\) smaller than the convergence radius of \(j(u)\).
The remote, subtracted part is analytic near \(a\).
On the local part expand \(j(r)\) in its convergent power series.
The base principal-value integral is, for \(0<u<\epsilon\),
\[
 \operatorname{PV}\int_0^\epsilon
 \frac{\sqrt r}{u-r}\,dr
 =-2\sqrt\epsilon+
 \sqrt u\log\frac{\sqrt\epsilon+\sqrt u}{\sqrt\epsilon-\sqrt u}.
\]
Its logarithm has a convergent odd power series in \(\sqrt u\);
multiplication by \(\sqrt u\) makes the result analytic in \(u\).
Polynomial division reduces each higher term
\(r^{n+1/2}/(u-r)\) to a polynomial in \(u\) plus \(u^n\)
times this base integral.  Uniform convergence on smaller
neighborhoods justifies summation and differentiation.
The subtraction term is independent of \(u\).
Thus \(s(a+u)\) has a real analytic right extension with value
\(S(a)\), as used in V48.  Since
\(a-\mu_R-S(a)=\delta\), the asserted form of \(A\) follows.
\end{proof}

\begin{theorem}[The two threshold laws]
\label{thm:v20v49-density}
For the parent density of V48,
\[
 \rho(a+u)=
 \begin{cases}
 \displaystyle \frac{j}{\delta^2}u^{1/2}
                   +u^{3/2}r_1(u),&\delta\neq0,\\
 \displaystyle \frac1{\pi^2j}u^{-1/2}
                   +u^{1/2}r_0(u),&\delta=0,
 \end{cases}
\]
where \(r_1,r_0\) are real analytic near zero.  In particular
the critical singularity is integrable and is not an atom.
\end{theorem}
\begin{proof}
Use the exact identity
\[
 \rho(a+u)=\frac{\sqrt u\,j(u)}
 {(\delta+u b(u))^2+\pi^2u j(u)^2}.
\]
If \(\delta\neq0\), the denominator has nonzero constant
term \(\delta^2\), so its reciprocal is analytic.
If \(\delta=0\), factor out \(u\); the remaining denominator
has nonzero constant term \(\pi^2j^2\).
Expanding the resulting analytic factors gives both formulas.
\end{proof}

\subsection{Control away from the endpoint}

\begin{lemma}[Integrable derivatives on the remote part]
\label{lem:v20v49-remote}
For any smooth function \(\vartheta\) equal to one near \(a\)
and zero outside a bounded interval, the function
\((1-\vartheta)\rho\), extended by zero below threshold, has
two integrable derivatives.  Its Fourier transform is \(O(t^{-2})\).
\end{lemma}
\begin{proof}
On compact subintervals of \((a,\infty)\) all factors in
V48's density are smooth and its denominator is strictly positive.
At infinity write \(J(w)=c+g(w)\), where
\(g^{(k)}(w)=O(w^{-2-k})\) for large \(w\).
The constant part of the subtracted real Cauchy transform is
\[
 c\log\frac{|x-a|}{|z_0-a|}.
\]
For completeness, bounds for the remainder and its first two
derivatives follow by splitting at \(w=x/2,3x/2\).
On the middle interval put \(w=xt\), \(1/2<t<3/2\);
subtract \(g(x)\) from the numerator in the principal-value
integral.  The resulting difference quotient and its first
two \(x\) derivatives are bounded using the displayed decay
of \(g\).  The principal value of the constant numerator on
this symmetric interval is zero.
On the other two intervals the denominator is bounded away
from zero in units of \(x\); differentiation, including the
moving endpoints, is bounded directly by the same decay.
The fixed bounded segment near \(a\) has integrable \(g\)
and denominators of order \(x\), so it causes no difficulty.
These estimates give
\[
 s(x)=O(\log x),\qquad s'(x)=O(x^{-1}),\qquad
 s''(x)=O(x^{-2}).
\]
Consequently \(A(x)=x+O(\log x)\), \(A'(x)=1+O(x^{-1})\),
and \(A''(x)=O(x^{-2})\).  Differentiating
\(\rho=J/(A^2+\pi^2J^2)\) gives
\[
 \rho(x)=O(x^{-2}),\quad \rho'(x)=O(x^{-3}),\quad
 \rho''(x)=O(x^{-4}).
\]
Thus the remote function and its first two derivatives are
integrable and vanish at both integration endpoints.
Two integrations by parts give the Fourier bound.
\end{proof}

\subsection{An endpoint Fourier estimate with explicit remainders}

\begin{lemma}[The two endpoint transforms]
\label{lem:v20v49-Fourier}
Suppose an integrable function \(f\) on \(u>0\) has the remote
regularity of Lemma \ref{lem:v20v49-remote}.
If near zero \(f(u)=B u^{1/2}+u^{3/2}r(u)\) with analytic
\(r\), then
\[
 \int_0^\infty e^{-itu}f(u)\,du
 =B\frac{\sqrt\pi}{2}e^{-3\pi i/4}t^{-3/2}+O(t^{-2}).
\]
If instead \(f(u)=B u^{-1/2}+u^{1/2}r(u)\), then
\[
 \int_0^\infty e^{-itu}f(u)\,du
 =B\sqrt\pi e^{-\pi i/4}t^{-1/2}+O(t^{-1}).
\]
\end{lemma}
\begin{proof}
Subtract \(B u^\alpha e^{-u}\), with \(\alpha=1/2\) or
\(-1/2\), from \(f\).  The transform of the subtracted
function is exactly
\[
 B\Gamma(\alpha+1)(1+it)^{-\alpha-1},
\]
by the convergent Laplace integral.  Its leading large-\(t\)
term is the one displayed, with an error
\(O(t^{-\alpha-2})\).
For \(\alpha=1/2\), the remaining function is \(O(u^{3/2})\),
its first derivative is \(O(u^{1/2})\), and its second derivative
is \(O(u^{-1/2})\) at zero.  The first two boundary terms
vanish and two integrations by parts give \(O(t^{-2})\).
For \(\alpha=-1/2\), the remainder is \(O(u^{1/2})\) and
its first derivative is integrable at zero; one integration
by parts gives \(O(t^{-1})\).
The remote assumptions justify the same operations at infinity.
\end{proof}

\subsection{Survival amplitudes and probabilities}

\begin{theorem}[Controlled long-time tails]
\label{thm:v20v49-tails}
Let \(\mathcal A(t)=\langle e_0,e^{-itH_R}e_0\rangle\).
For fixed parameters with \(\delta\neq0\),
\[
 \mathcal A(t)=
 \mathbf1_{\{\delta>0\}}Z_b e^{-itx_b}
 e^{-iat}\frac{j\sqrt\pi}{2\delta^2}
 e^{-3\pi i/4}t^{-3/2}+O(t^{-2}).
\]
At criticality \(\delta=0\),
\[
 \mathcal A(t)=e^{-iat}\frac{e^{-\pi i/4}}{\pi^{3/2}j}
 t^{-1/2}+O(t^{-1}).
\]
In the no-bound-state regimes, the survival probabilities are
\[
 |\mathcal A(t)|^2=
 \begin{cases}
 \displaystyle \frac{\pi j^2}{4\delta^4}t^{-3}
                        +O(t^{-7/2}),&\delta<0,\\
 \displaystyle \frac1{\pi^3j^2}t^{-1}
                        +O(t^{-3/2}),&\delta=0.
 \end{cases}
\]
\end{theorem}
\begin{proof}
V48 decomposes the spectral measure into its density plus
the subthreshold atom precisely when \(\delta>0\).
Apply the endpoint Fourier lemma to
\(f(u)=\rho(a+u)\), using the threshold expansions and
remote derivative bounds just proved.  Restore the phase
\(e^{-iat}\) and the atom.  In the cases without an atom,
squaring the amplitudes gives the error powers by the
cross-term estimate.  All leading constants are positive
in modulus because \(j>0\).
\end{proof}

For \(\delta>0\) the amplitude includes the nondecaying
bound-state oscillation, and the leading continuous term
interferes with it.  Its survival probability still tends
to \(Z_b^2\), as proved in V48.
The estimates above are for fixed parameters; their constants
are not uniform as \(\delta\to0\).  In particular the divergent
\(\delta^{-2}\) coefficient is not a valid way to take the
critical limit.

\begin{corollary}[No positive-rate exponential as the late-time law]
\label{cor:v20v49-nonexponential}
In the no-bound-state regimes the survival probability is not
asymptotic to \(C e^{-\Gamma t}\) for any \(C,\Gamma>0\).
\end{corollary}
\begin{proof}
The nonzero algebraic leading coefficients above imply that
the ratio to \(C e^{-\Gamma t}\) diverges as \(t\to\infty\).
\end{proof}

This does not preclude an intermediate-time exponential window
or a weak-coupling limit on a rescaled time axis.  Such claims
require different estimates and an explicit order of limits.

\subsection{Verification and programme decision}

The symbolic verifier
\[
 \texttt{scripts/verify\_sm\_v20\_v49\_threshold\_tails.py}
\]
checks both leading density coefficients and the gamma-function
constants.  The script SHA--256 is
\[
\texttt{1E949B9567BC6BDEB6AFB7E49004E26D7E265F1A59AA1E3C884818145E3B0B4D}.
\]
The artifact
\(\texttt{artifacts/sm\_v20\_v49\_threshold\_tails.json}\)
has SHA--256
\[
\texttt{AC850F7E451DF9862D68E887307E17D4761FF014F9E0BB24B1083FFEE11AB910}.
\]
The remainder estimates are analytic proofs in this note,
not conclusions from the symbolic leading-coefficient checks.

\begin{firewall}[Fixed-parameter projected dynamics]
\label{fire:v20v49}
These tails concern the ultraviolet-renormalized rest-parent
resonant block at fixed parameters.  They prove no uniform
weak-coupling law, many-body kinetic closure, unprojected
Yukawa theory, or SM--Einstein continuum.
\end{firewall}

\begin{decision}[V50 audit and subsequent acquisition]
\label{dec:v20v49-next}
V50 performs the compulsory companion audit and reviews V40--V49.
After that checkpoint, the next substantive target is a
weak-coupling limit with the renormalized parent energy kept
strictly above threshold.  It must derive a rescaled spectral
limit, distinguish its time range from the algebraic late-time
tails, and track the normalization fixed in V43.
\end{decision}

\begin{assessment}[V49 acquisition]
\label{ass:v20v49}
The continuous parent amplitude has a controlled \(t^{-3/2}\)
tail off criticality and a \(t^{-1/2}\) tail at criticality.
The corresponding decay probabilities are algebraic, with
explicit coefficients and Fourier remainder bounds.
This supplies the late-time comparison needed before any
weak-coupling exponential approximation is proposed.
\end{assessment}

\begin{record}[V49 append record]
\label{rec:v20v49}
The predecessor V48 has SHA--256
\texttt{4596C79E45C8CB4DC8DA5F9E2C2B3A569200151C22E5B66D69725542EF690039}.
Its complete prefix is retained.  Only the immediate active
SM predecessor is replaced after validation.
\end{record}



\section{V50: companion audit and review of V40--V49}
\label{sec:v20v50}

The required companion audit was performed on the current active
files, read-only.  The snapshots are:
\begin{itemize}
\item Main YM V47, 652916 bytes, SHA--256
\texttt{A578DC1E7F3460DD530DBB7B270A39A546DDC02BACCC2272CED3A22EBE2DC0B5}.
\item NS V26, 278283 bytes, SHA--256
\texttt{82C887F8C2C69E4F28FAD7C3DC8DE5B9099CB5A4BF431EBA1FFF3E91935978AD}.
\item Parallel YM V5, 54174 bytes, SHA--256
\texttt{306F1BB84CFA8A8D47EA569EC17E9545F9867C8D32B548EC9D9C444B4F3C0512}.
\end{itemize}
These are respectively the files
\texttt{Renewal\_Yang\_Mills\_Mass\_Gap\_Research\_Notes\_V47.tex},
\texttt{Renewal\_NS\_Stability\_Research\_Notes\_V26.tex}, and
\texttt{Renewal\_YM\_Parallel\_Research\_Notes\_V5.tex}.

Since V45, main YM has advanced from V43 to V47.  Its new analysis
includes level-uniform one-step covariance estimates and
square-summable composed transverse lifts, with conditional trace
estimates.  It explicitly separates those from the still-unproved
gauged Wilson response with fine-momentum derivatives.  NS and
parallel YM are unchanged by hash.

\begin{decision}[No cross-program theorem imported]
\label{dec:v20v50-audit}
The YM estimates concern spatial blocking and do not establish
convergence of a Yukawa parent spectral distribution on its
resonance scale.  No companion constant or theorem is imported.
The useful methodological agreement is to control the composed
quantity directly.  V51 therefore studies the normalized rescaled
spectral density rather than approximating separate terms of the
time-domain memory kernel.
\end{decision}

\subsection{The ten-version review}

V40 fixed the original direction of continuum triad densities and
kinetic--Einstein closure; V41 established finite-circle scaling.
V42 distinguished fixed-parent phase space from fixed-daughter
triad counting and proved a canonically normalized Yukawa benchmark.
V43 derived the canonical box vertex and found the omitted
standing-wave pair.  V44 derived exact CAR occupation identities,
corrected the bosonic interpretation of the Yukawa collision term,
and proved that equal occupation means need not determine
conversion curvature.  V45 audited the companions.

V46 then replaced an assumed instantaneous rate by an exact
one-parent memory equation and identified a logarithmic
ultraviolet divergence.  V47 constructed the corresponding
renormalized Hamiltonian by full norm-resolvent convergence.
V48 classified the parent's bound and continuous spectral
components; V49 proved controlled threshold tails.

\begin{assessment}[What the block actually closes]
\label{ass:v20v50-sweep}
Within the specified one-dimensional resonant one-parent model,
canonical normalization, removal of its ultraviolet cutoff,
the parent spectral measure, and fixed-coupling long-time
tails have analytic proofs.  The classical finite-mode results
and the canonical quantum model have not been identified as
equivalent full interacting theories.  V44's correlation
obstruction still prevents inferring a many-body kinetic closure.
No kinetic--Einstein continuum has been obtained.
\end{assessment}

\begin{decision}[V51 and the remaining bridge]
\label{dec:v20v50-next}
With renormalized parent energy a fixed positive distance above
threshold, scale the Yukawa coupling to zero and prove convergence
of the parent spectral density on the coupling-squared energy
scale.  If this limit is Cauchy, prove convergence in \(L^1\),
which controls survival amplitudes directly.  Keep this
one-parent exponential limit distinct from the still-open
many-body occupation hierarchy and from V49's fixed-coupling
late-time asymptotics.
\end{decision}

\begin{firewall}[No upgrade of scope in the review]
\label{fire:v20v50}
The one-parent operator is a projected resonant model.  Its
construction does not supply the full chiral Standard Model
Hilbert space, its interacting vacuum, its stress tensor, or
an Einstein continuum.  The review preserves those outstanding
requirements rather than redefining the endpoint.
\end{firewall}

\begin{record}[V50 append record]
\label{rec:v20v50}
The predecessor V49 has SHA--256
\texttt{291399A09318CC45C71750AC8E386BFFD174062DDEBED65AFBA9C413B995A2B4}.
Its complete prefix is retained.  Only the immediate active
SM file is replaced.  The next companion audit is V55 and
the next ten-version review is V60.
\end{record}



\section{V51: an exact weak-coupling spectral limit and exponential survival}
\label{sec:v20v51}

V49 rules out a positive-rate exponential as the fixed-coupling
late-time asymptotic law.  That statement leaves open a different
question: what happens when the coupling tends to zero while time
is scaled by its inverse square?  For the renormalized one-parent
model we can answer this directly from the complete spectral measure
of V48.  No random-phase or independent-occupation closure is used.

\subsection{A fixed energy above threshold}

Fix \(a=2m>0\), the vertex normalization parameter \(M>0\),
a real subtraction point \(z_0<a\), and a nonzero reference
coupling \(y_0\).  Define
\[
 J_0(w)=\frac{y_0^2}{4\pi M}\sqrt{1-a^2/w^2},\qquad w>a,
\]
and let \(S_0\) be its subtracted self-energy as in V47.
Let \(s_0(x)=\operatorname{Re}S_0(x+i0)\) on \(x>a\).
Choose a fixed renormalized parent energy \(E_0>a\) and set
\[
 y_\lambda=\lambda y_0,\qquad \mu_R(\lambda)=E_0,
 \qquad 0<\lambda\to0.
\]
Each \(H_\lambda\) below is the ultraviolet-renormalized
operator constructed in V47 with this coupling; the ultraviolet
limit is taken before the present weak-coupling limit.
Its parent resolvent denominator is
\[
 D_\lambda(z)=z-E_0-\lambda^2S_0(z).
\]

\begin{lemma}[No bound state for sufficiently small coupling]
\label{lem:v20v51-noatom}
For all sufficiently small \(\lambda\), the parent's spectral
measure is purely absolutely continuous and has total mass one.
Its density is
\[
 \rho_\lambda(x)=
 \frac{\lambda^2J_0(x)}
 {(x-E_0-\lambda^2s_0(x))^2+\pi^2\lambda^4J_0(x)^2},
 \qquad x>a.
\]
\end{lemma}
\begin{proof}
V48's critical parameter is
\(\mu_{\rm crit}(\lambda)=a-\lambda^2S_0(a)\).
The threshold self-energy is finite, and \(E_0-a>0\)
is fixed.  Hence \(E_0>\mu_{\rm crit}(\lambda)\) for
all sufficiently small \(\lambda\).
V48 then excludes a bound-state atom and a threshold atom,
and supplies the displayed density with integral one.
\end{proof}

\subsection{The rescaled density converges in total mass}

Extend the density by zero below threshold and introduce
\[
 r_\lambda(u)=\lambda^2\rho_\lambda(E_0+\lambda^2u),
 \qquad u\in\mathbb R.
\]
Let
\[
 b=s_0(E_0),\qquad J_*=J_0(E_0)>0,\qquad
 r_*(u)=\frac{J_*}{(u-b)^2+\pi^2J_*^2}.
\]

\begin{theorem}[Cauchy limit in \(L^1\)]
\label{thm:v20v51-L1}
The rescaled densities satisfy
\[
 \lim_{\lambda\to0}\int_{\mathbb R}
 |r_\lambda(u)-r_*(u)|\,du=0.
\]
Thus the rescaled parent energy distribution converges in total
variation to the Cauchy law with center \(b\) and scale \(\pi J_*\).
\end{theorem}
\begin{proof}
For every fixed \(u\), the point \(E_0+\lambda^2u\)
lies above threshold for all sufficiently small \(\lambda\).
Substitution into the exact density gives
\[
 r_\lambda(u)=
 \frac{J_0(E_0+\lambda^2u)}
 {[u-s_0(E_0+\lambda^2u)]^2+
     \pi^2J_0(E_0+\lambda^2u)^2}.
\]
Continuity of \(J_0,s_0\) near \(E_0>a\) proves
pointwise convergence to \(r_*\).
The change of variables preserves mass, so
\(\int r_\lambda=1\).  Also
\[
 \int_{\mathbb R}r_*(u)\,du
 =\left[\frac1\pi\arctan\frac{u-b}{\pi J_*}\right]_{-\infty}^\infty=1.
\]
Since \(\min(r_\lambda,r_*)\to r_*\) pointwise and is
dominated by the integrable \(r_*\), dominated convergence gives
\(\int\min(r_\lambda,r_*)\to1\).
Finally
\[
 \int|r_\lambda-r_*|
 =\int r_\lambda+\int r_*-2\int\min(r_\lambda,r_*)
 \longrightarrow0.
\]
This argument controls mass escaping the rescaled window;
pointwise convergence alone would not have done so.
\end{proof}

\subsection{Uniform absolute convergence of the survival amplitude}

Let
\[
 \mathcal A_\lambda(t)=
 \langle e_0,e^{-itH_\lambda}e_0\rangle,\qquad
 B_\lambda(\tau)=
 e^{iE_0\tau/\lambda^2}\mathcal A_\lambda(\tau/\lambda^2).
\]

\begin{theorem}[Weak-coupling exponential law]
\label{thm:v20v51-exponential}
One has
\[
 \sup_{\tau\in\mathbb R}
 \left|B_\lambda(\tau)-
       e^{-ib\tau-\pi J_*|\tau|}\right|\longrightarrow0.
\]
In particular, for rescaled forward time \(\tau\geq0\),
\[
 |\mathcal A_\lambda(\tau/\lambda^2)|^2
 \longrightarrow e^{-2\pi J_*\tau}
\]
uniformly in absolute error on the entire forward half-line.
\end{theorem}
\begin{proof}
The spectral theorem and the change of variables above give
\[
 B_\lambda(\tau)=\int_{\mathbb R}e^{-iu\tau}r_\lambda(u)\,du.
\]
The Fourier transform of \(r_*\) is
\(e^{-ib\tau-\pi J_*|\tau|}\).
For example, translate by \(b\), decompose
\[
 \frac{J_*}{v^2+\pi^2J_*^2}
 =\frac1{2\pi i}
 \left(\frac1{v-i\pi J_*}-\frac1{v+i\pi J_*}\right),
\]
and close the contour in the lower half-plane for \(\tau>0\)
and the upper half-plane for \(\tau<0\).  The residues give
the stated transform; at zero it equals its unit mass.
For every real \(\tau\), the difference of the transforms is
bounded by \(\|r_\lambda-r_*\|_1\), independently of \(\tau\).
Theorem \ref{thm:v20v51-L1} proves the first assertion.
Both amplitudes have modulus at most one, so their squared
moduli differ by at most twice that bound.
\end{proof}

For \(\tau\geq0\) the limiting parent amplitude is a scalar
contraction semigroup.  This is convergence of one matrix element,
not a claim that the full unitary group converges in operator norm
to an irreversible semigroup, and not a Lindblad theorem for
arbitrary many-body states.

\subsection{Normalization and energy shift}

When \(E_0=M>a\), the limiting probability rate in rescaled
time is exactly
\[
 2\pi J_0(M)=\frac{y_0^2}{2M}\sqrt{1-a^2/M^2}
 =\frac{y_0^2 k_*}{M^2}.
\]
In physical time this corresponds to the leading weak-coupling
coefficient \(\lambda^2 y_0^2k_*/M^2\), matching the
canonically normalized benchmark of V42--V43.
The phase \(b\) is the order-\(\lambda^2\) energy shift in
the fixed subtraction convention.

\begin{corollary}[A mass-centered renormalization convention]
\label{cor:v20v51-shift}
If instead \(\mu_R(\lambda)=E_0-\lambda^2b\), the same
argument gives the centered limit
\[
 e^{iE_0\tau/\lambda^2}\mathcal A_\lambda(\tau/\lambda^2)
 \longrightarrow e^{-\pi J_*|\tau|}.
\]
\end{corollary}
\begin{proof}
In the rescaled denominator replace
\(u-s_0(E_0+\lambda^2u)\) by
\(u+b-s_0(E_0+\lambda^2u)\), which converges to \(u\).
The density remains normalized and has no bound atom for
small coupling.  The same \(L^1\) argument and Fourier
transform prove the result.
\end{proof}

\subsection{Why the algebraic tails are not contradicted}

For every fixed nonzero small \(\lambda\), V49 still gives
a nonzero algebraic late-time tail.  The present theorem
controls absolute error as \(\lambda\to0\), not relative
error after the exponential has become extremely small.
The Fourier \(L^1\) bound contains no relative-error estimate
and no convergence rate in \(\lambda\).  It therefore permits
an algebraic tail to dominate the exponentially tiny prediction
at sufficiently late times while their absolute difference
remains small.  Nor is the present proof uniform when the
energy gap \(E_0-a\) tends to zero with \(\lambda\); the
critical-threshold scaling is a different problem.

\subsection{Verification and next dependency}

The symbolic verifier
\[
 \texttt{scripts/verify\_sm\_v20\_v51\_weak\_coupling.py}
\]
checks the exact rescaling algebra and normalization of the
limiting Cauchy density.  Its SHA--256 is
\[
\texttt{A5CAC5A362198916066E341D2750253420861FD7309BBEC29311D33EB4FC441C}.
\]
The artifact
\(\texttt{artifacts/sm\_v20\_v51\_weak\_coupling.json}\)
has SHA--256
\[
\texttt{BA3A4AD6E3A0503D1BDB7283A5A90FC416FE4C53A4DE78D7E9813B227DB5574F}.
\]
The convergence proof uses the exact spectral normalization;
it is not inferred from those symbolic checks.

\begin{firewall}[A one-parent limit, not occupation closure]
\label{fire:v20v51}
The theorem is an ultraviolet-renormalized, weak-coupling
survival limit for a fixed rest-parent energy above threshold.
It does not close the correlated CAR hierarchy of V44,
describe arbitrary matter occupations, establish a moving
parent's interacting boost law, or construct the full
Standard Model--Einstein continuum.
\end{firewall}

\begin{decision}[V52 acquisition]
\label{dec:v20v51-next}
Return to V44's unresolved many-channel structure by allowing
several parent levels coupled to the same pair continuum.
Derive the matrix self-energy and its weak-coupling spectral
limit on a degenerate parent subspace, identifying the
conditions for off-diagonal decay and nondecaying combinations.
This tests whether independent scalar rates can even be used
before a many-occupation kinetic--Einstein closure is attempted.
\end{decision}

\begin{assessment}[V51 acquisition]
\label{ass:v20v51}
The exponential weak-coupling law is now proved, rather than
postulated, within the renormalized one-parent block.
The proof uses \(L^1\) convergence of the exact rescaled
spectral densities and gives uniform absolute amplitude
control.  Its rate matches the canonical vertex, while the
fixed-coupling algebraic tails and the outstanding many-body
correlation problem remain intact.
\end{assessment}

\begin{record}[V51 append record]
\label{rec:v20v51}
The predecessor V50 has SHA--256
\texttt{49919AFC25EBAA98A48C6BC6EF2A5F8BE777DD8DCFD1E2F4467D54FF96771B5D}.
Its complete prefix is retained.  Only the immediate active
SM predecessor is replaced.  The next companion audit is V55.
\end{record}



\section{V52: shared decay channels, matrix survival, and exact dark combinations}
\label{sec:v20v52}

V51 proves a scalar survival limit.  V44, however, exhibits
cross-channel coherences, so assigning that scalar law separately
to each named parent would require justification.  We now give a
finite-multiplicity extension with a common momentum profile.
It exhibits an obstruction before any many-occupation closure:
the channel coupling can have a kernel.  All statements below
concern the projected one-excitation model, not the unprojected
Yukawa Hamiltonian or a completed Standard Model.

\subsection{Precisely specified common-profile model}

Use the parameters and functions \(a,E_0,J_0,S_0,s_0\) of V51,
with \(E_0>a\).  Write
\[
 \omega(p)=2\sqrt{m^2+p^2},\quad a=2m,\qquad
 v(p)=\frac{y_0p}{\sqrt{4\pi M}\sqrt{m^2+p^2}}.
\]
Thus the push-forward of \(|v(p)|^2\,dp\) by \(\omega\)
has density \(J_0\).  Let \(n,q\) be fixed positive integers,
let \(C:\mathbb C^n\to\mathbb C^q\) be a fixed complex matrix,
and put \(K=C^*C\).  The Hilbert space is
\[
 \mathcal H=\mathbb C^n\oplus L^2(\mathbb R,dp;\mathbb C^q).
\]
The finite-dimensional summand is the parent space, with inclusion
\(\iota u=(u,0)\) and projection \(P=\iota^*\).
For the same real cutoffs as V46 set \(v_\Lambda=\chi_\Lambda v\),
\[
 (V_\Lambda u)(p)=v_\Lambda(p)Cu,\qquad
 \Sigma_{0,\Lambda}(z)=\int_{\mathbb R}
             \frac{|v_\Lambda(p)|^2}{z-\omega(p)}\,dp .
\]
For \(\lambda>0\) the cutoff operator is
\[
 H_{\lambda,\Lambda}=
 \begin{pmatrix}
 E_0I_n-\lambda^2\Sigma_{0,\Lambda}(z_0)K&
 \lambda V_\Lambda^*\\
 \lambda V_\Lambda&\omega I_q
 \end{pmatrix}.
\]
Its domain is \(\mathbb C^n\oplus D(\omega I_q)\).
The counterterm is a matrix, including any off-diagonal entries of
\(K\); retaining only its diagonal would define a different model.
The common-profile assumption includes identical dispersions in all
\(q\) continuum channels.  It is not automatic for distinct physical
decay thresholds.

\begin{theorem}[Construction and matrix self-energy]
\label{thm:v20v52-construction}
For each fixed \(\lambda>0\), the cutoff operators converge in
norm resolvent to a self-adjoint operator \(H_\lambda\).
For nonreal \(z\), its compressed resolvent is
\[
 P(z-H_\lambda)^{-1}\iota
   =\big[(z-E_0)I_n-\lambda^2 S_0(z)K\big]^{-1}.
\]
Writing \(K=\sum_{\kappa\geq0}\kappa Q_\kappa\) for its distinct
eigenvalues and orthogonal spectral projections, the subspace
\(\iota\ker C\) is reducing and \(H_\lambda=E_0I\) there.
Each positive eigenvalue \(\kappa\), with multiplicity, supplies
the scalar model of V47 with coupling \(\lambda\sqrt{\kappa}\,y_0\).
Uncoupled continuum directions are free.
\end{theorem}
\begin{proof}
Take a singular-value decomposition of \(C\).
The associated fixed unitaries on parent and channel coordinates
preserve \(E_0I_n\) and \(\omega I_q\), respectively.
They turn each nonzero singular value \(\sqrt{\kappa}\)
into an independent parent--continuum block.  Its parent energy is
\(E_0-\lambda^2\kappa\Sigma_{0,\Lambda}(z_0)\), exactly the
subtraction convention of V47 with the stated coupling.
All remaining parent directions have energy \(E_0\), and the
remaining continuum directions have multiplication operator
\(\omega\).  There are finitely many blocks.  V47's norm-resolvent
convergence for each block therefore gives norm-resolvent
convergence of their orthogonal sum and a self-adjoint limit.
This also constructs the domains without applying the unbounded
uncut vertex directly to an arbitrary parent vector.

At cutoff, the Schur complement of \(z-\omega I_q\) is
\[
 (z-E_0)I_n-\lambda^2
       [\Sigma_{0,\Lambda}(z)-\Sigma_{0,\Lambda}(z_0)]K .
\]
Taking the blockwise limit proves the displayed formula.
The same fixed decomposition proves the reducing-subspace
statements.  Finally \(\ker K=\ker C\), since
\(\langle u,Ku\rangle=\|Cu\|^2\).
\end{proof}

\subsection{Matrix spectral measure and weak-coupling limit}

Set \(b=s_0(E_0)\) and \(J_*=J_0(E_0)>0\).
The parent spectral measure is the positive matrix measure
\(\mathsf M_\lambda(B)=P\,1_B(H_\lambda)\iota\), of total mass
\(I_n\).  Let \(Q_0\) denote the kernel projection, including
the convention \(Q_0=0\) when the kernel is trivial.

\begin{theorem}[A finite-matrix survival limit]
\label{thm:v20v52-limit}
For sufficiently small \(\lambda>0\),
\[
 \mathsf M_\lambda(dx)=Q_0\delta_{E_0}(dx)
            +\sum_{\kappa>0}Q_\kappa\rho_{\lambda,\kappa}(x)\,dx,
\]
where each scalar density has mass one and
\[
 \rho_{\lambda,\kappa}(x)=
 \frac{\lambda^2\kappa J_0(x)}
 {[x-E_0-\lambda^2\kappa s_0(x)]^2+
                       \pi^2\lambda^4\kappa^2J_0(x)^2},
 \qquad x>a.
\]
For the compressed, rotating parent propagator
\[
 T_\lambda(\tau)=e^{iE_0\tau/\lambda^2}
               P e^{-iH_\lambda\tau/\lambda^2}\iota
\]
one has
\[
 \sup_{\tau\in\mathbb R}
 \big\|T_\lambda(\tau)
       -\exp[-ibK\tau-\pi J_*K|\tau|]\big\|\longrightarrow0 .
\]
\end{theorem}
\begin{proof}
Each positive \(\kappa\) block obeys V48 with self-energy
\(\lambda^2\kappa S_0\).  Its bound-state criterion is excluded
for sufficiently small \(\lambda\), since \(E_0-a\) is fixed.
There are finitely many eigenvalues, so the same small-coupling
range works for all of them.  Their parent measures are the
displayed normalized densities.  The reducing kernel contributes
the exact atom, which can lie inside the continuum spectrum.

Under \(x=E_0+\lambda^2u\), the positive-block densities become
\[
 r_{\lambda,\kappa}(u)=
 \frac{\kappa J_0(E_0+\lambda^2u)}
 {[u-\kappa s_0(E_0+\lambda^2u)]^2+
               \pi^2\kappa^2J_0(E_0+\lambda^2u)^2},
\]
extended by zero below the rescaled threshold.  They converge
pointwise to the mass-one Cauchy densities
\[
 r_{*,\kappa}(u)=
 \frac{\kappa J_*}{(u-\kappa b)^2+\pi^2\kappa^2J_*^2}.
\]
The minimum-density argument of V51 gives
\(\|r_{\lambda,\kappa}-r_{*,\kappa}\|_1\to0\).
In particular the integral of the operator norm of the difference
of the absolutely continuous matrix densities is at most the
finite sum of these scalar errors.  The atoms \(Q_0\delta_0\)
agree exactly and must not be represented as Lebesgue densities.
Fourier transformation, with the same uniform bound as V51,
gives convergence to
\[
 Q_0+\sum_{\kappa>0}
      Q_\kappa e^{-i\kappa b\tau-\pi\kappa J_*|\tau|}.
\]
Functional calculus identifies this with the stated exponential.
\end{proof}

\begin{corollary}[Parent retention and the necessary basis choice]
\label{cor:v20v52-retention}
For a unit parent vector \(u\) and \(\tau\geq0\), the probability
of remaining somewhere in the parent subspace converges, uniformly
in absolute error, to
\[
 \|e^{-(ib+\pi J_*)K\tau}u\|^2
       =\sum_{\kappa\geq0}e^{-2\pi J_*\kappa\tau}\|Q_\kappa u\|^2 .
\]
Its limiting long-time value is \(\|Q_0u\|^2\).
In a fixed orthonormal parent basis, the limiting amplitude
evolves by independent scalar factors for all initial vectors
if and only if \(K\) is diagonal in that basis.
\end{corollary}
\begin{proof}
Both the compressed unitary and the limiting exponential are
contractions.  Their squared output norms differ by at most
twice their operator-norm difference on unit vectors.
Orthogonality of the projections proves the sum and its limit.
If \(K\) is diagonal the final assertion is immediate.
Conversely a diagonal amplitude propagator for all forward times
has a diagonal right derivative at zero.  That derivative is
\(-(ib+\pi J_*)K\), whose scalar prefactor is nonzero because
\(J_*>0\).  Hence \(K\) must be diagonal.
\end{proof}

This probability differs from return to the same initial vector,
which is \(|\langle u,T_\lambda(\tau)u\rangle|^2\).
No interchange of the weak-coupling and late-time limits is needed
for the corollary.  Separately, at every sufficiently small fixed
\(\lambda\), the integrable positive-block densities and the
Riemann--Lebesgue lemma give
\[
 P e^{-itH_\lambda}\iota=e^{-iE_0t}Q_0+o(1)
\]
in the finite-dimensional operator norm.  Thus the exact
fixed-coupling parent-retention limit also equals \(\|Q_0u\|^2\).

\subsection{An explicit interference obstruction}

Take two parents and one continuum channel with
\[
 C=\begin{pmatrix}1&1\end{pmatrix},\quad
 K=\begin{pmatrix}1&1\\1&1\end{pmatrix},\quad
 u_+=2^{-1/2}(1,1),\quad u_-=2^{-1/2}(1,-1).
\]
Direct multiplication gives \(Ku_+=2u_+\) and \(Cu_-=0\).
The symmetric combination has probability decay rate
\(4\pi J_*\) in rescaled time, while the antisymmetric
combination is an exact eigenvector at energy \(E_0\) for every
cutoff and in the renormalized limit.  For initial \(e_1=(1,0)\)
the limiting parent amplitudes are
\[
 e^{-(ib+\pi J_*)K\tau}e_1
 =\frac12
 \begin{pmatrix}
  1+e^{-2(ib+\pi J_*)\tau}\\
 -1+e^{-2(ib+\pi J_*)\tau}
 \end{pmatrix}.
\]
Parent retention is \((1+e^{-4\pi J_*\tau})/2\);
return to \(e_1\) is
\(\tfrac14|1+e^{-2(ib+\pi J_*)\tau}|^2\).
Their late-time limits are respectively \(1/2\) and \(1/4\).
Using only the two diagonal entries \(K_{11}=K_{22}=1\)
would predict complete loss from the parent space and is false.
The example is an algebraic interference mechanism, not a
prediction of a stable particle in the Standard Model.

\begin{assessment}[What has and has not advanced]
\label{ass:v20v52}
We have a cutoff-independent self-adjoint many-parent block,
its matrix self-energy, and a uniform weak-coupling parent
propagator limit under a stated common-profile hypothesis.
The decay matrix \(2\pi J_*C^*C\) retains coherences that
independent named-parent rates erase.  The exact kernel identifies
which combinations are nondecaying in this model.  None of these
statements closes V44's finite-density CAR hierarchy or constructs
gauge, chiral, stress-tensor, or Einstein continuum limits.
\end{assessment}

\begin{decision}[Next upstream test]
\label{dec:v20v52-next}
Exact degeneracy and a common profile permitted a fixed
singular-value decomposition.  Next admit an order-\(\lambda^2\)
Hermitian parent splitting \(B\) that need not commute with \(K\).
Determine whether the limit has generator
\(-i(B+bK)-\pi J_*K\), and whether \(\ker K\) stays nondecaying
only through its invariant part under \(B\).  This directly tests
the stability of the present obstruction instead of counting
uncoupled scalar laws as a many-body closure.
\end{decision}

\begin{record}[V52 append and verification record]
\label{rec:v20v52}
The predecessor V51 has SHA--256
\texttt{4699CE250C28AF38F1ECC04EFF956CEA38231A08E825607ACAF417F42D221B77}.
Its complete prefix is retained.  The construction is proved by
an exact finite orthogonal decomposition and V47's already proved
scalar norm-resolvent result; the limit uses normalized spectral
measures, not sampled fits.  The two-parent example is checked
directly by matrix multiplication in the text.  Source validation
checks balanced environments and delimiters, unique labels,
resolved references, and predecessor preservation; this is not
LaTeX compilation.  Only the immediate active SM predecessor is
replaced.  The next companion audit remains V55.
\end{record}


\section{V53: noncommuting parent splittings and the invariant dark subspace}
\label{sec:v20v53}

The kernel found in V52 is protected by exact degeneracy.
An order-\(\lambda^2\) parent energy splitting acts on the same
time scale as the decay.  We now allow such a splitting without
assuming that it commutes with the channel Gram matrix.
The proof uses an exact projected Duhamel equation; it does not
assume that the parent dynamics is Markovian before taking the limit.

\subsection{A bounded perturbation after ultraviolet removal}

Keep all the hypotheses and notation of V52, and let
\(B=B^*\) be a fixed \(n\)-by-\(n\) matrix.  Denote the operator
of V52 by \(H_{\lambda,0}\), and define
\[
 H_{\lambda,B}=H_{\lambda,0}+\lambda^2\iota BP,\qquad
 D(H_{\lambda,B})=D(H_{\lambda,0}).
\]
At finite cutoff add the same bounded term to
\(H_{\lambda,\Lambda}\).  In particular, the renormalized parent
matrix is \(E_0I_n+\lambda^2B\); the subtraction counterterm
remains \(-\lambda^2\Sigma_{0,\Lambda}(z_0)K\).

\begin{proposition}[The split model and its resolvent]
\label{prop:v20v53-resolvent}
The operator \(H_{\lambda,B}\) is self-adjoint.  It is the
norm-resolvent limit of the split cutoff operators, and
\[
 P(z-H_{\lambda,B})^{-1}\iota
 =\big[(z-E_0)I_n-\lambda^2B-\lambda^2S_0(z)K\big]^{-1}
 \quad (\operatorname{Im}z\ne0).
\]
\end{proposition}
\begin{proof}
The added term is bounded and self-adjoint, so bounded
self-adjoint perturbation preserves self-adjointness and the
domain.  For completeness, choose
\(|\operatorname{Im}z|>\lambda^2\|B\|\).
The resolvent of a perturbed cutoff operator is given by its
unperturbed resolvent times the convergent Neumann inverse for
the bounded perturbation.  V52's norm-resolvent convergence and
continuity of that inverse prove convergence at this \(z\).
The resolvent identity then gives norm-resolvent convergence
at every nonreal point.  The finite-cutoff Schur complement
adds exactly \(-\lambda^2B\) to the inverse compressed
resolvent.  Passing to the limit proves the formula.
Its inverse exists since, for any nonzero \(u\),
\[
 \operatorname{Im}\langle u,
 [(z-E_0)I_n-\lambda^2B-\lambda^2S_0(z)K]u\rangle
 =\operatorname{Im}z\left(\|u\|^2+
 \lambda^2\langle u,Ku\rangle
       \int_a^\infty\frac{J_0(w)}{|z-w|^2}\,dw\right),
\]
which has the strict sign of \(\operatorname{Im}z\).
\end{proof}

\subsection{An exact projected Duhamel equation}

For forward rescaled time define
\[
 T_{\lambda,B}(\tau)=e^{iE_0\tau/\lambda^2}
           P e^{-iH_{\lambda,B}\tau/\lambda^2}\iota,\qquad \tau\geq0.
\]
Let \(\gamma=\pi J_*>0\), \(b=s_0(E_0)\), and
\[
 A=-i(B+bK)-\gamma K,\qquad A_0=-(ib+\gamma)K.
\]

\begin{theorem}[Noncommuting weak-coupling limit]
\label{thm:v20v53-limit}
For every finite \(T\),
\[
 \sup_{0\leq\tau\leq T}\|T_{\lambda,B}(\tau)-e^{A\tau}\|
       \longrightarrow0\quad\hbox{as }\lambda\to0.
\]
More explicitly, with
\(\epsilon_\lambda=\sup_{\tau\geq0}
 \|T_{\lambda,0}(\tau)-e^{A_0\tau}\|\to0\) from V52,
the left side is bounded by
\[
 \epsilon_\lambda(1+\|B\|T)e^{\|B\|T}.
\]
No commutation of \(B\) and \(K\) is required.
\end{theorem}
\begin{proof}
The exact bounded-perturbation Duhamel identity, compressed
between \(P\) and \(\iota\), is
\[
 T_{\lambda,B}(\tau)=T_{\lambda,0}(\tau)
 -i\int_0^\tau T_{\lambda,0}(\tau-s)B
                          T_{\lambda,B}(s)\,ds .
\]
Indeed in physical time the perturbation is
\(\lambda^2\iota BP\), so inserting it in Duhamel factors
the integrand into the two compressed propagators.
Changing variables \(s=\lambda^2t\) cancels the prefactor,
and the scalar rotating phases multiply correctly.
This is an exact convolution identity even though neither
compressed propagator is generally a semigroup.

Finite-matrix variation of constants gives
\[
 e^{A\tau}=e^{A_0\tau}
 -i\int_0^\tau e^{A_0(\tau-s)}B e^{As}\,ds .
\]
Both \(e^{A_0t}\) and \(e^{At}\) are contractions for \(t\geq0\),
because their Hermitian generator parts equal \(-\gamma K\).
The compressed unitaries are also contractions.  Subtract the
two convolution identities and write the difference of integrands as
\[
 [T_{\lambda,0}(\tau-s)-e^{A_0(\tau-s)}]B T_{\lambda,B}(s)
 + e^{A_0(\tau-s)}B[T_{\lambda,B}(s)-e^{As}].
\]
The norm error at \(\tau\leq T\) is consequently at most
\[
 \epsilon_\lambda(1+\|B\|T)
       +\|B\|\int_0^\tau
                    \|T_{\lambda,B}(s)-e^{As}\|\,ds .
\]
The elementary integral Gronwall inequality gives the stated
bound and convergence.
\end{proof}

Unlike V52, this argument asserts uniform convergence only on
bounded rescaled time intervals.  Its bound grows with \(T\)
and does not justify exchanging \(\lambda\to0\) with
\(\tau\to\infty\).  For negative times the adjoint identity
determines the corresponding limit from the forward one.

\subsection{Which nondecaying combinations survive the splitting?}

Define
\[
 \mathcal W=\bigcap_{j=0}^{n-1}\ker(KB^j).
\]
The definition uses \(K\), not a choice of singular vectors.
It remains valid when \(K=0\) or when eigenvalues are repeated.

\begin{theorem}[Invariant kernel and decay on its complement]
\label{thm:v20v53-dark}
The space \(\mathcal W\) is the largest \(B\)-invariant subspace
contained in \(\ker K\).  It reduces \(B,K,A\).
On \(\mathcal W\), \(e^{A\tau}=e^{-iB\tau}\).
On \(\mathcal W^\perp\), every eigenvalue of \(A\) has strictly
negative real part.  Thus there are constants \(c,\eta>0\)
such that
\[
 \|e^{A\tau}|_{\mathcal W^\perp}\|\leq c e^{-\eta\tau}
 \quad(\tau\geq0)
\]
when this complement is nonzero.  In particular
\[
 \lim_{\tau\to\infty}\|e^{A\tau}u\|^2
                      =\|P_{\mathcal W}u\|^2.
\]
The embedded parent space \(\iota\mathcal W\) is also exactly
reducing for every \(H_{\lambda,B}\), where its Hamiltonian is
\(E_0I+\lambda^2 B|_{\mathcal W}\).
\end{theorem}
\begin{proof}
Cayley--Hamilton expresses \(B^n\) as a linear combination
of its lower powers.  Therefore \(u\in\mathcal W\) implies
\(Bu\in\mathcal W\), including the condition involving
\(KB^n u\).  Conversely any \(B\)-invariant subspace of
\(\ker K\) satisfies all the defining kernel conditions.
Since \(B\) is Hermitian, invariance of \(\mathcal W\)
implies invariance of its orthogonal complement.
Since \(K\) is Hermitian and vanishes on \(\mathcal W\),
it also preserves that complement.  These facts prove the
reducing statements and the expression on \(\mathcal W\).

If \(Au=\zeta u\), \(u\ne0\), then
\[
 \operatorname{Re}\zeta\,\|u\|^2
           =-\gamma\langle u,Ku\rangle\leq0 .
\]
Equality implies \(Ku=0\), by positivity, and hence
\(-iBu=\zeta u\).  Such a \(u\) is a \(B\)-eigenvector
in \(\ker K\), so it belongs to \(\mathcal W\).
There can be no such eigenvector in \(\mathcal W^\perp\).
All eigenvalues of the restriction therefore lie strictly
in the left half-plane.  In its finite Jordan normal form
each exponential is a decaying scalar exponential times a
polynomial.  Choosing \(\eta\) smaller than the minimum
strict spectral decay absorbs all these polynomials into
a finite constant \(c\).  Orthogonality of the reducing
decomposition proves the norm limit.

Finally \(C\mathcal W=0\), the counterterm \(K\) vanishes
on \(\mathcal W\), and \(B\) preserves it.  Already at
cutoff this is an uncoupled finite parent block.
Its orthogonal-sum decomposition persists under the
norm-resolvent construction, proving the last statement.
\end{proof}

This theorem characterizes the nondecaying part of the
limiting parent semigroup and an exact protected parent block.
It does not assert that every spectral question about the
complement of the full interacting Hamiltonian is settled.
The semigroup result alone supplies no uniform late-time
estimate for the exact split model.

\subsection{A splitting destroys the previous protected combination}

Use the bright/dark basis of the two-parent example in V52.
In that basis choose
\[
 K=\begin{pmatrix}2&0\\0&0\end{pmatrix},\qquad
 B=\begin{pmatrix}0&\delta\\\delta&0\end{pmatrix},
 \qquad \delta\in\mathbb R\setminus\{0\}.
\]
The old dark vector is \(u_d=(0,1)\).  But
\(KBu_d=(2\delta,0)\ne0\), so \(\mathcal W=\{0\}\).
The limiting generator is
\[
 A=\begin{pmatrix}
 -2\gamma-2ib&-i\delta\\-i\delta&0
 \end{pmatrix},\qquad
 \det(\zeta I-A)=\zeta^2+(2\gamma+2ib)\zeta+\delta^2.
\]
The theorem proves strict decay of both modes without choosing
a complex square-root branch.  The initially dark vector
illustrates why its zero instantaneous decay rate is insufficient.
If \(x(\tau)=e^{A\tau}u_d\), then
\[
 x_1(\tau)=-i\delta\tau+O(\tau^2),\qquad
 \frac{d}{d\tau}\|x(\tau)\|^2=-4\gamma|x_1(\tau)|^2.
\]
Integration gives the explicit short-time law
\[
 \|x(\tau)\|^2
       =1-\frac{4\gamma\delta^2}{3}\tau^3+O(\tau^4).
\]
This is a Taylor expansion of the limiting finite-dimensional
semigroup.  It is not a uniform short-physical-time expansion of
the ultraviolet-renormalized Hamiltonian.  It proves that checking
only \(Ku=0\) at the initial instant misses later leakage.

\begin{assessment}[Acquisition and remaining dependency]
\label{ass:v20v53}
Noncommuting order-\(\lambda^2\) parent splittings now have a
proved weak-coupling limit in the specified projected model.
The surviving protected space is the invariant part of the
coupling kernel, rather than the whole kernel.  This repairs
a potential overextension of V52.  Neither result controls
occupied fermion modes, Pauli blocking in a continuum limit,
or the connected correlations in V44.  Those are requirements
for matter kinetics and a conserved Einstein source.
\end{assessment}

\begin{decision}[Return to occupation correlations]
\label{dec:v20v53-next}
Before adding more one-excitation variants, return to V44's
many-occupation bottleneck.  Examine whether a factorized
initial state remains factorized under the exact resonant CAR
dynamics, derive the first generated connected correlations,
and test whether the available small-coupling estimates control
them on inverse-square time scales.  A negative test must identify
the missing continuum or mixing hypothesis explicitly; it must
not be relabeled a kinetic closure.
\end{decision}

\begin{record}[V53 append and verification record]
\label{rec:v20v53}
The predecessor V52 has SHA--256
\texttt{CE09DCA64A9C8D21186492FE81CE1D0C20D56C0ECCA4F4465D274E1456DD2069}.
Its complete prefix is retained.  All new estimates and the
two-dimensional example are proved in the text; no numerical
sampling is used to infer a theorem.  Source validation checks
predecessor preservation, environments, delimiters, unique labels,
and references.  It is not LaTeX compilation.
Only the immediate active SM predecessor is replaced.
The next companion audit remains V55.
\end{record}


\section{V54: correlations generated from product data and a kinetic-scale obstruction}
\label{sec:v20v54}

V44 rules out an exact occupation-only equation for arbitrary
correlated data.  A remaining possible argument is that independent
initial occupations and small coupling might suffice for the
factorized closure.  This section tests that argument directly.
We compute the first generated covariances for arbitrary
number-diagonal product data with finite oscillator support,
then give an exact resonant counterexample on inverse-square times.
The two-level dynamics itself was already available in V26;
the new result is the covariance and closure-error calculation,
including its nonvanishing time average on the proposed kinetic scale.

\subsection{General product data}

Use the single-channel operators and Hamiltonian of V44.
To avoid confusion with V53's splitting matrix, write
\[
 \mathcal B=N-NF-ND-FD
\]
for the blocking operator.  Let the initial state be the product of
a number-diagonal oscillator distribution with finite support and
independent diagonal fermion states.  Denote its data by
\[
 \mu=\langle N\rangle_0,\quad
 v=\langle(N-\mu)^2\rangle_0,\quad
 f=\langle F\rangle_0,\quad d=\langle D\rangle_0 .
\]
All three initial pair covariances vanish.  Put
\[
 U=(1-f)(1-d),\qquad V=fd,
\]
\[
 Q=\mu U-(\mu+1)V,\qquad
 L=(\mu^2+v-\mu)U-(\mu^2+v+\mu)V .
\]
These symbols refer only to initial data.

\begin{theorem}[Initial generation of connected occupations]
\label{thm:v20v54-generation}
For arbitrary fixed detuning and fixed \(g\), as \(t\to0\),
\[
 \begin{aligned}
 \operatorname{Cov}_t(N,F)
   &=|g|^2t^2[L-(\mu-f)Q]+O(t^4),\\
 \operatorname{Cov}_t(N,D)
   &=|g|^2t^2[L-(\mu-d)Q]+O(t^4),\\
 \operatorname{Cov}_t(F,D)
   &=|g|^2t^2(1-f-d)Q+O(t^4).
 \end{aligned}
\]
Consequently the error of V44's instantaneous occupation
factorization satisfies
\[
 \langle\mathcal B\rangle_t
 -\{n(t)[1-f(t)-d(t)]-f(t)d(t)\}
 =|g|^2t^2[(2\mu-1)Q-2L]+O(t^4).
\]
The remainders are local in time at fixed parameters; no
kinetic-time uniformity is claimed.
\end{theorem}
\begin{proof}
The conserved integers \(N+F,N+D\) decompose the supported
evolution into finitely many blocks.  A convertible block joins
\(|k,0,0\rangle\) and \(|k-1,1,1\rangle\), \(k\geq1\).
Singly occupied pairs are stationary for every number observable.
Starting from a diagonal state, the transition probability in a
convertible block is
\[
 \frac{k|g|^2}{k|g|^2+\Delta^2/4}
       \sin^2\!\left(t\sqrt{k|g|^2+\Delta^2/4}\right)
       =k|g|^2t^2+O(t^4).
\]
If the denominator vanishes, the probability is identically zero,
with the same limiting expansion.  Thus all number expectations
are even analytic functions of time on these finitely many blocks.

For a number function \(h(N,F,D)\), the coefficient of
\(|g|^2t^2\) in its expectation is the initial product average of
\[
 N(1-F)(1-D)[h(N-1,1,1)-h(N,0,0)]
 +(N+1)FD[h(N+1,0,0)-h(N,1,1)].
\]
The first expression is zero when \(N=0\); its apparent value
at \(N-1=-1\) is immaterial and can be assigned zero.
Applying this formula to \(N,F,D,FD,NF,ND\) yields respectively
the coefficients
\[
 -Q,\quad Q,\quad Q,\quad Q,\quad L,\quad L.
\]
For example the forward change of \(NF\) from an empty pair
is \(N-1\), and its reverse change from an occupied pair is
\(-N\).  Its coefficient is therefore
\(\langle N(N-1)\rangle_0U-\langle N(N+1)\rangle_0V=L\).
Subtracting the products of the evolving means proves the
three covariance formulas.  Their sum has coefficient
\(2L+(1-2\mu)Q\).  The exact covariance identity of V44,
with a minus sign on that sum, proves the last assertion.
\end{proof}

For empty fermion modes and a fixed oscillator number \(r\geq1\),
one has \(Q=r\) and \(L=r(r-1)\).  The three covariance
coefficients are respectively \(-r,-r,r\), and the blocking
error coefficient is \(r\).  Initially exact factorization
therefore fails immediately at second order.

\subsection{An exact product-state counterexample on long times}

Now impose exact resonance and take initial state
\(|r,0,0\rangle\), a product of three number states.
Let \(g_\lambda=\lambda g_0\) with \(g_0\ne0\).
In the two-dimensional invariant block, define
\[
 p_\lambda(t)=\sin^2(\lambda|g_0|\sqrt r\,t).
\]
This is the probability of \(|r-1,1,1\rangle\).
All the following formulas follow by averaging the two
number configurations, independently of their relative phase:
\[
 n=r-p_\lambda,\qquad f=d=p_\lambda,\qquad
 \langle NF\rangle=\langle ND\rangle=(r-1)p_\lambda,\qquad
 \langle FD\rangle=p_\lambda .
\]
In particular
\[
 \operatorname{Cov}(N,F)=\operatorname{Cov}(N,D)
      =-p_\lambda(1-p_\lambda),\qquad
 \operatorname{Cov}(F,D)=p_\lambda(1-p_\lambda).
\]
The exact blocking and its factorized approximation are
\[
 \langle\mathcal B\rangle=r-2rp_\lambda,\qquad
 \mathcal B_{\rm fac}
     =r-(2r+1)p_\lambda+p_\lambda^2,
\]
so their difference is exactly \(p_\lambda(1-p_\lambda)\).

\begin{theorem}[Small coupling alone does not propagate occupation factorization]
\label{thm:v20v54-obstruction}
For every fixed \(\tau>0\), the blocking error at
\(t=\tau/\lambda^2\) has no limit as \(\lambda\to0\).
It has subsequences with values zero and \(1/4\).
For every \(T>0\),
\[
 \lim_{\lambda\to0}\int_0^T
 \left|\langle\mathcal B\rangle_{\tau/\lambda^2}
             -\mathcal B_{\rm fac}(\tau/\lambda^2)\right|\,d\tau
 =\frac{T}{8}.
\]
Thus neither pointwise vanishing nor \(L^1\) vanishing of the
factorization error on rescaled time follows from product
initial data and weak coupling alone.
\end{theorem}
\begin{proof}
Write \(\alpha=|g_0|\sqrt r>0\).  At rescaled time
\(p_\lambda=\sin^2(\alpha\tau/\lambda)\).
The sequences
\(\lambda_k=\alpha\tau/(k\pi)\) and
\(\widetilde\lambda_k=\alpha\tau/(k\pi+\pi/4)\)
tend to zero and give respectively \(p_\lambda=0\) and
\(p_\lambda=1/2\).  The asserted error values follow.
Moreover the error is nonnegative and
\[
 p_\lambda(1-p_\lambda)
       =\frac{1-\cos(4\alpha\tau/\lambda)}8.
\]
Its integral is exactly
\[
 \frac{T}{8}
 -\frac{\lambda}{32\alpha}\sin(4\alpha T/\lambda),
\]
which proves the limit.
\end{proof}

This example lies inside a fixed finite resonant block.
It does not disprove a kinetic theorem that requires an
infinite-volume limit, a continuous detuning spectrum, or
a specified many-mode scaling.  It disproves the implication
from the two stated assumptions alone.  Nor does it contradict
V51: that theorem removes the ultraviolet cutoff in a continuum
one-parent model before taking a weak-coupling limit, whereas
this block retains a single coherent resonant pair.

\subsection{What has to be estimated in the many-channel problem}

Restoring mode labels is essential before using any of these
quantities as a kinetic or stress-tensor source.  For the
disjoint-pair model of V44, write
\[
 z_i=g_i\langle X_i\rangle,\quad
 E_i=\langle\mathcal B_i\rangle-
      \{n(1-f_i-d_i)-f_id_i\},\quad
 Y_{ij}=\langle P_i^\dagger P_j\rangle .
\]
The exact identity already proved there becomes
\[
 \dot z_i=i\Delta_i z_i
 +i|g_i|^2\{n(1-f_i-d_i)-f_id_i\}
 +i|g_i|^2E_i-i\sum_{j\ne i}g_i\overline g_jY_{ij}.
\]
Thus the omitted forcing of the coherence equation is
\[
 \mathcal R_i=|g_i|^2E_i-
                \sum_{j\ne i}g_i\overline g_jY_{ij}.
\]
A proof of a proposed closure must control the effect of this
forcing after the detuning propagation and the further time
integration giving \(f_i\), in the topology and weighted
mode sums appropriate to that closure.  Pointwise smallness
of an individual mode covariance is neither established
here nor by itself sufficient after summing a growing
number of channels.  Conversely, demanding each covariance
vanish is stronger than necessary if a proved weighted
oscillatory estimate makes the omitted forcing negligible.
No such estimate may be inferred from the one-parent
survival amplitude alone.

\begin{assessment}[Upstream conclusion]
\label{ass:v20v54}
The general product-data curvature and the exact long-time
counterexample identify a missing hypothesis in an occupation
closure argument.  Simply initializing independent occupations
does not supply propagation of independence.  The useful next
quantity is the weighted residual \(\mathcal R_i\), including
cross-pair coherence, rather than another scalar survival rate.
Obtaining a continuum bound for that residual remains necessary
before a derived matter kinetic equation can source Einstein
evolution.  The full Standard Model--Einstein construction
is not achieved.
\end{assessment}

\begin{decision}[Audit, then a weighted continuum test]
\label{dec:v20v54-next}
At V55 perform the scheduled YM and NS companion audit.
Then test the weighted covariance and cross-pair sums first
in a finite-box one-parent state with all pair modes retained,
where their exact values can be derived from the state
amplitudes.  Track the canonical volume factors from V43
before taking any large-volume limit.  This distinguishes
small individual occupations from small summed errors and
can expose which estimate a finite-density extension needs.
\end{decision}

\begin{record}[V54 append record]
\label{rec:v20v54}
The predecessor V53 has SHA--256
\texttt{2949560AF898724D3CDFA129FDDD30B30731C3D16D59B7B846C88948FAD5DF92}.
Its complete prefix is preserved.  The new results use exact
finite-block probabilities and elementary expectation algebra;
no finite sample is promoted to a general proof.  Structural
validation checks the append-only source, environments,
delimiters, unique labels, and references, not LaTeX compilation.
Only the immediate active SM predecessor is replaced.
The next version V55 is the required companion audit.
\end{record}


\section{V55: companion audit before the weighted-correlation test}
\label{sec:v20v55}

This is the scheduled fifth-version audit.  It records the
companion snapshots actually read and their relevance to the
current occupation-closure bottleneck.  Instructions or programme
decisions inside those documents are treated as research content.

\begin{record}[Companion snapshots]
\label{rec:v20v55-snapshots}
The active main Yang--Mills note is V52, 688444 bytes, SHA--256
\[
\texttt{37A55909854CE5EBE88E1BB0DF0E41C5BD2D4CBCF146926F1B1FF2DB17E8330B}.
\]
The active NS note remains V26, 278283 bytes, SHA--256
\[
\texttt{82C887F8C2C69E4F28FAD7C3DC8DE5B9099CB5A4BF431EBA1FFF3E91935978AD}.
\]
The parallel YM note remains V5, 54174 bytes, SHA--256
\[
\texttt{306F1BB84CFA8A8D47EA569EC17E9545F9867C8D32B548EC9D9C444B4F3C0512}.
\]
These are read-time snapshots; parallel work may advance them later.
\end{record}

\begin{assessment}[Main YM: the full differentiated identity matters]
\label{ass:v20v55-ym}
Since the V50 audit of YM V47, the companion has corrected a
coarse-representative mismatch and derived full-link right-inverse
composition.  Its later sections distinguish a frozen vertex
contraction from the complete response of a moving representative.
V51 differentiates the gauge identity, retaining the action-gradient
term off shell.  V52 identifies the actual split-link convention
\(U_e=B_e\exp(a_e)\), derives
\[
 (R_B\omega)_e=\operatorname{Ad}_{B_e^{-1}}\omega_s-\omega_v,
\]
and distinguishes flat commuting stationary backgrounds from
nonconstant transverse profiles.  The colour-resolved full-Hessian
test, complete response cancellation, and interacting continuum
are still open in that snapshot.

No gauge identity from that lattice calculation is imported as
a theorem about the present Yukawa covariance sums.  The useful
methodological lesson is concrete: a small or vanishing isolated
term does not establish the full identity after omitted terms
are restored.  In the SM calculation the corresponding task is
to retain both the number-covariance and cross-pair-coherence
parts of V54's exact residual.
\end{assessment}

\begin{assessment}[NS and parallel YM]
\label{ass:v20v55-other}
NS V26 retains its rank-one Oseen derivative-norm analysis and
does not prove global regularity.  It supplies no dispersive
quantum pair-correlation estimate.  Parallel YM V5 retains its
obstruction to a volume-uniform global extensive-charge tail
under tensorization.  Its fixed-window alternative reinforces
the need to specify the volume scaling of a target quantity,
but it does not bound the SM residual either.  No new analytic
estimate is transferred from either unchanged source.
\end{assessment}

\begin{decision}[Proceed immediately to V56]
\label{dec:v20v55-next}
Compute exact one-parent, many-pair occupation covariances and
pair coherences at finite box size.  Bound their weighted sums
separately, using the canonical \(L^{-1/2}\) vertex scaling.
If a coherent sum persists while individual mode errors vanish,
retain it as an explicit memory contribution; do not infer
independent collision rates.  The next companion audit and
ten-version review are V60.
\end{decision}

\begin{record}[V55 append record]
\label{rec:v20v55}
The predecessor V54 has SHA--256
\texttt{3B76B1313C9DE549EA16832ACAED03EF4C124753CAE6B3D2C765FB3C50BA40E5}.
Its complete prefix is retained.  Only the immediate active SM
predecessor is replaced; companions and frozen archives are
preserved.  This audit does not close any continuum gate.
\end{record}


\section{V56: small mode covariances do not control the coherent channel sum}
\label{sec:v20v56}

Following the V55 audit, we test the full residual identified
in V54 at finite box size.  The state space remains the
one-parent sector, but all the allowed disjoint pair modes
are retained.  This is an exactly solvable diagnostic of the
weighted sums that a finite-density kinetic argument must
eventually control, not that finite-density theorem itself.

\subsection{Exact occupation algebra in the full one-parent block}

For \(r\) disjoint pairs use V44's operators \(P_i^\dagger,X_i\).
Let \(|H\rangle\) contain one parent and no pairs, and
\(|i\rangle\) contain no parent and precisely pair \(i\).
Consider any normalized vector in this invariant block,
\[
 |\psi\rangle=A|H\rangle+\sum_{i=1}^r b_i|i\rangle,\qquad
 P=|A|^2,\quad p_i=|b_i|^2,\quad P+\sum_i p_i=1.
\]
Here the scalar \(P\) is a probability, not a projection or
the pair annihilator \(P_i\).  The couplings are arbitrary
complex numbers \(g_i\).

\begin{theorem}[Exact covariances and residual]
\label{thm:v20v56-residual}
In this state,
\[
 n=P,\quad f_i=d_i=p_i,\quad
 \langle NF_i\rangle=\langle ND_i\rangle=0,\quad
 \langle F_iD_i\rangle=p_i,
\]
\[
 E_i=\langle\mathcal B_i\rangle-
       \{n(1-f_i-d_i)-f_id_i\}
       =(2P-1)p_i+p_i^2,\qquad
 Y_{ij}=\langle P_i^\dagger P_j\rangle=\overline b_i b_j .
\]
Define
\[
 d_g=\sum_i|g_i|^2p_i,\qquad s_g=\sum_i\overline g_i b_i,
 \qquad
 \mathcal R_i=|g_i|^2E_i-\sum_{j\ne i}g_i\overline g_jY_{ij}.
\]
Then the summed omitted forcing in V54 obeys the exact identity
\[
 \sum_i\mathcal R_i
   =2P d_g+\sum_i|g_i|^2p_i^2-|s_g|^2.
\]
In particular, with \(g_{\max}^2=\max_i|g_i|^2\),
\[
 \sum_i |g_i|^2|E_i|\leq g_{\max}^2(1-P),\qquad
 0\leq 2P d_g+\sum_i|g_i|^2p_i^2
                         \leq2g_{\max}^2(1-P).
\]
\end{theorem}
\begin{proof}
Number observables are diagonal on the displayed basis.
Parent and pair occupation are mutually exclusive, while the
two fermions in pair \(i\) occur together.  This gives the
first line and \(\langle\mathcal B_i\rangle=P-p_i\).
Subtracting \(P(1-2p_i)-p_i^2\) gives \(E_i\).
The even pair operator \(P_i^\dagger P_j\) maps \(|j\rangle\)
to \(|i\rangle\); the CAR convention of V44 therefore gives
the stated \(Y_{ij}\), including \(i=j\).
Consequently
\[
 \sum_{i\ne j}g_i\overline g_jY_{ij}
                =|s_g|^2-d_g.
\]
Combining this with the formula for \(E_i\) proves the residual.

Since \(0\leq P\leq1-p_i\),
\(|2P-1+p_i|\leq1-p_i\).  Hence
\(|E_i|\leq p_i(1-p_i)\leq p_i\), proving the first bound.
For the second, all terms are nonnegative and
\[
 2P d_g+\sum_i|g_i|^2p_i^2
 \leq g_{\max}^2[2P(1-P)+(1-P)^2]
 =g_{\max}^2(1-P)(1+P).
\]
The result follows from \(P\leq1\).
\end{proof}

Thus if \(g_{\max}\to0\), the total number-factorization
contribution vanishes uniformly over normalized one-parent
states.  The cross-pair sum has no comparable conclusion:
its leading part is the squared coherent overlap \(|s_g|^2\).
The theorem also gives the precise uniform estimate
\[
 \left|\sum_i\mathcal R_i+|s_g|^2\right|
                   \leq2g_{\max}^2(1-P).
\]
It is the summed residual, not a claim about each
\(\mathcal R_i\) separately.

\subsection{A dynamically attained counterexample to a small-entry argument}

Assume \(G^2=\sum_i|g_i|^2>0\).  For a fixed \(P\in[0,1]\),
the normalized pair amplitudes
\[
 b_i=\sqrt{1-P}\,\frac{g_i}{G}
\]
give \(|s_g|^2=(1-P)G^2\).
Hence this overlap need not vanish when each \(g_i\) does.
This is not merely a choice of algebraically allowed states:
it occurs under a resonant Hamiltonian with equal pair energies.

\begin{proposition}[Uniformly small entries with a finite collective term]
\label{prop:v20v56-collective}
Let all pair energies equal the parent energy, let
\(g_i=G/\sqrt r\) with fixed \(G>0\), and start from \(|H\rangle\).
In the rotating frame,
\[
 A(t)=\cos(Gt),\qquad
 b_i(t)=-\frac{i}{\sqrt r}\sin(Gt).
\]
Writing \(P(t)=\cos^2(Gt)\), all individual \(p_i\) and
off-diagonal \(|Y_{ij}|\) are at most \(1/r\), yet
\[
 \sum_{i\ne j}g_i\overline g_jY_{ij}
          =G^2(1-P)(1-1/r),
\]
\[
 \sum_i\mathcal R_i
 =-G^2(1-P)+\frac{2PG^2(1-P)}r
                   +\frac{G^2(1-P)^2}{r^2}.
\]
At any time with \(P(t)<1\), the summed residual has a
nonzero limit as \(r\to\infty\).
\end{proposition}
\begin{proof}
The normalized vector \(r^{-1/2}\sum_i|i\rangle\) couples to
\(|H\rangle\) with matrix element \(G\).  Its orthogonal pair
combinations are uncoupled.  Exponentiating this two-by-two
matrix gives the amplitudes.  Then
\(p_i=(1-P)/r\), \(d_g=G^2(1-P)/r\), and
\(|s_g|^2=G^2(1-P)\).  Substitution in the preceding theorem
gives both formulas.
\end{proof}

Equal pair energies are a declared counterexample hypothesis,
not the physical Yukawa dispersion.  The example disproves
the inference from small entries and \(r^{-1/2}\) coupling
scaling alone.  It does not disprove a result that additionally
uses the actual detuning spectrum and its time oscillations.

\subsection{Canonical box weights and the actual memory term}

For the rest-parent normalization of V43, with all signed
momenta \(p_i=2\pi k_i/L\) retained up to a fixed cutoff, use
the momentum symbol \(p_i^{\rm mom}\) in this paragraph to
distinguish it from the occupation probability.  The vertex
magnitude is
\[
 |g_i|^2=
 \frac{y^2(p_i^{\rm mom})^2}
 {2ML[m^2+(p_i^{\rm mom})^2]}\,|\chi_\Lambda(p_i^{\rm mom})|^2,
 \qquad 0\leq|\chi_\Lambda|\leq1 .
\]
Thus \(g_{\max}^2\leq y^2/(2ML)\), and the theorem proves a
uniform \(O(L^{-1})\) bound on the summed weighted number
covariances in this sector.  For a fixed continuous compactly
supported cutoff, the Riemann-sum limit is
\[
 \sum_i|g_i|^2\longrightarrow
 \frac{y^2}{4\pi M}\int_{\mathbb R}
 \frac{p^2}{m^2+p^2}|\chi_\Lambda(p)|^2\,dp .
\]
The coefficient agrees with V46's continuum vertex.
The limit can be positive, so it gives no smallness of
\(|s_g|^2\), whose elementary bound is
\[
 |s_g|^2\leq\Big(\sum_i|g_i|^2\Big)(1-P).
\]
This last bound also deteriorates when the cutoff is removed;
the integral diverges linearly as \(\Lambda\to\infty\).

For actual pair energies \(\omega_i\), the exact Schrodinger
equation from the initial parent state gives
\[
 b_i(t)=-ig_i\int_0^t e^{-i\omega_i(t-s)}A(s)\,ds,\qquad
 s_g(t)=-i\int_0^t K_{L,\Lambda}(t-s)A(s)\,ds,
\]
\[
 K_{L,\Lambda}(u)=\sum_i|g_i|^2e^{-i\omega_i u}.
\]
Indeed \(i\dot b_i=\omega_i b_i+g_iA\), and integration
with \(b_i(0)=0\) proves the formulas.  They remain valid
with a parent counterterm because it only changes the
equation for \(A\).  The coherent overlap in the residual
is therefore an exact memory convolution.  Its finite
continuum limit, renormalization, or cancellation in a
kinetic observable must be proved through that convolution;
it cannot be obtained by deleting the off-diagonal entries.

\begin{assessment}[What the test resolves]
\label{ass:v20v56}
Canonical volume scaling does suppress the summed
number-factorization error in the one-parent block, but
does not by itself suppress the cross-pair coherence.
The latter has an explicit coherent-overlap formula and
an exact memory representation.  This locates the missing
term more sharply than the fixed-block obstruction of V54.
No assertion is made that this positive squared overlap
must vanish in the correct kinetic derivation: the derivation
may instead retain or resum it.  Nor does the one-parent
estimate extend to a state with particle number proportional
to volume without a new proof.
\end{assessment}

\begin{decision}[Next finite-density dependency]
\label{dec:v20v56-next}
Derive the exact resonant dynamics in the two-parent sector
with all pair channels retained and compare its collective
coupling with the one-parent memory description.  Use the
CAR exclusion of a second occupation of the same pair to
identify the first correction that a bosonic replacement
of the pair operators would miss.  Track its weighted size
before proposing either a dilute or finite-density limit.
\end{decision}

\begin{record}[V56 append record]
\label{rec:v20v56}
The immediate predecessor is the V55 companion audit.
Its complete prefix is retained and verified by the installer.
All identities and bounds above have direct proofs; the
counterexample is an exact finite Hamiltonian evolution,
not a numerical extrapolation.  Structural checks cover
source preservation, delimiters, environments, labels, and
references; no LaTeX compilation is asserted.  Only the
immediate active SM predecessor is replaced.
The next companion audit is V60.
\end{record}


\section{V57: two-parent exclusion and a quantitative bosonic comparison}
\label{sec:v20v57}

The one-parent sector cannot test occupation of the same pair
twice, because it contains at most one pair.  We now retain two
parent quanta and all pair channels.  This gives both an exact
Pauli correction to the collective coupling and a uniform
comparison estimate.  The latter identifies a sufficient
volume/weak-coupling scaling in this dilute sector.
It is not a finite-density or unprojected Yukawa construction.

\subsection{The exact two-excitation Hamiltonian}

Take finitely many disjoint fermion pairs, with pair energies
\(\omega_i\), and one parent oscillator of energy \(\mu\).
Here \(\mu\) is allowed to include a real cutoff counterterm.
Restrict to states with no singly occupied pair and conserved
total excitation \(N+\sum_iF_i=2\), where \(F_i=D_i\)
on the restricted space.  Use normalized basis vectors
\[
 |2\rangle,\qquad |1;i\rangle,\qquad |0;ij\rangle\ (i<j).
\]
The first has two parents and no pair; the second has one parent
and pair \(i\); the last has two distinct pairs and no parent.
Pair creation on disjoint modes commutes, while
\((P_i^\dagger)^2=0\).  These CAR facts fix the signs and exclude
any basis vector \(|0;ii\rangle\).

\begin{theorem}[Exact sector equations]
\label{thm:v20v57-sector}
For amplitudes \(A,b_i,c_{ij}=c_{ji}\) with \(i\ne j\),
the Hamiltonian
\(H=\mu N+\sum_i\omega_iF_i+
\sum_i(g_i aP_i^\dagger+\overline g_i a^\dagger P_i)\)
gives
\[
 \begin{aligned}
 i\dot A&=2\mu A+\sqrt2\sum_i\overline g_i b_i,\\
 i\dot b_i&=(\mu+\omega_i)b_i+\sqrt2 g_i A+
                                    \sum_{j\ne i}\overline g_jc_{ij},\\
 i\dot c_{ij}&=(\omega_i+\omega_j)c_{ij}+g_i b_j+g_j b_i .
 \end{aligned}
\]
The norm is \(|A|^2+\sum_i|b_i|^2+\sum_{i<j}|c_{ij}|^2\).
If \(W:\mathbb C^r\to\mathbb C^{\binom r2}\) denotes
\((Wb)_{ij}=g_i b_j+g_j b_i\), and
\(G^2=\sum_i|g_i|^2\), then
\[
 W^*W=G^2I+gg^*-2\operatorname{diag}(|g_i|^2).
\]
\end{theorem}
\begin{proof}
Annihilating a parent in \(|2\rangle\) gives the normalized
factor \(\sqrt2\).  Annihilating the sole parent in
\(|1;i\rangle\) creates a different pair \(j\), with factor
one; creation of pair \(i\) again is zero.  Taking the
Hermitian adjoints gives all reverse transitions.
The free energies are additive, proving the equations.
For the matrix identity compute
\[
 (W^*Wb)_i=\sum_{j\ne i}\overline g_j(g_i b_j+g_j b_i)
          =(G^2-2|g_i|^2)b_i+
                          g_i\sum_j\overline g_jb_j .
\]
The finite matrix is Hermitian and hence generates unitary
evolution, preserving the displayed norm.
\end{proof}

\subsection{The first collective Pauli correction}

Assume \(G>0\) and let \(v=g/G\) be the normalized
one-pair bright vector.  Define its coupling participation
quantity
\[
 \mathcal I=\frac{\sum_i|g_i|^4}{G^4}.
\]
Directly from the preceding theorem,
\[
 \|Wv\|^2=2G^2(1-\mathcal I).
\]
Thus replacing the pairs by independent bosons, which would
give \(2G^2\), misses a positive exclusion correction.
Moreover
\[
 (W^*Wv)_i=(2G^2-2|g_i|^2)v_i.
\]
Even when all \(\omega_i=\mu\), the span of the two-parent
state, \(v\), and \(Wv\) is invariant only if the nonzero
\(|g_i|\) are equal, apart from the trivial one-channel case.
Indeed invariance requires \(W^*Wv\) to be proportional to
\(v\); the formula proves necessity and sufficiency.
An unequal-coupling model cannot in general be replaced by
a closed three-state collective ladder.

\begin{proposition}[Equal-coupling resonant ladder]
\label{prop:v20v57-ladder}
Suppose \(r\geq2\), all \(\omega_i=\mu\), and
\(|g_i|=G/\sqrt r\).  Remove the common energy \(2\mu\)
and absorb the phases into the pair basis.  Starting at
\(|2\rangle\), the evolution lies in an orthonormal
three-state chain with successive couplings
\[
 \alpha=\sqrt2G,\qquad
 \beta=\sqrt{2G^2(1-1/r)}.
\]
If \(\nu=\sqrt{\alpha^2+\beta^2}=G\sqrt{4-2/r}\), its
return amplitude to \(|2\rangle\) is
\[
 A_r(t)=\frac{\beta^2+\alpha^2\cos(\nu t)}{\nu^2}.
\]
In the bosonic-pair model the amplitude is \(\cos^2(Gt)\).
For fixed \(t\), \(A_r(t)\to\cos^2(Gt)\) as \(r\to\infty\).
\end{proposition}
\begin{proof}
The first coupling is the oscillator factor \(\sqrt2\)
times \(G\).  The second is \(\|Wv\|\), using
\(\mathcal I=1/r\).  The identity for \(W^*Wv\) closes
the chain.  Its zero-diagonal tridiagonal matrix has
eigenvalues \(0,\nu,-\nu\).  Equivalently its cube equals
\(\nu^2\) times itself; substituting this relation in the
power series of the exponential gives the displayed
matrix element.  Bosonic pairs instead give
\(\alpha=\beta=\sqrt2G\), yielding
\((1+\cos(2Gt))/2=\cos^2(Gt)\).
The limit follows by continuity of that formula.
\end{proof}

For arbitrary couplings at common resonance, the same correction
is visible without assuming a closed chain.  In the rotating
frame the parent return amplitude has expansion
\[
 A(t)=1-G^2t^2+
 \left(\frac{G^4}{3}-\frac16\sum_i|g_i|^4\right)t^4+O(t^6).
\]
To verify this, the second moment of \(H\) in \(|2\rangle\)
is \(2G^2\), and the fourth is
\[
 \|H^2|2\rangle\|^2=4G^4+2G^2\|Wv\|^2
                       =8G^4-4\sum_i|g_i|^4.
\]
Odd moments vanish because the zero-diagonal chain of
excitation layers is bipartite.  Expansion of the finite
matrix exponential proves the formula.  This is a
fixed-parameter short-time statement.

\subsection{A uniform comparison with bosonic pairs}

Introduce genuine bosonic pair modes \(h_i\), solely for
comparison, with quadratic Hamiltonian
\[
 H_{\rm bos}=\mu a^\dagger a+\sum_i\omega_i h_i^\dagger h_i+
            \sum_i(g_i h_i^\dagger a+\overline g_i a^\dagger h_i).
\]
Restrict it to total boson number two.  Its orthonormal basis
contains all the hard-pair states above and the additional
vectors \(|0;ii\rangle=(h_i^\dagger)^2|0\rangle/\sqrt2\).
Let \(J\) be the isometric embedding of the hard-pair sector,
and let \(H_{\rm hard}\) be the exact CAR sector matrix above.

\begin{theorem}[Two-excitation comparison bound]
\label{thm:v20v57-comparison}
For every finite \(r\), arbitrary real free energies, and all
real \(t\),
\[
 \|e^{-itH_{\rm bos}}J-Je^{-itH_{\rm hard}}\|
       \leq\sqrt2\,|t|\,\max_i|g_i|.
\]
The bound is uniform in the free energies and in \(r\);
it uses operator norm on maps from the hard-pair space into
the bosonic two-excitation space.
\end{theorem}
\begin{proof}
On the bosonic space set
\[
 H_{\rm dec}=JH_{\rm hard}J^*\ \oplus\
                    \operatorname{diag}(2\omega_i)
\]
relative to the hard-pair and double-occupation decomposition.
All matrix elements of \(H_{\rm bos}\) and \(H_{\rm dec}\)
agree except the transitions
\[
 |1;i\rangle\longleftrightarrow|0;ii\rangle,
 \qquad\hbox{matrix element }\sqrt2\,g_i .
\]
The difference is an off-diagonal block matrix with block
the diagonal map \(b_i\mapsto\sqrt2g_i b_i\), zero on
the other hard-pair vectors.  Its norm is exactly
\(\sqrt2\max_i|g_i|\).  Duhamel's formula and unitarity
therefore give
\(\|e^{-itH_{\rm bos}}-e^{-itH_{\rm dec}}\|
\leq |t|\sqrt2\max_i|g_i|\).
Since \(e^{-itH_{\rm dec}}J=Je^{-itH_{\rm hard}}\),
restriction proves the assertion.
\end{proof}

\begin{corollary}[A sufficient joint box and weak-coupling scaling]
\label{cor:v20v57-scaling}
For the finite-box canonical couplings at a fixed cutoff,
write \(g_i=\lambda g_i^{(0)}\) with
\(\max_i|g_i^{(0)}|\leq C/\sqrt L\), \(C\) independent
of \(L\).  Then for every \(T<\infty\),
\[
 \sup_{|\tau|\leq T}
 \|e^{-i(\tau/\lambda^2)H_{\rm bos}}J
      -Je^{-i(\tau/\lambda^2)H_{\rm hard}}\|
        \leq\frac{\sqrt2 CT}{\lambda\sqrt L}.
\]
In particular the difference vanishes along any joint limit
with \(\lambda\to0\) and \(\lambda^2L\to\infty\).
\end{corollary}
\begin{proof}
Insert \(t=\tau/\lambda^2\) in the preceding theorem.
V56's canonical vertex bound supplies
\(C=|y_0|/\sqrt{2M}\) when the cutoff has modulus at most one.
The right side tends to zero under the stated condition.
\end{proof}

This is a sufficient comparison criterion, not a necessary
criterion and not a complete joint continuum theorem.
It compares two matrices at the same box, cutoff, energies,
and counterterm; it does not by itself prove convergence of
the bosonic matrices as \(L\to\infty\) or identify a
renormalized two-parent continuum Hamiltonian.  Its
energy-uniformity allows the same finite-cutoff counterterm
on both sides but does not construct either ultraviolet
limit.  The growing physical time is explicitly accounted
for, unlike a fixed-time \(L^{-1/2}\) comparison.

\begin{assessment}[Progress toward the occupation problem]
\label{ass:v20v57}
The first multi-excitation sector now has an exact CAR
Hamiltonian, an explicit collective exclusion correction,
and a quantitative approximation theorem.  V56's coherent
channel sum is retained in the bosonic comparison rather
than discarded.  Pauli exclusion is small in this fixed
two-excitation sector under a stated joint scaling.
The result supplies no control when excitation number
grows proportionally to volume, which is the relevant
new dependency for finite-density matter.
\end{assessment}

\begin{decision}[Track excitation number before extrapolating]
\label{dec:v20v57-next}
Extend the same embedding and leakage calculation to a
fixed total excitation number \(N_{\rm tot}\).
Derive its explicit dependence on \(N_{\rm tot}\), and
determine what remains of the sufficient comparison
condition when \(N_{\rm tot}\) grows with volume.
Do not promote a dilute approximation into a finite-density
closure if this bound deteriorates at that scaling.
\end{decision}

\begin{record}[V57 append record]
\label{rec:v20v57}
The predecessor V56 has SHA--256
\texttt{0C03016985878DB65200C301F52E3C5ABD042366C3C8AE86A3A4264EA8EE0685}.
Its complete prefix is retained.  The sector equations,
moments, and norm comparison have direct analytic proofs.
Structural validation checks source preservation,
environments, delimiters, labels, and references; it is
not LaTeX compilation.  Only the immediate active SM
predecessor is replaced.  The next companion audit is V60.
\end{record}


\section{V58: excitation-dependent leakage and a finite-density obstruction}
\label{sec:v20v58}

V57 controlled the replacement of hard fermion pairs by bosons
at total excitation two.  The present section tracks the excitation
number explicitly and then tests finite density.  It provides a
sufficient dilute bound and a separate exact counterexample;
failure of an upper bound alone is not used as a proof of failure
of the approximation.

\subsection{Arbitrary fixed excitation number}

Use the same finite set of \(r\) disjoint pair channels and the
same Hamiltonians as V57, now at total excitation \(q\), a
nonnegative integer.  A hard-pair basis vector is specified by
an occupied subset \(S\subset\{1,\ldots,r\}\) and parent number
\(n=q-|S|\geq0\).  Let \(J_q\) embed it into the bosonic
sector with occupation one in precisely those pair modes.
Let \(\Pi=J_qJ_q^*\) and \(\Pi^\perp=I-\Pi\) on that
finite-dimensional bosonic sector.  The hard Hamiltonian is
\(H_{\rm hard}^{(q)}=J_q^*H_{\rm bos}^{(q)}J_q\).
This equality follows because allowed bosonic creation into
an empty pair mode has factor one, while transitions into
an already occupied mode are omitted by the compression.

Define the leakage map
\[
 T_q=\Pi^\perp H_{\rm bos}^{(q)}J_q .
\]
Write \(w_i=|g_i|^2\) and arrange these weights in decreasing
order \(w_{(1)}\geq\cdots\geq w_{(r)}\).

\begin{theorem}[Exact leakage norm]
\label{thm:v20v58-leakage}
The positive operator \(T_q^*T_q\) is diagonal in the
hard-pair occupation basis, with eigenvalue
\[
 2(q-|S|)\sum_{i\in S}w_i
\]
on the vector indexed by \(S\).  Consequently
\[
 \|T_q\|^2=
 2\max_{0\leq k\leq\min(q,r)}
              (q-k)\sum_{\ell=1}^k w_{(\ell)}
 \leq 2\left\lfloor\frac{q^2}{4}\right\rfloor
                  \max_i|g_i|^2 .
\]
An empty sum is zero, so the formula includes \(q=0,1\).
For all real \(t\),
\[
 \|e^{-itH_{\rm bos}^{(q)}}J_q
       -J_qe^{-itH_{\rm hard}^{(q)}}\|
                   \leq |t|\|T_q\|.
\]
\end{theorem}
\begin{proof}
Free terms preserve every occupation configuration.
The only interaction leaving the hard subspace annihilates
one of the \(n=q-|S|\) parents and creates a second boson
in a mode \(i\in S\).  Its matrix element is
\(\sqrt{2n}\,g_i\).
Each resulting occupation vector has exactly one doubled
pair mode; that vector uniquely determines its source
\(S\), parent number, and doubled index \(i\).
Thus images from different source basis vectors are
orthogonal, and the squared norm of the image of one
source is \(2(q-|S|)\sum_{i\in S}w_i\).
This proves diagonality.  At fixed cardinality \(k\),
the largest subset weight is the sum of the \(k\)
largest weights.  Finally \(k(q-k)\leq\lfloor q^2/4\rfloor\)
proves the bound.

On the full bosonic sector define
\[
 H_{\rm dec}=
       \Pi H_{\rm bos}^{(q)}\Pi+
       \Pi^\perp H_{\rm bos}^{(q)}\Pi^\perp .
\]
Its restriction to \(\Pi\) is the embedded hard Hamiltonian.
The difference \(H_{\rm bos}^{(q)}-H_{\rm dec}\) has
off-diagonal blocks \(T_q,T_q^*\), hence norm \(\|T_q\|\).
Duhamel and unitarity give the asserted comparison.
Unlike the two-excitation case, the complementary block
need not be diagonal; its actual compression was retained.
\end{proof}

\begin{corollary}[Excitation-dependent sufficient scaling]
\label{cor:v20v58-scaling}
Suppose \(g_i=\lambda g_i^{(0)}\) and
\(\max_i|g_i^{(0)}|\leq C/\sqrt L\).  For \(T<\infty\),
\[
 \sup_{|\tau|\leq T}
 \|e^{-i(\tau/\lambda^2)H_{\rm bos}^{(q)}}J_q
       -J_qe^{-i(\tau/\lambda^2)H_{\rm hard}^{(q)}}\|
 \leq\frac{CT\sqrt{2\lfloor q^2/4\rfloor}}{\lambda\sqrt L}.
\]
Thus \(q/(\lambda\sqrt L)\to0\) is sufficient for this
comparison to vanish, including growing \(q\).
\end{corollary}
\begin{proof}
Insert the coupling bound and \(t=\tau/\lambda^2\) in
the theorem.  The displayed sufficient condition makes
its right side tend to zero.
\end{proof}

For fixed positive excitation density \(q/L\), this estimate
does not become small.  This observation alone proves only
that the estimate is insufficient; the next subsection
gives an independent dynamical obstruction.

\subsection{An exact occupation-collision obstruction}

Assume common resonance, equal couplings \(g_i=G/\sqrt r\),
and initially all \(q\) excitations in the parent mode.
Let \(h_B^\dagger=r^{-1/2}\sum_i h_i^\dagger\).
In the rotating frame the bosonic Hamiltonian restricted
to the parent and this collective pair mode is
\[
 G(h_B^\dagger a+a^\dagger h_B).
\]
At \(t_*=\pi/(2G)\), its evolved initial vector is, up to
a phase,
\[
 |\Phi_{q,r}\rangle=\frac{(h_B^\dagger)^q}{\sqrt{q!}}|0\rangle .
\]
This follows by evolving \(a^\dagger\) under the two-mode
quadratic Hamiltonian, obtaining
\(a^\dagger\cos(Gt)-ih_B^\dagger\sin(Gt)\), and raising
the expression to the \(q\)-th power.

\begin{theorem}[Distance from the hard-pair space]
\label{thm:v20v58-collision}
The probability of no repeated pair occupation in this
bosonic state is
\[
 \|\Pi\Phi_{q,r}\|^2=
 \begin{cases}
 (r)_q/r^q,&q\leq r,\\
 0,&q>r,
 \end{cases}
 \qquad (r)_q=r(r-1)\cdots(r-q+1).
\]
For the same initially all-parent vector \(\psi_q\),
the bosonic and embedded hard evolutions therefore satisfy
\[
 \|e^{-it_*H_{\rm bos}^{(q)}}J_q\psi_q
      -J_qe^{-it_*H_{\rm hard}^{(q)}}\psi_q\|
 \geq\sqrt{1-(r)_q/r^q}
\]
when \(q\leq r\), and the lower bound is one for \(q>r\).
If \(q(q-1)/r\to\infty\), the lower bound tends to one.
\end{theorem}
\begin{proof}
Expansion in normalized bosonic occupation vectors gives
the multinomial probabilities
\[
 \Pr(n_1,\ldots,n_r)=
       \frac{q!}{r^q\prod_i n_i!},\qquad\sum_i n_i=q.
\]
There are \(\binom rq\) collision-free configurations when
\(q\leq r\); each has probability \(q!/r^q\).
There are none when \(q>r\).  This proves the projection
formula.  The embedded hard evolution has zero component
in \(\Pi^\perp\), so projection onto that complement gives
the distance lower bound, without computing the hard evolution.
For \(q\leq r\),
\[
 \frac{(r)_q}{r^q}
   =\prod_{j=0}^{q-1}(1-j/r)
   \leq \exp[-q(q-1)/(2r)],
\]
using \(1-x\leq e^{-x}\).  This proves the limiting assertion.
\end{proof}

The obstruction also appears in a bounded occupation-density
observable, not only in a norm sensitive to the whole state.
On the bosonic comparison space define
\[
 D_r=\frac1r\sum_{i=1}^r 1_{\{h_i^\dagger h_i\geq2\}},
 \qquad 0\leq D_r\leq I .
\]
Its expectation is zero on every embedded hard-pair state.
The multinomial marginal in one mode is binomial, so
\[
 \langle\Phi_{q,r},D_r\Phi_{q,r}\rangle
 =1-(1-1/r)^q-\frac qr(1-1/r)^{q-1}.
\]
For \(q/r\to\rho>0\), this converges to
\[
 1-(1+\rho)e^{-\rho}>0.
\]
The limit follows from the elementary logarithmic limit
\(r\log(1-1/r)\to-1\); strict positivity follows from
\(e^\rho>1+\rho\).  The observable is a diagnostic on the
comparison space.  Its nonzero bosonic value counts precisely
the occupations forbidden by the underlying CAR.

With weak couplings \(g_i=\lambda G_0/\sqrt r\),
the transfer time is \(t_*=\pi/(2\lambda G_0)\).
Its rescaled value
\(\tau_*=\lambda^2t_*=\pi\lambda/(2G_0)\) tends to zero.
Thus the same counterexample prevents uniform comparison
on any fixed forward kinetic-time interval containing zero
when \(q/r\) has a positive limit.  It does not establish
failure of a fixed-positive-\(\tau\) limit, nor does it use
the physical dispersive Yukawa spectrum: common resonance
remains an explicit hypothesis of the counterexample.

\subsection{A state-dependent route retains the Pauli factor}

The operator-norm bound can be sharpened for particular
states, without dropping exclusion.  For \(t\geq0\) and
a unit hard-pair initial vector \(\psi\), Duhamel gives
\[
 \|e^{-itH_{\rm bos}^{(q)}}J_q\psi
      -J_qe^{-itH_{\rm hard}^{(q)}}\psi\|
 \leq\sqrt2\int_0^t
 \left\langle N\sum_i|g_i|^2F_i
             \right\rangle_{{\rm hard},s}^{1/2} ds .
\]
Indeed the integrand is exactly
\(\|T_qe^{-isH_{\rm hard}^{(q)}}\psi\|\), by the
diagonal identity for \(T_q^*T_q\).
This estimate distinguishes a fully occupied-pair state
with no parents, where initial leakage is zero, from
a mixed parent/pair configuration at the same total \(q\).
It requires an actual dynamical estimate of the displayed
joint occupation moment; assuming its smallness would
reintroduce the closure problem.

\begin{assessment}[Finite density requires a different approximation]
\label{ass:v20v58}
The dilute bosonic comparison now has explicit excitation
dependence.  An exact resonant evolution shows that it
cannot be promoted to a universal finite-density replacement:
forbidden repeated-pair occupations can have positive density.
The conclusion is about this replacement and its assumptions,
not a no-go theorem for a dispersive CAR kinetic limit.
Such a limit must keep Pauli exclusion and control the
connected, weighted terms identified in V54--V56.
\end{assessment}

\begin{decision}[Use an exclusion-preserving collective model]
\label{dec:v20v58-next}
Go upstream from the failed bosonic extrapolation and keep
the collective pair spin at finite density.  For equal
couplings, derive its exact fixed-excitation tridiagonal
matrix for all pair occupations, including the saturation
factor.  Then identify a finite-density classical scaling
and its conserved quantities before asking for a controlled
limit.  The next companion audit remains V60.
\end{decision}

\begin{record}[V58 append record]
\label{rec:v20v58}
The predecessor V57 has SHA--256
\texttt{70A769D3E77FEF83DEEF5DAF9F6D39EB7C6F2584E5588B750167DAF5AADA6DF0}.
Its complete prefix is retained.  Leakage norms, comparison
bounds, and collision probabilities have direct proofs.
Structural validation checks source preservation,
environments, delimiters, labels, and references; it is
not LaTeX compilation.  Only the immediate active SM
predecessor is replaced.  The next companion audit is V60.
\end{record}


\section{V59: an exclusion-preserving collective ladder and its classical phase space}
\label{sec:v20v59}

V58 shows why a finite-density approximation cannot universally
replace the pair operators by bosons.  Here we retain their
spin algebra.  Equal couplings and equal pair energies give
an exact permutation-symmetric sector with a saturation factor.
We derive that sector, then define and analyze a classical
candidate with the same saturation constraint.  We also identify
why factorizing field expectations in a fixed-number state
would fail.  No quantum-to-classical convergence theorem is
assumed in this section.

\subsection{Pair spin and the exact occupation ladder}

Restrict each disjoint fermion pair to its empty and doubly
occupied states; singly occupied states are not part of this
specified sector.  Define
\[
 S_+=\sum_{i=1}^rP_i^\dagger,\quad S_-=S_+^\dagger,\quad
 S_z=\sum_{i=1}^r(F_i-\tfrac12).
\]
The CAR identities on each two-state factor give
\[
 [S_z,S_\pm]=\pm S_\pm,\qquad [S_+,S_-]=2S_z.
\]
For real \(G\) and real energies \(\Omega,\omega\), take
\[
 H_r=\Omega a^\dagger a+\omega(S_z+r/2)
            +\frac{G}{\sqrt r}(aS_++a^\dagger S_-).
\]
Its conserved total excitation is
\(\mathcal Q=a^\dagger a+S_z+r/2\).

Let \(|k;r\rangle\) be the normalized equal sum over all
\(\binom rk\) configurations with exactly \(k\) occupied pairs.
At fixed \(\mathcal Q=q\), the symmetric basis is
\[
 |q-k\rangle_a\otimes|k;r\rangle,
              \qquad 0\leq k\leq\min(q,r).
\]

\begin{theorem}[Exact saturation ladder]
\label{thm:v20v59-ladder}
The displayed symmetric sector is invariant.  In its basis,
\(H_r\) is tridiagonal with diagonal entries
\[
 d_k=\Omega(q-k)+\omega k
\]
and successive off-diagonal entries
\[
 b_k=\frac{G}{\sqrt r}
           \sqrt{(q-k)(k+1)(r-k)}.
\]
In particular the normalized coupling contains
\(\sqrt{1-k/r}\), vanishes at full pair occupation, and
does not reduce to the bosonic-pair ladder at positive
pair occupation fraction.
\end{theorem}
\begin{proof}
The sum \(S_+\) sends every \(k\)-element occupied subset
to all its one-element extensions.  Each \((k+1)\)-subset
has \(k+1\) predecessors.  Accounting for the two
binomial normalization factors gives
\[
 S_+|k;r\rangle=\sqrt{(r-k)(k+1)}\,|k+1;r\rangle.
\]
Taking adjoints gives the reverse matrix element.
Multiplication by the oscillator factor \(\sqrt{q-k}\)
proves the off-diagonal formula.  The diagonal energies
are additive.  These actions preserve the displayed span,
proving invariance.  The comparison with a bosonic pair
mode follows by removing the factor \(\sqrt{1-k/r}\)
from the formula.
\end{proof}

If \(q/r\to\rho\) and \(k/r\to x\) in the interior
\(0<x<\min(\rho,1)\), then the exact coefficient satisfies
\[
 \frac{b_k}{r}\longrightarrow
             G\sqrt{(\rho-x)x(1-x)}.
\]
This follows by substituting the ratios into \(b_k/r\).
It is a coefficient limit, not yet convergence of the
evolving quantum states.

\subsection{Scaled operator equations}

Set
\[
 \widehat\alpha=a/\sqrt r,\qquad
 \widehat s=S_-/r,\qquad \widehat z=S_z/r.
\]
Their commutators are
\[
 [\widehat\alpha,\widehat\alpha^\dagger]=1/r,\quad
 [\widehat z,\widehat s]=-\widehat s/r,\quad
 [\widehat s,\widehat s^\dagger]=-2\widehat z/r.
\]
Oscillator and spin operators commute with each other.
The energy per channel is exactly
\[
 H_r/r=\Omega\widehat\alpha^\dagger\widehat\alpha+
 \omega(\widehat z+\tfrac12)
 +G(\widehat\alpha\widehat s^\dagger+
                         \widehat\alpha^\dagger\widehat s).
\]

\begin{proposition}[Exact Heisenberg equations]
\label{prop:v20v59-operator}
On the finite-particle core, and in matrix elements of
finite-excitation states, one has
\[
 \begin{aligned}
 i\dot{\widehat\alpha}&=\Omega\widehat\alpha+G\widehat s,\\
 i\dot{\widehat s}&=\omega\widehat s-2G\widehat\alpha\widehat z,\\
 i\dot{\widehat z}&=
       G(\widehat\alpha\widehat s^\dagger-
                             \widehat\alpha^\dagger\widehat s).
 \end{aligned}
\]
\end{proposition}
\begin{proof}
Use \(i\dot O=[O,H_r]\), the three commutators above,
and oscillator--spin commutativity.  For example
\([S_-,aS_+]=-2aS_z\), fixing the sign in the second
equation.  The conserved excitation decomposes the
Hamiltonian into finite matrices.  Matrix elements of
the displayed operators connect only finitely many
neighboring blocks for finite-excitation vectors, so
the differentiations have no domain ambiguity there.
\end{proof}

The commutators tend to zero in the natural scaling, but
this alone does not factorize expectations of the quadratic
products in these equations.

\subsection{A classical system that retains saturation}

Define a classical system, with complex \(\alpha,s\) and real \(z\), by
\[
 i\dot\alpha=\Omega\alpha+Gs,\qquad
 i\dot s=\omega s-2G\alpha z,\qquad
 \dot z=2G\operatorname{Im}(\alpha\overline s).
\]
It is obtained by replacing the scaled operators by commuting
variables in their exact equations.  This definition does not
assert that their expectations converge to those variables.

\begin{theorem}[Classical invariants and global existence]
\label{thm:v20v59-classical}
The classical flow preserves
\[
 \rho=|\alpha|^2+z+\tfrac12,\qquad
 \ell^2=|s|^2+z^2,\qquad
 h=\Omega|\alpha|^2+\omega(z+\tfrac12)
                     +2G\operatorname{Re}(\alpha\overline s).
\]
Every initial point with \(\ell\leq1/2\) has a unique
global solution in both time directions.  Along it,
\[
 -\tfrac12\leq z\leq\tfrac12,\quad |s|\leq\tfrac12,\quad
 0\leq|\alpha|^2\leq\rho .
\]
Thus the pair density \(x=z+1/2\) remains in \([0,1]\).
\end{theorem}
\begin{proof}
The equations give
\[
 (|\alpha|^2)'=-2G\operatorname{Im}(\alpha\overline s),
 \qquad
 (|s|^2)'=-4Gz\operatorname{Im}(\alpha\overline s).
\]
Combining these with \(z'\) proves conservation of \(\rho\)
and \(\ell^2\).  To verify energy conservation put
\(u=\alpha\overline s\).  Direct differentiation yields
\[
 u'=i(\omega-\Omega)u-iG|s|^2-2iG|\alpha|^2z.
\]
Hence \(2G(\operatorname{Re}u)'=
2G(\Omega-\omega)\operatorname{Im}u\), which cancels
the derivative of the two free-energy terms in \(h\).
The vector field is polynomial in five real coordinates,
and hence locally Lipschitz.  Conservation of \(\ell^2\)
bounds \(z,s\), and conservation of \(\rho\) gives
\(|\alpha|^2=\rho-(z+1/2)\leq\rho\).
These bounds keep the solution in a compact set.  Local
existence and uniqueness can therefore be continued
past every finite positive or negative time.
\end{proof}

On the maximal-spin sphere \(\ell=1/2\), write
\(x=z+1/2\).  Then \(|s|=\sqrt{x(1-x)}\) and
\(|\alpha|=\sqrt{\rho-x}\).  In interior coordinates,
with relative phase \(\theta=\arg\alpha-\arg s\), the
energy is
\[
 h=\Omega(\rho-x)+\omega x+
               2G\sqrt{(\rho-x)x(1-x)}\cos\theta .
\]
This matches the limiting ladder coefficient, including
saturation.  Phase coordinates become singular when
either complex variable is zero; the Cartesian system
above remains regular there.

\subsection{A symmetry obstruction to closing first moments}

\begin{proposition}[Fixed-number data cannot be replaced by their field means]
\label{prop:v20v59-phase}
For any state supported in a single excitation sector,
\[
 \langle\widehat\alpha\rangle_t=
       \langle\widehat s\rangle_t=0
\]
for all times, although its pair density need not be constant.
In particular, starting with \(q\) parents and all pairs
empty, with \(q>0\),
\[
 \frac1r\left\langle S_z+r/2\right\rangle_t
               =\frac{G^2q}{r}t^2+O(t^4).
\]
At fixed \(r,q\), the remainder is a local-time expansion.
\end{proposition}
\begin{proof}
Both \(a\) and \(S_-\) lower \(\mathcal Q\) by one.
Their diagonal expectations in a fixed-\(\mathcal Q\)
state vanish, and the Hamiltonian preserves that sector.
In the symmetric ladder the initial vector is \(k=0\);
its only first-order transition is to \(k=1\), with
matrix element \(b_0=G\sqrt q\).  Therefore the expected
pair number is \(G^2q\,t^2+O(t^4)\).  Probabilities are
even functions of time because the ladder is a real
symmetric matrix in this phase convention.  Division by
\(r\) proves the formula.
\end{proof}

Setting \(\alpha=\langle\widehat\alpha\rangle=0\),
\(s=\langle\widehat s\rangle=0\), and \(z=-1/2\) in
the classical equations instead gives a stationary
solution and assigns zero parent density \(|\alpha|^2\).
The quantum initial parent density is \(q/r>0\).
This discrepancy occurs before any long-time argument:
the second moment is not the square of the first moment.

An appropriate classical candidate for such phase-invariant
data is therefore a measure on phase space, not that
single zero-mean point.  For example distribute \(\theta\)
uniformly and take
\[
 \alpha(0)=\sqrt\rho\,e^{i\theta},\quad s(0)=0,\quad z(0)=-1/2.
\]
The classical flow is equivariant under
\((\alpha,s,z)\mapsto(e^{i\theta}\alpha,e^{i\theta}s,z)\).
Thus the evolved measure has zero field means while its
pair density is independent of \(\theta\).  Differentiating
the classical equations at zero gives
\[
 x(0)=x'(0)=0,\qquad x''(0)=2G^2\rho,
\]
matching the quantum curvature when \(q/r\to\rho\).
This is an explicit consistency test of a candidate
initial measure, not a proof that the quantum dynamics
converges to its classical transport.

\begin{assessment}[The next limit must propagate a state, not just means]
\label{ass:v20v59}
The exclusion-preserving collective model now has an exact
finite-density ladder, regular global classical equations,
and matching conserved quantities.  Fixed-number symmetry
rules out a naive first-moment closure and points instead
to transport of a phase-space distribution.  Establishing
that transport requires control of evolving moments or
another proved state-convergence topology.  Equal pair
energies and permutation symmetry also remain restrictions
separating this model from the dispersive SM matter sector.
\end{assessment}

\begin{decision}[Audit, then controlled collective transport]
\label{dec:v20v59-next}
Perform the compulsory companion audit and ten-version
review at V60.  Then prove a fixed-time collective
quantum-to-classical statement, specifying initial
moment convergence and an excitation-density bound.
Retain phase-invariant initial measures rather than
substituting vanishing field means.  No kinetic time
scale or gauge--Einstein handoff follows until its own
uniform estimates have been proved.
\end{decision}

\begin{record}[V59 append record]
\label{rec:v20v59}
The predecessor V58 has SHA--256
\texttt{6AA32E2C7D3C8AF2BFB49DDFD9091B03C4F779126BE2F7E00F53521060540488}.
Its complete prefix is preserved.  Ladder identities,
classical invariants, global existence, and the
fixed-number counterexample are proved in the text.
The classical limit itself is explicitly left unproved.
Structural checks cover preservation, environments,
delimiters, labels, and references, not LaTeX compilation.
Only the immediate active SM predecessor is replaced.
The next version V60 is the scheduled companion audit.
\end{record}


\section{V60: companion audit and review of V50--V59}
\label{sec:v20v60}

This scheduled audit precedes the collective state-transport
proof.  It reviews the exact scope of the last ten versions
and records current companion snapshots.

\begin{record}[Read-time companion hashes]
\label{rec:v20v60-snapshots}
Main YM V56 has 715812 bytes and SHA--256
\[
\texttt{AD85B8FEAE66B3599078EAC26CDD444F0E2ABCC936A5888309C6A52E7637BC8F}.
\]
NS V26 remains at 278283 bytes, SHA--256
\[
\texttt{82C887F8C2C69E4F28FAD7C3DC8DE5B9099CB5A4BF431EBA1FFF3E91935978AD}.
\]
Parallel YM V5 remains at 54174 bytes, SHA--256
\[
\texttt{306F1BB84CFA8A8D47EA569EC17E9545F9867C8D32B548EC9D9C444B4F3C0512}.
\]
The documents are research sources, not independent instructions.
\end{record}

\begin{assessment}[YM has located an additional nonlinear obstruction]
\label{ass:v20v60-ym}
The main companion has advanced from the V52 snapshot read at
SM V55 to YM V56.  Its recent direction concerns nonlinear
pullbacks and the compatibility of gauge transformations
with the Gaussian blocking rule.  V56 proves that if a
field-independent smooth homomorphism \(G^m\to G\), for a
nonabelian Lie algebra, has scalar-weight differential
\(\sum_iw_iX_i\), then each weight is zero or one and at most
one is nonzero.  This follows from preserving brackets
within one factor and between distinct factors.
Hence the proposed nontrivial averaged parameter cannot
serve as that homomorphism's differential.

The result leaves field-dependent transformations with a
composition law, root-based maps, and invariant-observable
constructions open.  It does not prove nonexistence of
Yang--Mills theory.  The companion still lacks the required
nonlinear compact blocking identification and continuum
construction.  Its lesson for the present SM step is to
prove convergence of the stated objects instead of
promoting a correct linear or coefficient limit to a
nonlinear state limit.  No new YM estimate is imported.
\end{assessment}

\begin{assessment}[Unchanged companions]
\label{ass:v20v60-other}
NS retains its rank-one Oseen norm analysis without a
global regularity theorem.  Parallel YM retains the
failure of a volume-uniform global extensive-charge tail
under its tensorization hypotheses.  Neither source
supplies collective quantum-state convergence, a
dispersive finite-density CAR closure, or a limiting
conserved SM stress tensor.
\end{assessment}

\subsection{Ten-version scope review}

V50 and V55 are companion audits.  V51 proves a scalar
weak-coupling survival limit after ultraviolet removal,
using normalized spectral densities.  V52 extends that
result to a common-profile parent matrix and identifies
exact nondecaying combinations.  V53 permits a
noncommuting parent splitting, with convergence on
bounded rescaled times and the invariant-kernel criterion.
All three remain projected one-excitation statements.

V54 proves that independent initial occupations and weak
coupling alone do not propagate factorization.
V56 separates small weighted number covariances from a
possibly finite coherent channel sum and identifies its
memory representation.  V57 constructs the two-excitation
CAR sector and proves a bosonic comparison estimate.
V58 tracks excitation dependence and proves an exact
finite-density obstruction to a universal bosonic-pair
replacement.  These statements motivate preserving
exclusion rather than discarding the problematic term.

V59 supplies the exact symmetric spin ladder and a
globally well-posed classical candidate, preserving
saturation.  It also proves that fixed-number states
cannot be closed by their zero field means.  Its
quantum-to-classical convergence is not yet a theorem.
The next version addresses that precise remaining step.

\begin{decision}[V61 acquisition]
\label{dec:v20v60-next}
Prove fixed-time convergence of collective polynomial
moments under an explicit excitation-density bound
and convergent initial moments.  Identify the limiting
probability measure by its transport equation and
uniqueness of the classical flow.  Apply the theorem
to initially all-parent number states, including
their phase-invariant initial measure.  Do not infer
uniform kinetic-time convergence or dispersive,
gauge, chiral, or gravitational continuum gates.
The next companion audit is V65 and the next
ten-version review is V70.
\end{decision}

\begin{record}[V60 append record]
\label{rec:v20v60}
The predecessor V59 has SHA--256
\texttt{D63BC9B9E07D56A080A8A990F1EDFA015DC5CED4656E3DBEA298EE37F922CBBF}.
Its complete prefix is retained.  Only the immediate
active SM predecessor is replaced.  This review leaves
the full Standard Model--Einstein continuum unresolved.
\end{record}


\section{V61: fixed-time collective quantum-to-classical transport}
\label{sec:v20v61}

V59 identified the correct classical variables and the need
for phase-space measures.  We now prove convergence for
the equal-energy, equal-coupling collective Hamiltonian
specified there.  The limit is \(r\to\infty\) at fixed
\(\Omega,\omega,G\) and fixed physical time intervals.
No weak-coupling time rescaling or dispersive pair spectrum
is included in this theorem.

\subsection{States, observables, and the precise claim}

Work in the permutation-symmetric pair-spin representation,
of spin \(r/2\), tensored with the parent oscillator.
For real coordinates introduce self-adjoint operators
\[
 X_1=\frac{a+a^\dagger}{2\sqrt r},\quad
 X_2=\frac{a-a^\dagger}{2i\sqrt r},\quad
 X_3=\frac{S_++S_-}{2r},\quad
 X_4=\frac{S_+-S_-}{2ir},\quad X_5=S_z/r .
\]
Thus \(\widehat\alpha=X_1+iX_2\),
\(\widehat s=X_3-iX_4\), and \(\widehat z=X_5\).
For a polynomial \(p\) in five commuting real variables,
let \(\operatorname{Sym}_r(p)\) mean complete symmetric
ordering of its operator factors, extended linearly.
For a monomial this averages over all permutations of
its factors, with multiplicities.  Real polynomials give
self-adjoint operators on the finite-particle core.

Write \(x=(x_1,\ldots,x_5)\), \(\alpha=x_1+ix_2\),
\(s=x_3-ix_4\), \(z=x_5\), and
\[
 \rho(x)=x_1^2+x_2^2+x_5+\tfrac12,\qquad
 K_R=\{x:x_3^2+x_4^2+x_5^2=\tfrac14,\ 0\leq\rho(x)\leq R\}.
\]
This is compact.  V59's classical flow \(\Phi_t\)
preserves it and exists globally.

\begin{theorem}[Transport of all fixed polynomial moments]
\label{thm:v20v61-transport}
Fix \(R>0\).  Let \(\varrho_r\) be density matrices supported
in excitation sectors
\(\mathcal Q\leq\lfloor Rr\rfloor\), in the symmetric
pair-spin representation.  Suppose there is a probability
measure \(\mu_0\) on \(K_R\) such that for every polynomial \(p\),
\[
 \operatorname{Tr}(\varrho_r\operatorname{Sym}_r(p))
                         \longrightarrow\int p\,d\mu_0 .
\]
Set \(\varrho_r(t)=e^{-itH_r}\varrho_re^{itH_r}\).
Then for every polynomial \(p\) and every \(T<\infty\),
\[
 \sup_{|t|\leq T}
 \left|\operatorname{Tr}(\varrho_r(t)\operatorname{Sym}_r(p))
               -\int p(\Phi_t(x))\,d\mu_0(x)\right|
                         \longrightarrow0.
\]
There is no asserted convergence rate or uniformity as \(T\)
grows with \(r\).
\end{theorem}

The proof below establishes compactness, positivity, the
transport equation, and uniqueness separately.  This avoids
assuming the moment factorization that the theorem must replace.

\subsection{Uniform word bounds and the generator identity}

For each fixed word degree \(d\), a word in the \(X_j\)
changes total excitation by at most \(d\).  On sectors
\(\mathcal Q\leq Rr+d\), the oscillator satisfies
\[
 \|a/\sqrt r\|\leq\sqrt{R+d/r},\qquad
 \|a^\dagger/\sqrt r\|\leq\sqrt{R+(d+1)/r},
\]
where these are norms as maps from the indicated subspace
into the full Hilbert space.  The scaled spin operators
have bounds independent of \(r\), since the spin is \(r/2\).
Successive application therefore bounds every fixed word,
and its expectation in \(\varrho_r(t)\), uniformly in \(r,t\).
Excitation conservation makes the same bounds available
at all times.  No oscillator compression is inserted into
the equations of motion.

Every coordinate commutator is \(1/r\) times a constant
or a linear spin coordinate: for example
\[
 [X_1,X_2]=i/(2r),\qquad [X_3,X_4]=iX_5/r,
\]
with cyclic versions for the spin and zero mixed
oscillator--spin commutators.
Exchanging adjacent factors in any fixed word hence changes
its restricted norm by \(O(1/r)\).  In particular products,
adjoints, and reorderings of fixed polynomials agree with
their commutative counterparts to this order on the
excitation-bounded states.

Let \(F\) be the real polynomial vector field defined by
V59's classical equations.  Their exact operator equations
give \(\dot X_j=\operatorname{Sym}_r(F_j)\): quadratic
products there are between commuting oscillator and
spin coordinates.  Differentiating a fixed word by the
product rule and using the adjacent-exchange estimate
therefore gives, for every fixed polynomial \(p\),
\[
 \frac{d}{dt}\operatorname{Tr}
      (\varrho_r(t)\operatorname{Sym}_r(p))
 =\operatorname{Tr}
      (\varrho_r(t)\operatorname{Sym}_r(F\cdot\nabla p))
      +e_{r,p}(t),\qquad
 |e_{r,p}(t)|\leq C_{p,R}/r .
\]
Here the constant may depend on the fixed Hamiltonian
parameters and polynomial degree.  There are only
finitely many adjacent exchanges for each polynomial,
and the word bounds control all the resulting terms.
This proves the identity with its error, rather than
assuming a closed expectation equation.

\subsection{Every subsequential moment limit is a probability measure}

The word bounds and generator identity give uniform
boundedness and equicontinuity of every polynomial moment
on \([-T,T]\).  Arzela--Ascoli and a diagonal subsequence
over the countable monomials produce uniform limits
\(L_t(p)\).  They are normalized, and positive on squares:
\[
 L_t(1)=1,\qquad L_t(\overline p p)\geq0 .
\]
Indeed the product-reordering estimate identifies the
limit with that of the nonnegative expectation
\(\operatorname{Sym}_r(p)^*\operatorname{Sym}_r(p)\).

Here is a bounded-coordinate argument for representation,
including the support needed below.  Choose one fixed
constant \(M>\sqrt{R+1}+1\).  Applied to the vectors
\(\operatorname{Sym}_r(p)\psi\), which change the
excitation bound by only the fixed degree of \(p\),
the coordinate estimates give in the limit
\[
 L_t(x_j^2|p|^2)\leq M^2 L_t(|p|^2),\qquad j=1,\ldots,5 .
\]
For mixed density matrices the same inequality follows
by summing their vector decompositions.  Form the Hilbert
space completion of polynomials with inner product
\(L_t(\overline p q)\), quotienting its zero vectors.
Multiplication by each \(x_j\) is well defined, bounded
by \(M\), and symmetric, hence self-adjoint.  These five
bounded operators commute.  Their joint spectral theorem,
evaluated in the vector represented by the constant one,
gives a probability measure \(\mu_t\) on \([-M,M]^5\)
whose polynomial integrals are \(L_t\).
Polynomials are uniformly dense in continuous functions
on this box, so this measure is unique.

The exact Casimir and excitation identities are
\[
 X_3^2+X_4^2+X_5^2=\tfrac14+\frac1{2r},\qquad
 X_1^2+X_2^2+X_5+\tfrac12=\mathcal Q/r+\frac1{2r}.
\]
The first is the spin-\(r/2\) Casimir divided by \(r^2\);
the second follows from the oscillator commutator.
Apply each identity between
\(\operatorname{Sym}_r(p)^*\) and \(\operatorname{Sym}_r(p)\).
The excitation of this tested vector is at most
\(\lfloor Rr\rfloor+\deg p\) and is nonnegative.
Passing to the limit with the reordering estimate yields
\[
 \begin{aligned}
 L_t\big(|p|^2(x_3^2+x_4^2+x_5^2-\tfrac14)\big)&=0,\\
 L_t(|p|^2\rho)&\geq0,\qquad
 L_t(|p|^2(R-\rho))\geq0 .
 \end{aligned}
\]
These imply support in \(K_R\).  To see this directly,
approximate continuous localized test functions on the
box by polynomials.  If a displayed inequality failed
pointwise on a set of positive measure, a test supported
where its factor is strictly negative would contradict
the integrated inequality.  The equality forces its
factor to vanish almost everywhere, by the same argument
for both signs.  Thus \(\mu_t(K_R)=1\).

\subsection{The limiting equation and uniqueness}

Integrating the generator identity and passing to the
uniform moment limits gives
\[
 \int p\,d\mu_t-\int p\,d\mu_0
       =\int_0^t\int F\cdot\nabla p\,d\mu_s\,ds .
\]
Weak continuity of \(\mu_t\) follows from continuity of
moments and uniform polynomial approximation on the
common compact box.  The equation extends to smooth
tests and smooth time-dependent tests on a neighborhood
of \(K_R\).  One justification is simultaneous uniform
approximation of a function and its first derivatives
on a containing box by polynomial approximants; the
multivariable Bernstein approximants have this property
for smooth functions.  Their derivative formulas are
finite differences of the sampled function, which
converge uniformly to its derivatives.  The bounded
support and bounded \(F\) permit passage to the integrals.

For completeness, choose a smooth compactly supported
extension of \(F\) that agrees with it on a neighborhood
of \(K_R\).  Its global smooth flow agrees with \(\Phi_t\)
on \(K_R\), because that compact set is invariant.
For a smooth terminal test \(f\) and fixed terminal
time \(t_1\), use
\[
 \phi(s,x)=f(\Phi_{t_1-s}(x))
\]
near \(K_R\), defined outside by the extended flow.
It satisfies \(\partial_s\phi+F\cdot\nabla\phi=0\)
on \(K_R\).  The time-dependent weak equation gives
\[
 \int f\,d\mu_{t_1}=\int f(\Phi_{t_1}(x))\,d\mu_0(x).
\]
Smooth-test uniqueness on compact support implies
\(\mu_t=(\Phi_t)_*\mu_0\).  The same argument applies
on negative time intervals.

Every convergent subsequence has therefore produced the
same moment curves.  If the full sequence failed to
converge uniformly for a particular polynomial, an
equicontinuous subsequence separated from that curve
would have a further convergent subsequence, contradicting
this uniqueness.  This finishes the proof of
Theorem \ref{thm:v20v61-transport}.

\subsection{Application to all-parent number states}

\begin{corollary}[Phase-invariant initial data are covered]
\label{cor:v20v61-number}
Suppose \(q_r/r\to\rho_*\geq0\) and choose \(R>\rho_*\).
For sufficiently large \(r\), start with \(q_r\) parent
quanta and the symmetric empty-pair state.  The initial
measure in the theorem is the uniform phase measure
\[
 \alpha=\sqrt{\rho_*}e^{i\theta},\quad
 s=0,\quad z=-\tfrac12,\qquad
 \theta\ \hbox{uniform on }[0,2\pi).
\]
All fixed polynomial moments converge to its transported
moments, uniformly on bounded physical time intervals.
In particular the quantum pair density converges to the
pair density of the classical solution from
\((\sqrt{\rho_*},0,-1/2)\).
\end{corollary}
\begin{proof}
For oscillator monomials in normal order,
\[
 \left\langle
 \frac{(a^\dagger)^k a^\ell}{r^{(k+\ell)/2}}
 \right\rangle
 =\begin{cases}(q_r)_k/r^k,&k=\ell,\\0,&k\ne\ell,\end{cases}
\]
where the falling factorial is zero if \(k>q_r\).
These tend to \(\rho_*^k\) in the equal-power case,
exactly the circle-measure moments.  Reordering to
symmetric form costs \(O(1/r)\) at fixed degree by
the established word bounds.

In the empty-pair state \(X_5=-1/2\) exactly and
\(\langle X_3^2\rangle=\langle X_4^2\rangle=1/(4r)\).
Every fixed word containing a transverse spin coordinate
has vanishing limiting expectation: reorder one such
factor to the right, at cost \(O(1/r)\), and apply
Cauchy--Schwarz with its \(O(r^{-1/2})\) norm on the
empty state and the uniform bound on the remaining word.
The other spin moments equal powers of \(-1/2\).
Initial oscillator--spin independence completes the
initial moment calculation.  The theorem applies.
Finally phase equivariance from V59 makes \(z(t)\)
independent of the initial phase, proving the
pair-density assertion.
\end{proof}

\begin{assessment}[A proved collective state limit, with its scope fixed]
\label{ass:v20v61}
The phase-space transport proposed in V59 is now proved
for excitation-bounded symmetric collective states with
convergent initial moments.  It retains Pauli saturation
and covers number states despite their vanishing field
means.  The proof gives neither a quantitative rate nor
uniformity on growing kinetic times.  Equal pair energies,
one collective parent mode, and permutation symmetry
still exclude the dispersive finite-density SM problem.
No limiting stress tensor or Einstein continuum follows
from this collective transport theorem alone.
\end{assessment}

\begin{decision}[Restore more than one pair energy]
\label{dec:v20v61-next}
Replace the single symmetric pair group by finitely
many groups with distinct pair energies and couplings,
each of positive limiting size fraction.  Prove the
corresponding coupled-spin transport and retain each
group's saturation constraint.  This tests the first
energy-resolution step before any passage to a
continuous pair spectrum or a kinetic time scale.
\end{decision}

\begin{record}[V61 append record]
\label{rec:v20v61}
The immediate predecessor is the V60 companion audit.
Its complete prefix is preserved and checked by the
installer.  The convergence argument uses uniform word
bounds, positive limiting moments, the joint spectral
theorem, and uniqueness of compactly supported classical
transport.  Numerical samples are not used as proof.
Structural validation checks source preservation,
environments, delimiters, labels, and references; it
is not LaTeX compilation.  Only the immediate active
SM predecessor is replaced.  The next audit is V65.
\end{record}


\section{V62: finitely many pair energies and coupled-spin transport}
\label{sec:v20v62}

V61 proved collective transport with a single pair energy.
This section restores finitely many energies and couplings
while preserving permutation symmetry only within each group.
The groups share one parent oscillator and are not independent.
We track their size fractions explicitly, prove the corresponding
transport statement, and derive a memory equation whose
saturation factor remains dynamical.

\subsection{Microscopic grouping and normalization}

Fix \(m\geq1\).  Group \(\nu\) contains \(r_\nu\) disjoint
fermion pairs with common energy \(\omega_\nu\) and real
coupling coefficient \(G_\nu\).  Let
\[
 r=\sum_{\nu=1}^m r_\nu,\qquad
 \theta_{\nu,r}=r_\nu/r\longrightarrow\theta_\nu>0,
 \qquad\sum_{\nu=1}^m\theta_\nu=1 .
\]
The number \(m\), the energies, and the \(G_\nu\) are fixed
as \(r\to\infty\).  On the tensor product of the symmetric
spin-\(r_\nu/2\) representations and the parent oscillator set
\[
 H_r=\Omega a^\dagger a+
 \sum_\nu\omega_\nu(S_{\nu,z}+r_\nu/2)
 +\frac1{\sqrt r}\sum_\nu G_\nu
          (aS_{\nu,+}+a^\dagger S_{\nu,-}).
\]
The normalization is by total \(r\), not by \(r_\nu\).
The conserved excitation is
\[
 \mathcal Q=a^\dagger a+\sum_\nu(S_{\nu,z}+r_\nu/2).
\]
In the fixed-\(\mathcal Q=q\) sector, let \(k_\nu\) be
the occupied-pair number in group \(\nu\).
The normalized basis has
\(0\leq k_\nu\leq r_\nu\) and \(\sum_\nu k_\nu\leq q\).
Its diagonal energy and transition coefficients are
\[
 d_{\boldsymbol k}=\Omega(q-\sum_\nu k_\nu)
                            +\sum_\nu\omega_\nu k_\nu,
\]
\[
 b_{\boldsymbol k,\nu}=\frac{G_\nu}{\sqrt r}
 \sqrt{(q-\sum_\eta k_\eta)(k_\nu+1)(r_\nu-k_\nu)} .
\]
These follow by applying V59's normalized subset-counting
identity within group \(\nu\), with the common parent
annihilation factor.  They prove exact invariance of
this finite grouped sector and retain separate saturation
at every \(k_\nu=r_\nu\).

\subsection{Exact scaled equations and their classical candidate}

Define
\[
 \widehat\alpha=a/\sqrt r,\qquad
 \widehat s_\nu=S_{\nu,-}/r_\nu,\qquad
 \widehat z_\nu=S_{\nu,z}/r_\nu .
\]
The same commutator calculation as V59, now with the
declared group normalization, gives exactly
\[
 \begin{aligned}
 i\dot{\widehat\alpha}
    &=\Omega\widehat\alpha+
             \sum_\nu\theta_{\nu,r}G_\nu\widehat s_\nu,\\
 i\dot{\widehat s}_\nu
    &=\omega_\nu\widehat s_\nu
                         -2G_\nu\widehat\alpha\widehat z_\nu,\\
 i\dot{\widehat z}_\nu
    &=G_\nu(\widehat\alpha\widehat s_\nu^\dagger
                         -\widehat\alpha^\dagger\widehat s_\nu).
 \end{aligned}
\]
For example the first coupling has coefficient
\(G_\nu S_{\nu,-}/r=\theta_{\nu,r}G_\nu\widehat s_\nu\),
whereas the spin equation's \(r_\nu\) cancels out.
Missing this distinction would change the limiting
excitation and energy balances.

The limiting classical system is therefore
\[
 i\dot\alpha=\Omega\alpha+\sum_\nu\theta_\nu G_\nu s_\nu,
 \qquad
 i\dot s_\nu=\omega_\nu s_\nu-2G_\nu\alpha z_\nu,
 \qquad
 \dot z_\nu=2G_\nu\operatorname{Im}(\alpha\overline s_\nu).
\]

\begin{theorem}[Groupwise saturation and global flow]
\label{thm:v20v62-flow}
The classical system preserves
\[
 \rho=|\alpha|^2+\sum_\nu\theta_\nu(z_\nu+\tfrac12),
 \qquad \ell_\nu^2=|s_\nu|^2+z_\nu^2,
\]
\[
 h=\Omega|\alpha|^2+
       \sum_\nu\theta_\nu\omega_\nu(z_\nu+\tfrac12)
       +2\sum_\nu\theta_\nu G_\nu
                           \operatorname{Re}(\alpha\overline s_\nu).
\]
If every \(\ell_\nu\leq1/2\), solutions exist uniquely
for all real times, with \(0\leq z_\nu+1/2\leq1\)
and \(0\leq|\alpha|^2\leq\rho\).
\end{theorem}
\begin{proof}
Differentiation of \(|\alpha|^2\) gives
\(-2\sum_\nu\theta_\nu G_\nu
\operatorname{Im}(\alpha\overline s_\nu)\),
which cancels the weighted sum of the \(z_\nu\) derivatives.
Each spin-length cancellation is local to its group,
as in V59.

For the energy, put \(u_\nu=\alpha\overline s_\nu\).
Then
\[
 \dot u_\nu=i(\omega_\nu-\Omega)u_\nu
       -i\sum_\eta\theta_\eta G_\eta s_\eta\overline s_\nu
       -2iG_\nu|\alpha|^2z_\nu .
\]
The first terms in the weighted real part cancel the
free-energy derivatives.  The last terms are purely
imaginary.  The middle terms sum, after multiplication
by \(\theta_\nu G_\nu\), to
\[
 -i\left|\sum_\nu\theta_\nu G_\nu s_\nu\right|^2,
\]
whose real part is zero.  Thus \(h\) is conserved;
this cancellation includes interactions between groups.
Spin-length bounds and conserved \(\rho\) bound every
coordinate.  The polynomial vector field is locally
Lipschitz, so boundedness gives global continuation and
uniqueness, exactly as in V59.
\end{proof}

\subsection{A fixed-time grouped transport theorem}

Use real oscillator coordinates as in V61, and for
each group the three real coordinates of
\((s_\nu,z_\nu)\), normalized by \(r_\nu\).
Let \(\operatorname{Sym}_{\boldsymbol r}(p)\) denote
complete symmetric ordering of polynomials in these
\(2+3m\) real coordinates.  Define the compact set
\[
 K_R^{(m)}=\{(\alpha,(s_\nu,z_\nu)_\nu):
       |s_\nu|^2+z_\nu^2=\tfrac14\ \hbox{for all }\nu,\
                      0\leq\rho\leq R\}.
\]
Let \(\Phi_t^{(m)}\) be the flow just proved.

\begin{theorem}[Quantum transport with finitely many energies]
\label{thm:v20v62-transport}
Let \(\varrho_r\) be density matrices in the grouped
symmetric representation, supported in
\(\mathcal Q\leq\lfloor Rr\rfloor\).  Suppose all initial
symmetrically ordered polynomial moments converge to
those of a probability measure \(\mu_0\) on \(K_R^{(m)}\).
Then for every fixed polynomial \(p\) and \(T<\infty\),
\[
 \sup_{|t|\leq T}\left|
 \operatorname{Tr}\!\left(e^{-itH_r}\varrho_re^{itH_r}
                         \operatorname{Sym}_{\boldsymbol r}(p)\right)
   -\int p(\Phi_t^{(m)}(x))\,d\mu_0(x)\right|
                          \longrightarrow0 .
\]
\end{theorem}
\begin{proof}
We give the changes to V61's proof explicitly.
A fixed coordinate word changes excitation by at most
its degree.  Oscillator word bounds therefore still
use \(Rr+O(1)\), while each normalized spin coordinate
is uniformly bounded.  Different groups commute.
The within-group commutators are \(1/r_\nu\) times
linear spin coordinates.  Since every limiting fraction
is positive, all \(r_\nu\to\infty\); fixed-word
reordering errors are bounded by
\(C_p\max_\nu r_\nu^{-1}\).

The exact equations have vector field \(F_{\boldsymbol\theta_r}\).
Replacing its fractions by the limiting ones adds a
fixed-word error at most
\(C_p\sum_\nu|\theta_{\nu,r}-\theta_\nu|\).
Consequently the generator identity has error bounded by
\[
 C_{p,R}\left(\frac1r+\max_\nu\frac1{r_\nu}
                    +\sum_\nu|\theta_{\nu,r}-\theta_\nu|\right)
                         \longrightarrow0 .
\]
This gives moment equicontinuity and subsequential limits.
Positivity on squares, bounded coordinate multiplication,
and the joint spectral representation follow word for
word from V61 in \(2+3m\) dimensions.

For support, use separately the exact group Casimirs
\[
 X_{\nu,1}^2+X_{\nu,2}^2+X_{\nu,3}^2
                   =\tfrac14+\frac1{2r_\nu}.
\]
The excitation relation is
\[
 X_1^2+X_2^2+\sum_\nu\theta_{\nu,r}
                      (X_{\nu,3}+\tfrac12)
                   =\mathcal Q/r+\frac1{2r}.
\]
Sandwich these identities between fixed polynomial
operators and pass to the limit.  The group fractions
converge and all tested spin coordinates are bounded,
so the support inequalities have precisely the limiting
\(\rho\).  The resulting measure is supported in
\(K_R^{(m)}\).

The integrated generator identity is the weak transport
equation for \(F_{\boldsymbol\theta}\).  The preceding
global-flow theorem, compact support, smooth test
approximation, and backward characteristic tests give
uniqueness as in V61.  Hence every subsequential limit
is \((\Phi_t^{(m)})_*\mu_0\).  Equicontinuity then yields
uniform convergence of each polynomial moment on the
stated time interval.
\end{proof}

This proof fixes \(m\) before taking \(r\to\infty\).
Its constants and observables can depend on \(m\),
and its commutator estimate requires every group
population to tend to infinity.  It proves no uniform
simultaneous limit with infinitely many sparsely
populated energy groups.

For initially \(q_r\) parents and all groups empty,
with \(q_r/r\to\rho_*\) and \(R>\rho_*\), the initial
measure is again the uniform phase circle
\[
 \alpha=\sqrt{\rho_*}e^{i\theta},\quad
 s_\nu=0,\quad z_\nu=-1/2\quad\hbox{for every }\nu.
\]
The oscillator factorial-moment calculation of V61
is unchanged, and each empty group has transverse
variances \(1/(4r_\nu)\).  Independence initially
therefore proves the required joint moment convergence.
Phase equivariance makes every group density
\(z_\nu(t)+1/2\) independent of the initial phase.
The theorem proves convergence of these densities
and their fixed polynomial correlations; it does
not assume factorization of evolving quantum states.

\subsection{Exact classical memory with dynamic saturation}

Solving the linear-in-\(s_\nu\) equation at a given
\(\alpha,z_\nu\) trajectory gives
\[
 s_\nu(t)=e^{-i\omega_\nu t}s_\nu(0)
       +2iG_\nu\int_0^t e^{-i\omega_\nu(t-u)}
                                  \alpha(u)z_\nu(u)\,du.
\]
Substitution into the parent equation yields
\[
 \begin{split}
 \dot\alpha(t)={}&-i\Omega\alpha(t)
       -i\sum_\nu\theta_\nu G_\nu
                         e^{-i\omega_\nu t}s_\nu(0)\\
 &+2\sum_\nu\theta_\nu G_\nu^2
       \int_0^t e^{-i\omega_\nu(t-u)}
                                  \alpha(u)z_\nu(u)\,du .
 \end{split}
\]
These identities follow by the integrating factor
\(e^{i\omega_\nu t}\), so have no approximation error.
In the formal empty-pair linearization \(z_\nu=-1/2\),
the memory kernel is
\[
 K(t)=\sum_\nu\theta_\nu G_\nu^2e^{-i\omega_\nu t},
\]
with the negative convolution sign familiar from the
one-parent memory equation.  At finite density the
factors \(z_\nu(u)\) evolve and must remain inside the
integral.  Freezing them at \(-1/2\) is a linearization,
not a consequence of the transport theorem.

\begin{assessment}[Energy resolution without losing exclusion]
\label{ass:v20v62}
The fixed-time quantum-to-classical theorem now allows
finitely many pair energies and couplings with the
correct group weights.  The common parent still couples
all groups, and their exact classical memory retains
dynamic saturation.  This advances energy resolution
but leaves continuous spectra, kinetic times, ultraviolet
removal, spatially varying matter, and a conserved
Einstein source unproved.
\end{assessment}

\begin{decision}[Continuous parameter limit at fixed time]
\label{dec:v20v62-next}
Construct the classical continuum of spin variables
over a compact set of pair energies and couplings.
Prove well-posedness and a quadrature-stability estimate
under explicit initial regularity.  Combine it with
the present theorem only in a stated sequential order:
large populations in finitely many groups first, then
refinement of the group measure.  Do not claim a
simultaneous quantum limit without uniform estimates.
\end{decision}

\begin{record}[V62 append record]
\label{rec:v20v62}
The predecessor V61 has SHA--256
\texttt{9E366C5E882E7D3BCF7D7268735CE2F268ECB1DA19F42EE5A23BB4693D121EA2}.
Its complete prefix is preserved.  The new proof
tracks group fractions, commutator errors, Casimirs,
and transport uniqueness explicitly.  Structural
checks cover preservation, environments, delimiters,
labels, and references; they are not LaTeX compilation.
Only the immediate active SM predecessor is replaced.
The next companion audit is V65.
\end{record}


\section{V63: a compact continuous pair spectrum and quadrature stability}
\label{sec:v20v63}

V62 fixes the number of populated energy groups before taking
their quantum large-population limit.  We now refine the
resulting classical group measure.  The construction is over
a compact parameter set and retains a spin at each parameter.
It is a continuum of pair energies, not a spatial quantum
Standard Model or an ultraviolet limit.

\subsection{The continuum equations}

Let \(K\subset\mathbb R^2\) be compact, with parameters
\(\xi=(\omega,g)\), and let \(\nu\) be a Borel probability
measure on \(K\).  Write \(g_*=\sup_K|g|\).
For a common complex parent variable \(\alpha(t)\) and
continuous spin fields \(s(t,\xi)\in\mathbb C\),
\(z(t,\xi)\in\mathbb R\), consider
\[
 \begin{aligned}
 i\dot\alpha&=\Omega\alpha+\int_K g\,s(t,\xi)\,d\nu(\xi),\\
 i\partial_t s(t,\xi)&=\omega s(t,\xi)-2g\alpha(t)z(t,\xi),\\
 \partial_t z(t,\xi)&=2g\operatorname{Im}
                                   (\alpha(t)\overline{s(t,\xi)}).
 \end{aligned}
\]
The initial data are \(\alpha_0\in\mathbb C\) and continuous
\(s_0,z_0\) satisfying \(|s_0(\xi)|^2+z_0(\xi)^2\leq1/4\)
at every \(\xi\).  Set \(w=(\operatorname{Re}s,\operatorname{Im}s,z)\).

\begin{theorem}[Global continuum evolution]
\label{thm:v20v63-global}
There is a unique global solution, continuously differentiable
in time with values in \(\mathbb C\times C(K;\mathbb R^3)\).
It preserves, pointwise in \(\xi\), \(|s|^2+z^2\), and preserves
\[
 \rho_\nu=|\alpha|^2+\int_K(z+\tfrac12)\,d\nu,
\]
\[
 h_\nu=\Omega|\alpha|^2+
          \int_K\{\omega(z+\tfrac12)
                         +2g\operatorname{Re}(\alpha\overline s)\}\,d\nu .
\]
In particular \(0\leq z+1/2\leq1\) and
\[
 |\alpha(t)|\leq A_*:=\sqrt{|\alpha_0|^2+1}
\]
for every real \(t\), uniformly over the choice of \(\nu\).
\end{theorem}
\begin{proof}
The right side is a locally Lipschitz map on the stated
Banach space: multiplication is continuous in the supremum
norm, parameter coefficients are bounded, and integration
against a probability measure is bounded.  Picard iteration
therefore gives a unique local solution.  The spin-length
and excitation cancellations of V62 hold pointwise and
under the integral, respectively.  The energy calculation
is also identical after replacing finite sums by integrals:
the cross contribution is
\(-i|\int_K g s\,d\nu|^2\), with zero real part.
Continuity on compact \(K\) justifies differentiation
under the integral on every local time interval.
These conserved quantities give
\(|w|\leq1/2\) and \(|\alpha|^2\leq\rho_\nu(0)
\leq|\alpha_0|^2+1\).  The solution stays in a bounded
set on which the polynomial vector field is Lipschitz.
Continuation in both time directions proves global existence.
\end{proof}

\subsection{Regularity in the energy and coupling parameters}

The spin equation is a rotation equation
\[
 \partial_t w=M(\alpha,\xi)w,\qquad
 M(\alpha,\xi)=
 \begin{pmatrix}
 0&\omega&-2g\operatorname{Im}\alpha\\
 -\omega&0&2g\operatorname{Re}\alpha\\
 2g\operatorname{Im}\alpha&-2g\operatorname{Re}\alpha&0
 \end{pmatrix}.
\]
It is real skew-symmetric.  Its propagators are orthogonal
even when \(\alpha\) varies with time.

\begin{lemma}[Parameter Lipschitz bound]
\label{lem:v20v63-lipschitz}
If \(w_0\) is Lipschitz on \(K\) with constant \(L_0\)
for Euclidean distance, then
\[
 \operatorname{Lip}_K w(t,\cdot)
              \leq L_0+(\tfrac12+A_*)|t|.
\]
Consequently for \(|t|\leq T\),
\[
 \operatorname{Lip}_K(g\,s(t,\cdot))
 \leq B_T:=\tfrac12+
              g_*[L_0+(\tfrac12+A_*)T].
\]
\end{lemma}
\begin{proof}
Compare the spin solutions at \(\xi=(\omega,g)\) and
\(\eta=(\widetilde\omega,\widetilde g)\), driven by the
same \(\alpha(t)\).  Variation of constants with the
orthogonal propagator at \(\xi\) bounds their difference by
\[
 |w_0(\xi)-w_0(\eta)|+
 \int_0^{|t|}\tfrac12
       (|\omega-\widetilde\omega|+
                           2A_*|g-\widetilde g|)\,du .
\]
The matrix-norm bound follows by separating the frequency
and coupling entries of \(M\); the multiplying spin has
norm at most \(1/2\).  Each coordinate difference is
bounded by \(|\xi-\eta|\), proving the first assertion.
For the second, split
\(g s(\xi)-\widetilde g s(\eta)\) into a coupling
difference and a spin difference, using \(|s|\leq1/2\).
\end{proof}

\subsection{Stability under refinement of the measure}

For probability measures on \(K\), define
\[
 W_1(\nu,\widetilde\nu)
  =\inf_{\pi}\int_{K\times K}|\xi-\eta|\,d\pi(\xi,\eta),
\]
where \(\pi\) ranges over their couplings.  For any
complex Lipschitz function \(f\),
\[
 \left|\int f\,d\nu-\int f\,d\widetilde\nu\right|
                    \leq\operatorname{Lip}(f)W_1(\nu,\widetilde\nu).
\]
Indeed integrate \(f(\xi)-f(\eta)\) against a coupling,
use the Lipschitz bound, and take the infimum.

\begin{theorem}[Uniform fixed-time quadrature estimate]
\label{thm:v20v63-stability}
Let the two continuum equations use measures
\(\nu,\widetilde\nu\) and the same initial
\(\alpha_0,w_0\), with \(w_0\) Lipschitz as above.
Their solutions satisfy
\[
 \sup_{|t|\leq T}
 \left\{|\alpha_\nu(t)-\alpha_{\widetilde\nu}(t)|
       +\sup_{\xi\in K}|w_\nu(t,\xi)-w_{\widetilde\nu}(t,\xi)|\right\}
 \leq B_T T e^{g_*T}W_1(\nu,\widetilde\nu).
\]
The constants are independent of the number of atoms of either
measure.
\end{theorem}
\begin{proof}
It suffices to prove the forward estimate; reversing time
gives the same bounds.  Let \(a(t)\) be the parent difference
in modulus and \(b(t)\) the supremum spin difference.
Use variation of constants for the parent with its scalar
free rotation.  Split the source difference into a field
difference integrated against \(\widetilde\nu\) and the
fixed field \(g s_\nu\) integrated against the measure
difference.  The preceding lemma gives
\[
 a(t)\leq g_*\int_0^t b(u)\,du+
                           B_Tt\,W_1(\nu,\widetilde\nu).
\]
For the spins at the same \(\xi\), the difference of
the two matrices \(M\) has norm at most \(2|g|a(t)\).
Multiply by the spin norm at most \(1/2\), and use the
orthogonal propagator to obtain
\[
 b(t)\leq g_*\int_0^t a(u)\,du.
\]
Adding and applying the integral Gronwall bound on
\([0,T]\) proves the assertion.
\end{proof}

For an atomic measure
\(\nu_m=\sum_{\ell=1}^{m}\theta_\ell\delta_{\xi_\ell}\)
with positive weights, the restrictions of the field
solution to its nodes solve exactly V62's grouped
classical system.  Defining the spin field at other
parameters with the same parent solution is only a
convenient continuous extension for comparison; those
extra spins do not enter the atomic parent source.

Atomic refinements with \(W_1(\nu_m,\nu)\to0\) always
exist on compact \(K\).  Partition \(K\) into finitely
many Borel sets of diameter at most \(\varepsilon\),
discard sets of zero measure, select a node in each
remaining set, and use its measure as the weight.
Transporting each set to its node yields
\(W_1(\nu_m,\nu)\leq\varepsilon\).
The stability theorem therefore supplies convergence
of the grouped classical dynamics to the continuum
dynamics, uniformly on every fixed time interval.

\subsection{A precise sequential quantum consequence}

\begin{corollary}[Initially empty pairs and macroscopic parent number]
\label{cor:v20v63-sequential}
Fix \(\rho_*\geq0\) and atomic refinements \(\nu_m\to\nu\)
as above.  For each fixed \(m\), take V62's grouped
quantum systems with \(r_\ell/r\to\theta_\ell\),
initial parent number \(q_r/r\to\rho_*\), and all pair
groups empty.  Let \(\alpha_\nu,z_\nu\) solve the
continuum equations with
\[
 \alpha_\nu(0)=\sqrt{\rho_*},\quad
 s_\nu(0,\xi)=0,\quad z_\nu(0,\xi)=-1/2.
\]
For every \(T<\infty\), the quantum parent-density
expectation converges in the sequential order
\(r\to\infty\) first, \(m\to\infty\) second, uniformly
for \(|t|\leq T\), to \(|\alpha_\nu(t)|^2\).
For every Lipschitz function \(\psi\) on \(K\), the
same sequential convergence holds for the observable
\[
 \frac1r\sum_{\ell=1}^m\psi(\xi_\ell)
                  (S_{\ell,z}+r_\ell/2)
\]
toward
\[
 \int_K\psi(\xi)(z_\nu(t,\xi)+\tfrac12)\,d\nu(\xi).
\]
\end{corollary}
\begin{proof}
At each fixed \(m\), V62 proves uniform fixed-time
convergence of these polynomial expectations to the
atomic classical solution.  Its initial phase measure
does not alter the parent or group densities, by phase
equivariance.  The quantum fractions in the smeared
observable converge to their declared positive weights.
The stability theorem then controls the difference
between the atomic and continuum solutions.

For the smeared observable, split that difference
into a field error and quadrature of the continuum
field.  The first is bounded by
\(\|\psi\|_\infty\sup_\xi|z_{\nu_m}-z_\nu|\).
The second is bounded by
\[
 \operatorname{Lip}\{\psi(z_\nu+\tfrac12)\}
                               W_1(\nu_m,\nu),
\]
uniformly on \([-T,T]\), using the parameter Lipschitz
bound.  Both tend to zero.  The parent-density
difference is at most \(2A_*\) times the parent
amplitude difference.  Taking the two limits in the
stated order proves the result.
\end{proof}

The sequential assertion can equivalently be written
as a limit of the inner \(\limsup_{r\to\infty}\) of
the uniform expectation error, followed by \(m\to\infty\).
It does not allow an arbitrary simultaneous choice
\(m=m(r)\): V62 has no rate uniform in growing group
number.  It also begins with grouped quantum models,
not a proof that grouping the actual dispersive
finite-box Hamiltonian has a negligible error.

\begin{assessment}[A bounded-spectrum continuum, with exclusions retained]
\label{ass:v20v63}
There is now a globally well-posed classical continuum
over compact pair parameters, with pointwise Pauli
saturation, conserved excitation and energy, and an
explicit quadrature estimate.  It is the sequential
limit of the specified grouped quantum density
observables at fixed times.  Spatial locality,
ultraviolet removal, a simultaneous ungrouped quantum
limit, kinetic times, and a stress-tensor/Einstein
handoff are separate unresolved requirements.
\end{assessment}

\begin{decision}[Match the parameter measure to canonical box normalization]
\label{dec:v20v63-next}
At a fixed momentum cutoff, identify the parameter
measure, the group coupling, and the physical-volume
density corresponding to V43's canonical Yukawa box
vertex.  Track the conversion between normalization
per pair mode and per spatial volume.  Determine
which constants grow with the cutoff before attempting
an ultraviolet statement.  The next companion audit
remains V65.
\end{decision}

\begin{record}[V63 append record]
\label{rec:v20v63}
The predecessor V62 has SHA--256
\texttt{4BBC3AC3649C8753E2664AF138822DF894ACEDBF5FF6FF9A93125FECFF175882}.
Its complete prefix is retained.  The new estimates
use Banach-space local existence, conserved bounds,
orthogonal spin propagators, and an explicit transport
coupling estimate.  Structural checks cover source
preservation, environments, delimiters, labels, and
references; they are not LaTeX compilation.
Only the immediate active SM predecessor is replaced.
The next companion audit is V65.
\end{record}


\section{V64: physical-volume normalization and the ultraviolet source obstruction}
\label{sec:v20v64}

The probability measure in V63 counts fractions of pair modes.
It is not the density of states per spatial volume.  We now
match the two conventions at a fixed momentum cutoff and
derive the cutoff dependence of the parent source.  A
finite-energy counterexample then shows why the compact-set
existence proof cannot simply be declared uniform in the cutoff.

\subsection{Canonical box normalization}

Fix \(m>0\), vertex parameter \(M>0\), and real \(y\ne0\).
For the rest-parent channel set
\[
 E(p)=\sqrt{m^2+p^2},\qquad \omega(p)=2E(p),\qquad
 h(p)=\frac{yp}{\sqrt{2M}E(p)}.
\]
A fixed pair-phase convention has been chosen so that the
box vertex is \(g_{L,p}=h(p)/\sqrt L\).
Its magnitude agrees with V43 and V56.  Retain all signed
momenta \(p=2\pi k/L\) with \(|p|\leq\Lambda\), including
the uncoupled zero mode.  Then
\[
 r_{L,\Lambda}=2\left\lfloor\frac{\Lambda L}{2\pi}\right\rfloor+1,
 \qquad
 \frac{r_{L,\Lambda}}L\longrightarrow
                         d_\Lambda=\frac{\Lambda}{\pi}.
\]
The empirical probability measure of these momenta tends,
at fixed \(\Lambda>0\), to
\[
 d\nu_\Lambda(p)=\frac{dp}{2\Lambda},\qquad |p|\leq\Lambda .
\]
These are ordinary Riemann-sum statements.

To express the vertex in V62's form \(G/\sqrt r\),
the finite-box group coefficient must be
\[
 G_{L,\Lambda}(p)=\sqrt{r_{L,\Lambda}/L}\,h(p)
       \longrightarrow G_\Lambda(p)=\sqrt{d_\Lambda}\,h(p).
\]
Thus V63's parameter measure is the push-forward of
\(\nu_\Lambda\) by
\(p\mapsto(\omega(p),G_\Lambda(p))\).
The equality of these normalizations is exact at the
coefficient level; a dynamical error estimate for replacing
the ungrouped box by energy groups is still a separate task.

\begin{proposition}[Equations and conservation per spatial volume]
\label{prop:v20v64-volume}
Write \(\beta=\sqrt{d_\Lambda}\alpha\) for the
parent amplitude per square root of spatial volume.
The cutoff classical equations become
\[
 \begin{aligned}
 i\dot\beta&=\Omega_\Lambda\beta+
           \int_{-\Lambda}^{\Lambda}h(p)s(t,p)\,\frac{dp}{2\pi},\\
 i\partial_t s&=\omega(p)s-2h(p)\beta z,\\
 \partial_t z&=2h(p)\operatorname{Im}(\beta\overline s).
 \end{aligned}
\]
Here \(\Omega_\Lambda\) is any real parent coefficient,
including a declared cutoff counterterm.
They preserve the excitation density and energy density
\[
 \mathcal N_\Lambda=|\beta|^2+
          \int_{-\Lambda}^{\Lambda}(z+\tfrac12)\,\frac{dp}{2\pi},
\]
\[
 \mathcal E_\Lambda=\Omega_\Lambda|\beta|^2+
  \int_{-\Lambda}^{\Lambda}
  \{\omega(p)(z+\tfrac12)+
              2h(p)\operatorname{Re}(\beta\overline s)\}\,\frac{dp}{2\pi}.
\]
\end{proposition}
\begin{proof}
Substitute \(G_\Lambda=\sqrt{d_\Lambda}h\) and
\(\alpha=\beta/\sqrt{d_\Lambda}\) into V63's equations.
The parent source uses
\(d_\Lambda\,dp/(2\Lambda)=dp/(2\pi)\), and the spin
coupling uses \(G_\Lambda\alpha=h\beta\).
Multiply V63's conserved quantities by \(d_\Lambda\)
to obtain the displayed densities and their conservation.
The same formulas follow by direct differentiation.
\end{proof}

At the box level \(\widehat\beta=a/\sqrt L\), so
\(\langle a^\dagger a\rangle/L
=(r_{L,\Lambda}/L)\langle\widehat\alpha^\dagger
\widehat\alpha\rangle\).  Fixed density per volume and
fixed excitation per pair mode therefore differ as the
cutoff varies: \(d_\Lambda\to\infty\).

\subsection{Two explicit ultraviolet weights}

Define
\[
 B_\Lambda=\int_{-\Lambda}^{\Lambda}h(p)^2\,\frac{dp}{2\pi},
 \qquad
 I_\Lambda=\int_{-\Lambda}^{\Lambda}
                       \frac{h(p)^2}{\omega(p)}\,\frac{dp}{2\pi}.
\]
Elementary integration gives
\[
 B_\Lambda=\frac{y^2}{2\pi M}
                 [\Lambda-m\arctan(\Lambda/m)],
\]
\[
 I_\Lambda=\frac{y^2}{4\pi M}
 \left[\operatorname{arsinh}(\Lambda/m)
                 -\frac{\Lambda}{\sqrt{\Lambda^2+m^2}}\right].
\]
For the first, use \(p^2/(m^2+p^2)=1-m^2/(m^2+p^2)\).
For the second, differentiate
\(\operatorname{arsinh}(p/m)-p/\sqrt{p^2+m^2}\),
whose derivative is \(p^2/(p^2+m^2)^{3/2}\).
Consequently \(B_\Lambda\) grows linearly and
\[
 I_\Lambda=\frac{y^2}{4\pi M}
             [\log(2\Lambda/m)-1+o(1)] .
\]
The common-profile coupling bound of V63 also grows:
\(\sup|G_\Lambda|\) is of order \(\sqrt\Lambda\).
Thus its constants are not ultraviolet-uniform after
matching the physical normalization.

Let \(x=z+1/2\).  The spin constraint implies
\[
 0\leq x\leq1,\qquad |s|^2\leq x(1-x)\leq x.
\]
Hence excitation conservation gives
\(\int|s|^2\,dp/(2\pi)\leq\mathcal N_\Lambda\).
But Cauchy--Schwarz only bounds the source by
\[
 \left|\int h s\,\frac{dp}{2\pi}\right|
       \leq\sqrt{B_\Lambda\mathcal N_\Lambda}.
\]
If instead one controls the positive free pair energy
\(\mathcal E_{\rm pair}=\int\omega x\,dp/(2\pi)\), the
weighted estimate improves to
\[
 \left|\int h s\,\frac{dp}{2\pi}\right|
       \leq\sqrt{I_\Lambda\mathcal E_{\rm pair}}.
\]
Even this bound grows as \(\sqrt{\log\Lambda}\).
The next theorem shows that this latter growth reflects
actual allowed data, rather than only a loose estimate.

\begin{theorem}[Bounded energy does not uniformly bound the instantaneous source]
\label{thm:v20v64-source}
Fix \(\mathcal E_*>0\).  For all sufficiently large
\(\Lambda\), there are continuous initial spin fields
on \([-\Lambda,\Lambda]\) with maximal spin length,
parent amplitude \(\beta(0)=0\), and
\[
 0\leq\mathcal E_\Lambda(0)=\mathcal E_{\rm pair}(0)
          \leq\mathcal E_*,
 \qquad
 \mathcal N_\Lambda(0)\leq\mathcal E_*/(2m),
\]
but
\[
 |\dot\beta(0)|=\sqrt{\mathcal E_* I_\Lambda/2}
                                      \longrightarrow\infty .
\]
This holds for any choice of the real coefficient
\(\Omega_\Lambda\).
\end{theorem}
\begin{proof}
Set
\[
 c_\Lambda=\sqrt{\mathcal E_*/(2I_\Lambda)},\qquad
 s_\Lambda(p)=c_\Lambda h(p)/\omega(p).
\]
The function \(h/\omega\) is bounded on the real line,
and \(c_\Lambda\to0\).  Thus \(|s_\Lambda|\leq1/2\)
for large \(\Lambda\).  Define the continuous functions
\[
 x_\Lambda(p)=
       \frac{1-\sqrt{1-4s_\Lambda(p)^2}}2,\qquad
 z_\Lambda(p)=x_\Lambda(p)-\tfrac12 .
\]
They satisfy
\(|s_\Lambda|^2+z_\Lambda^2=1/4\) and
\[
 s_\Lambda^2\leq x_\Lambda\leq2s_\Lambda^2.
\]
It follows that
\[
 \int\omega x_\Lambda\,\frac{dp}{2\pi}
    \leq2c_\Lambda^2 I_\Lambda=\mathcal E_* .
\]
Since \(\omega\geq2m\), the excitation-density bound
also follows.  With \(\beta(0)=0\), both parent and
interaction energy vanish, whatever \(\Omega_\Lambda\).
The initial parent equation gives
\[
 |\dot\beta(0)|=
 \left|\int h s_\Lambda\,\frac{dp}{2\pi}\right|
      =c_\Lambda I_\Lambda
      =\sqrt{\mathcal E_*I_\Lambda/2}.
\]
This diverges logarithmically under the square root,
proving all assertions.
\end{proof}

These are cutoff-dependent initial data in a uniformly
energy-bounded class.  The result rules out a uniform
instantaneous-source bound from that energy bound alone.
It does not rule out evolution with an unbounded generator,
convergence for more restricted compatible initial data,
or a renormalized mild evolution.  In particular it
does not prove that a continuum theory fails to exist.

\subsection{Where the previous scalar subtraction reappears}

At subtraction point zero, below the pair threshold,
the canonical linear self-energy is
\[
 \Sigma_\Lambda(0)=
     \int_{-\Lambda}^{\Lambda}
       \frac{h(p)^2}{-\omega(p)}\,\frac{dp}{2\pi}
                       =-I_\Lambda .
\]
V47's one-parent subtraction therefore corresponds to
\(\Omega_\Lambda=\Omega_R+I_\Lambda\).
This identity matches the physical normalization exactly.
It does not, by itself, establish that the same prescription
constructs the nonlinear finite-density spin evolution.

At finite cutoff, define the dressed spin variable
\[
 \eta(p)=s(p)+\frac{h(p)}{\omega(p)}\beta .
\]
With the preceding counterterm, direct substitution yields
\[
 i\dot\beta=\Omega_R\beta+\int h\eta\,\frac{dp}{2\pi},
\]
\[
 i\partial_t\eta=
       \omega\eta-2h\beta x+\frac{h}{\omega}i\dot\beta,
                 \qquad x=z+\tfrac12 .
\]
The first equality cancels \(I_\Lambda\beta\) exactly.
For the second, use
\(i\partial_t s=\omega s+h\beta-2h\beta x\)
and differentiate the defining correction.
These are finite-cutoff identities, not uncut equations
until their domains and integrals have been controlled.
They isolate the nonlinear saturated forcing
\(-2h\beta x\) from the linear ultraviolet dressing.

\begin{assessment}[A concrete ultraviolet dependency]
\label{ass:v20v64}
The per-mode and per-volume normalizations now match the
canonical Yukawa vertex.  The source obstruction explains
why the compact-spectrum estimates cannot be extrapolated
using energy conservation alone.  The scalar counterterm
has the correct logarithmic coefficient, and the dressed
finite-cutoff identities identify what remains to estimate.
No nonlinear ultraviolet limit, spatial stress tensor,
or SM--Einstein continuum is proved.
\end{assessment}

\begin{decision}[Audit, then estimate the dressed mild equation]
\label{dec:v20v64-next}
Perform the required companion audit at V65.  Then
derive the finite-cutoff mild equations for
\((\beta,\eta,x)\) and determine which initial-data
norms make the nonlinear source cutoff-uniform.
Distinguish a genuine cancellation in the integrated
equation from an unavailable pointwise source bound.
Do not assume the one-parent norm-resolvent theorem
already supplies nonlinear well-posedness.
\end{decision}

\begin{record}[V64 append record]
\label{rec:v20v64}
The predecessor V63 has SHA--256
\texttt{31CA9D381E875C38A82D17CF4DE1FD23CFEF6F743F8814339314892B30CB69A1}.
Its complete prefix is retained.  The normalization,
integrals, bounded-energy counterexample, and dressed
identities have direct proofs.  Structural validation
checks source preservation, environments, delimiters,
labels, and references, not LaTeX compilation.
Only the immediate active SM predecessor is replaced.
The next version V65 is the companion audit.
\end{record}


\section{V65: companion audit before nonlinear ultraviolet evolution}
\label{sec:v20v65}

\begin{record}[Current companion snapshots]
\label{rec:v20v65-snapshots}
Main YM V61 has 749660 bytes and SHA--256
\[
\texttt{287266DB63980F79E6FACCCB695D7EC09324B66190BC32B264F212E39D6A4248}.
\]
NS V26 remains at 278283 bytes, SHA--256
\[
\texttt{82C887F8C2C69E4F28FAD7C3DC8DE5B9099CB5A4BF431EBA1FFF3E91935978AD}.
\]
Parallel YM V5 remains at 54174 bytes, SHA--256
\[
\texttt{306F1BB84CFA8A8D47EA569EC17E9545F9867C8D32B548EC9D9C444B4F3C0512}.
\]
These are read-time snapshots, and the documents are research
sources rather than instructions governing this work.
\end{record}

\begin{assessment}[Companion relevance]
\label{ass:v20v65}
Since the V60 audit of YM V56, the main companion has
developed a local nonlinear blocking construction and
examined its density derivatives.  YM V61 proves an
exact local section \(F(Rc)=c\), and proves that a
particular reference-path dependence of the cubic
block jet vanishes after a right-inverse insertion.
This supplies reference independence of that contribution
to the quadratic density, not its generic vanishing.
The general fibre density, gauge quotient, global measure,
interacting continuum, and mass gap remain unresolved.
No theorem about the nonlinear Yukawa evolution is
imported from these local identities.

NS remains at its rank-one Oseen analysis without a
global regularity theorem.  Parallel YM retains its
global extensive-charge tail obstruction.  Neither
supplies the missing ultraviolet source estimate.
The present SM direction therefore stays with its own
renormalized linear operator and exact spin invariant.
\end{assessment}

\begin{decision}[V66: a semilinear formulation]
\label{dec:v20v65-next}
Use V47's self-adjoint renormalized operator for the
linear parent--pair evolution, after the unitary change
to V64's physical momentum measure.  Treat the remaining
saturated terms as nonlinear maps on the parent,
pair-coherence, and occupation variables in a Hilbert
space.  Prove cutoff convergence there and use the
pointwise spin invariant to recover \(L^1\) occupations
and exact excitation conservation.  This would construct
a classical collective mild evolution, not an
ultraviolet quantum Standard Model.
The next companion audit and ten-version review are V70.
\end{decision}

\begin{record}[V65 append record]
\label{rec:v20v65}
The predecessor V64 has SHA--256
\texttt{EA35FC17F0D379854B79996D2ED3933A1DFE76BF8DEE4022AA6108D6EAB2ED14}.
Its complete prefix is retained.  Only the immediate
active SM predecessor is replaced.  The full continuum
goal remains unresolved.
\end{record}


\section{V66: global renormalized nonlinear mild evolution with Pauli saturation}
\label{sec:v20v66}

V64 excludes a uniform instantaneous-source bound from
energy control alone.  It does not exclude mild evolution
generated by an unbounded self-adjoint operator.  We now
use exactly that distinction: V47 handles the singular
linear parent--pair coupling, and the residual saturated
terms are locally Lipschitz on a Hilbert space.  A
pointwise invariant recovers \(L^1\) occupation convergence,
which a bare \(L^2\) estimate would not provide.

\subsection{Linear operator and nonlinear state space}

Use V64's functions \(h,\omega\) and measure
\(d\pi=dp/(2\pi)\).  Let
\[
 \mathcal H=\mathbb C\oplus L^2(\mathbb R,d\pi;\mathbb C).
\]
For real cutoffs \(0\leq\chi_\Lambda\leq1\), compactly
supported and tending pointwise to one, set
\[
 h_\Lambda=\chi_\Lambda h,\quad
 I_\Lambda=\int\frac{h_\Lambda^2}{\omega}\,d\pi,\quad
 \Omega_\Lambda=\Omega_R+I_\Lambda .
\]
We use cutoffs equal to one on expanding intervals,
as in V47; sharp interval cutoffs are also allowed
by the same dominated-tail proof.  Define
\[
 H_\Lambda=
 \begin{pmatrix}
 \Omega_\Lambda&\langle h_\Lambda,\cdot\rangle\\
 h_\Lambda&\omega
 \end{pmatrix}.
\]
Multiplication of the continuum coordinate by
\((2\pi)^{-1/2}\) maps this problem unitarily to V47,
with subtraction point zero and \(v=h/\sqrt{2\pi}\).
Hence \(H_\Lambda\) converges in norm resolvent to
the self-adjoint operator \(H_R\) constructed there.

The real Hilbert state space for the nonlinear problem is
\[
 \mathcal X=\mathcal H\oplus L^2(\mathbb R,d\pi;\mathbb R),
 \qquad Y=(\beta,s,x).
\]
Use the unitary group \(e^{-itH_R}\) on \((\beta,s)\)
and the identity on \(x\), denoting their direct sum by
\(\mathcal U_R(t)\); define \(\mathcal U_\Lambda(t)\)
in the same way.  Let \(H_*=\|h\|_\infty<\infty\) and set
\[
 \mathcal F_\Lambda(\beta,s,x)=
       (0,\ 2ih_\Lambda\beta x,\
                       2h_\Lambda\operatorname{Im}(\beta\overline s)),
\]
with \(\mathcal F\) the same formula using \(h\).
The last component is real.  All products belong to
the stated \(L^2\) spaces because \(\beta\) is a scalar.
On every ball of \(\mathcal X\) these maps have a
common Lipschitz constant, independent of \(\Lambda\).
This follows by expanding differences of each bilinear
product and using
\(\|h_\Lambda f\|_2\leq H_*\|f\|_2\).
Moreover \(\mathcal F_\Lambda(Y)\to\mathcal F(Y)\)
for every \(Y\), by dominated convergence.

The uncut equation is defined by its mild form
\[
 Y(t)=\mathcal U_R(t)Y_0+
              \int_0^t\mathcal U_R(t-u)\mathcal F(Y(u))\,du .
\]
The integral is a Hilbert-space integral.  This definition
does not require that \(\int h s\,d\pi\) exist pointwise
or that \(Y(t)\) belong to the domain of the generator.

\subsection{Local convergence of the semilinear equations}

\begin{lemma}[Common local existence and cutoff convergence]
\label{lem:v20v66-local}
For every \(Y_0\in\mathcal X\), the uncut and cutoff
mild equations have unique local solutions with a
common existence interval depending only on
\(\|Y_0\|_{\mathcal X}\) and \(H_*\).
On that interval,
\[
 \sup_{|t|\leq T}\|Y_\Lambda(t)-Y(t)\|_{\mathcal X}
                                 \longrightarrow0 .
\]
The same conclusion holds on any interval on which
both solutions have a common a priori norm bound.
\end{lemma}
\begin{proof}
The propagators have norm one and the nonlinearities
have common ball bounds and Lipschitz constants.
Picard contraction on continuous Hilbert-space paths
therefore supplies a common local interval and uniqueness.

Self-adjoint norm-resolvent convergence implies
\(\mathcal U_\Lambda(t)Z\to\mathcal U_R(t)Z\),
uniformly on compact time intervals for each fixed \(Z\).
One spectral-calculus justification is to truncate to
a compact energy interval using a continuous spectral
cutoff.  On that interval the functions
\(e^{-itE}\) times the cutoff form a compact family
of continuous functions vanishing at infinity, so
finite nets give uniform strong convergence.
For a fixed vector, the omitted spectral tails are
uniformly small by strong functional calculus and
normalization.  The identity component is unchanged.
Norm-one bounds then make this convergence uniform
for \(Z\) in a compact subset of \(\mathcal X\).

Likewise the common Lipschitz bounds make
\(\mathcal F_\Lambda\to\mathcal F\) uniform on compact
subsets.  Subtract the two mild equations.  Terms
evaluated on the fixed limiting trajectory tend
uniformly to zero, including the propagator differences
on the compact range of \(\mathcal F(Y(u))\).
The remaining term is bounded by the common Lipschitz
constant times the integral of
\(\|Y_\Lambda(u)-Y(u)\|\).
Gronwall proves convergence.  The same argument
iterated on bounded solution intervals proves
the last assertion.
\end{proof}

\subsection{Physical data and the exact spin invariant}

Assume now that
\[
 \beta_0\in\mathbb C,\quad x_0\in L^1(d\pi),\quad
 0\leq x_0\leq1,\quad
 |s_0|^2\leq x_0(1-x_0)\quad\hbox{almost everywhere}.
\]
Then \(x_0,s_0\in L^2\).  Define the nonnegative
integrable spin deficit
\[
 d_0(p)=x_0(p)-x_0(p)^2-|s_0(p)|^2,
 \qquad
 N_0=|\beta_0|^2+\int x_0\,d\pi<\infty .
\]

\begin{lemma}[Cutoff invariants and a uniform bound]
\label{lem:v20v66-cutoff}
For these initial data, every cutoff solution is global,
and for every time it satisfies almost everywhere
\[
 x_\Lambda-x_\Lambda^2-|s_\Lambda|^2=d_0,\qquad
 0\leq x_\Lambda\leq1,\qquad
 |s_\Lambda|^2\leq x_\Lambda(1-x_\Lambda).
\]
It also satisfies
\[
 |\beta_\Lambda(t)|^2+\int x_\Lambda(t)\,d\pi=N_0,
 \qquad
 \|Y_\Lambda(t)\|_{\mathcal X}^2\leq2N_0 .
\]
\end{lemma}
\begin{proof}
At finite cutoff the component equations are
\[
 i\dot\beta_\Lambda=\Omega_\Lambda\beta_\Lambda+
                                  \int h_\Lambda s_\Lambda\,d\pi,
\]
\[
 i\partial_t s_\Lambda=\omega s_\Lambda+
                           h_\Lambda\beta_\Lambda(1-2x_\Lambda),
 \qquad
 \partial_t x_\Lambda=
                     2h_\Lambda\operatorname{Im}(\beta_\Lambda\overline s_\Lambda).
\]
On the compact support of \(h_\Lambda\), the frequency
is bounded; the mild equations have the corresponding
pointwise absolutely continuous representatives.
Outside that support \(s_\Lambda\) rotates freely and
\(x_\Lambda=x_0\).  Thus the pointwise rotation argument
applies also to merely measurable physical initial data.
It gives
\(\partial_t(x_\Lambda-x_\Lambda^2-|s_\Lambda|^2)=0\).
Equivalently the spin length
\(|s_\Lambda|^2+(x_\Lambda-1/2)^2=1/4-d_0\)
is preserved, proving the constraint.

The derivative of \(|\beta_\Lambda|^2\) cancels
\(\int\partial_t x_\Lambda\,d\pi\).
This integral is legitimate: only a compact momentum
set changes, and \(h_\Lambda s_\Lambda\) is integrable
by Cauchy--Schwarz.  Initial \(L^1\) occupation is
preserved outside that set.  Hence excitation is
conserved on each local interval.
Finally \(\|x_\Lambda\|_2^2\leq\int x_\Lambda\) and
\(\|s_\Lambda\|_2^2\leq\int x_\Lambda\), so
\[
 |\beta_\Lambda|^2+\|s_\Lambda\|_2^2+\|x_\Lambda\|_2^2
              \leq2N_0.
\]
The local continuation criterion now gives global
existence in both time directions.
\end{proof}

\begin{theorem}[Global ultraviolet limit and exact occupation conservation]
\label{thm:v20v66-global}
For every physical initial datum above, the uncut mild
equation has a unique global solution \(Y\in C(\mathbb R;\mathcal X)\).
On every bounded time interval,
\[
 \sup_{|t|\leq T}\|Y_\Lambda(t)-Y(t)\|_{\mathcal X}\to0,
 \qquad
 \sup_{|t|\leq T}\|x_\Lambda(t)-x(t)\|_1\to0 .
\]
For every time the uncut solution satisfies
\[
 x-x^2-|s|^2=d_0\quad\hbox{almost everywhere},\quad
 0\leq x\leq1,\quad |s|^2\leq x(1-x),
\]
\[
 |\beta(t)|^2+\int x(t)\,d\pi=N_0.
\]
The solution is independent of the allowed cutoff
profile, with the subtraction convention held fixed.
\end{theorem}
\begin{proof}
On the common local interval, the preceding local
convergence lemma gives convergence in \(\mathcal X\).
At each fixed time an almost-everywhere convergent
subsequence of the two \(L^2\) coordinates passes
the pointwise invariant and constraint to the limit.
Thus \(x=x^2+|s|^2+d_0\) almost everywhere.
In particular \(x\in L^1\), since its right side is
integrable.  Subtracting this identity from the cutoff
one gives the quantitative bound
\[
 \begin{split}
 \|x_\Lambda-x\|_1\leq&
   (\|x_\Lambda\|_2+\|x\|_2)\|x_\Lambda-x\|_2\\
 &+(\|s_\Lambda\|_2+\|s\|_2)\|s_\Lambda-s\|_2 .
 \end{split}
\]
The fixed deficit \(d_0\) cancels.  The local
Hilbert-space convergence therefore proves uniform
\(L^1\) convergence on the same interval.
Passing to the cutoff conservation law yields the
exact uncut excitation law and the bound \(2N_0\).

These bounds permit repeated local continuation with
a time step depending only on \(N_0,H_*\); consequently
the solution and convergence extend to every bounded
time interval.  The same invariant argument applies
at each step.  Uniqueness is the local Lipschitz
uniqueness in \(\mathcal X\), extended globally.
Different allowed cutoffs have the same linear limit
\(H_R\) by V47 and the same nonlinear limit
\(\mathcal F\), so uniqueness proves profile independence.
\end{proof}

This argument closes a specific loophole in a possible
compactness proof: \(L^2\) convergence of occupations
alone would permit loss of their \(L^1\) mass.
The fixed pointwise deficit identity supplies the
additional information needed for exact conservation.

\subsection{Scope of the constructed evolution}

The theorem removes the momentum cutoff for this
classical collective saturated model in a precise mild
topology.  It includes the initially empty-pair data
\(x_0=s_0=0\), with arbitrary finite \(\beta_0\).
For such data the initial vector need not belong to
the linear operator domain, so differentiability of
the parent amplitude or an instantaneous source
formula is not asserted.  This is consistent with
V64's source obstruction.

No conservation of a renormalized energy functional
has been proved here.  The finite-cutoff expression
contains counterterms, and taking its limit requires
separate domain or form control.  Nor has the sequence
of microscopic grouped quantum models been shown to
converge uniformly as this ultraviolet cutoff is
removed.  The present construction concerns a
nonlinear classical Hilbert-space equation, with
one parent mode and the specified pair continuum.
It supplies no spatially local stress tensor, chiral
gauge sector, or Einstein evolution.

\begin{decision}[Next domain and energy question]
\label{dec:v20v66-next}
Determine the renormalized energy or quadratic-form
domain compatible with this nonlinear flow.  First
write the finite-cutoff conserved energy in terms
of the linear renormalized form and the spin deficit,
and test which physical initial data have a finite
limit.  Do not infer a conserved gravitational source
from excitation conservation alone.
\end{decision}

\begin{record}[V66 append record]
\label{rec:v20v66}
The immediate predecessor is the V65 companion audit.
Its complete prefix is retained and verified by the
installer.  The ultraviolet argument combines V47's
linear norm-resolvent construction with a new
locally Lipschitz nonlinear equation, uniform physical
invariants, and an \(L^1\) recovery estimate.
Structural checks cover source preservation,
environments, delimiters, labels, and references;
they are not LaTeX compilation.
Only the immediate active SM predecessor is replaced.
The next companion audit and ten-version review are V70.
\end{record}


\section{V67: the dressed energy domain and conservation under the nonlinear flow}
\label{sec:v20v67}

V66 constructed global mild evolution without asserting finite
energy.  We now identify a renormalized energy domain and prove
conservation there.  The argument uses physical cutoff recovery
data, lower semicontinuity, and reversibility of the unique mild
flow.  It does not differentiate an undefined instantaneous source.

\subsection{The linear quadratic form}

Keep \(d\pi=dp/(2\pi)\), \(\omega\geq2m>0\), and \(h\)
from V64--V66.  Set
\[
 k=h/\omega,\qquad \eta=s+k\beta,\qquad
 \|f\|_\omega^2=\int\omega|f|^2\,d\pi .
\]
Here \(k\in L^2(d\pi)\), but
\(\|k\|_\omega^2=\int h^2/\omega\,d\pi=\infty\).
The first assertion follows from \(k(p)=O(|p|^{-1})\);
the second is V64's logarithmic divergence.

\begin{theorem}[Closed dressed form of the renormalized operator]
\label{thm:v20v67-form}
On \(\mathcal H=\mathbb C\oplus L^2(d\pi)\), define
\[
 \mathcal D_{\rm form}=
       \{(\beta,s):s+k\beta\in L^2(\omega\,d\pi)\},\qquad
 q_R(\beta,s)=\Omega_R|\beta|^2+\|s+k\beta\|_\omega^2 .
\]
This is a densely defined, closed, lower-semibounded
quadratic form.  Its associated self-adjoint operator
is precisely \(H_R\) from V66.
Its operator domain and action can be written as
\[
 \eta\in D(\omega),\qquad
 H_R(\beta,s)=
       \big(\Omega_R\beta+\langle k,\omega\eta\rangle,\
                                                  \omega\eta\big).
\]
\end{theorem}
\begin{proof}
The map \((\beta,s)\mapsto(\beta,s+k\beta)\) is a bounded
invertible map of \(\mathcal H\).
The space \(\mathbb C\oplus L^2(\omega\,d\pi)\) is dense
and complete in its weighted norm.  Adding a sufficiently
large multiple of \(|\beta|^2+\|s\|_2^2\) to \(q_R\)
produces a norm equivalent to
\(|\beta|^2+\|\eta\|_\omega^2\), since \(\omega\geq2m\)
and \(k\in L^2\).  This proves density, closedness, and
the lower bound \(q_R\geq-\max(0,-\Omega_R)|\beta|^2\).

Polarize the form.  Testing against vectors \((0,f)\)
forces the second operator component to equal \(\omega\eta\)
and requires \(\eta\in D(\omega)\).
Testing against \((b,-kb)\), whose dressed coordinate is
zero, then forces the first component to be
\(\Omega_R\beta+\langle k,\omega\eta\rangle\).
These conditions are also sufficient in the polarized
identity, giving the stated domain and action.

To identify this operator with V47's construction, solve
its resolvent equation at nonreal \(z\).
The continuum component gives
\[
 s=(z-\omega)^{-1}f+(z-\omega)^{-1}h\,\beta .
\]
Inserting it into the parent component gives the
subtracted denominator
\[
 z-\Omega_R-
      \int h^2\left(\frac1{z-\omega}+\frac1\omega\right)d\pi.
\]
All these terms are convergent: the dressed coefficient
\((z-\omega)^{-1}h+k\) has one extra power of decay.
The other resolvent blocks are the same weighted
resolvent vectors and row functionals as V47, after
the measure change of V66.  Thus the full resolvents,
and hence the self-adjoint operators, agree.
\end{proof}

\subsection{The physical energy domain}

For V66's physical data let
\[
 d=x-x^2-|s|^2\geq0 .
\]
Define the finite-energy class by the additional conditions
\[
 \eta=s+k\beta\in L^2(\omega\,d\pi),\qquad
 x\in L^2(\omega\,d\pi),\qquad
 d\in L^1(\omega\,d\pi).
\]
The renormalized energy is
\[
 \mathscr E_R(\beta,s,x)=
    \Omega_R|\beta|^2+\|\eta\|_\omega^2+
                             \|x\|_\omega^2+\int\omega d\,d\pi.
\]
It is bounded below on an excitation surface
\(N=|\beta|^2+\int x\,d\pi\) by
\(-\max(0,-\Omega_R)N\).

To explain this formula, use the sharp cutoff
\(\chi_\Lambda=1_{\{|p|\leq\Lambda\}}\).
For any cutoff state with \(s,x\) supported in this interval,
the finite-cutoff conserved energy is exactly
\[
 \begin{split}
 \mathscr E_\Lambda
 &= (\Omega_R+I_\Lambda)|\beta|^2+
             \int\omega x\,d\pi+
                      2\operatorname{Re}\int h_\Lambda\beta\overline s\,d\pi\\
 &=\Omega_R|\beta|^2+
       \|s+(h_\Lambda/\omega)\beta\|_\omega^2+
                           \|x\|_\omega^2+\int\omega d\,d\pi .
 \end{split}
\]
Indeed \(x=|s|^2+x^2+d\); completing the square in
the \((\beta,s)\) terms uses
\(\int h_\Lambda^2/\omega\,d\pi=I_\Lambda\).
All integrals in this cutoff identity are finite.

\begin{lemma}[Physical recovery and varying initial data]
\label{lem:v20v67-recovery}
For every physical finite-energy state \(Y=(\beta,s,x)\),
define
\[
 Y^{[\Lambda]}=(\beta,\chi_\Lambda s,\chi_\Lambda x).
\]
Its deficit is \(d^{[\Lambda]}=\chi_\Lambda d\), its
dressed cutoff coordinate is
\(\eta^{[\Lambda]}=\chi_\Lambda\eta\), and
\[
 Y^{[\Lambda]}\to Y\ \hbox{in }\mathcal X,\qquad
 \|x^{[\Lambda]}-x\|_1\to0,\qquad
 \mathscr E_\Lambda(Y^{[\Lambda]})\to\mathscr E_R(Y).
\]
The corresponding cutoff solutions converge to V66's
uncut solution in \(C([-T,T];\mathcal X)\) and in
\(C([-T,T];L^1)\) for occupations, for each finite \(T\).
\end{lemma}
\begin{proof}
Indicator truncation preserves the spin constraint,
sets the outside spins to their empty state, and gives
the stated deficit and dressed coordinate exactly.
Tail convergence in each indicated space, and monotone
convergence of the nonnegative weighted energy terms,
prove recovery.

The local proof of V66 applies also to convergent initial
data: subtracting mild equations adds only
\(\|Y^{[\Lambda]}-Y\|_{\mathcal X}\) to its Gronwall
estimate.  Excitation bounds for these data are uniformly
at most \(N\), so the extension is global on compact
time intervals.  For occupation convergence, V66's
invariant-identity estimate now has the additional
term \(\|d^{[\Lambda]}-d\|_1\), which tends to zero.
This proves the claimed convergences.
\end{proof}

Using the same untruncated dressed state at every cutoff
would not give this recovery: its \(s\) may fail to have
finite bare weighted norm even at finite coupling cutoff.
The explicit physical truncation is therefore part of
the proof, not an optional change of notation.

\subsection{Conservation without a pointwise source formula}

\begin{theorem}[Energy invariance of the renormalized nonlinear flow]
\label{thm:v20v67-conservation}
V66's global physical mild flow preserves the finite-energy
class and satisfies
\[
 \mathscr E_R(Y(t))=\mathscr E_R(Y(0))
                  \quad\hbox{for every }t\in\mathbb R .
\]
The deficit is the fixed function \(d_0\).
Moreover \((\eta(t),x(t))\) is continuous with values in
\(L^2(\omega\,d\pi)\oplus L^2(\omega\,d\pi)\).
\end{theorem}
\begin{proof}
Take the recovery sequence from the lemma.
Each cutoff solution remains supported in the cutoff
interval in its \(s,x\) coordinates: outside it the
initial state is empty and the coupling is zero.
Its finite energy is conserved by the finite-cutoff
equations, whose momentum coefficients are bounded
on that interval.  Its deficit remains
\(\chi_\Lambda d_0\).

At a fixed time, the lemma gives strong convergence
of \(\beta_\Lambda,s_\Lambda,x_\Lambda\) in the
unweighted spaces.  Since \(k_\Lambda\to k\) in
\(L^2\), it also gives
\(\eta_\Lambda=s_\Lambda+k_\Lambda\beta_\Lambda\to\eta\)
in unweighted \(L^2\).
The completed-square energy identity and bounded
\(\beta_\Lambda\) give a uniform bound on the two
weighted squared norms.  Weighted lower semicontinuity
(or subsequence Fatou) and
\(\int\omega\chi_\Lambda d_0\,d\pi\to\int\omega d_0\,d\pi\)
therefore imply
\[
 \mathscr E_R(Y(t))\leq\mathscr E_R(Y(0)).
\]
The possibly negative scalar term causes no problem:
it converges separately, and all remaining terms
are nonnegative.  This also proves that \(Y(t)\)
belongs to the finite-energy class.

Apply this same inequality with initial datum \(Y(t)\)
and elapsed time \(-t\).  The autonomous mild equation
is uniquely solvable in both directions by V66, so this
backward solution reaches \(Y(0)\).  Hence
\(\mathscr E_R(Y(0))\leq\mathscr E_R(Y(t))\).
The two inequalities prove equality without
differentiating the uncut energy.

On any compact time interval the two weighted coordinates
are bounded by energy conservation, excitation conservation,
and the fixed weighted deficit.  Unweighted continuity
then implies their weighted weak continuity: test first
against compactly supported momentum functions and use
the weighted bound to approximate an arbitrary weighted
test.  Their combined squared norm is
\[
 \|\eta(t)\|_\omega^2+\|x(t)\|_\omega^2
 =\mathscr E_R(Y(0))-\Omega_R|\beta(t)|^2
                                  -\int\omega d_0\,d\pi,
\]
which is continuous.  Weak continuity together with
continuity of the Hilbert norm proves strong continuity
of the pair \((\eta,x)\).
\end{proof}

The same argument also gives, for the recovery solutions,
strong convergence of \((\eta_\Lambda(t),x_\Lambda(t))\)
in the two weighted spaces at each fixed time.
Indeed the weak limits are identified by unweighted
convergence, and conservation plus convergence of the
initial energies and deficit integrals gives convergence
of the combined weighted norms.  No uniform-in-time
weighted convergence rate is asserted.

\subsection{Nonempty finite-energy data and an important exclusion}

The bare all-parent datum \(s=x=0\), \(\beta\ne0\),
is allowed in V66 but has
\(\eta=k\beta\notin L^2(\omega\,d\pi)\).
Its energy as defined here is infinite.  Thus global
Hilbert-space evolution and finite renormalized energy
are distinct assertions.

The finite-energy class nevertheless contains states
with nonzero parent amplitude.  Fix any \(\beta\).
For sufficiently large \(P\), the tail
\(|k(p)\beta|\) is smaller than \(1/2\) for \(|p|\geq P\).
Choose a continuous cutoff of this tail and set \(s=0\)
near zero and \(s=-k\beta\) for \(|p|\geq2P\), keeping
\(|s|\leq1/2\) throughout.  Put
\[
 x=\frac{1-\sqrt{1-4|s|^2}}2,\qquad d=0.
\]
Then the spin length is maximal, \(\eta\) is supported
in a compact interval, and in the tail \(x=O(|p|^{-2})\).
Thus \(x\in L^1\), \(s\in L^2\), and
\(\int\omega x^2\,d\pi<\infty\); all energy conditions
hold.  Its bare pair energy \(\int\omega x\,d\pi\)
generally diverges, while the completed-square
renormalized energy is finite.  This example explains
why separating the dressed form from the bare
positive pair energy is essential.

\begin{assessment}[Energy has been constructed only for this model]
\label{ass:v20v67}
The renormalized saturated classical continuum now has
a closed linear form, a nonempty physical finite-energy
class, and an exactly conserved nonlinear energy there.
The proof is compatible with V64's source obstruction
and does not require an instantaneous uncut source.
This is a global energy for the one-parent-mode
collective model.  It is not a local stress tensor,
nor a microscopic quantum energy construction or an
Einstein source.
\end{assessment}

\begin{decision}[Test the locality requirement before a gravity handoff]
\label{dec:v20v67-next}
Identify which spatial observables are actually represented
by the fixed rest-parent, opposite-momentum pair sector.
Determine whether its conserved global energy can supply
the local energy and momentum densities required by an
Einstein equation, and isolate the missing parent-momentum
channels if it cannot.  Do not introduce a gravitational
closure by assigning a local stress tensor to a global
energy without a derivation.
\end{decision}

\begin{record}[V67 append record]
\label{rec:v20v67}
The predecessor V66 has SHA--256
\texttt{E6A46F6DC67942974FACC24AAEBADDBA9F396ABDC2B3D5FE98147CFEEB5E918F}.
Its complete prefix is retained.  The energy proof
uses exact cutoff identities, physical recovery,
lower semicontinuity, and reversibility of the
already constructed mild flow.  Structural checks
cover source preservation, environments, delimiters,
labels, and references; they are not LaTeX compilation.
Only the immediate active SM predecessor is replaced.
The next companion audit and ten-version review are V70.
\end{record}


\section{V68: the locality obstruction and overlapping pair channels}
\label{sec:v20v68}

V67 supplies a conserved global energy for the rest-parent
collective model.  An Einstein source requires local energy,
momentum, and metric response, so the next question is what
spatial information that model actually contains.  We analyze
the finite-box sector before any cutoff limit.  The conclusions
are kinematic selection rules and exact CAR identities, not
a claim that homogeneous matter cannot gravitate.

\subsection{All retained configurations have zero total momentum}

On a circle of length \(L\), momenta lie in
\((2\pi/L)\mathbb Z\).  Let \(a_q,c_p,d_u\) denote
parent, fermion, and antifermion modes.  The normal-ordered
total momentum is
\[
 \mathsf P=\sum_q q\,a_q^\dagger a_q+
          \sum_p p\,c_p^\dagger c_p+
          \sum_u u\,d_u^\dagger d_u .
\]
The collective sector used in V43--V67 contains only
the parent \(a_0\) and pair creations
\(c_p^\dagger d_{-p}^\dagger\).  Every such creation
has total momentum zero.  Hence on its embedded
finite-particle subspace, \(\mathsf P=0\).
Let \(P_{\rm coll}\) be its orthogonal projection.
This subspace is smaller than the full zero-total-momentum
sector.

\begin{theorem}[Compression removes nonzero momentum transfers]
\label{thm:v20v68-transfer}
Suppose an observable component \(A_k\) obeys
\[
 U(a)A_kU(a)^*=e^{ika}A_k,\qquad U(a)=e^{ia\mathsf P}.
\]
For \(k\ne0\), its matrix elements between collective
states vanish:
\[
 P_{\rm coll}A_kP_{\rm coll}=0 .
\]
The assertion holds for bounded operators, or on any
common finite-particle domain where the stated matrix
elements exist.  Every density matrix supported in
the collective sector is translation invariant; its
one-point expectations of translation-covariant local
observables are spatially constant when defined.
\end{theorem}
\begin{proof}
Translations act as the identity on the collective
subspace.  Sandwiching the covariance identity between
two vectors in that subspace gives
\(\langle u,A_kv\rangle=e^{ika}\langle u,A_kv\rangle\).
For a nonzero circle Fourier momentum one can choose
\(a\) with \(e^{ika}\ne1\), proving the vanishing.
The density matrix is unchanged by conjugation with
translations, so the covariance law of a local
observable gives constant one-point expectations.
\end{proof}

Thus a nonconstant lapse or another spatial source,
coupled through nonzero Fourier components of a local
density, has no first-order matrix elements inside
this projection.  This does not say those components
are zero in the full theory.  They connect to excluded
momentum sectors.

\subsection{Local response is not reconstructed by multiplying compressed observables}

For any projection \(P\) and composable operators \(A,B\),
\[
 PABP=PAP\,PBP+PA(I-P)BP .
\]
This follows by inserting \(I=P+(I-P)\).
The second term is the intermediate excursion omitted
by multiplying the separately compressed operators.

For an explicit finite-particle example, in the
one-boson space let \(|p\rangle=a_p^\dagger|0\rangle\),
and take \(P=|0_{\rm mom}\rangle\langle0_{\rm mom}|\),
where \(|0_{\rm mom}\rangle=a_0^\dagger|0\rangle\)
is the zero-momentum one-particle state, not the vacuum.
The Fourier number-density operators
\[
 n_k=\sum_p a_{p+k}^\dagger a_p
\]
satisfy \(n_k|p\rangle=|p+k\rangle\).  For \(k\ne0\),
\[
 Pn_kP=Pn_{-k}P=0,\qquad Pn_{-k}n_kP=P .
\]
Only modes \(0,k,-k\) needed for these matrix elements
need be retained in a finite cutoff example.
This is a number-density example of the compression
problem, not an identification of number density
with the relativistic stress tensor.

The same projection identity applies to a response
with intermediate propagation, such as
\(PA_{-k}e^{-itH}A_kP\), when defined.
Global energy conservation within \(P\mathcal H\)
does not determine the missing intermediate matrix
elements.  In particular, it cannot by itself
reconstruct all local stress responses.

\subsection{Restoring parent momenta changes the pair algebra}

For arbitrary momenta define
\[
 P_{p,u}^\dagger=c_p^\dagger d_u^\dagger,\qquad
 P_{p,u}=d_uc_p .
\]
Unlike the rest-parent matched set \(u=-p\), the
full set contains pairs sharing a fermion mode.

\begin{theorem}[Exact overlapping-pair commutator]
\label{thm:v20v68-overlap}
With canonical anticommutation between all fermion
and antifermion modes,
\[
 [P_{p,u},P_{v,w}^\dagger]
 =\delta_{pv}\delta_{uw}
       -\delta_{pv}d_w^\dagger d_u
       -\delta_{uw}c_v^\dagger c_p .
\]
In particular, if \(p=v\) but \(u\ne w\), the
commutator is \(-d_w^\dagger d_u\), not zero.
Moreover
\[
 P_{p,u}^\dagger P_{p,w}^\dagger=0
 \quad\hbox{for }u\ne w .
\]
Thus distinct overlapping pairs cannot be represented
by independent two-state spins.
\end{theorem}
\begin{proof}
Expand
\[
 d_uc_pc_v^\dagger d_w^\dagger
   =\delta_{pv}d_ud_w^\dagger-d_uc_v^\dagger c_pd_w^\dagger .
\]
Use \(d_ud_w^\dagger=\delta_{uw}-d_w^\dagger d_u\);
anticommute the remaining fermion factors to put the
second creation pair on the left.  The result is
\[
 \delta_{pv}\delta_{uw}-\delta_{pv}d_w^\dagger d_u
 -\delta_{uw}c_v^\dagger c_p+
                  c_v^\dagger d_w^\dagger d_uc_p .
\]
Subtracting the reversed product proves the commutator.
The second identity follows from \((c_p^\dagger)^2=0\).
Independent spin creation operators on distinct
two-state factors would have a nonzero product on
their joint empty state, contradicting this identity.
\end{proof}

The momentum-complete conversion vertex has the form
\[
 H_I=\sum_{q,p}
 \left(g_{q,p}a_qc_p^\dagger d_{q-p}^\dagger+
       \overline g_{q,p}a_q^\dagger d_{q-p}c_p\right)
\]
at a finite mode cutoff, with specified coefficients.
Every term conserves total momentum, but changing
\(q\) permits the same \(c_p\) to appear in many
pairs.  The independent disjoint-pair spin algebra
of the collective model therefore cannot be carried
over unchanged.

There is also a direct failure of invariance of the
matched collective sector under generic such vertices.
Choose distinct \(p,v\), with modes and nonzero
coefficient \(g_{p-v,p}\) retained.  Starting from
\[
 c_p^\dagger d_{-p}^\dagger
 c_v^\dagger d_{-v}^\dagger|0\rangle ,
\]
the inverse-conversion monomial
\(a_{p-v}^\dagger d_{-v}c_p\) produces, up to a nonzero
CAR sign,
\[
 a_{p-v}^\dagger c_v^\dagger d_{-p}^\dagger|0\rangle .
\]
The output still has total momentum
\((p-v)+v-p=0\), but contains a moving parent and
an unmatched pair.  It lies outside the collective
sector.  Total-momentum conservation does not
justify retaining only zero-momentum parent modes.

\subsection{A scalar energy value does not fix metric response}

Even in a free finite-box example, global energy and
zero mean momentum do not determine pressure response.
Let a fermion of mass \(m\) have
\[
 E_n(L)=\sqrt{m^2+(2\pi n/L)^2}.
\]
At a reference length \(L_0\), write \(p=2\pi n/L_0\ne0\)
and \(E_p=\sqrt{m^2+p^2}\).  Choose \(0<\varepsilon<m\).
The two diagonal density matrices
\[
 \rho_0=\frac{\varepsilon}{m}|0_{\rm mom}\rangle
          \langle0_{\rm mom}|+
              (1-\varepsilon/m)|{\rm vac}\rangle\langle{\rm vac}|,
\]
\[
 \rho_p=\frac{\varepsilon}{2E_p}
       (|p\rangle\langle p|+|-p\rangle\langle-p|)
              +(1-\varepsilon/E_p)|{\rm vac}\rangle\langle{\rm vac}|
\]
are positive, normalized, and have the same energy
\(\varepsilon\) and zero mean momentum at \(L_0\).
Keep their probabilities fixed when varying \(L\).
The mechanical response of each occupied level is
\[
 -L\frac{dE_n(L)}{dL}\bigg|_{L=L_0}=\frac{p^2}{E_p}.
\]
It vanishes for the zero-momentum level.  Thus the
expectations of this length-response observable are
respectively zero and
\(\varepsilon p^2/E_p^2>0\).
This derivative is computed directly from the free
dispersion.  It demonstrates the missing metric-response
information without assigning a stress tensor to
V67's interacting collective model.

\begin{assessment}[A precise limit on the proposed gravity handoff]
\label{ass:v20v68}
The retained sector has homogeneous one-point data,
omits intermediate nonzero momentum transfers, and
is not invariant under a generic momentum-complete
conversion Hamiltonian once two pairs are present.
Restoring the missing channels introduces overlapping
pair commutators and shared-mode exclusion, beyond
the independent-spin equations.  V67's energy theorem
remains valid in its stated sector, but it cannot
alone provide a local Einstein source or its metric
response.  Homogeneous gravitational models are not
ruled out; their pressure and covariant balance
still require an actual derivation.
\end{assessment}

\begin{decision}[Return to the momentum-complete finite-cutoff hierarchy]
\label{dec:v20v68-next}
For a specified finite set of parent and fermion modes,
derive the exact evolution of the fermion one-particle
density matrices and pair coherences with all allowed
parent momenta.  Use the overlapping-pair commutator
instead of independent spins, and identify the joint
correlations required to preserve local momentum
transfer.  This is the upstream algebraic dependency
before attempting a spatial quantum closure or a
stress-tensor construction.
\end{decision}

\begin{record}[V68 append record]
\label{rec:v20v68}
The predecessor V67 has SHA--256
\texttt{1ED7BD7C2FF7C25DB6D1E8CD3688259DD70D4776AC9A69D5010AD0BCD011D8BD}.
Its complete prefix is retained.  The new statements
are proved by translation covariance, projection
identities, CAR, and a free spectral derivative.
Archived keyword matches concerned different
cofinal-locality and block-covariance statements;
they are not repeated here.
Structural checks cover source preservation,
environments, delimiters, labels, and references;
they are not LaTeX compilation.
Only the immediate active SM predecessor is replaced.
The next companion audit and ten-version review are V70.
\end{record}


\section{V69: the momentum-resolved density and conversion hierarchy}
\label{sec:v20v69}

V68 showed that restoring spatial momentum transfers destroys
the independent-pair spin description.  We now derive the
exact replacement at finite mode cutoff.  Fermion density
matrices, rather than only their diagonal occupations, are
retained.  The resulting system exposes the mixed correlations
that a spatial approximation must estimate; it is not declared
closed.

\subsection{Specified finite-mode model}

Let parent modes have labels \(a,b\), momenta \(q_a\), and
energies \(\Omega_a\).  Fermion modes have labels \(i,j\),
momenta \(p_i\), and energies \(E_i\); antifermion labels
are \(r,s\), with momenta \(u_r\) and energies \(\widetilde E_r\).
All three sets are finite.  Set \(P_{ir}^\dagger=c_i^\dagger d_r^\dagger\),
\(P_{ir}=d_rc_i\), and
\[
 H=H_0+\sum_{a,i,r}
       (g_{a;ir}a_aP_{ir}^\dagger+
                         \overline g_{a;ir}a_a^\dagger P_{ir}),
\]
\[
 H_0=\sum_a\Omega_a a_a^\dagger a_a+
       \sum_iE_i c_i^\dagger c_i+
       \sum_r\widetilde E_r d_r^\dagger d_r,\qquad
 g_{a;ir}=0\ \hbox{unless }q_a=p_i+u_r.
\]
Coefficients are otherwise arbitrary fixed complex numbers.
No dispersion or renormalization estimate is inferred from
this finite algebra.

The excitation
\[
 \mathcal N=\sum_a a_a^\dagger a_a+
       \tfrac12\sum_i c_i^\dagger c_i+
       \tfrac12\sum_r d_r^\dagger d_r
\]
is conserved.  Every fixed eigenvalue of \(\mathcal N\)
has a finite-dimensional Fock subspace: there are finitely
many fermion configurations and a bounded total parent
number.  We work in any finite sum of these invariant
subspaces.  All operators and differentiations used below
then have unambiguous finite matrix elements and unitary
evolution.

Define the expectations, with the displayed index convention,
\[
 B_{ab}=\langle a_b^\dagger a_a\rangle,\quad
 C_{ij}=\langle c_j^\dagger c_i\rangle,\quad
 D_{rs}=\langle d_s^\dagger d_r\rangle,\quad
 T_{a;ir}=\langle a_a c_i^\dagger d_r^\dagger\rangle .
\]
The tensor \(T\) is defined for all its indices, not only
where the interaction coefficient is nonzero.  Off-diagonal
density equations require conversion observables carrying
nonzero momentum transfer.

\subsection{Exact one-particle density equations}

\begin{theorem}[Density matrices retain the conversion coherences]
\label{thm:v20v69-densities}
The exact equations are
\[
 \begin{aligned}
 i\dot C_{ij}&=(E_i-E_j)C_{ij}
       +\sum_{a,r}
       (g_{a;ir}T_{a;jr}-\overline g_{a;jr}\,\overline T_{a;ir}),\\
 i\dot D_{rs}&=(\widetilde E_r-\widetilde E_s)D_{rs}
       +\sum_{a,i}
       (g_{a;ir}T_{a;is}-\overline g_{a;is}\,\overline T_{a;ir}),\\
 i\dot B_{ab}&=(\Omega_a-\Omega_b)B_{ab}
       -\sum_{i,r}g_{b;ir}T_{a;ir}
       +\sum_{i,r}\overline g_{a;ir}\,\overline T_{b;ir}.
 \end{aligned}
\]
\end{theorem}
\begin{proof}
Use \(i\dot O=[O,H]\).  The CAR give
\[
 [c_j^\dagger c_i,P_{vr}^\dagger]
             =\delta_{iv}P_{jr}^\dagger,\qquad
 [c_j^\dagger c_i,P_{vr}]
             =-\delta_{jv}P_{ir},
\]
and the analogous two identities for antifermions.
The CCR give
\([a_b^\dagger a_a,a_e]=-\delta_{be}a_a\) and
\([a_b^\dagger a_a,a_e^\dagger]=\delta_{ae}a_b^\dagger\).
The free commutators supply the stated energy differences.
Taking expectations and using adjoints for the reverse
conversion monomials proves each equation.
\end{proof}

\subsection{The next equation and its precise missing terms}

For compactness set
\[
 Q_{ir;js}=\langle P_{ir}^\dagger P_{js}\rangle,\quad
 M^c_{ab;ji}=\langle a_b^\dagger a_a c_i^\dagger c_j\rangle,\quad
 M^d_{ab;sr}=\langle a_b^\dagger a_a d_r^\dagger d_s\rangle,
\]
and \(\Delta_{a;ir}=E_i+\widetilde E_r-\Omega_a\).

\begin{theorem}[Exact conversion-coherence hierarchy]
\label{thm:v20v69-conversion}
For all \(a,i,r\),
\[
 \begin{split}
 i\dot T_{a;ir}={}&-\Delta_{a;ir}T_{a;ir}
       +\sum_{j,s}\overline g_{a;js}Q_{ir;js}
       -\sum_b\overline g_{b;ir}B_{ab}\\
 &+\sum_{b,s}\overline g_{b;is}M^d_{ab;sr}
       +\sum_{b,j}\overline g_{b;jr}M^c_{ab;ji}.
 \end{split}
\]
\end{theorem}
\begin{proof}
Write \(X_{a;ir}=a_aP_{ir}^\dagger\).
Pair creations commute, including the case in which
their product vanishes by exclusion.  Hence
\([X_{a;ir},X_{b;js}]=0\).
V68's overlapping-pair identity gives
\[
 \begin{split}
 [X_{a;ir},X_{b;js}^\dagger]
 ={}&\delta_{ab}P_{ir}^\dagger P_{js}\\
 &-a_b^\dagger a_a
     (\delta_{ji}\delta_{sr}
            -\delta_{ji}d_r^\dagger d_s
            -\delta_{sr}c_i^\dagger c_j).
 \end{split}
\]
Indeed \(a_a a_b^\dagger=a_b^\dagger a_a+\delta_{ab}\);
the remaining difference of pair products is the
negative of \([P_{js},P_{ir}^\dagger]\).
The free commutator is \(-\Delta_{a;ir}X_{a;ir}\).
Multiplying the interaction commutator by
\(\overline g_{b;js}\), summing, and taking expectations
proves the formula.
\end{proof}

The four-body term and mixed moments can be separated
from their products of one-particle data without assuming
that the differences are small.  Define exactly
\[
 \begin{aligned}
 K^{cd}_{ir;js}&=Q_{ir;js}-C_{ji}D_{sr},\\
 K^{ad}_{ab;sr}&=M^d_{ab;sr}-B_{ab}D_{sr},\\
 K^{ac}_{ab;ji}&=M^c_{ab;ji}-B_{ab}C_{ji}.
 \end{aligned}
\]
These are defects relative to the displayed factorization,
not a claim that it is the full Wick rule for every
possible quasifree or anomalous state.
The exact conversion equation becomes
\[
 \begin{split}
 i\dot T_{a;ir}={}&-\Delta_{a;ir}T_{a;ir}
  +\sum_{j,s}\overline g_{a;js}C_{ji}D_{sr}
  -\sum_b\overline g_{b;ir}B_{ab}\\
 &+\sum_{b,s}\overline g_{b;is}B_{ab}D_{sr}
  +\sum_{b,j}\overline g_{b;jr}B_{ab}C_{ji}
  +\mathcal R_{a;ir},
 \end{split}
\]
where
\[
 \mathcal R_{a;ir}
 =\sum_{j,s}\overline g_{a;js}K^{cd}_{ir;js}
  +\sum_{b,s}\overline g_{b;is}K^{ad}_{ab;sr}
  +\sum_{b,j}\overline g_{b;jr}K^{ac}_{ab;ji}.
\]
Thus the spatial closure problem has three explicitly
weighted defect tensors, including off-diagonal shared-mode
terms.  Setting them to zero is an additional approximation,
not an algebraic consequence of retaining density matrices.

\subsection{Exact consistency checks and conservation laws}

\begin{proposition}[Charges, momentum, and matrix Pauli bounds]
\label{prop:v20v69-conservation}
The exact hierarchy has
\[
 \frac{d}{dt}\operatorname{Tr}C
 =\frac{d}{dt}\operatorname{Tr}D
 =-\frac{d}{dt}\operatorname{Tr}B
 =2\operatorname{Im}\sum_{a,i,r}g_{a;ir}T_{a;ir}.
\]
It preserves fermion charge \(\operatorname{Tr}C-\operatorname{Tr}D\),
total excitation \(\operatorname{Tr}B+
(\operatorname{Tr}C+\operatorname{Tr}D)/2\), and mean
total momentum
\[
 \sum_aq_aB_{aa}+\sum_ip_iC_{ii}+\sum_ru_rD_{rr}.
\]
Moreover at every time
\[
 B\geq0,\qquad 0\leq C\leq I,\qquad 0\leq D\leq I .
\]
\end{proposition}
\begin{proof}
Take diagonal traces of the density equations.  The
energy differences vanish, and the interaction terms
give the displayed imaginary parts.  Weighting instead
by the respective momenta gives a sum of
\[
 2(-q_a+p_i+u_r)\operatorname{Im}(g_{a;ir}T_{a;ir}),
\]
which is zero on the support of \(g\).
Finally for any coefficient vector, a quadratic form
of \(C\) is an expectation of a fermion annihilation
combination's adjoint times itself (with the conjugate
coefficient convention dictated by \(C_{ij}\)).
It is nonnegative, and the CAR bound the combination's
occupation by the squared coefficient norm.  This gives
\(0\leq C\leq I\); the same proof gives the \(D\) bounds,
while the CCR give positivity of \(B\).
Unitary state evolution preserves these facts.
\end{proof}

For a single parent and one pair, the conversion
equation reduces to
\[
 i\dot T=-\Delta T+
     \overline g\{\langle FD\rangle-\langle N\rangle+
                         \langle ND\rangle+\langle NF\rangle\}
           =-\Delta T-\overline g\langle\mathcal B\rangle,
\]
matching V44's sign and spontaneous/blocking terms.
For a state in a fixed total-excitation sector the
bare pair mean \(\langle P_{ir}\rangle\) can vanish
while \(T_{a;ir}\) is nonzero: \(X_{a;ir}\) preserves
total excitation.  Replacing the neutral conversion
tensor by products of charged field means would
therefore lose actual dynamics.

\subsection{Momentum transfer is kept at the level of observables}

The density entry \(c_j^\dagger c_i\) carries momentum
\(p_j-p_i\).  In its equation, a term with nonzero
\(g_{a;ir}\) multiplies
\(T_{a;jr}\), whose operator has momentum
\[
 p_j+u_r-q_a=p_j-p_i.
\]
The other term has the same transfer after taking
the adjoint.  This verifies the covariance of the
off-diagonal equation directly.  The antifermion
and parent equations obey the analogous rule.
Consequently discarding all \(T\) entries except
those with \(q_a=p_i+u_r\) would remove the observables
needed for nonzero-transfer density evolution,
even though only such momentum-conserving triples
occur in the Hamiltonian itself.

The finite-cutoff model conserves global momentum
and energy, but these identities alone are not
a local continuity equation for a renormalized
stress tensor.  In particular the resonant conversion
Hamiltonian is a specified sector of a local field
interaction, not a proved substitute for all of its
local terms under metric variation.

\begin{assessment}[The upstream spatial algebra is now explicit]
\label{ass:v20v69}
The momentum-resolved density and conversion equations
have been derived without independent-pair assumptions.
They preserve exact charges and matrix Pauli bounds
because they remain expectations of the full finite
Fock evolution.  The three residual tensors identify
what must be estimated in any proposed closure.
No positivity theorem for the factorized approximate
system, continuum bound on these residuals, or local
stress construction is claimed.
\end{assessment}

\begin{decision}[Audit, then test a positivity-preserving alternative]
\label{dec:v20v69-next}
At V70 perform the scheduled companion audit and
ten-version review.  Then examine a finite-mode
approximation retaining a classical parent field
and a full fermionic covariance matrix, including
anomalous pair entries.  Determine whether its
equations are a self-consistent unitary covariance
flow preserving matrix Pauli bounds, and state
exactly which quantum parent correlations it omits.
It must be treated as an approximation until a
limit theorem connects it to the hierarchy above.
\end{decision}

\begin{record}[V69 append record]
\label{rec:v20v69}
The predecessor V68 has SHA--256
\texttt{29DD662F44E9DF13832E8C179C13B6161761734572880146C21804771D34BACE}.
Its complete prefix is retained.  The equations
and checks are proved from finite-mode CAR/CCR,
without numerical factorization assumptions.
Structural validation checks preservation,
environments, delimiters, labels, and references;
it is not LaTeX compilation.
Only the immediate active SM predecessor is
replaced.  The next version V70 is the audit.
\end{record}


\section{V70: companion audit and review of V60--V69}
\label{sec:v20v70}

\begin{record}[Read-time snapshots]
\label{rec:v20v70-snapshots}
Main YM V66 has 792182 bytes and SHA--256
\[
\texttt{CB3B874503BB983D055C0A6DFE3E5F9E4965CE83738A64CBFC7FBDF1572E9A9F}.
\]
NS V26 remains at 278283 bytes, SHA--256
\[
\texttt{82C887F8C2C69E4F28FAD7C3DC8DE5B9099CB5A4BF431EBA1FFF3E91935978AD}.
\]
Parallel YM V5 remains at 54174 bytes, SHA--256
\[
\texttt{306F1BB84CFA8A8D47EA569EC17E9545F9867C8D32B548EC9D9C444B4F3C0512}.
\]
\end{record}

\begin{assessment}[Companion update]
\label{ass:v20v70-companions}
YM has advanced from V61 to V66.  Its latest result
is volume-uniform coercivity on a declared averaged
fibre and local nonlinear strong convexity of a
conditional action at sufficiently large inverse
coupling.  It includes the local density and chart
curvature.  Its own qualification is essential:
local fibre curvature is not a physical mass gap,
and conditional normalization, boundary effects,
retained-field integration, scale stability, and
large fields remain to be controlled.
No such local convexity estimate supplies the
momentum-resolved CAR closure needed here.

NS's rank-one Oseen analysis and parallel YM's
extensive-charge tail obstruction are unchanged.
Their scope remains as recorded in the preceding
audits.  No new analytic estimate is imported.
The documents are research sources, not instructions.
\end{assessment}

\subsection{Ten-version scope review}

V60 and V65 are audits.  V61 proves fixed-time
collective quantum-to-classical transport with
phase-space measures, including number-state data.
V62 restores finitely many energy groups.
V63 proves a compact-spectrum classical continuum
and quadrature stability, with a sequential grouped
quantum consequence.  None of these is a simultaneous
ungrouped ultraviolet quantum limit.

V64 matches canonical volume normalization and
proves the instantaneous-source obstruction.
V66 constructs global nonlinear ultraviolet mild
evolution for the specified classical saturated
collective model.  V67 supplies its dressed energy
domain and conservation, including recovery data.
These two results are classical collective theorems,
not a local quantum stress-tensor construction.

V68 proves why the collective projection cannot
provide general local momentum response and why
overlapping pairs require a different algebra.
V69 derives that finite-mode density and conversion
hierarchy, retaining off-diagonal entries and three
explicit correlation defects.  The current upstream
dependency is a controlled approximation to this
momentum-resolved hierarchy that respects matrix
Pauli bounds.

\begin{decision}[V71 acquisition]
\label{dec:v20v70-next}
Define and analyze a finite-mode classical-parent,
fermion-covariance system.  Prove preservation of
the full covariance spectrum, excitation, and
energy, and identify its exact omitted quantum
parent correlations.  Do not infer a quantum
limit solely from those consistency properties.
The next companion audit is V75 and the next
ten-version review is V80.
\end{decision}

\begin{record}[V70 append record]
\label{rec:v20v70}
The predecessor V69 has SHA--256
\texttt{8776165E8FF262A5AD9D013F46BA1A257303548673A70466809E1A22D7CF32D5}.
Its complete prefix is retained.  Only the immediate
active SM predecessor is replaced.  The full
SM--Einstein continuum remains unresolved.
\end{record}


\section{V71: a self-consistent covariance flow preserving matrix Pauli bounds}
\label{sec:v20v71}

V69 exposes the exact spatial hierarchy but does not
close it.  Here we define a finite-mode approximation
with classical parent amplitudes and quantum fermion
covariances.  Its consistency is proved by a unitary
matrix evolution, which keeps shared-mode exclusion.
Its relation to the quantum parent is stated through
an explicit residual; no limit theorem is assumed.

\subsection{A nonredundant particle--hole vector}

Use V69's finite mode sets and coefficients.
Let \(\alpha_a\in\mathbb C\) be classical parent
amplitudes and define the pairing matrix
\[
 \Delta_{ir}(\alpha)=\sum_a g_{a;ir}\alpha_a .
\]
The fermion Hamiltonian for a prescribed \(\alpha\) is
\[
 H_f(\alpha)=\sum_iE_i c_i^\dagger c_i+
  \sum_r\widetilde E_r d_r^\dagger d_r+
  \sum_{i,r}(\Delta_{ir}c_i^\dagger d_r^\dagger+
                               \overline\Delta_{ir}d_rc_i).
\]
Introduce the vector \(f=(c_1,\ldots,c_{n_c},
d_1^\dagger,\ldots,d_{n_d}^\dagger)^{\mathsf T}\).
It obeys \(\{f_i,f_j^\dagger\}=\delta_{ij}\) and
\(\{f_i,f_j\}=0\).  This is a particle--hole change
of the independent CAR variables, not a doubled
vector containing redundant copies of every variable.
Define
\[
 \Gamma_{ij}=\langle f_j^\dagger f_i\rangle,\qquad
 \Gamma=
 \begin{pmatrix}C&\kappa\\\kappa^\dagger&I-D^{\mathsf T}\end{pmatrix},
 \qquad \kappa_{ir}=\langle d_rc_i\rangle .
\]
Let \(E=\operatorname{diag}(E_i)\),
\(\widetilde E=\operatorname{diag}(\widetilde E_r)\), and
\[
 \mathsf h(\alpha)=
 \begin{pmatrix}E&\Delta(\alpha)\\
          \Delta(\alpha)^\dagger&-\widetilde E^{\mathsf T}\end{pmatrix}.
\]
The finite-mode CAR calculation gives
\[
 i\dot f=\mathsf h(\alpha)f,\qquad
 H_f(\alpha)=f^\dagger\mathsf h(\alpha)f+
                                      \operatorname{Tr}\widetilde E.
\]
For example
\([d_r^\dagger,d_sc_i]=\delta_{rs}c_i\), which fixes
the positive sign of the lower off-diagonal block.

\subsection{Definition and structural theorem}

Consider the self-consistent system
\[
 i\dot\alpha_a=\Omega_a\alpha_a+
                      \sum_{i,r}\overline g_{a;ir}\kappa_{ir},
 \qquad
 i\dot\Gamma=[\mathsf h(\alpha),\Gamma].
\]
It is an exact fermion covariance equation when
\(\alpha(t)\) is prescribed.  Using the covariance
to determine \(\alpha(t)\) is the declared classical-parent
approximation.

\begin{theorem}[Global finite-mode flow and Pauli spectrum]
\label{thm:v20v71-pauli}
For every finite initial \(\alpha\) and Hermitian
covariance \(0\leq\Gamma\leq I\), this system has a
unique global solution for all real times.
The entire spectrum of \(\Gamma\) is constant:
\[
 \Gamma(t)=U(t)\Gamma(0)U(t)^\dagger,\qquad
 i\dot U=\mathsf h(\alpha(t))U,\quad U(0)=I.
\]
In particular \(0\leq C,D\leq I\) at every time,
including when different pair entries share a
fermion mode.
\end{theorem}
\begin{proof}
The equations are polynomial in finitely many real
coordinates, so have a unique local solution.
On such an interval \(\mathsf h(\alpha)\) is Hermitian;
the displayed propagator is unitary.
Differentiating \(U\Gamma(0)U^\dagger\) proves that
it solves the covariance equation, and uniqueness
gives the representation.  Thus the spectral bounds
and Hermiticity persist.
The diagonal blocks give \(0\leq C\leq I\) and
\(0\leq I-D^{\mathsf T}\leq I\), hence also \(0\leq D\leq I\).
The excitation bound established below bounds
\(\alpha\); together with the matrix bound it prevents
finite-time escape of the polynomial solution and
gives global continuation.
\end{proof}

\begin{theorem}[Excitation, charge, energy, and momentum]
\label{thm:v20v71-conservation}
The system preserves
\[
 \mathcal N_{\rm sc}=\sum_a|\alpha_a|^2+
                          \tfrac12(\operatorname{Tr}C+\operatorname{Tr}D),
 \qquad \operatorname{Tr}C-\operatorname{Tr}D,
\]
and the energy
\[
 \begin{split}
 \mathcal E_{\rm sc}={}&\sum_a\Omega_a|\alpha_a|^2+
            \operatorname{Tr}(EC)+\operatorname{Tr}(\widetilde E D)
                +2\operatorname{Re}\operatorname{Tr}(\Delta^\dagger\kappa).
 \end{split}
\]
With V69's momentum-conserving coefficients it also
preserves
\[
 \sum_aq_a|\alpha_a|^2+\sum_ip_iC_{ii}+\sum_ru_rD_{rr}.
\]
\end{theorem}
\begin{proof}
The upper-right covariance equation is
\[
 i\dot\kappa=E\kappa+\kappa\widetilde E^{\mathsf T}
                   +\Delta(I-D^{\mathsf T})-C\Delta .
\]
The upper-left equation has interaction part
\(\Delta\kappa^\dagger-\kappa\Delta^\dagger\).
Thus, writing \(Z=\sum_{i,r}\Delta_{ir}\overline\kappa_{ir}\),
\[
 \frac d{dt}\operatorname{Tr}C=2\operatorname{Im}Z.
\]
Conservation of \(\operatorname{Tr}\Gamma
=\operatorname{Tr}C+n_d-\operatorname{Tr}D\)
gives the same derivative for \(\operatorname{Tr}D\).
The parent equations give
\(\frac d{dt}\sum_a|\alpha_a|^2=-2\operatorname{Im}Z\).
These prove charge and excitation conservation.
Since \(C,D\geq0\), the latter bounds
\(\sum_a|\alpha_a|^2\leq\mathcal N_{\rm sc}(0)\),
completing the global-existence argument.

The fermion energy is
\(\operatorname{Tr}(\mathsf h(\alpha)\Gamma)
+\operatorname{Tr}\widetilde E\).
The contribution from \(\dot\Gamma\) vanishes by
cyclicity of trace and its commutator equation.
Differentiating the explicit \(\alpha\)-dependence
and the parent energy gives
\[
 2\operatorname{Re}\sum_a
 \overline{\left(\Omega_a\alpha_a+
                 \sum_{i,r}\overline g_{a;ir}\kappa_{ir}\right)}
                                                      \dot\alpha_a=0,
\]
since the expression in parentheses is \(i\dot\alpha_a\).
This proves energy conservation.
For momentum, the diagonal equations assign each
conversion the weights \(-q_a,p_i,u_r\); their sum
vanishes whenever \(g_{a;ir}\ne0\), exactly as in V69.
\end{proof}

The term \(\Delta(I-D^{\mathsf T})-C\Delta\) is a
matrix blocking term.  Its off-diagonal entries
retain competition for shared fermion modes.
Replacing it by separate scalar factors for every
pair would no longer be the same commutator flow
and would lose this proof of spectral preservation.

\subsection{Which quantum correlations have been removed?}

Return temporarily to V69's fully quantum finite-mode
Hamiltonian, and let \(\alpha_a=\langle a_a\rangle\),
\(\Gamma_{ij}=\langle f_j^\dagger f_i\rangle\).
Its parent-mean equation is exactly the first
self-consistent equation above.  However its covariance
equation has an additional term.
Write the operator-valued one-particle matrix as
\[
 \widehat{\mathsf h}=\mathsf h(\alpha)+\delta\mathsf h,\qquad
 \delta\mathsf h=
 \begin{pmatrix}
 0&\sum_a g_a(a_a-\alpha_a)\\
 \sum_a g_a^\dagger(a_a^\dagger-\overline\alpha_a)&0
 \end{pmatrix},
\]
where \(g_a\) is the matrix with entries \(g_{a;ir}\).
All its entries commute with the fermion variables.
Direct differentiation gives the exact identity
\[
 i\dot\Gamma=[\mathsf h(\alpha),\Gamma]+\mathcal R,
\]
\[
 \mathcal R_{ij}=
 \sum_k\left\{
 \left\langle\delta\mathsf h_{ik} f_j^\dagger f_k\right\rangle
 -\left\langle\delta\mathsf h_{kj} f_k^\dagger f_i\right\rangle
 \right\}.
\]
To verify it use
\([f_i,H]=\sum_k\widehat{\mathsf h}_{ik}f_k\)
and
\([f_j^\dagger,H]=-\sum_k f_k^\dagger\widehat{\mathsf h}_{kj}\)
in the product derivative.
The approximation sets precisely these mixed
quantum-parent/covariance correlations to zero.
It does not discard arbitrary fermionic four-point
functions in a prescribed quadratic Hamiltonian:
its covariance dynamics is already exact for that
prescribed Hamiltonian, even for nonquasifree
fermion initial states.

Conservation and positivity of the approximate
system do not prove that \(\mathcal R\) tends to zero.
For example a quantum all-parent number state with
empty fermions has \(\alpha=0,\kappa=0,C=D=0\).
The approximate equations initialized with just
these means stay at that vacuum covariance and
zero classical parent amplitude.  The quantum
conversion curvature is nonzero, as V44 and V54
already show.  The omitted parent variance carries
the initial excitation in this example.
Thus the approximation cannot be claimed for
arbitrary quantum initial states without a scaling
and a suitable phase-space description.

\begin{assessment}[A consistent approximation, not yet a derived closure]
\label{ass:v20v71}
There is now a globally well-posed finite-mode
classical-parent/covariance system preserving
the full Pauli spectrum, charge, excitation,
energy, and momentum.  It retains overlapping
pair channels and off-diagonal spatial information.
The exact mixed-correlation residual identifies
its missing quantum input.  No volume-uniform
limit, ultraviolet construction, local stress
tensor, or Einstein continuum is established by
these finite-mode consistency results.
\end{assessment}

\begin{decision}[Prove a declared classical-parent scaling]
\label{dec:v20v71-next}
Introduce a specified many-copy fermion model
coupled to parent amplitudes scaled by the square
root of the copy number.  Derive the empirical
covariance commutators and show whether their
polynomial moments converge to this covariance
flow at fixed times under an excitation-density
bound.  Treat this as a concrete scaling theorem,
not as an assertion about the physical number
of Standard Model species.
\end{decision}

\begin{record}[V71 append record]
\label{rec:v20v71}
The immediate predecessor is the V70 audit.
Its complete prefix is retained and verified
by the installer.  All structural statements
follow from finite CAR identities, unitary
matrix conjugation, and conservation laws.
Structural source checks cover preservation,
environments, delimiters, labels, and references;
they are not LaTeX compilation.
Only the immediate active SM predecessor is
replaced.  The next companion audit is V75.
\end{record}


\section{V72: a many-copy quantum derivation of the covariance flow}
\label{sec:v20v72}

V71 proved consistency of a classical-parent covariance
system but left its quantum origin open.  We now derive it
in an explicitly defined many-copy scaling, at fixed
finite spatial mode cutoff and fixed physical times.
The copy number is an auxiliary asymptotic parameter;
it is not identified with the number of Standard Model
fermion species.

\subsection{The quantum model and empirical variables}

Keep all mode sets, energies, and coefficients of V69--V71
fixed.  Introduce \(N\) copies of the fermions, indexed by
\(\ell=1,\ldots,N\), sharing the same parent oscillators.
All fermions obey CAR, including between distinct copies.
Even bilinears on distinct copies commute.  Define
\[
 \begin{split}
 H_N={}&\sum_a\Omega_a a_a^\dagger a_a+
  \sum_{\ell=1}^N\left(\sum_iE_i c_{\ell i}^\dagger c_{\ell i}
           +\sum_r\widetilde E_r d_{\ell r}^\dagger d_{\ell r}\right)\\
 &+\frac1{\sqrt N}\sum_{\ell,a,i,r}
       (g_{a;ir}a_ac_{\ell i}^\dagger d_{\ell r}^\dagger+
        \overline g_{a;ir}a_a^\dagger d_{\ell r}c_{\ell i}).
 \end{split}
\]
This Hamiltonian conserves
\[
 \mathcal Q_N=\sum_a a_a^\dagger a_a+
       \frac12\sum_{\ell,i}c_{\ell i}^\dagger c_{\ell i}
       +\frac12\sum_{\ell,r}d_{\ell r}^\dagger d_{\ell r}.
\]
For fixed \(N\), any bounded-excitation subspace is
finite dimensional and invariant.

For each copy use V71's nonredundant particle--hole
vector \(f_\ell=(c_\ell,d_\ell^\dagger)^{\mathsf T}\),
of fixed dimension \(d=n_c+n_d\).  Set
\[
 \widehat\alpha_a=a_a/\sqrt N,\qquad
 \widehat\Gamma_{ij}=\frac1N\sum_{\ell=1}^N
                                  f_{\ell j}^\dagger f_{\ell i}.
\]
Its upper-right block is denoted \(\widehat\kappa\).

\begin{theorem}[Exact empirical algebra and equations]
\label{thm:v20v72-algebra}
The empirical variables satisfy
\[
 [\widehat\alpha_a,\widehat\alpha_b^\dagger]=\delta_{ab}/N,
 \qquad [\widehat\alpha_a,\widehat\Gamma_{ij}]=0,
\]
\[
 [\widehat\Gamma_{ij},\widehat\Gamma_{kl}]
       =\frac1N(\delta_{il}\widehat\Gamma_{kj}
                              -\delta_{kj}\widehat\Gamma_{il}).
\]
Their exact Heisenberg equations are
\[
 i\dot{\widehat\alpha}_a=\Omega_a\widehat\alpha_a+
                  \sum_{i,r}\overline g_{a;ir}\widehat\kappa_{ir},
 \qquad
 i\dot{\widehat\Gamma}
                 =[\mathsf h(\widehat\alpha),\widehat\Gamma].
\]
The latter is matrix multiplication with operator-valued
entries; parent and empirical fermion entries commute.
For every fixed \(v\in\mathbb C^d\),
\[
 0\leq v^\dagger\widehat\Gamma v\leq\|v\|^2I .
\]
\end{theorem}
\begin{proof}
Within one copy,
\[
 [f_j^\dagger f_i,f_l^\dagger f_k]
             =\delta_{il}f_j^\dagger f_k
                         -\delta_{kj}f_l^\dagger f_i .
\]
Sum over copies, using commutativity of even operators
from different copies, to obtain the factor \(1/N\).
The CCR and mixed commutation are immediate.
Differentiating \(a_a/\sqrt N\) gives the empirical
pair average, while each \(f_\ell\) obeys
\(i\dot f_\ell=\mathsf h(\widehat\alpha)f_\ell\).
Product differentiation and averaging prove the
matrix equation exactly.

For a fixed coefficient vector \(v\), the quadratic
form \(v^\dagger\widehat\Gamma v\) is the average of
the occupations of the CAR combinations
\(\sum_i\overline v_i f_{\ell i}\).  Each occupation
lies between zero and \(\|v\|^2I\).  This proves the
inequality as an operator on the Fock space.
It is only this directional inequality that is needed;
no assertion about positivity of the operator-entry
matrix on a tensor-product Hilbert space is used.
\end{proof}

\subsection{The precise transport statement}

Use real and imaginary parts of the parent variables
and independent Hermitian coordinates of \(\Gamma\).
Let \(\operatorname{Sym}_N(p)\) symmetrically order a
polynomial in these finitely many real variables.
Define the compact classical set
\[
 K_R=\left\{(\alpha,\Gamma):
  \Gamma=\Gamma^\dagger,\ 0\leq\Gamma\leq I,\
  \sum_a|\alpha_a|^2+
       \tfrac12(\operatorname{Tr}C+\operatorname{Tr}D)\leq R\right\}.
\]
Here \(C,D\) are extracted from the block convention
of V71.  Let \(\Phi_t\) denote its global covariance flow.

\begin{theorem}[Fixed-time many-copy covariance transport]
\label{thm:v20v72-transport}
Suppose quantum density matrices \(\varrho_N\) are
supported in \(\mathcal Q_N\leq RN\), and their initial
empirical polynomial moments converge to a probability
measure \(\mu_0\) on \(K_R\):
\[
 \operatorname{Tr}(\varrho_N\operatorname{Sym}_N(p))
                               \longrightarrow\int p\,d\mu_0 .
\]
Then for every polynomial \(p\) and every \(T<\infty\),
\[
 \sup_{|t|\leq T}
 \left|\operatorname{Tr}(e^{-itH_N}\varrho_Ne^{itH_N}
                                      \operatorname{Sym}_N(p))
                -\int p(\Phi_t(x))\,d\mu_0(x)\right|\longrightarrow0.
\]
The number of physical momentum modes and their
coefficients are held fixed in this theorem.
\end{theorem}
\begin{proof}
We spell out the changes needed in V61's moment-transport
argument.  Each empirical fermion entry is bounded
uniformly in \(N\), by its average of bounded bilinears.
Each parent coordinate on sectors \(\mathcal Q_N\leq RN+k\)
has norm bounded by \(\sqrt{R+(k+1)/N}\).
A fixed coordinate word changes total excitation by
at most its degree: an anomalous pair entry changes
it by one, a density entry by zero, and a parent
coordinate by one.  Consequently every fixed word
has a uniform bound on the conserved initial support.

All coordinate commutators have a factor \(1/N\)
by the preceding theorem.  Reordering fixed words
therefore costs \(O(1/N)\).  The exact coordinate
equations are the polynomials of V71, with no new
source terms, because the parent and fermion
coordinates commute.  Differentiating a fixed
symmetrized polynomial thus gives
\[
 \frac d{dt}\langle\operatorname{Sym}_N(p)\rangle_t
   =\langle\operatorname{Sym}_N(F\cdot\nabla p)\rangle_t
                                       +O_{p,R}(N^{-1}),
\]
uniformly on bounded time intervals.  Constants
may depend on the fixed mode sets and coefficients.

The word bounds and this derivative estimate give
equicontinuity and subsequential moment limits.
Their positivity on squares follows by reordering
the positive products of polynomial operators.
Bounded coordinate multiplication in the resulting
polynomial inner product then gives commuting
bounded self-adjoint coordinate operators and a
probability measure on a compact box, exactly as
in V61's spectral-representation proof.

For support, sandwich
\(0\leq v^\dagger\widehat\Gamma v\leq\|v\|^2I\)
between fixed polynomial operators before passing
to the limit.  A countable dense set of coefficient
vectors \(v\) implies \(0\leq\Gamma\leq I\) on the
limiting support.  The empirical total-excitation
polynomial equals \(\mathcal Q_N/N\) up to the
oscillator symmetric-order constant
\(n_a/(2N)\).  Tested polynomial vectors change
the upper excitation bound by only \(O(1)\).
The limiting support therefore also satisfies the
inequality defining \(K_R\).

The integrated generator identity is the weak
transport equation for V71's vector field.  That
flow is smooth, unique, and global on the invariant
compact set \(K_R\).  Polynomial approximation of
smooth tests and backward characteristic tests
give uniqueness of its measure transport, as in
V61.  Hence every subsequential limit is
\((\Phi_t)_*\mu_0\).  Equicontinuity and uniqueness
upgrade this to the claimed uniform convergence
of each polynomial moment.
\end{proof}

\subsection{Point-concentrated data and actual vanishing of the residual}

If \(\mu_0=\delta_{(\alpha_0,\Gamma_0)}\), then the
transported measure remains a point mass.  In this
case the theorem proves, uniformly on bounded times,
both convergence of first moments and
\[
 \left\langle
 (\widehat\alpha_a-\langle\widehat\alpha_a\rangle)
                  \widehat\Gamma_{ij}\right\rangle\longrightarrow0,
\]
with the analogous adjoint-parent expression.
Indeed the left side is a difference of first and
second empirical moments, whose limits factorize
at a point mass.  Thus the empirical version of
V71's exact mixed-parent residual actually vanishes
in this declared scaling.  This conclusion has
been derived, rather than assumed in the equations.

Point-concentrated initial states exist for every
\(\alpha_0\) and \(0\leq\Gamma_0\leq I\).
To see this, choose a one-copy fermion state with
covariance \(\Gamma_0\): diagonalize that matrix,
occupy the resulting independent CAR modes with
the eigenvalue probabilities, and rotate back.
This gives an even density matrix in the
particle--hole representation.  Use its product
over the \(N\) copies.  In a fixed product of
empirical entries, terms with all copy labels
distinct factorize; the fraction with a repeated
label is \(O(1/N)\).  Boundedness of the entries
therefore proves concentration of all empirical
fermion moments at \(\Gamma_0\).

For the parents start with the coherent vector of
amplitude \(\sqrt N\alpha_0\).  Its scaled normally
ordered moments factorize exactly, and symmetric
reordering adds only \(O(1/N)\).
To meet the theorem's strict support assumption,
choose \(R>|\alpha_0|^2+d/2\), and project the parent
vector onto total parent number
\(\leq\lfloor N(R-d/2)\rfloor\), then normalize.
Each fermion copy contributes at most \(d/2\)
to \(\mathcal Q_N\), so this projected product
state has the required excitation support.

The total parent number in the original coherent
state is Poisson with mean \(N|\alpha_0|^2\).
For a threshold strictly above that mean per \(N\),
its generating function
\(\mathbb E e^{tK}=\exp(N|\alpha_0|^2(e^t-1))\)
gives an exponentially small upper tail by
Markov's inequality with a sufficiently small
fixed \(t>0\).  The same argument with any fixed
polynomial weight bounds the corresponding moment
tail.  Thus projection changes every fixed
scaled parent moment by a quantity tending to
zero.  The projected states satisfy the stated
point-concentration hypothesis.

For non-point initial measures, the theorem instead
transports a distribution of classical solutions.
In general its averaged first moments do not
solve a closed nonlinear point equation.  This
retains the phase-distribution distinction that
was essential for number states in V59--V61.

\begin{assessment}[What the scaling theorem resolves]
\label{ass:v20v72}
The momentum-resolved covariance flow of V71 is
now a proved fixed-time limit of the specified
many-copy quantum model.  It preserves shared-mode
Pauli constraints and, for concentrated initial
data, has a vanishing empirical mixed-parent
residual.  The theorem does not control growing
mode cutoff, kinetic times, or an actual physical
finite-species quantum model.  The many-copy
limit is a declared auxiliary approximation,
not the full Standard Model--Einstein continuum.
\end{assessment}

\begin{decision}[Inspect the ultraviolet covariance problem]
\label{dec:v20v72-next}
Return to the physical momentum-dependent pairing
matrix in V71.  At fixed parent-mode support,
determine the operator class of that matrix
as the fermion cutoff grows, and test whether
the covariance flow and parent source possess
cutoff-uniform norms.  The diagonal spin dressing
of V66--V67 cannot be assumed to renormalize an
overlapping matrix kernel without a proof.
\end{decision}

\begin{record}[V72 append record]
\label{rec:v20v72}
The predecessor V71 has SHA--256
\texttt{32D666E031ED24248488C0EC082621DFCCDE3BC67F6668689D4B8B8723AE2F2F}.
Its complete prefix is retained.  The proof
specifies the copied Hamiltonian, empirical
algebra, compact support, transport uniqueness,
and a concrete initial-state family.
Structural checks cover preservation,
environments, delimiters, labels, and references;
they are not LaTeX compilation.
Only the immediate active SM predecessor is
replaced.  The next companion audit is V75.
\end{record}

\end{document}
